\documentclass[paper=a4, fontsize=11pt]{jlreq} % LuaLaTeX または upLaTeX でのコンパイルを推奨

% --- パッケージ群 ---
\usepackage[utf8]{inputenc}
\usepackage[T1]{fontenc}
\usepackage[margin=20mm]{geometry} % 上下左右の余白を20mmに設定
\usepackage{amsmath, amssymb, amsthm, amsfonts} % 数学の基本
\usepackage{physics} % \ket{0}, \pdv{x} などが簡単に打てる
\usepackage{bm} % 太字ベクトル \bm{v}
\usepackage{graphicx}
\usepackage{hyperref} % リンク・目次用
\usepackage{cite} % 引用用
\usepackage{etoolbox} % 環境の終わりに自動でマークを入れるために追加
\usepackage{ytableau} % ★追加: Young図形の描画に必須
\usepackage{tikz}     % ★追加: TikZの読み込み
\usepackage{mathtools}
\usepackage[notref,notcite]{showkeys}
\usepackage{xcolor}
\usepackage{mdframed}
\usetikzlibrary{arrows.meta, decorations.pathreplacing}% ★追加: TikZの波括弧や矢印のスタイルに必須

% 本文用の環境（色は目に優しいDarkBlueなどを指定）
\newenvironment{story}{%
  \begin{mdframed}[
    backgroundcolor=red!5, % redを5%の濃さで出力（ほんのり薄い赤）
    linewidth=0pt,         % 枠線はなし（スッキリさせるため）
    innertopmargin=10pt,   % 上下の余白
    innerbottommargin=10pt,
    innerrightmargin=10pt, % 左右の余白
    innerleftmargin=10pt
  ]%
}{%
  \end{mdframed}%
}

\newtheoremstyle{break}% スタイル名
  {1em}% 上部のスペース（約1行分空ける）
  {1em}% 下部のスペース（約1行分空ける）
  {\normalfont}% 本文のフォント（日本語なので斜体にしない）
  {}% インデント量
  {\bfseries}% 見出しのフォント（太字）
  {}% 見出し後の句読点（不要なら空に）
  {\newline}% ★ここで改行を指定！★
  {}% 見出しのカスタム指定（デフォルト）

% --- 定理環境の設定 ---
\theoremstyle{break}
\newtheorem{theorem}{定理}[section]
\newtheorem{definition}[theorem]{定義}
\newtheorem{proposition}[theorem]{命題}
\newtheorem{lemma}[theorem]{補題}
\newtheorem{corollary}[theorem]{系}
\newtheorem{example}[theorem]{例}
\newtheorem{problem}[theorem]{問題} % 自作問題用


% --- 環境の終わりにマークをつける設定 ---
\newcommand{\myendmark}{\null\hfill$\diamondsuit$} % 既存のコマンドと被らないよう \myendmark に変更
\AtEndEnvironment{theorem}{\myendmark}
\AtEndEnvironment{definition}{\myendmark}
\AtEndEnvironment{proposition}{\myendmark}
\AtEndEnvironment{lemma}{\myendmark}
\AtEndEnvironment{corollary}{\myendmark}
\AtEndEnvironment{example}{\myendmark}
\AtEndEnvironment{problem}{\myendmark}

% --- Q&A（自問自答）用の専用スタイル ---
\newtheoremstyle{qa}% スタイル名
  {1ex}% 上部のスペース
  {1ex}% 下部のスペース
  {\normalfont}% 本文のフォント（日本語が崩れないようnormalfontを指定）
  {}% インデント量
  {\bfseries}% 見出しのフォント（太字）
  {}% 見出し後の句読点
  {\newline}% 見出しの後に改行を入れる
  {\thmname{#1}\thmnote{：#3}}% カスタム指定（例：「Q：なぜ〜？」と出力）

\theoremstyle{qa}
\newtheorem*{Q}{Q}

\title{パンルヴェ方程式勉強ノート}
\author{松浦光佑}
\date{\today}

\begin{document}
\maketitle
\tableofcontents % 自動で目次を作成

% 各章の読み込み
\input{sections/01_intro.tex}
% !TEX root = ../main.tex

\section{フックス型微分方程式}

\subsection{2階線型常微分方程式}
%#region
\begin{equation}\label{eq:2-ode}
\frac{d^{2}y}{dx^{2}} + p_{1}(x)\frac{dy}{dx} + p_{2}(x)y = 0   
\end{equation}
を考える。$x\in\mathbb{C}$とする。よって解$y$は解析的である。
また、$p_{1}(x),p_{2}(x)$は$x$の有理関数とする。よって\eqref{eq:2-ode}は広くても$\mathbb{P}^{1}(\mathbb{C})$上の微分方程式である。
%#endregion
%#region
\begin{definition}[$x=\infty$における微分方程式]
    \eqref{eq:2-ode}で$x=\frac{1}{u}$とする。$u=0$の近傍で考えたものを\textbf{$x=\infty$における微分方程式}という。
    （\cite[p.49]{book:okamoto} を参照）
\end{definition}
%#endregion
%#region
ある有理関数$q_{1},q_{2}$を用いて
\[
\frac{d^{2}y}{du^{2}} + q_{1}(u)\frac{dy}{du} + q_{2}(u)y = 0
\]
のように書いたとき、\eqref{eq:2-ode}に連鎖律 $\frac{d}{dx} = -u^2 \frac{d}{du}$ を用いたものとの係数比較により、
\[
q_{1}(u) = \frac{2}{u}-\frac{1}{u^{2}}p_{1}\left(\frac{1}{u}\right),\, q_2(u)=\frac{1}{u^{4}}p_{2}\left(\frac{1}{u}\right)
\]
という関係がわかる。
%#endregion
%#region
$p_{1}, p_{2}$ を $x=x_{0}$ で正則とする。このとき、コーシーの存在定理（Cauchy's existence theorem）から、任意の $(y_{0}, y'_{0}) \in \mathbb{C}^{2}$ に対して、初期条件
$$
y(x_{0}) = y_{0}, \quad \frac{dy}{dx}(x_{0}) = y'_{0}
$$
を満たす微分方程式 \eqref{eq:2-ode} の解 $y = \varphi(x; x_{0}, y_{0}, y'_{0})$ がただ一つ存在する。

ここで、$u = x - x_{0}$ とおく。$p_1, p_2$ および解 $\varphi$ は $x=x_0$ で正則であるから、以下のようにテイラー展開できる。
$$
\varphi(x; x_{0}, y_{0}, y'_{0}) = \sum_{n=0}^{\infty} c_{n} u^{n} \quad (c_{0} = y_{0}, \, c_{1} = y'_{0})
$$
$$
p_{1}(x) = \sum_{n=0}^{\infty} a_{n} u^{n}, \quad p_{2}(x) = \sum_{n=0}^{\infty} b_{n} u^{n}
$$
これらを \eqref{eq:2-ode} に代入すると、
$$
\sum_{n=2}^{\infty} c_{n} n(n-1) u^{n-2} + \left( \sum_{n=0}^{\infty} a_{n} u^{n} \right) \left( \sum_{n=1}^{\infty} c_{n} n u^{n-1} \right) + \left( \sum_{n=0}^{\infty} b_{n} u^{n} \right) \left( \sum_{n=0}^{\infty} c_{n} u^{n} \right) = 0
$$
となる。この式の $u$ の各べきの係数を比較することで、$c_n$ を漸化式として具体的に順次決定することができる。

\begin{definition}[解析関数の特異点の厳密な定義]\label{def:singular_point_strict}
    ある領域 $D$ で正則な関数（微分方程式の解の分枝）$f(x)$ について、その境界上の点 $\zeta \in \partial D$ が特異点であるとは、$\zeta$ を中心とするいかなる小さな開円板 $U$ をとっても、$f(x)$ を $D$ から $U$ へ解析接続できないこと、すなわち「共通部分 $D \cap U$ において $f(x)$ と一致する $U$ 上の正則関数が存在しないこと」をいう。
\end{definition}

\begin{Q}[【核心】「限界まで」という曖昧な表現を、「いかなる近傍への解析接続も不可能」と厳密化することで得られる論理的な利点は何か？]
\textbf{A.} 特異点であることの定義が「条件を満たす正則関数が存在しない」という明確な否定命題になるため、今回のように「$\xi$ を含む開円板上で正則なもの（解析接続）が構成できる」ことを示せば、直ちに「特異点の定義を満たさない」と結論づける論理（対偶や背理法）が数学的な隙なく回るようになるからである。
\end{Q}

\begin{Q}[【核心】なぜ特異点を単なる「最初の収束円の外側」ではなく、「解析接続の限界（境界）」として定義する必要があるのか？]
\textbf{A.} 単に一つの収束円の外側というだけでは、別の点を中心に展開し直す（解析接続する）ことでその先の点でも正則になり得るため、「いかなる局所的な展開を用いても絶対に到達・通過できない行き止まり」を特異点と定義して初めて、解の真の存在範囲と性質を議論できるからである。
\end{Q}

また、得られた級数 $\sum_{n=0}^{\infty} c_{n} u^{n}$ の収束半径を $r$ とすると、命題\ref{prop:singular_point_on_convergence_circle}により、収束円周 $|x-x_0|=r$ 上には解 $\varphi$ の特異点が少なくとも1つ存在する。

\begin{Q}[【補足】解の冪級数の収束半径 $r$ は、具体的にどのくらい大きくなることが保証されているか？]
\textbf{A.} 初期値 $x_0$ から、「係数関数 $p_1(x), p_2(x)$ の特異点」までの距離の最小値以上になることが保証されている。線形常微分方程式においては、解の特異点は係数の特異点以外には決して新たに生じない（特異点が固定されている）という強力な性質があるためである。
\end{Q}

\begin{Q}[【補足】「ある展開中心での冪級数の収束円上の特異点」は、大域的な関数にとっての「真の特異点」と常に一致するのか？]
\textbf{A.} 常に一致するとは限らない。収束円上の特異点はあくまで「その点を中心としたテイラー展開が破綻する限界点」に過ぎず、別の点を経由して解析接続すれば正則に振る舞うこともある。しかし、線形常微分方程式においては前述の通り特異点が係数によって固定されるため、収束限界点は「係数の特異点」あるいは「多価関数の分岐点」といった真の特異点構造に由来するものと一致する。
\end{Q}
%#endregion
%#region
\begin{proposition}\label{prop:linear_ode_singularities_fixed}
    $x=\xi$ を $y=\varphi(x;x_{0};y_{0},y_{0}')$ の特異点とすれば、$\xi$ は $p_{1}(x)$ または $p_{2}(x)$ の特異点である。
    （\cite[p.50 命題2.1]{book:okamoto} を参照）
\end{proposition}

$n$階でも成り立つので以下$n$階で示す。

\begin{proof}
    対偶を示す。すなわち「$x=\xi$ が係数 $p_{1}(x),\dots,p_{n}(x)$ の正則点であれば、解 $y=\varphi(x;x_{0};y_{0},\dots,y_{0}^{(n-1)})$ は $x=\xi$ において正則である（特異点を持たない）」ことを示す。

    $\xi$ に向かう解析接続の経路上で、$\xi$ に十分近い点 $x=x_{1}$ をとる。具体的には、$x_{1}$ から「係数 $p_1(x),\dots,p_n(x)$ の最も近い特異点」までの距離を $R$ としたとき、$|x_1 - \xi| < R$ となるように $x_1$ を選ぶ（$\xi$ は係数の正則点なので、$\xi$ に十分近づけば必ずこのような $x_1$ が取れる）。

    系\ref{cor:linear_maximal_disk}より、$x_1$ での解の冪級数展開は、少なくとも半径 $R$ の開円板上で正則となる。
    $|x_1 - \xi| < R$ であるから、この開円板は $x=\xi$ を完全に含んでいる。
    したがって、解 $y=\varphi(x;x_{0};y_{0},\dots,y_{0}^{(n-1)})$ は $x=\xi$ を含む開円板上まで解析接続可能となり、$x=\xi$ で正則となる。
    ゆえに定義\ref{def:singular_point_strict}より、$x=\xi$ は解の特異点とはなり得ない。
\end{proof}

\begin{Q}[【補足】なぜ線形微分方程式では、解が勝手に（係数にない場所で）特異点を持つことはないのか？]
\textbf{A.} コーシーの存在定理により、係数が正則である限り、解は局所的に必ず冪級数として一意に構成でき、その収束半径は次の「係数の特異点」にぶつかるまで解析接続で伸ばし続けられるからである。
\end{Q}

\begin{Q}[【重要】この命題が、後のパンルヴェ方程式の学習においてなぜ重要になるのか？]
\textbf{A.} 線形方程式では特異点の場所が初期値によらず係数のみで決まる（固定特異点）のに対し、パンルヴェ方程式のような「非線形」方程式では、初期値によって解の特異点の場所が動く（動く特異点）という決定的な違いを対比させるための基礎となるからである。
\end{Q}
%#endregion
%#region
\begin{definition}[線形微分方程式の特異点]\label{def:linear_equation_singularity}
    \textbf{線形微分方程式の特異点}とは、係数関数 $p_{1}(x), p_{2}(x)$ の特異点のことをいう。
    また、$x=\infty$ での特異点は、$x=\frac{1}{u}$ と変数変換したときの $u=0$ における係数の特異点として定義する。
\end{definition}

\begin{Q}[【核心】命題\ref{prop:linear_ode_singularities_fixed}の逆、すなわち「方程式の係数が特異点を持つならば、解も必ずそこで特異点を持つ」は成り立つか？]
\textbf{A.} 成り立たない。以下の例のように、方程式の係数としては特異点であっても、すべての解がそこで正則に振る舞う「見掛けの特異点（apparent singularity）」が存在し得るからである。
\end{Q}

\begin{example}[見掛けの特異点の例]\label{ex:apparent_singularity}
    次の1階線形微分方程式を考える。
    \begin{equation}\label{eq:apparent_ex}
    \frac{dy}{dx}=\frac{y}{x}
    \end{equation}
    この方程式は係数関数が $1/x$ であるため、$x=0$ に特異点をもつ。
    しかし、この微分方程式の一般解は任意の初期値（積分定数）$C \in \mathbb{C}$ に対して $y=Cx$ となり、$x=0$ において正則である。
    したがって、$x=0$ は方程式の特異点であるが、解の特異点とはならない。
\end{example}

\begin{definition}
    このように、線形微分方程式の特異点であって、任意の初期値に対して、解の正則点となる点を\textbf{みかけの特異点}という。
\end{definition}

%#endregion
%#region

$\Xi$を\eqref{eq:2-ode}の特異点の集合とし、$\mathbb{X}=\mathbb{P}^{1}(\mathbb{C})\setminus\Xi$とする。解の基本系を$\vec{\varphi}(x)$とする。

\begin{definition}
    \[
    W(\vec{\varphi})=
    \left[
    \begin{matrix}
        \varphi_{1}(x) & \varphi_{2}(x) \\
        \varphi_{1}'(x) & \varphi_{2}'(x) \\
    \end{matrix}
    \right]
    \]
    を$\vec{\varphi}(x)$の\textbf{ロンスキー行列}といい、
    \[
    w(\vec{\varphi})=\text{det}(W(\vec{\varphi}))
    \]
    を$\vec{\varphi}(x)$の\textbf{ロンスキアン}という
\end{definition}

\begin{proposition}\label{prop:wronskian-ode}
    \eqref{eq:2-ode}の解の基本系$\vec{\varphi}(x)$のロンスキアンは
    \[
    \frac{dw}{dx}+p_{1}(x)w=0
    \]
    なる$1$階微分方程式を満たす。
    （\cite[p.51 命題2.2]{book:okamoto} を参照）
\end{proposition}

この命題は$n$階でも成り立つ。

\begin{proof}
    定理\ref{thm:derivative-of-determinant}から、
    \[
    \frac{d}{dx}(w(\vec{\varphi})(x))=\sum_{i=1}^{n}w_{i}(\vec{\varphi})(x)
    \]
    である。ここで、$w_{i}(\vec{\varphi})(x)$は$w(\vec{\varphi})(x)$の$i$行目を微分した行列の行列式である。
    ここで$i=1,\cdots,n-1$に対して、$i$行目と$i+1$行目が一致するので、
    \[
    w_{i}(\vec{\varphi})(x)=0
    \]
    よって
    \[
    \frac{d}{dx}(w(\vec{\varphi})(x))=w_{n}(\vec{\varphi})(x)
    \]
    だが、$w_{n}(\vec{\varphi})(x)$の$n$行目について
    \[
    (\varphi_{1}^{(n)}(X),\cdots,\varphi_{n}^{(n)}(x))
    \]
    となっているが、満たす微分方程式によって
    \[
    (-p_{1}(x)\varphi_{1}^{(n-1)}(x)-\cdots-p_{n}(x)\varphi_{1}(x),\cdots,-p_{1}(x)\varphi_{n}^{(n-1)}(x)-\cdots-p_{n}(x)\varphi_{n}(x))
    \]
    とあらわせる。detの多重線形性と合わせると$n-2$階微分以下は0になるので、
    \[
    w_{n}(\vec{\varphi})(x)=-p_{1}(x)w(\vec{\varphi})(x)
    \]
    である。よって、
    \[
    \frac{d}{dx}(w(\vec{\varphi})(x))+p_{1}(x)w(\vec{\varphi})(x)=0
    \]
\end{proof}

\begin{proposition}
$\mathbb{X}$内の$x_{0}$を始点とする2つの閉じた道が互いにホモトープならば、
\[
\Gamma(\gamma)=\Gamma(\gamma ')
\]
が成り立つ
証明は定理\ref{thm:monodromy}を参照
\end{proposition}

\begin{Q}[一価性定理の仮定は満たしてるの？]
\textbf{A.} 係数が正則なら定理\ref{thm:linear_ode_existence_and_continuation}からすべての曲線にそって解析接続可能である。\\
\end{Q}

\begin{Q}[ホモトープって何？]
\textbf{A.} 曲線が連続変形できること。(定義\ref{def:homotope}を参照)\\
\end{Q}

\begin{Q}[解を解析接続したものも解になるの？]
\textbf{A.} なる。(定理\ref{thm:ode_solution_continuation}を参照)\\
\end{Q}
%TODO(基本群等について書いてもいい)

\begin{definition}[モノドロミー行列]
    $\Gamma(\gamma)$を\textbf{$\gamma$に対する$\vec{\varphi}(x)$の回路行列}、あるいは、\textbf{$\gamma$に関する$\vec{\varphi}(x)$のモノドロミー行列}という。
\end{definition}

$\Gamma(\gamma_{2}\cdot\gamma_{1})=\Gamma(\gamma_{2})\Gamma(\gamma_{1})$や、$\Gamma(\gamma^{-1})=\Gamma(\gamma)^{-1}$がわかるので、モノドロミー行列は通常の積に関して群構造をもち、その群を$\mathcal{G}(\vec{\varphi}(x))$とかく。\\
また基本群$\pi_1(x_{0};\mathbb{X})$から$\mathcal{G}(\vec{\varphi}(x))$への準同型
\[
\rho([\gamma])=\Gamma(\gamma)
\]
があり、これは$\pi_1(x_{0};\mathbb{X})$の$\mathbb{C}^{2}$上の表現となっている。

\begin{Q}[$\pi_1(x_{0};\mathbb{X})$の$\mathbb{C}^{2}$上の表現って？]
\textbf{A.} 定義\ref{dfn:representation}を参照。\\
\end{Q}

\begin{definition}
    $\mathcal{G}(\vec{\varphi}(x))$を、\eqref{eq:2-ode}の解の基本形$\vec{\varphi}(x)$の\textbf{モノドロミー群}という。
\end{definition}

ここでもう一つの解の基本系 $\vec{\psi}(x)$ をとっておくと、可逆行列 $C$ を用いて
\[
\vec{\psi}(x) = \vec{\varphi}(x)C
\]
となる。また、解析接続と線型操作は交換できるので、
\[
\vec{\psi}^{\gamma}(x) = \vec{\varphi}^{\gamma}(x)C
\]
よって、$\vec{\psi}(x)$ のモノドロミー行列は、
\begin{equation}\label{eq:monodromy_change_of_basis}
\vec{\psi}^{\gamma}(x) = \vec{\varphi}^{\gamma}(x)C= \vec{\varphi}(x)\Gamma(\gamma)C = \vec{\psi}(x)C^{-1}\Gamma(\gamma)C
\end{equation}
より、$C^{-1}\Gamma(\gamma)C$ となる。
よって、
\[
\mathcal{G}(\vec{\psi}(x))=C^{-1}\mathcal{G}(\vec{\varphi}(x))C
\]
より、表現の同値類、$\hat{\rho}:\pi(x_{0};\mathbb{X})\to GL(\mathbb{C})/\sim$を考えることができる。

\begin{definition}[モノドロミー]\label{def:monodromy}
     $\hat{\rho} : \pi_{1}(x_{0};\mathbb{X})\to GL(\mathbb{C})/\sim ; [\gamma] \mapsto [\Gamma(\gamma)]$ を\textbf{モノドロミー}という。
\end{definition}
%TODO (例を書いてもいい)
%#endregion

\subsection{確定特異点}
%#region
2階微分方程式\eqref{eq:2-ode}を考える。$x=\xi$を\eqref{eq:2-ode}の特異点とする。\eqref{eq:2-ode}の特異点ではない点$x_{0}$を取り、$x_{0}$を始点・終点として、$\xi$の周りを1周する閉曲線$\gamma$をとる。

\begin{Q}[そんな閉曲線$\gamma$は常にとれるの？]
\textbf{A.} 今は$p_{1}, p_{2}$を有理関数としているため特異点は離散的である。したがって、他の特異点を巻き込まないように$\xi$の周りを回る道をとることができる。
\end{Q}

\eqref{eq:2-ode}の解の基本系$\vec{\varphi}(x)$の$\gamma$に沿った解析接続を考える。$u=x-\xi$とする。このとき、解が
\[
\vec{\varphi}(x)=(\varphi_{1}(x),\varphi_{2}(x))=\left(u^{\alpha}\sum_{n=0}^{\infty}a_{n}u^{n},u^{\beta}\sum_{n=0}^{\infty}b_{n}u^{n}\right)
\]
とかけるとする。

\begin{Q}[冪級数部分の収束円盤内に$x_{0}$をとることはできる？]
\textbf{A.} できる。特異点は離散的なので、$\xi$から最も近い他の特異点までの距離より小さい収束半径を持つ円盤を考え、その内部に$x_{0}$を設定すればよい。
\end{Q}

このとき、$\gamma$に沿った解の解析接続$\vec{\varphi}^{\gamma}(x)$は、$u \to ue^{2\pi i}$ の変化（偏角が $2\pi$ 増える）を考慮して
\[
\vec{\varphi}^{\gamma}(x)=\left(u^{\alpha}e^{2\pi i\alpha}\sum_{n=0}^{\infty}a_{n}u^{n},u^{\beta}e^{2\pi i\beta}\sum_{n=0}^{\infty}b_{n}u^{n}\right)
\]
となり、行列表記を用いると
\[
\vec{\varphi}^{\gamma}(x)=\vec{\varphi}(x)
\begin{pmatrix}
    e^{2\pi i\alpha} & 0 \\
    0 & e^{2\pi i\beta} \\
\end{pmatrix}
=
\begin{pmatrix}
    e(\alpha) & 0 \\
    0 & e(\beta) \\
\end{pmatrix}
=:\vec{\varphi}(x)\Gamma(\gamma)
\]
となる。\\
また、任意の解の基本系$\vec{\phi}(x)$について、\eqref{eq:2-ode}は線型であるから、ある可逆行列$C$を用いて$\vec{\phi}(x)=\vec{\varphi}(x)C$とかける。このとき、モノドロミー行列は\eqref{eq:monodromy_change_of_basis}から$C^{-1}\Gamma(\gamma) C$とかける。

以下では、具体的な解の表示を仮定せずに一般論を進める。
\eqref{eq:2-ode}の解の基本系を$\vec{\varphi}(x)$とし、閉曲線$\gamma$に沿った解析接続のモノドロミー行列を$\Gamma(\gamma)$とする。その固有値を $e^{2\pi i \alpha}, e^{2\pi i \beta}$ とおく（ここでは記述の簡略化のため $e(\alpha)=e^{2\pi i \alpha}$ のように書くこととする）。

\begin{Q}[$\alpha, \beta$ の取り方には不定性がある？]
\textbf{A.} ある。任意の整数 $n \in \mathbb{Z}$ に対して $e(\alpha) = e(\alpha+n)$ となるため、$\alpha, \beta$ の値には整数分の不定性が生じる。しかし、ここではある特定の $\alpha, \beta$ を一つ固定したとして議論を進める。
\end{Q}

\begin{Q}[なぜ $\alpha-\beta \in \mathbb{Z}$ かどうかで場合分けするのか？]
\textbf{A.} モノドロミー行列 $\Gamma(\gamma)$ が対角化可能か、それともジョルダン細胞（非対角成分に1が残る形）になるかの分岐点だからである。特性指数の差が整数の場合、解に $\log$（対数項）が現れる可能性がある。
\end{Q}

$\alpha-\beta \in \mathbb{Z}$ の場合と $\alpha-\beta \not\in \mathbb{Z}$ の場合で分けて考える。

\begin{itemize}
    \item[\textbf{[1]}] $\alpha-\beta \notin \mathbb{Z}$ のとき\\
    このとき2つの固有値は異なるので $\Gamma(\gamma)$ は対角化可能である。ある正則行列$P$を用いて
    \[
    P^{-1}\Gamma(\gamma)P=
    \begin{pmatrix}
        e(\alpha) & 0 \\
        0 & e(\beta) \\
    \end{pmatrix}
    \]
    とできる。よって、\eqref{eq:monodromy_change_of_basis}により $\vec{\psi}(x)=\vec{\varphi}(x)P$ なる解の基本系に取り直すことで、モノドロミー行列は最初から対角行列であるとしてよい。
    ここで、行列関数 $\Phi(x)$ を
    \[
    \Phi(x)=
    \begin{pmatrix}
        (x-\xi)^{\alpha} & 0 \\
        0 & (x-\xi)^{\beta} \\
    \end{pmatrix}
    \]
    で定める。$\Phi(x)$ が $\gamma$ に沿って1周すると $\Phi^{\gamma}(x) = \Phi(x) \begin{pmatrix} e(\alpha) & 0 \\ 0 & e(\beta) \end{pmatrix}$ となる。
    ここで $\vec{z}(x)=\vec{\psi}(x)\Phi(x)^{-1}$ とおくと、
    \[
    \vec{z}^{\gamma}(x) = \vec{\psi}^{\gamma}(x)\left(\Phi^{\gamma}(x)\right)^{-1} 
    = \vec{\psi}(x)\begin{pmatrix} e(\alpha) & 0 \\ 0 & e(\beta) \end{pmatrix} \begin{pmatrix} e(-\alpha) & 0 \\ 0 & e(-\beta) \end{pmatrix}\Phi(x)^{-1} 
    = \vec{\psi}(x)\Phi(x)^{-1} = \vec{z}(x)
    \]
    となり、$\vec{z}(x)$ は $x=\xi$ の周りで1価であることがわかる。

    \item[\textbf{[2]}] $\alpha-\beta \in \mathbb{Z}$ のとき\\
    $\Gamma(\gamma)$ が対角化可能であれば、[1]のときと同様にして $\vec{z}(x)$ は1価となる。
    一方、$\Gamma(\gamma)$ が対角化不可能で、ある正則行列$P$を用いて
    \[
    P^{-1}\Gamma(\gamma)P=
    \begin{pmatrix}
        e(\alpha) & 1 \\
        0 & e(\alpha) \\
    \end{pmatrix}
    \]
    となるジョルダン細胞のケースを考える。[1]と同様に解の基本系 $\vec{\psi}(x)$ を取り直すことで、モノドロミー行列がこの形であるとしてよい。このとき、
    \[
    \Phi(x)=
    \begin{pmatrix}
        (x-\xi)^{\alpha} & \frac{1}{2\pi i e(\alpha)}(x-\xi)^{\alpha}\log(x-\xi) \\
        0 & (x-\xi)^{\alpha} \\
    \end{pmatrix}
    \]
    と置く。$\log(x-\xi)$ は1周すると $\log(x-\xi) + 2\pi i$ となるため、
    \[
    \Phi^{\gamma}(x)= \Phi(x)
    \begin{pmatrix}
        e(\alpha) & 1 \\
        0 & e(\alpha) \\
    \end{pmatrix}
    \]
    を満たすことが確認できる。したがって、[1]と同様に $\vec{z}(x)=\vec{\psi}(x)\Phi(x)^{-1}$ と置くことで、逆行列との積が打ち消し合い $\vec{z}^{\gamma}(x) = \vec{z}(x)$ となるため、$\vec{z}(x)$ が1価であることがわかる。
\end{itemize}

以上より、以下の極めて重要な命題が成り立つ。

\begin{proposition}\label{prop:single_valued_z}
    \eqref{eq:2-ode}の任意の解の基本系 $\vec{\varphi}(x)$ に対して、ある適当な2次の行列関数 $\Phi(x)$ を選ぶことで、
    \[
    \vec{z}(x)=\vec{\varphi}(x)\Phi(x)^{-1} \quad (\text{すなわち } \vec{\varphi}(x) = \vec{z}(x)\Phi(x))
    \]
    で定められる $\vec{z}(x)$ が、$x=\xi$ の周りで1価の解析関数になるようにできる。\\
    また、$\Gamma(\gamma)$が対角化可能であるという条件のもとで、上の$\vec{\psi}(x)$に対しては$\Phi(x)$を対角行列になるようにとることができる。
    （\cite[p.59 命題2.5]{book:okamoto} を参照）
\end{proposition}

命題\ref{prop:single_valued_z}の$\vec{z}(x)$は定理\ref{thm:laurent_expansion}より、$x=\xi$の周りで、
\begin{equation}\label{eq:power_series_z}
\vec{z}(x)=\sum_{n=-\infty}^{\infty}\vec{a}_{n}u^{n}
\end{equation}
とローラン級数に展開される。


\begin{definition}[微分方程式の確定特異点]
    \eqref{eq:power_series_z}において、ある整数$j_{0}$があって、$j\leq j_{0}$となるすべてのjに対し、
    \[
    \vec{a}_{j} = \vec{0}
    \]
    となる時$\xi$は\eqref{eq:2-ode}の\textbf{確定特異点}という。確定特異点でない時\textbf{不確定特異点}という。
    （\cite[p.59 定義2.4]{book:okamoto} を参照）
\end{definition}

よって微分方程式の特異点$x=\xi$が微分方程式の意味で確定特異点でない時、命題\ref{prop:single_valued_z}の$\vec{z}(x)$は$x=\xi$を真性特異点としてもつ。\\

微分方程式だけでなく解析関数の分枝にも確定特異点は定義されている。\\
$\theta_{1},\theta_{2}\in\mathbb{R}$と$r>0$に対して、
\[
S(\theta_{1},\theta_{2}) = \{u\,|\, 0<|x-\xi|<r, \theta_{1} < \arg u < \theta_{2} \}
\]
とおく。ここで$u=x-\xi$である。

\begin{definition}[解析的な解の確定特異点]
    $x=\xi$で特異点$\omega$をもつ解析関数の分枝$\varphi(x)$に対して、$S(\theta_{1},\theta_{2})$で$\varphi(x)$が定義されるように$r$を十分小さくとる。
    ある$A>0$がとれて$\theta_{1}<\theta_{2}$に対し、
    \[
    \lim_{\substack{u\to 0 \\ u\in S(\theta_{1},\theta_{2})}}|u|^{A}|\varphi(x)| = 0
    \]
    となるとき、$\varphi(x)$は$x=\xi$に\textbf{確定特異点}を持つという。
    （\cite[p.60 定義2.5]{book:okamoto} を参照）
\end{definition}

\begin{Q}[なぜ特異点をわざわざ $\omega$ と書き分けているの？]
\textbf{A.} $\varphi(x)$ が対数関数や累乗根のような多価関数であり、その「分枝（branch）」を厳密に区別するためである。
\begin{itemize}
    \item \textbf{厳密な理由（原著の立場）:} 複素幾何学（リーマン面の理論）の観点では、ベース空間の $x=\xi$ という1点の上にも、分枝の選び方によって異なる複数の特異点 $\omega_1, \omega_2, \dots$ が重なって存在し得る。
    \item \textbf{本ノートでの扱い:} 確定特異点の定義における極限評価 $\lim |u|^A|\varphi(x)|=0$ では、ベース空間での距離 $|u| = |x-\xi|$ のみが計算に関与する。「特定の分枝を選び、扇形領域内で近づく」という前提さえあれば解析的な論理は完結するため、本ノートでは無駄を省き、原則として $x=\xi$ と表記を統一する。
\end{itemize}
\end{Q}

\begin{example}\label{ex:regular_singular_points}
    $\alpha\in\mathbb{C}$に対して、以下の関数の特異点の性質を調べる。
    
    \textbf{(1)} $x^{\alpha}$ は $x=0$ に確定特異点を持つ。
    \begin{Q}[なぜ $x=0$ は確定特異点といえるのか？]
    \textbf{A.} 定義に従い $\lim_{x\to 0} |x|^A |x^\alpha| = 0$ となる $A>0$ が存在するか調べる。
    $x = r e^{i\theta}$ とおくと、
    \[
    |x^{\alpha}| = |e^{\alpha(\log r + i\theta)}| = e^{\text{Re}(\alpha)\log r - \text{Im}(\alpha)\theta} = e^{-\text{Im}(\alpha)\theta} r^{\text{Re}(\alpha)}
    \]
    となる。扇形領域内で $\theta$ は有界であるから、$e^{-\text{Im}(\alpha)\theta}$ はある定数 $B(\theta)$ で抑えられる。
    ここで $|x|^A |x^\alpha| \propto r^{A + \text{Re}(\alpha)}$ となるため、$A + \text{Re}(\alpha) > 0$、すなわち $A > -\text{Re}(\alpha)$ となるように十分大きな $A>0$ をとれば、極限は $0$ になる。
    \end{Q}

    \textbf{(2)} $x^{\alpha}\log x$ は $x=0$ に確定特異点を持つ。
    \begin{Q}[$\log x$ があっても確定特異点になるのはなぜ？]
    \textbf{A.} 対数関数の発散は、多項式（べき乗）の減衰よりも遅いためである。
    (1)と同様に考えると $|x^{\alpha}\log x| \propto r^{\text{Re}(\alpha)}|\log r + i\theta|$ となる。
    極限 $\lim_{r\to 0} r^{A + \text{Re}(\alpha)} |\log r + i\theta|$ を評価する際、任意の $\epsilon > 0$ に対して $\lim_{r\to 0} r^\epsilon \log r = 0$ であるから、(1)と同じく $A > -\text{Re}(\alpha)$ となるように $A$ をとれば、多項式部分が $\log$ を押し潰して極限は $0$ になる。
    \end{Q}

    \textbf{(3)} $x^{\alpha}e^{x}$ の特異点 $x=\infty$ は確定特異点ではない（不確定特異点である）。
    \begin{Q}[なぜ $x=\infty$ で条件を満たさないのか？]
    \textbf{A.} 指数関数の爆発スピードが、いかなる多項式の減衰スピードも上回ってしまうためである。
    $u = 1/x$ とおいて $u \to 0$ の極限を考える。例えば正の実軸に沿って $u \to +0$ と近づく場合、
    \[
    |u|^A |x^{\alpha}e^{x}| = |u|^{A - \text{Re}(\alpha)} e^{\frac{1}{u}}
    \]
    となる。$u \to +0$ のとき、指数関数 $e^{1/u}$ の発散は多項式 $|u|^{A-\text{Re}(\alpha)}$ の減衰よりも圧倒的に速いため、$A$ をどれだけ大きく（正に）とっても極限は $0$ にならない。したがって確定特異点ではない（真性特異点である）。
    \end{Q}
    
    （\cite[p.60 例2.2]{book:okamoto} を参照）
\end{example}

また、次の極めて重要な命題が成り立つ。

\begin{proposition}\label{prop:equivalence_regular_singular}
    線型微分方程式\eqref{eq:2-ode}の特異点$x=\xi$が確定特異点であるための必要十分条件は、\eqref{eq:2-ode}の任意の解の基本系$\vec{\phi}(x)$の各成分$\phi_{1}(x),\phi_{2}(x)$が$x=\xi$上に確定特異点を持つことである。
    （\cite[p.61 命題2.6]{book:okamoto} を参照）
\end{proposition}

\begin{Q}[代数的な定義（ローラン展開）と解析的な定義（増大度）が一致するのはなぜ？]
\textbf{A.} 複素関数論における「孤立特異点の分類定理（リーマンの定理）」が背後で強力に効いているためである。証明の核心は前節で示した定理\ref{thm:growth_and_pole}（極と増大度の同値性）にある。この定理を見通しよく使うために、まず以下の補題を用意する。
\end{Q}

\begin{lemma}[多項式増大度の和と積]\label{lem:polynomial_growth_closure}
関数 $f(x), g(x)$ が $x=\xi$ （$u = x-\xi = 0$）の近傍で「多項式増大度」を持つとする。すなわち、ある定数 $A, B \ge 0$ と $C_1, C_2 > 0$ が存在して、十分小さな穴あき近傍 $0 < |u| < r$ で
\[
|f(x)| \le C_1|u|^{-A}, \quad |g(x)| \le C_2|u|^{-B}
\]
を満たすとする。このとき、$f(x)+g(x)$ および $f(x)g(x)$ もまた $x=\xi$ の近傍で多項式増大度を持つ。
具体的には、ある定数 $C_3, C_4 > 0$ に対して以下が成り立つ。
\begin{enumerate}
    \item $|f(x)+g(x)| \le C_3|u|^{-\max(A, B)}$
    \item $|f(x)g(x)| \le C_4|u|^{-(A+B)}$
\end{enumerate}
\end{lemma}

\begin{proof}
和については三角不等式より $|f(x)+g(x)| \le |f(x)| + |g(x)| \le C_1|u|^{-A} + C_2|u|^{-B}$ であり、$|u|<1$ の近傍では次数の低い（マイナスに大きい）方が支配的になるため、定数を調整すれば $|u|^{-\max(A, B)}$ で抑えられる。積についてはそのまま $|f(x)g(x)| \le C_1 C_2 |u|^{-(A+B)}$ となる。
\end{proof}

\begin{proof}[命題\ref{prop:equivalence_regular_singular}の証明]
\begin{itemize}
    \item \textbf{【$\Rightarrow$ の証明（ODE確定 $\implies$ 解析的確定）】}\\
    ODEの確定特異点の定義より、命題\ref{prop:single_valued_z}で定まる1価関数 $z_{k}(x)$ のローラン展開は有限の負べきで打ち切られる（すなわち高々極である）。よって定理\ref{thm:growth_and_pole}より、$z_k(x)$ は $x=\xi$ の近傍で多項式増大度を持つ。\\
    また、基本解行列 $\Phi(x)$ の各成分 $\Phi_{kj}(x)$ は $u^{\alpha}$ や $u^{\alpha}\log u$ の線型結合であるため、例題\ref{ex:regular_singular_points}で見たようにこれらも多項式増大度を持つ。\\
    ここで、解の各成分 $\phi_i(x)$ は
    \[
    \phi_{i}(x) = z_{1}(x)\Phi_{1i}(x) + z_{2}(x)\Phi_{2i}(x)
    \]
    と表される。これは「多項式増大度を持つ関数の積と和」であるから、補題\ref{lem:polynomial_growth_closure}より $\phi_i(x)$ 自身も多項式増大度を持つ。すなわち、解析的な意味で確定特異点であることが示された。

    \item \textbf{【$\Leftarrow$ の証明（解析的確定 $\implies$ ODE確定）】}\\
    仮定より、解 $\phi_i(x)$ は $x=\xi$ の近傍で多項式増大度を持つ。\\
    これをODEの確定特異点の定義（$z_k(x)$ のローラン展開が有限の負べきで打ち切られること）に結びつけるには、定理\ref{thm:growth_and_pole}より、$z_k(x)$ が多項式増大度を持つことを示せば十分である。\\
    命題\ref{prop:single_valued_z}より $\vec{z}(x) = \vec{\phi}(x)\Phi(x)^{-1}$ であるから、各成分は
    \[
    z_{i}(x) = \phi_{1}(x)\Phi_{1i}^{-1}(x) + \phi_{2}(x)\Phi_{2i}^{-1}(x)
    \]
    と書ける。ここで、逆行列 $\Phi(x)^{-1}$ の各成分 $\Phi_{ji}^{-1}(x)$ は $u^{-\alpha}$ や $u^{-\alpha}\log u$ の線型結合で表される（$\log u$ の逆数などは出現しない）ため、これらもやはり多項式増大度を持つ。\\
    したがって、補題\ref{lem:polynomial_growth_closure}より、その積と和で表される $z_i(x)$ も多項式増大度を持つ。\\
    $z_i(x)$ は $x=\xi$ の周りで1価の解析関数であり、かつ多項式増大度を持つため、定理\ref{thm:growth_and_pole}によりそのローラン展開は有限の負べきで必ず打ち切られる。これはまさにODEの確定特異点の定義そのものである。
\end{itemize}
\end{proof}

\begin{Q}[【補足】$\Phi(x)^{-1}$ を計算するとき、$\log u$ の逆数（$1/\log u$）のような扱いにくい項は出てこないの？]
\textbf{A.} 全く出てこない。行列の逆行列計算の美しい性質により、$\log u$ は分子に留まる。
$\alpha - \beta \in \mathbb{Z}$ でジョルダン細胞になる場合（最も複雑なケース）、$\Phi(x)$ は
\[
\Phi(x) = 
\begin{pmatrix}
    u^\alpha & K u^\alpha \log u \\
    0 & u^\alpha
\end{pmatrix}
\]
という形であった（$K$ は定数）。この行列の行列式は $\det \Phi(x) = u^{2\alpha}$ である。逆行列 $\Phi(x)^{-1}$ を真面目に計算すると、
\[
\Phi(x)^{-1} = \frac{1}{u^{2\alpha}}
\begin{pmatrix}
    u^\alpha & -K u^\alpha \log u \\
    0 & u^\alpha
\end{pmatrix}
=
\begin{pmatrix}
    u^{-\alpha} & -K u^{-\alpha} \log u \\
    0 & u^{-\alpha}
\end{pmatrix}
\]
となる。見ての通り、$\Phi(x)^{-1}$ の成分に現れるのは $u^{-\alpha}$ と $u^{-\alpha} \log u$ のみであり、$1/\log u$ は出現しない。これらは例題\ref{ex:regular_singular_points}で確認した通り、すべて多項式オーダー $|u|^{-M}$ で安全に上から抑えることができる。
\end{Q}

\begin{Q}[【補足】「ローラン展開の負べきが有限個（極）」だと、なぜ $|u|^{-N}$ で抑えられると直感的に言えるの？]
\textbf{A.} 原点付近（$u \to 0$）では、最も次数の低い（マイナスに大きい）項が圧倒的に支配的になるからである。
$z_k(x)$ が $x=\xi$ ($u=0$) に $N$ 位の極を持つとすると、ローラン展開は
\[
z_k(x) = \frac{a_{-N}}{u^N} + \frac{a_{-N+1}}{u^{N-1}} + \cdots + a_0 + a_1 u + \cdots
\]
となる。両辺に $u^N$ をかけると、
\[
u^N z_k(x) = a_{-N} + a_{-N+1}u + \cdots 
\]
となり、これは $u \to 0$ で正則（有限の値 $a_{-N}$ に収束）である。したがって、$u$ が十分小さい領域では $|u^N z_k(x)| \le C$ (ある定数) と有界になり、両辺を割ることで $|z_k(x)| \le C|u|^{-N}$ という多項式オーダーの評価が直ちに得られる。無限に続く真性特異点（例：$e^{1/u}$）ではこの議論が破綻する。（※厳密な同値性の証明は定理\ref{thm:growth_and_pole}を参照）
\end{Q}
%#endregion
%#region
\begin{definition}[ガウスの微分方程式]
    複素パラメータ $a, b, c$ に対して、
    \begin{equation}\label{eq:gauss_hypergeometric}
    x(1-x)\frac{d^{2}y}{dx^{2}}+(c-(a+b+1)x)\frac{dy}{dx}-aby = 0
    \end{equation}
    と書かれる2階線形微分方程式を\textbf{ガウスの微分方程式（超幾何微分方程式）}という。
\end{definition}

\begin{lemma}[ガウスの超幾何関数]\label{lem:hypergeometric}
    \[
        F(a,b,c;x)=\sum_{j=0}^{\infty}\frac{(a)_{j}(b)_{j}}{(c)_{j}j!}x^{j} \qquad \left( \text{ただし } (a)_{j}=a(a+1)(a+2)\dots(a+j-1) \right)
    \]
    の収束半径は1である。これは\textbf{ガウスの超幾何関数}と呼ばれる。
\end{lemma}

\begin{proof}
    命題\ref{prop:dalembert_radius}より収束半径は
    \[
    \lim_{j\to\infty}\left| \frac{(a)_{j}(b)_{j}}{(c)_{j}j!} \frac{(c)_{j+1}(j+1)!}{(a)_{j+1}(b)_{j+1}} \right| = \lim_{j\to\infty}\left| \frac{(c+j)(j+1)}{(a+j)(b+j)} \right| = 1
    \]
    となる。
\end{proof}

\begin{Q}[【核心】ガウスの超幾何関数の収束半径が「1」になることは、ガウスの微分方程式の係数の形からどのように「予言」できるか？]
\textbf{A.} ガウスの微分方程式を標準形 $\frac{d^2y}{dx^2} + p_1(x)\frac{dy}{dx} + p_2(x)y = 0$ で表すと、分母に $x(1-x)$ が来るため、方程式の特異点は $x=0$ と $x=1$（および $\infty$）であることがわかる。超幾何関数 $F(a,b,c;x)$ は原点 $x=0$ まわりでの解であるため、そこから「最も近い次の特異点 $x=1$」までの距離がちょうど1であり、これが収束半径の限界と一致するからである。
\end{Q}
%#endregion
%#region
\begin{lemma}
    \eqref{eq:gauss_hypergeometric}は$x=0,1,\infty$で特異点をもつ。
\end{lemma}

\begin{proof}
    $x=0, 1$ が特異点であることは、方程式の両辺を $x(1-x)$ で割ることで係数の分母がゼロになることからすぐにわかる。

    無限遠点 $x=\infty$ については、$x=\frac{1}{u}$ を代入して $u=0$ での振る舞いを調べる。微分の連鎖律を用いて \eqref{eq:gauss_hypergeometric} を書き換えると、
    \begin{align*}
    \frac{1}{u}\left(1-\frac{1}{u}\right)\left(u^{4}\frac{d^{2}y}{du^{2}}+2u^{3}\frac{dy}{du}\right)+\left(c-(a+b+1)\frac{1}{u}\right)\left(-u^{2}\frac{dy}{du}\right)-aby &= 0 \\
    u^{2}(u-1)\frac{d^{2}y}{du^{2}}+2u(u-1)\frac{dy}{du}-u(cu-(a+b+1))\frac{dy}{du}-aby &= 0 \\
    u^{2}(u-1)\frac{d^{2}y}{du^{2}}+u((2-c)u+(a+b-1))\frac{dy}{du}-aby &= 0
    \end{align*}
    となる。両辺を $u^{2}(u-1)$ で割ると、各係数が $u=0$ を極（特異点）として持つことがわかる。つまり $x=\infty$ も特異点である。

    また、これら以外の点では係数は明らかに正則であるため、特異点の集合は $\{0,1,\infty\}$ である。
\end{proof}

\begin{Q}[【核心】なぜ複素平面だけでなく、無限遠点 $\infty$ を含めた「リーマン球面」全体で特異点の個数を数える必要があるのか？]
\textbf{A.} リーマン球面全体で見渡したとき、ガウスの微分方程式は「特異点をちょうど3つ $\{0, 1, \infty\}$ だけ持つ」という極めて対称的で純粋な構造を持っていることがわかり、実はこの「特異点が3つ」という幾何学的な条件だけで方程式の形が（変数変換を除いて）一意に決定されてしまうという、複素微分方程式の根源的な性質（リーマンのP方程式）に繋がるからである。
\end{Q}
%#endregion
%#region
\begin{lemma}\label{lem:hypergeometric_solution}
    $F(a,b,c;x)$ は \eqref{eq:gauss_hypergeometric} の解である。
\end{lemma}

\begin{proof}
    級数解 $y = \sum_{j=0}^{\infty}\frac{(a)_{j}(b)_{j}}{(c)_{j}j!}x^{j}$ の微分はそれぞれ、
    \[
    y' = \sum_{j=1}^{\infty}\frac{(a)_{j}(b)_{j}}{(c)_{j}(j-1)!}x^{j-1}, \quad
    y'' = \sum_{j=2}^{\infty}\frac{(a)_{j}(b)_{j}}{(c)_{j}(j-2)!}x^{j-2}
    \]
    となる。これらを \eqref{eq:gauss_hypergeometric} の左辺に代入して展開すると、
    \begin{align*}
    (\text{左辺}) &= x(1-x)\sum_{j=2}^{\infty}\frac{(a)_{j}(b)_{j}}{(c)_{j}(j-2)!}x^{j-2} + (c-(a+b+1)x)\sum_{j=1}^{\infty}\frac{(a)_{j}(b)_{j}}{(c)_{j}(j-1)!}x^{j-1} - ab\sum_{j=0}^{\infty}\frac{(a)_{j}(b)_{j}}{(c)_{j}j!}x^{j} \\
    &= \sum_{j=1}^{\infty}\frac{(a)_{j+1}(b)_{j+1}}{(c)_{j+1}(j-1)!}x^{j} - \sum_{j=2}^{\infty}\frac{(a)_{j}(b)_{j}}{(c)_{j}(j-2)!}x^{j} \\
    &\quad + c\sum_{j=0}^{\infty}\frac{(a)_{j+1}(b)_{j+1}}{(c)_{j+1}j!}x^{j} - (a+b+1)\sum_{j=1}^{\infty}\frac{(a)_{j}(b)_{j}}{(c)_{j}(j-1)!}x^{j} - ab\sum_{j=0}^{\infty}\frac{(a)_{j}(b)_{j}}{(c)_{j}j!}x^{j}
    \end{align*}
    となる（$x y''$ と $c y'$ の項ではインデックスを $j \to j+1$ とずらした）。
    定数項（$x^0$ の係数）は $c\frac{ab}{c} - ab = 0$ となる。
    $x^j \ (j \ge 1)$ の係数を集め、$\frac{(a)_{j}(b)_{j}}{(c)_{j}j!}$ をくくり出すと、
    \begin{align*}
    & \frac{(a)_{j+1}(b)_{j+1}}{(c)_{j+1}j!} (j + c) - \frac{(a)_{j}(b)_{j}}{(c)_{j}j!} \left( j(j-1) + (a+b+1)j + ab \right) \\
    &= \frac{(a)_{j}(b)_{j}}{(c)_{j}j!} \left[ \frac{(a+j)(b+j)}{c+j}(j+c) - (j^2 + (a+b)j + ab) \right] \\
    &= \frac{(a)_{j}(b)_{j}}{(c)_{j}j!} \left[ (a+j)(b+j) - (a+j)(b+j) \right] = 0
    \end{align*}
    となり、すべての次数の係数が $0$ となる。したがって $F(a,b,c;x)$ は方程式を満たす。
\end{proof}

\begin{Q}[【核心】ガウスの超幾何関数という複雑な級数は、どこから天下り的にやってきたのか？]
\textbf{A.} 天下りではなく、未定係数法による必然である。微分方程式の解を $y = \sum A_j x^j$ と仮定して代入し、各次数の係数がゼロになるように整理すると、$A_{j+1}$ と $A_j$ の間に $A_{j+1} = \frac{(a+j)(b+j)}{(c+j)(j+1)}A_j$ という2項間漸化式が必ず導かれる。これを $A_0 = 1$ から順次解き下した結果が、ポッホハマー記号を用いた $F(a,b,c;x)$ の係数に他ならない。
\end{Q}
%#endregion
%#region
\begin{definition}[複素領域における解の「分枝」の1次独立性とロンスキアン]\label{def:linear_independence_complex}
    複素領域における2階線形微分方程式について、特異点を含まないある単連結な領域（例えば小さな開円板やスリット領域）$D$ 上で固定された\textbf{解の分枝のペア $(\varphi_1(x), \varphi_2(x))$} が\textbf{1次独立（線形独立）}であるとは、任意の複素定数 $C_1, C_2$ に対して $D$ 上で
    \[ 
    C_1\varphi_1(x) + C_2\varphi_2(x) \equiv 0 \implies C_1 = 0 \text{ かつ } C_2 = 0 
    \]
    が成り立つことをいう。解析的には、領域 $D$ 上でロンスキアン $W(\varphi_1, \varphi_2)(x) \not\equiv 0$ であることと同値である。
    
    （※大域的な多価関数そのものが無条件に独立であるかを問うのではなく、あくまで「局所領域 $D$ 上で切り出された特定の分枝の組み合わせ」に対して独立性が定義・判定されることに注意する。）
\end{definition}

\begin{Q}[【核心】局所的な領域 $D$ での判定だけで、多価関数という複雑な対象の「全域での独立性」を保証できるのはなぜか？]
\textbf{A.} 複素解析の「一致の定理」と微分方程式の「アーベルの公式」が、局所と大域を鉄の論理で結んでいるからである。
ある1つの領域 $D$ でロンスキアンが非ゼロ（独立）ならば、特異点を回って別の分枝に移っても、それは正則なモノドロミー行列を掛ける「基底の変換」に過ぎず、独立性は全域で不変である。逆に、もし最初から従属なペアを選んでいれば、領域 $D$ での判定で確実に $0$ となり、偽陽性（誤判定）は絶対に起こらない。
\end{Q}

\begin{Q}[【核心】2つの多価関数を比較する際、なぜ「解析接続の経路を同期させる（シートを揃える）」ことが絶対ルールなのか？]
\textbf{A.} リーマン面という螺旋階段において、同じ階層（シート）で評価しなければ数学的に無意味だからである。
同じ複素座標 $x$ であっても、多価関数には「0周目の $x$」や「1周目の $x$」といった異なる実体が存在する。これらを混同して比較すると、本来独立なはずの解がたまたま同じ向きを向いて従属に見えてしまう「分枝の罠」に嵌まる。これを防ぐため、必ず共通の基準点から「全く同じ経路」で解析接続された値同士を戦わせる必要がある。
\end{Q}

\begin{Q}[【核心】特異点を回って別の分枝に移動したとき、解を結びつけていた定数 $C_1, C_2$ が変化してしまう心配はないのか？]
\textbf{A.} 全くない。$C_1, C_2$ は解析接続の影響を受けない「普遍的な複素定数」だからである。
2つの解を同じ経路で同期させて解析接続する限り、多価性によって関数自身の値が $\lambda$ 倍に変化しても、等式の両辺でその因子が約分されるため、線形関係を規定する定数は全分枝において完全に同一のまま保たれる。
\end{Q}
%#endregion
%#region
\begin{example}[$x=0$ まわりの局所解]\label{ex:gauss_solutions_zero}
    $c\notin\mathbb{Z}$ ならば、ガウスの微分方程式 \eqref{eq:gauss_hypergeometric} は $x=0$ の近傍において、
    \[
    \varphi_{1}(x)=F(a,b,c;x),\quad \varphi_{2}(x)=x^{1-c}F(a-c+1,b-c+1,2-c;x)
    \]
    という1次独立な2つの解をもつ。
\end{example}

\begin{proof}
    定理と命題を用いて、以下の3つのステップに分けて証明する。

    \begin{description}
        \item[【ステップ1】$\varphi_2(x)$ が解であること]\mbox{}\\
        $\varphi_{1}(x)$ が解であることは既に補題\ref{lem:hypergeometric_solution}で示した。$\varphi_{2}(x)$ について、$y=x^{1-c}z$ とおいて \eqref{eq:gauss_hypergeometric} に代入し、微分の連鎖律を適用すると、
        \begin{gather*}
        x(1-x)\left\{ c(c-1)x^{-1-c}z + 2(1-c)x^{-c}\frac{dz}{dx} + x^{1-c}\frac{d^{2}z}{dx^{2}} \right\} \\
        + (c-(a+b+1)x)\left\{ (1-c)x^{-c}z + x^{1-c}\frac{dz}{dx} \right\} - abx^{1-c}z = 0
        \end{gather*}
        となる。超幾何微分方程式の標準形に合わせるため、この両辺に $x^{c-1}$ を掛けて整理すると、
        \begin{gather*}
        x(1-x)\frac{d^{2}z}{dx^{2}} + \left\{ 2(1-c)(1-x) + (c-(a+b+1)x) \right\}\frac{dz}{dx} \\
        + \left\{ c(c-1)x^{-1}(1-x) + (c-(a+b+1)x)(1-c)x^{-1} - ab \right\}z = 0 \\
        \intertext{$z', z$ の係数をそれぞれ展開して因数分解すると、}
        x(1-x)\frac{d^{2}z}{dx^{2}} + \left\{ (2-c) - ((a-c+1)+(b-c+1)+1)x \right\}\frac{dz}{dx} - (a-c+1)(b-c+1)z = 0
        \end{gather*}
        が得られる。これはパラメータを $(a-c+1), (b-c+1), (2-c)$ としたガウスの超幾何微分方程式そのものである。
        したがって、補題\ref{lem:hypergeometric_solution}よりその解として $z=F(a-c+1,b-c+1,2-c;x)$ がとれる。よって、$y=x^{1-c}F(a-c+1,b-c+1,2-c;x)$ は \eqref{eq:gauss_hypergeometric} の解である。

        
        \item[【ステップ2】2つの解が1次独立であること]\mbox{}\\
        まず、特異点 $x=0$ を含まない単連結な領域 $D$ 上で $\log x$ の分枝を1つ固定し、$\varphi_2(x)$ を $D$ 上の一価正則な関数として確定させる。
        このとき、$\varphi_1, \varphi_2$ が $D$ 上で1次独立でないと仮定すると、ある定数 $(C_{1},C_{2})\in\mathbb{C}^{2}\setminus\{(0,0)\}$ に対して、
        \[
        C_{1}\varphi_{1}(x)+C_{2}\varphi_{2}(x)\equiv 0
        \]
        が成り立つ。命題\ref{prop:independence_branch_invariance}より、この等式全体を $x=0$ の周りに反時計回りに1周する共通の経路に沿って解析接続したとき、元の領域 $D$ に戻ってきて得られる新しい分枝同士の間でも、同じ定数 $(C_1, C_2)$ による線形関係が成立する。

        1周解析接続した結果、$\varphi_1(x)$ は一価関数であるため元の分枝と完全に一致するが、$\varphi_2(x)$ は $x^{1-c} = \exp((1-c)\log x)$ の因子を持つため、元の分枝から $e^{2\pi i(1-c)}$ 倍された値となる。したがって $D$ 上で、
        \[
        C_{1}\varphi_{1}(x)+C_{2}e^{2\pi i(1-c)}\varphi_{2}(x)\equiv 0
        \]
        となる。この新旧2つの等式の差をとると $C_2(e^{2\pi i(1-c)}-1)\varphi_2(x) \equiv 0$ となるが、$c \notin \mathbb{Z}$ より $1-c$ は整数ではないため $e^{2\pi i(1-c)} \neq 1$ である。よって $C_2=0$ となり、最初の式から $C_1=0$ も導かれるため仮定に矛盾する。ゆえに、選び出した分枝の組によらず $\varphi_1, \varphi_2$ は1次独立である。
        
        \item[【ステップ3】解がこれら2つで全てであること]\mbox{}\\
        もし、上に挙げた2つ以外にこれらと1次独立な解が存在すると仮定すると、解空間の次元が3次元以上となってしまう。しかしこれは、定理\ref{thm:solution_space_dimension}で保証された「解空間が2次元ベクトル空間をなす」という事実に矛盾する。よって、局所解はこの2つのみで完全に尽くされている。
    \end{description}
\end{proof}

\begin{Q}[【核心】切り込み（スリット）を越えて解析接続を行うと、分枝の定義域を超えてしまい論理が破綻しないか？]
\textbf{A.} 破綻しない。なぜなら、切り込みの上で方程式を評価するのではなく、始点の「安全な領域 $D$」上の等式を出発点とし、1周回って「再び領域 $D$ に戻ってきた新しい分枝」と元の分枝を比較しているからである。一致の定理により定数の関係は保たれるため、同じ $D$ 上での比較が正当化される。
\end{Q}

\begin{Q}[【核心】$\varphi_{2}(x)$ は $x=0$ で多価性や発散を伴うが、微分方程式の「解」として問題ないか？]
\textbf{A.} 全く問題ない。$x=0$ は方程式自身の「特異点」であるため、解が特異性を持つことは自然である。むしろ特異点の周りでは、このような $x^\rho \times (\text{正則関数})$ の形（フロベニウスの解）を想定して基本解を探すのが定石である。
\end{Q}

\begin{Q}[【核心】$x=0$ の真上ではコーシーの定理が使えないのに、なぜ解空間が「2次元」であり、この2つの解で「全て」だと断言できるのか？]
\textbf{A.} $x=0$ を除外した「穴あき近傍（$0 < |x| < r$）」で考えているからである。正則点で確立された「解空間は2次元」という代数的な事実は穴あき近傍でも変わらないため、そこで1次独立な2つの解を見つけさえすれば、いかなる解もその線形結合で完全に表現できる。
\end{Q}

\begin{Q}[【核心】特異点を回って解が「多価関数」になる現象は、「解空間は2次元」という枠組みと矛盾しないか？]
\textbf{A.} 矛盾しない。方程式の係数が一価であるため、特異点を1周して得た「新しい分枝」も再び解となる。解空間が局所的に2次元である以上、いかに多価に分岐しようとも、その新しい分枝は必ず元の2つの基底の線形結合で表され、多価性は「2次元空間内の線形変換（モノドロミー行列）」として吸収される。
\end{Q}

\begin{Q}[【実践】$c \notin \mathbb{Z}$ のとき、2つの解の1次独立性が計算せずとも直ちにわかるのはなぜか？]
\textbf{A.} 特異点 $x \to 0$ における漸近挙動が、一方は $\varphi_1(x) \to 1$、もう一方は $\varphi_2(x) \sim x^{1-c}$ と決定的に異なるからである。先頭のオーダーが異なる関数同士が定数倍で結びつくことは絶対に不可能であるため、これだけで1次独立性が確定する。
\end{Q}
%#endregion
%#region
これまでの議論から、$\{\varphi_{1}(x), \varphi_{2}(x)\}$ は方程式 \eqref{eq:gauss_hypergeometric} の基本系をなし、それぞれ $x=0$ に確定特異点をもつことがわかっている。

\begin{Q}[命題\ref{prop:equivalence_regular_singular}を適用するには「\textbf{任意の}解の基本系」が $x=0$ に確定特異点をもつことを示す必要がある。なぜ特定の基本系 $\{\varphi_{1}(x), \varphi_{2}(x)\}$ の性質から「任意」へと議論を拡張できるのか？]
\textbf{A.} 微分方程式の線型性と、定義\ref{def:linear_independence_complex}における確定特異点（増大度）の条件が、線型結合によって保たれる性質だからである。

具体的には以下の論理による：
\begin{enumerate}
    \item $\{\varphi_{1}(x), \varphi_{2}(x)\}$ が基本系であるため、\eqref{eq:gauss_hypergeometric} の任意の解は $c_{1}\varphi_{1}(x) + c_{2}\varphi_{2}(x)$ という定数係数の線型結合で一意に表される。
    \item 確定特異点をもつ（すなわち高々多項式オーダーの増大度をもつ）関数の線型結合もまた、同じ特異点において確定特異点をもつ。
    \item ゆえに、\eqref{eq:gauss_hypergeometric} の「任意の解の基本系」を構成する各成分も、すべて $x=0$ に確定特異点をもつことが保証される。
\end{enumerate}
\end{Q}

上記の理由により、\eqref{eq:gauss_hypergeometric} の任意の解の基本系は $x=0$ を確定特異点にもつ。よって、命題\ref{prop:equivalence_regular_singular}より、$x=0$ は微分方程式 \eqref{eq:gauss_hypergeometric} の確定特異点である。また、$x=1,x=\infty$も確定特異点であることが知られている。
%#endregion
%#region
\begin{definition}
    \eqref{eq:2-ode}の特異点$x=\xi$がすべて確定特異点である時、\eqref{eq:2-ode}を\textbf{フックス型微分方程式}という。
\end{definition}
%#endregion
%#region
\begin{example}
    ガウスの超幾何微分方程式 \eqref{eq:gauss_hypergeometric} はフックス型微分方程式である。
\end{example}

\begin{proof}
    今までの議論から、\eqref{eq:gauss_hypergeometric} の特異点は $x=0,1,\infty$ のみであり、それらはすべて確定特異点であることがわかっているので、定義からフックス型微分方程式である。
\end{proof}

\begin{Q}[ガウスの超幾何微分方程式のように「リーマン球面上に確定特異点をちょうど3つだけ持つ」フックス型方程式は、パンルヴェ方程式などの非線形理論へ向かう上でどのような特別な立ち位置にあるか？]
\textbf{A.} 「変形」の余地を持たない、最も基本的で**「剛（rigid）」**な方程式としての立ち位置です。

リーマン球面上の任意の3点は、一次分数変換（モビウス変換）によって常に $0, 1, \infty$ に移すことができます。そのため、確定特異点が3つのフックス型方程式は、各特異点における特性指数（局所的な振る舞い）を指定するだけで**方程式全体が一意に決定されてしまいます**（リーマンのP関数）。
ここから特異点を4つ（例えば $0, 1, \infty, t$）に増やし、その位置 $t$ を動かしても方程式の全体的な性質（モノドロミー）が変わらないように要請する（モノドロミー保存変形）ことで、初めてパンルヴェ方程式のような非線形微分方程式が舞台に上がってきます。
\end{Q}
%#endregion
%#region
\begin{proposition}
    \eqref{eq:2-ode}の特異点$x=\xi$が確定特異点になる必要十分条件は$(x-\xi)^{i}p_{i}(x)$が$x=\xi$で正則になることである。
\end{proposition}

\begin{proof}
    命題\ref{prop:fuchs_theorem_n}の$n=2$の場合である。
\end{proof}
%#endregion
%#region
上の命題で冪級数の構成部分を再考する。
特異点 $x=\xi$ の周りでの局所座標を $u = x - \xi$ とおく。
\eqref{eq:2-ode}の両辺に $u^{2}$ をかけると、$\frac{d}{dx} = \frac{d}{du}$ であるから、
\[
u^{2}\frac{d^{2}y}{du^{2}} + (u p_{1}(x)) u\frac{dy}{du} + (u^{2} p_{2}(x)) y = 0
\]
となる。

\begin{Q}[【核心】なぜわざわざ $u^2$ を掛け、$u p_1(x)$ と $u^2 p_2(x)$ が正則であるという状況を考えるのか？]
\textbf{A.} この条件こそが、$x=\xi$ が方程式の「確定特異点（regular singular point）」であることの定義だからである。フックス（Fuchs）の定理により、特異点であってもこの条件を満たせば、解の少なくとも1つは $u^s \sum c_n u^n$ という冪級数（フロベニウス級数）の形で正則な解を持つことが保証されるため、この前提が不可欠となる。
\end{Q}

上の仮定より、$u p_{1}(x), u^{2} p_{2}(x)$ は $u=0$ において正則なので、それぞれ $u$ の関数として $q(u), r(u)$ と置き直すと、方程式は次のように書ける。
\begin{equation}\label{eq:ode_u}
u^{2}\frac{d^{2}y}{du^{2}} + q(u) u\frac{dy}{du} + r(u) y = 0
\end{equation}
ここで、解 $\varphi(u)$ が
\[
\varphi(u) = u^{s}\sum_{n=0}^{\infty}c_{n}u^{n} \quad (c_{0}\neq 0)
\]
と書けるとする。$q(u), r(u)$ もテイラー展開して
\[
q(u) = \sum_{n=0}^{\infty}q_{n}u^{n}, \quad r(u) = \sum_{n=0}^{\infty}r_{n}u^{n}
\]
とし、これらを\eqref{eq:ode_u}に代入して $u^{s+n}$ の係数を整理すると、
\[
\sum_{n=0}^{\infty} u^{s+n} \left[ (s+n)(s+n-1)c_{n} + \sum_{k=0}^{n} (s+k)c_{k}q_{n-k} + \sum_{k=0}^{n} c_{k}r_{n-k} \right] = 0
\]
を得る。$n=0$（すなわち最低次 $u^s$ の係数）を見ると、
\[
(s(s-1) + sq_{0} + r_{0})c_{0} = 0
\]
となる。仮定より $c_{0} \neq 0$ であるため、次の方程式が得られる。
\begin{equation}\label{eq:indicial}
f(s) := s(s-1) + sq_{0} + r_{0} = 0
\end{equation}

\begin{definition}
\eqref{eq:indicial}を、$x=\xi$ における元の微分方程式\eqref{eq:2-ode}の\textbf{決定方程式}（または特性方程式）という。また、この2次方程式の2つの根を\textbf{決定指数}（特性指数）という。
\end{definition}

決定指数のうち、実部が大きい方（あるいは任意の1つ）を $s_{1}$ とし、もう1つを $s_{2}$ とする。
一般の $u^{s+n}$ の係数が $0$ になる条件から、以下の漸化式が得られる。
\begin{equation}\label{eq:recurrence}
f(s+n)c_{n} + \sum_{k=0}^{n-1} \left( (s+k)q_{n-k} + r_{n-k} \right) c_{k} = 0
\end{equation}
右側の和を既知の項 $R_n(c_{0}, c_{1}, \cdots, c_{n-1})$ とおけば、$f(s+n)c_{n} = -R_n$ となる。
したがって、$f(s+n) \neq 0$ であれば、$c_{n}$ は $c_{0}$ から順に一意に定まる。

決定指数の差 $s_{1} - s_{2}$ が整数でない場合、任意の自然数 $n$ に対して $f(s_{1}+n) \neq 0$ かつ $f(s_{2}+n) \neq 0$ となる。よって $s=s_{1}$ と $s=s_{2}$ のそれぞれから解 $\varphi_{1}(x), \varphi_{2}(x)$ が構成できる。これらは最低次の項（$u^{s_1}$ と $u^{s_2}$）の冪が異なるため、互いに1次独立な解となる。

\begin{Q}[【核心】決定指数の差が整数のとき、なぜ解に対数項 $\log u$ が加わることがあるのか？]
\textbf{A.} $s_{1} - s_{2} = J \in \mathbb{Z}_{\ge 0}$ と書ける場合、小さい方の根 $s=s_{2}$ を用いて漸化式\eqref{eq:recurrence}を解こうとすると、$n=J$ のステップで $f(s_2+J) = f(s_1) = 0$ となる。
このとき、漸化式は $0 \cdot c_{J} = -R_J$ となり、もし $R_J \neq 0$ であれば、これを満たす係数 $c_{J}$ は存在せず、解の構成が破綻してしまうからである。この破綻を吸収し独立な解を2つ見つけるためには、フロベニウス級数の形を拡張し、必然的に $\log u$ を含む項を導入しなければならない。
\end{Q}
%#endregion
%#region
\eqref{eq:2-ode}において、無限遠点 $x=\infty$ での振る舞いを調べるため、局所座標 $x = \frac{1}{u}$ を導入する。
連鎖律より $\frac{dy}{dx} = -u^{2}\frac{dy}{du}$、$\frac{d^{2}y}{dx^{2}} = 2u^{3}\frac{dy}{du} + u^{4}\frac{d^{2}y}{du^{2}}$ となるため、これを\eqref{eq:2-ode}に代入して $u^{4}$ で割ると、次の方程式を得る。
\begin{equation}\label{eq:ode_infinity}
\frac{d^{2}y}{du^{2}} + \left(\frac{2}{u} - \frac{1}{u^{2}}p_{1}\left(\frac{1}{u}\right)\right)\frac{dy}{du} + \frac{1}{u^{4}}p_{2}\left(\frac{1}{u}\right)y = 0
\end{equation}

\begin{Q}[【核心】なぜ $x=\infty$ の特異性を調べるのに、わざわざ $x=1/u$ という変換を行って $u=0$ での振る舞いを見るのか？]
\textbf{A.} 複素解析において、無限遠点 $\infty$ は特別な場所ではなく、リーマン球面 $\hat{\mathbb{C}} = \mathbb{C} \cup \{\infty\}$ 上の1点に過ぎないからである。しかし、我々が持っているフロベニウスの方法などの「局所的な解析手法」は原点（有限の点）の周りで構築されているため、チャート変換 $x=1/u$ を用いて $\infty$ を原点 $u=0$ に引き戻すことで、有限の特異点と全く同じ土俵で解析できるようになる。
\end{Q}

\eqref{eq:ode_infinity}に対して、有限の確定特異点の条件（フックスの定理の前提）を適用することで、以下の命題が成り立つ。

\begin{proposition}\label{prop:regular_infinity}
$x$ の有理関数を係数とする\eqref{eq:2-ode}に対して、無限遠点 $x=\infty$ が確定特異点になるための必要十分条件は、$\frac{1}{u}p_{1}\left(\frac{1}{u}\right)$ および $\frac{1}{u^{2}}p_{2}\left(\frac{1}{u}\right)$ が $u=0$ で正則になることである。
（※これは言い換えれば、$x \to \infty$ において $p_{1}(x) = O(x^{-1})$ かつ $p_{2}(x) = O(x^{-2})$ となることと同値である。）
\end{proposition}
%#endregion
%#region
次に、特異点において決定指数（特性指数）の差が整数となる例外的なケースを考察する。
通常、決定指数の差が整数の場合は対数項 $\log u$ が解に出現するが、漸化式における障害となる項 $R_J$ が偶然 $0$ になる場合、対数項を持たない正則な解が存在する。

\begin{definition}[非対数的特異点 / 見かけの特異点]\label{def:non_logarithmic}
確定特異点 $x=\xi$ における2つの決定指数 $s_{1}, s_{2}$ （$s_{1} \ge s_{2}$）の差が整数 $s_{1} - s_{2} = J \in \mathbb{Z}_{\ge 0}$ であり、かつ小さい方の根 $s=s_{2}$ に対する漸化式の第 $J$ ステップにおいて障害項が消滅する（すなわち $R_{J}(c_{0}, \cdots, c_{J-1}) = 0$ が成り立つ）とき、この特異点を\textbf{非対数的特異点}（または\textbf{見かけの特異点}）と呼ぶ。
\end{definition}

非対数的特異点においては、$n=J$ での漸化式が $0 \cdot c_{J} = 0$ となるため、$c_{J}$ は矛盾を引き起こすことなく、任意の複素定数 $h \in \mathbb{C}$ として自由に選んで計算を進めることができる。

\begin{Q}[【核心】小さい方の根 $s_1$ から出発して $c_J = h$ と置いた解は、なぜ大きい方の根 $s_2$ の解 $\varphi_{2}(x)$ を $h$ 倍して足し合わせた形になるのか？]
\textbf{A.} $c_J = h$ としたとき、$n > J$ における漸化式は $c_0, \cdots, c_{J-1}, h$ に依存する項 $R_n$ を持ち、次のように書ける。
\[
f(s_1+n)c_n + R_n(c_0, \cdots, c_{J-1}, h) = 0
\]
ここで、$R_n$ は各係数に対して線形であるため、次のように分離できる。
\[
R_n(c_0, \cdots, c_{J-1}, h) = R_n(c_0, \cdots, c_{J-1}, 0) + h R_n(0, \cdots, 0, 1)
\]
この分離により、以降の係数 $c_n$ も「$h=0$ としたときの係数」と「$h$ に比例する係数」の和に完全に分かれる。
後半の $h$ に比例する部分は、初項を $c_J=1$ として $u^{s_2+J} = u^{s_1}$ から始まる級数であり、これはまさに大きい根 $s_2$ に対応する解 $\varphi_{2}(x)$ そのものである。
したがって、任意の定数 $h$ に対する全体の解は、
\[
\varphi(x) = \varphi_{1}(x) + h\varphi_{2}(x)
\]
（ただし $\varphi_{1}(x)$ は $h=0$ とした場合の解）という形に自然に分離されるのである。
\end{Q}
決定指数の差が整数となる場合について、方程式を正規化してさらに深く考察する。
特異点 $x=\xi$ における小さい方の決定指数を $s_{1}$、大きい方の決定指数を $s_{2}$ とし、その差を $J = s_{2} - s_{1} \in \mathbb{Z}_{\ge 0}$ とする。
未知関数 $y$ を
\[
y = (x-\xi)^{s_{1}} w
\]
と変換して\eqref{eq:2-ode}を $w$ についての微分方程式に書き換えると、変換後の微分方程式の $x=\xi$ における決定指数は $s_{1}-s_{1}=0$ および $s_{2}-s_{1}=J$ となる。
この整数 $J \ge 0$ の値によって、変換後の特異点の性質は以下のように分類される。

\begin{itemize}
    \item \textbf{$J=0$ のとき:} 決定指数は $0, 0$ となり重根を持つため、必ず対数項が発生し、非対数的特異点にはなり得ない。
    \item \textbf{$J=1$ のとき:} 決定指数は $0, 1$ となる。これが非対数的特異点である場合、2つの独立な解はともに $x=\xi$ で正則となる。フックスの定理の逆から、このとき $x=\xi$ は変換後の微分方程式にとって特異点ではなく、単なる正則点（ordinary point）となる。
    \item \textbf{$J \ge 2$ のとき:} 変換後の微分方程式は $x=\xi$ で依然として確定特異点を持つ。これが非対数的特異点である場合、決定指数が $0, J$ であるため、変換後の任意の解は $x=\xi$ で正則なテイラー展開を持つ。
\end{itemize}

\begin{Q}[【核心】$J=1$ と $J \ge 2$ では、どちらも「2つの解が正則」になるのに、なぜ $J=1$ は「正則点」で、$J \ge 2$ は「見かけの特異点」になるのか？]
\textbf{A.} その違いは、解から微分方程式の係数を逆算したときの「ロンスキアン（Wronskian）」の値に現れる。
微分方程式 $w'' + P(x)w' + Q(x)w = 0$ の係数 $P, Q$ は、2つの解 $w_{1}, w_{2}$ のロンスキアン $W = w_{1}w_{2}' - w_{1}'w_{2}$ を分母に持つ。
決定指数が $0$ と $J$ のとき、$w_{1}$ は定数項から始まり（$w_{1}(\xi) \neq 0$）、$w_{2}$ は $(x-\xi)^{J}$ から始まる。
\begin{itemize}
    \item \textbf{$J=1$ の場合：} $w_{2}'(\xi) \neq 0$ となるため、$x=\xi$ において $W(\xi) \neq 0$ となる。分母が $0$ にならないため、係数 $P(x), Q(x)$ も正則となり、そもそも微分方程式にとって特異点ではない（正則点である）。
    \item \textbf{$J \ge 2$ の場合：} $w_{2}(\xi) = 0$ だけでなく $w_{2}'(\xi) = 0$ も成り立つため、$x=\xi$ において $W(\xi) = 0$ となってしまう。分母が $0$ になるため係数 $P(x), Q(x)$ は極を持ち、方程式の特異性が残る（見かけの特異点となる）。
\end{itemize}
\end{Q}

\begin{Q}[【核心】$W(\xi)=0$ になるとしても、$P(x) = -\frac{w_{1}w_{2}'' - w_{1}''w_{2}}{W}$ の分子も同時に $0$ になって、極がキャンセルされる（除去可能な特異点になる）可能性はないのか？]
\textbf{A.} 結論から言えば、分子も $0$ になるが、零点の位数が異なるためキャンセルしきれず、必ず極が残る。
局所座標を $u = x-\xi$ とする。2つの解のテイラー展開の最低次を見ると、
\begin{align*}
w_{1}(u) &= c_{0} + c_{1}u + c_{2}u^{2} + \cdots \quad (c_{0} \neq 0) \\
w_{2}(u) &= u^{J}(d_{0} + d_{1}u + \cdots) \quad (d_{0} \neq 0)
\end{align*}
と書ける。分母のロンスキアン $W = w_{1}w_{2}' - w_{1}'w_{2}$ の最低次の項を計算すると、
\[
W(u) \approx c_{0} \cdot (J d_{0} u^{J-1}) - c_{1} \cdot (d_{0} u^{J}) \approx J c_{0} d_{0} u^{J-1}
\]
となり、$W$ は位数 $J-1$ の零点を持つ。
一方、$P(x)$ の分子 $N_{P} = w_{1}w_{2}'' - w_{1}''w_{2}$ の最低次の項を計算すると、
\[
N_{P}(u) \approx c_{0} \cdot J(J-1)d_{0}u^{J-2} - 2c_{2} \cdot d_{0}u^{J} \approx J(J-1)c_{0}d_{0}u^{J-2}
\]
となり、分子は位数 $J-2$ の零点を持つ。
したがって、極の位数は引き算されて、
\[
P(u) = -\frac{N_{P}(u)}{W(u)} \approx -\frac{J(J-1)c_{0}d_{0}u^{J-2}}{J c_{0} d_{0} u^{J-1}} = -\frac{J-1}{u}
\]
となる。$J \ge 2$ のとき $J-1 \ge 1$ であるため、キャンセルしきれずに必ず $u=0$ に $1$ 位の極（留数 $1-J$）が残るのである。

（※ちなみに、このとき $u P(u) \to 1-J$ となることから、決定方程式 $t(t-1) + p_{0}t + q_{0} = 0$ に $p_{0} = 1-J, q_{0} = 0$ を代入すると $t(t-J) = 0$ となり、決定指数が $0$ と $J$ になる事実とも見事に辻褄が合う。）
\end{Q}

\begin{example}
微分方程式
\begin{equation}\label{eq:euler_example}
x^{2}\frac{d^{2}y}{dx^{2}} + (1-2s)x\frac{dy}{dx} + (s^{2}-1)y = 0
\end{equation}
において、$x=0$ は非対数的特異点である。
\end{example}

\begin{proof}
方程式\eqref{eq:euler_example}の両辺を $x^2$ で割ると、$p_{1}(x) = \frac{1-2s}{x}, p_{2}(x) = \frac{s^{2}-1}{x^{2}}$ となる。
$x p_{1}(x) = 1-2s$ および $x^{2} p_{2}(x) = s^{2}-1$ はともに定数（正則）であるため、$x=0$ は確定特異点である。
また、これらは定数関数であるため、展開係数は $q_{0} = 1-2s, r_{0} = s^{2}-1$ であり、$n \ge 1$ に対しては $q_{n} = 0, r_{n} = 0$ となる。

決定方程式は
\[
f(t) = t(t-1) + (1-2s)t + (s^{2}-1) = t^{2} - 2st + s^{2} - 1 = (t-s-1)(t-s+1) = 0
\]
より、決定指数は $t = s+1, s-1$ である。
小さい方を $s_{1} = s-1$、大きい方を $s_{2} = s+1$ とすると、その差は $J = s_{2} - s_{1} = 2$ である。

ここで、小さい方の根 $t=s_{1}$ から出発する漸化式を考える。一般の漸化式において、障害となる項は
\[
R_{n} = \sum_{k=0}^{n-1} \left( (t+k)q_{n-k} + r_{n-k} \right) c_{k}
\]
と書けるが、今 $n \ge 1$ に対して $q_{n} = r_{n} = 0$ であるため、和の中身は常に $0$ となる。
したがって、任意の $n \ge 1$ に対して $R_{n} \equiv 0$ である。
特に $n = J = 2$ においても $R_{2}(c_{0}, c_{1}) = 0$ が自動的に成り立つため、$x=0$ は非対数的特異点である。
\end{proof}

\begin{Q}[【核心】この例題において、$R_n \equiv 0$ となるのは偶然か？]
\textbf{A.} 偶然ではない。この方程式は、各項の次数が揃っている「オイラーの微分方程式（Euler's equidimensional equation）」である。オイラー方程式は $y = x^{\lambda}$ という単項式の代入だけで厳密解が求まるという性質を持つため、そもそも無限級数（フロベニウス級数）で解を構成する必要がない。級数の高次の項が不要であることの代数的な現れが、「$n \ge 1$ で $q_n=r_n=0$ となり、$R_n \equiv 0$ となる」という構造そのものなのである。
\end{Q}
%#endregion

\subsection{不確定特異点}

\subsection{高階フックス型線型常微分方程式}
%#region
$n$階のフックス型微分方程式
\begin{equation}\label{eq:linear-ode-n}
    \frac{d^{n}y}{dx^{n}}+\sum_{j=1}^{n} p_{j}(x)\frac{d^{n-j}y}{dx^{n-j}} = 0   
\end{equation}
について考える。

\begin{definition}
    \eqref{eq:linear-ode-n}の特異点$x=\xi$が\textbf{確定特異点}であるとは、全ての解$\varphi(x)$が$x=\xi$で正則であるか、確定特異点を持つことである
\end{definition}

\begin{proposition}[フックスの定理]\label{prop:fuchs_theorem_n}
    $n$階線型常微分方程式\eqref{eq:linear-ode-n}
    の特異点 $x=\xi$ が確定特異点になる（すなわち任意の解が $x=\xi$ において確定増大度をもつ）ための必要十分条件は、各 $i=1,\dots,n$ について $(x-\xi)^{i}p_{i}(x)$ が $x=\xi$ で正則になることである。
    これは$n$階でも成り立つ。
    （\cite[定理13.4 p.128]{book:takano}）
\end{proposition}

\begin{Q}[$n$階のフックスの定理（命題\ref{prop:fuchs_theorem_n}）の証明は、どのような数学的アイデアに基づいているか？]
\textbf{A.} 十分性（係数の正則性 $\implies$ 確定特異点）は「フロベニウスの方法」による解の直接構成で示し、必要性（確定特異点 $\implies$ 係数の正則性）は「ロンスキアンの増大度評価」を用いて示します。

証明の骨格は以下の通り：
\begin{itemize}
    \item \textbf{十分性の証明（解の構成）：}\\
    $(x-\xi)^{i}p_{i}(x)$ が正則であると仮定する。方程式 \eqref{eq:linear-ode-n} の両辺に $(x-\xi)^n$ をかけると、一般化されたオイラーの微分方程式（$x^n y^{(n)} + \cdots = 0$ の形）の摂動系とみなせる。このとき、解をフロベニウス級数 $y = (x-\xi)^{\rho} \sum_{k=0}^{\infty} c_k (x-\xi)^k$ の形で仮定し、決定方程式（indicial equation）を解くことで、$n$個の線型独立な解（対数項を含む場合もある）を具体的に構成できる。これらはすべて高々多項式オーダーの増大度を持つため、確定特異点となる。

    \item \textbf{必要性の証明（係数の逆算）：}\\
    任意の解の基本系 $y_1, \dots, y_n$ がすべて確定増大度をもつと仮定する。方程式の係数 $p_i(x)$ は、基本系から作られるロンスキアン $W(x) = \det(y_j^{(k-1)})$ を用いてクラメルの公式から逆算できる（例えば $p_1(x) = -W'(x)/W(x)$）。解が確定増大度を持つことから各導関数の増大度も評価でき、結果として得られる $p_i(x)$ の極の位数が「高々 $i$ であること（すなわち $(x-\xi)^i p_i(x)$ が正則であること）」が導かれる。
\end{itemize}
\end{Q}

\begin{Q}[フックスの定理の証明において、十分性（解の構成）に比べて、必要性（ロンスキアンを用いた逆算）の証明が圧倒的に「面倒」なのはなぜか？]
\textbf{A.} 最大の理由は、複素領域特有の「モノドロミー（解の多価性）」の処理と、分母に現れる行列式の「零点の評価」が極めて繊細な作業を要求するから。

それぞれの証明のアプローチの違いから、難易度に以下のような差が生じる：
\begin{itemize}
    \item \textbf{十分性（構成的で楽）：}\\
    係数の形（極の位数）が既に与えられているため、フロベニウス級数（べき級数＋対数項）を代入し、決定方程式と漸化式を機械的に解くだけで済む。計算は長いものの、論理はアルゴリズム的で平坦である。
    
    \item \textbf{必要性（技巧的で面倒）：}\\
    「すべての解が確定増大度を持つ」という解析的な仮定だけから、係数 $p_i(x)$ の極の位数を評価しなければならない。係数はロンスキアンを用いた分数の形で表されますが、特異点近傍で分母が $0$ に近づく際の挙動を正確に抑え込む必要がある。さらに、解が特異点の周りを回ったときの多価性を処理するため、「モノドロミー行列のジョルダン標準形」を用いて、評価しやすい特別な基本系（標準基本系）をわざわざ選び直すという、高度に技巧的なステップを踏むことになる。
\end{itemize}
\end{Q}

一般の $n$ 階線形微分方程式
\begin{equation}\label{eq:n-ode}
\frac{d^{n}y}{dx^{n}} + \sum_{j=1}^{n} q_{j}(x)\frac{d^{n-j}y}{dx^{n-j}} = 0
\end{equation}
について、$x=\infty$ での振る舞いを調べるため、局所座標 $x = \frac{1}{u}$ を導入する。

\begin{proposition}\label{prop:regular_infinity_n}
方程式\eqref{eq:n-ode}に対して、特異点 $x=\infty$ が確定特異点になるための必要十分条件は、各 $j = 1, 2, \dots, n$ について、$x^{j}q_{j}(x)$ が $x=\infty$ において正則になることである。
（すなわち、$x \to \infty$ において $q_{j}(x) = O(x^{-j})$ となることと同値である。）
\end{proposition}

\begin{Q}[【核心】$x=\frac{1}{u}$ と変換したとき、連鎖律によって微分方程式は低階の項が複雑に混ざり合った形になるはずである。それにもかかわらず、なぜ $x^{j}q_{j}(x)$ が正則という、各項が独立したシンプルな条件に帰着するのか？]
\textbf{A.} 変換によって生じる $u$ の冪乗の「最高階部分の寄与」に注目すれば明らかになる。
連鎖律 $\frac{d}{dx} = -u^{2}\frac{d}{du}$ を繰り返し適用すると、$k$ 階微分の最高階部分は
\[
\frac{d^{k}y}{dx^{k}} = (-1)^{k} u^{2k} \frac{d^{k}y}{du^{k}} + (\text{低階の微分項})
\]
となる。これを\eqref{eq:n-ode}に代入すると、全体の最高階項は $(-1)^{n} u^{2n} \frac{d^{n}y}{du^{n}}$ となる。
新しい微分方程式を標準形（最高階の係数を $1$）にするため、全体を $(-1)^{n} u^{2n}$ で割る。このとき、第 $j$ 項（元の $\frac{d^{n-j}y}{dx^{n-j}}$ に対応する部分）の最高階の係数には、
\[
\frac{(-1)^{n-j} u^{2(n-j)}}{(-1)^{n} u^{2n}} \cdot q_{j}\left(\frac{1}{u}\right) = (-1)^{-j} u^{-2j} q_{j}\left(\frac{1}{u}\right)
\]
がかかることになる。
$u=0$ が確定特異点になるためには（フックスの定理より）、この $\frac{d^{n-j}y}{du^{n-j}}$ の係数が $u=0$ で高々 $j$ 位の極（つまり $u^{-j}$ の特異性）に収まらなければならない。
すなわち、$u^{-2j} q_{j}(\frac{1}{u}) = O(u^{-j})$ が要求される。両辺に $u^{2j}$ を掛ければ $q_{j}(\frac{1}{u}) = O(u^{j})$ となり、これを元の $x$ に戻せば $q_{j}(x) = O(x^{-j})$ 、つまり $x^{j}q_{j}(x)$ が $x=\infty$ で正則であることが導かれる。
連鎖律による低階の複雑な混ざり合いは、常にこの最も厳しい条件（最高階からの寄与）に吸収されるため、結果として各係数ごとの美しい条件式だけが残るのである。
\end{Q}
%#endregion
%#region
$n$ 階線形微分方程式\eqref{eq:linear-ode-n}において、特異点 $x=\xi$ の周りの局所座標を $u = x-\xi$ とする。
方程式の両辺に $u^{n}$ をかけると、
\[
u^{n}\frac{d^{n}y}{du^{n}} + \sum_{j=1}^{n} u^{n-j} \left( u^{j}p_{j}(u+\xi) \right) \frac{d^{n-j}y}{du^{n-j}} = 0
\]
となる。ここで、$x=\xi$ が確定特異点であるという仮定から、$u^{j}p_{j}(u+\xi)$ は $u=0$ で正則となるため、これを $q_{j}(u)$ と置き直す。

解をフロベニウス級数 $\varphi(u) = u^{s} \sum_{k=0}^{\infty} c_{k} u^{k} \quad (c_{0} \neq 0)$ と仮定し、上式に代入して最低次 $u^{s}$ の係数を比較する。
微分演算子について $u^{k} \frac{d^{k}}{du^{k}} (u^{s}) = s(s-1)\cdots(s-k+1)u^{s}$ となる性質を用いると、$u^{s}$ の係数が消えるための条件として次の方程式が得られる。

\begin{equation}\label{eq:indicial-n}
f(s) = s(s-1)\cdots(s-n+1) + \sum_{j=1}^{n-1} q_{j}(0) s(s-1)\cdots(s-n+j+1) + q_{n}(0) = 0
\end{equation}
（ただし、$q_{j}(0)$ は $(x-\xi)^{j}p_{j}(x)$ の $x=\xi$ における極限値である。）

\begin{definition}
\eqref{eq:linear-ode-n}に対し、多項式 $f(s)$ を特異点 $x=\xi$ における\textbf{決定多項式}（または特性多項式）という。さらに、$f(s)=0$ を $x=\xi$ における\textbf{決定方程式}（または特性方程式）と呼び、その根 $s$ を\textbf{決定指数}（または特性指数）という。
\end{definition}

\begin{Q}[【核心】決定方程式 $f(s)=0$ の各項が、$s^{n}$ のような単純な冪乗ではなく、$s(s-1)\cdots(s-n+1)$ のような「下降階乗冪（Falling factorial）」の形をとるのはなぜか？]
\textbf{A.} これは、特異点近傍での主要な振る舞いが「指数関数 $e^{sx}$」ではなく「冪関数 $u^{s}$」によって支配されているからである。
定数係数の線形微分方程式では解を $e^{sx}$ と仮定するため、微分するたびに $s$ が前に出て $s^{n}$ が現れる。しかし確定特異点においては、方程式が本質的にオイラーの微分方程式（各項が $x^{k} y^{(k)}$ の形を持つ）に帰着される。冪関数 $u^{s}$ を $k$ 回微分して $u^{k}$ を掛けると、$s^{k}$ ではなく $s(s-1)\cdots(s-k+1)$ という係数が飛び出してくる。この微分の代数的な構造が、決定方程式に下降階乗冪を強制しているのである。
\end{Q}

決定方程式の根（決定指数）の一つを $s$ とする。
解が $u=0$ の周りで
\[
\varphi(u) = u^{s}\sum_{k=0}^{\infty} c_{k}u^{k} \quad (c_{0}\neq 0)
\]
のように展開されると仮定し、方程式に代入して $u^{s+k}$ の係数を整理すると、一般に次のような漸化式が得られる。
\begin{equation}\label{eq:recurrence-n}
f(s+k)c_{k} + R_{k}(c_{0}, \cdots, c_{k-1}) = 0 \quad (k \ge 1)
\end{equation}
したがって、$f(s+k) = 0$ となる正の整数 $k$ が存在しない限り、$c_{k}$ は $c_{0}$ から順番に矛盾なく一意に決定できる。

次に、この特異点 $x=\xi$ における $n$ 個の決定指数 $s_{1}, s_{2}, \dots, s_{n}$ の和について考察する。
方程式\eqref{eq:indicial-n}の決定多項式 $f(s)$ に対して「解と係数の関係」を適用することで、以下の美しい関係が導かれる。
\begin{equation}\label{eq:sum-of-exponents}
\sum_{j=1}^{n} s_{j} = \frac{1}{2}n(n-1) - q_{1}(0)
\end{equation}

\begin{Q}[【核心】決定多項式 $f(s)$ から、なぜこの特性指数の和の式が導出されるのか？]
\textbf{A.} $n$ 次方程式 $f(s)=0$ において、根の和 $\sum s_{j}$ は「$s^{n-1}$ の係数にマイナスをつけたもの」となる（解と係数の関係）。
決定多項式\eqref{eq:indicial-n}の各項の展開を考えると、
\begin{align*}
f(s) &= s(s-1)\cdots(s-n+1) + q_{1}(0)s(s-1)\cdots(s-n+2) + \cdots \\
&= \left( s^{n} - (0+1+2+\cdots+n-1)s^{n-1} + \cdots \right) + q_{1}(0)\left( s^{n-1} + \cdots \right) + \cdots \\
&= s^{n} + \left( -\frac{1}{2}n(n-1) + q_{1}(0) \right)s^{n-1} + (\text{低次の項})
\end{align*}
となる。したがって、$s^{n-1}$ の係数の符号を反転させることで、ただちに $\sum s_{j} = \frac{1}{2}n(n-1) - q_{1}(0)$ が得られる。
\end{Q}

\begin{Q}[【核心】なぜわざわざ、特異点における特性指数の「和」を計算したのか？これにはどのような大域的な意味があるのか？]
\textbf{A.} この局所的な和こそが、リーマン球面 $\hat{\mathbb{C}}$ 全体における極めて強力な拘束条件「フックスの関係式（Fuchs' relation）」を導くためのパーツだからである。
各特異点での $q_{1}(0)$ は、微分方程式の第2項の係数 $p_{1}(x)$ のその特異点における留数（Residue）に相当する。複素解析の留数定理により、リーマン球面上のすべての特異点（$\infty$ を含む）の留数を足し合わせると必ず $0$ になる。
つまり、すべての特異点について\eqref{eq:sum-of-exponents}を足し合わせると $q_{1}(0)$ の寄与が完全にキャンセルされ、「全特異点の特性指数の総和は、空間のトポロジー（特異点の数と階数 $n$）だけで決まる一定の整数になる」という大域的な定理が成立する。パンルヴェ方程式などのモノドロミー保存変形を考える際、この関係式は系の自由度を落とす絶対的な土台として機能する。
\end{Q}
%#endregion
%#region
ここで
\[
f(u;s)=s(s-1)\cdots(s-n+1) + \sum_{j=1}^{n-1} q_{j}(u) s(s-1)\cdots(s-n+j+1) + q_{n}(u)
\]
とおく。
\begin{proposition}
    \eqref{eq:linear-ode-n}は
    \[
    f(u;D)y=0\quad D=u\frac{d}{du}
    \]
    と表せる。
\end{proposition}

\begin{proof}
補題\ref{lem:falling_factorial}から従う。
\end{proof}

\begin{lemma}\label{lem:falling_factorial}
    \[
    u^{k}\frac{d^{k}}{du^{k}}=D(D-1)(D-2)\cdots(D-k+1)
    \]
\end{lemma}

\begin{proof}
$k=1$のとき、$u\frac{d}{du} = D$ であるため、成立する。
$k>1$のとき、帰納法を用いる。$u^{k}\frac{d^{k}}{du^{k}} = u \cdot u^{k-1} \frac{d^{k-1}}{du^{k-1}} \cdot \frac{d}{du}$ と書ける。帰納法の仮定から、$u^{k-1} \frac{d^{k-1}}{du^{k-1}} = D(D-1)\cdots(D-k+2)$ であるため、
\[
u^{k}\frac{d^{k}}{du^{k}} = u \cdot D(D-1)\cdots(D-k+2) \cdot \frac{d}{du}
\]
となる。ここで、$(D-a)\frac{d}{du}=\frac{d}{du}(D-a-1)$ という関係を用いると、
\[
u^{k}\frac{d^{k}}{du^{k}} = u \cdot D(D-1)\cdots(D-k+2) \cdot \frac{d}{du} = u\frac{d}{du}(D-1)\cdots(D-k+2)(D-k+1) = D(D-1)(D-2)\cdots(D-k+1)
\]
となり、主張が成立する。
\end{proof}

\begin{definition}[フロベニウスの方法]
このようにして、冪級数解を求める方法を\textbf{フロベニウスの方法}（Frobenius method）と呼ぶ。
\end{definition}

無限遠点 $x=\infty$ のまわりでの微分方程式の表現について、オイラー作用素（Euler operator）を用いてより代数的に洗練された表示を考える。
方程式\eqref{eq:linear-ode-n}に $x^{n}$ を掛け、$q_{j}\left(\frac{1}{x}\right) = x^{j}p_{j}(x)$ とおくと、
\begin{equation}\label{eq:ode-euler-x}
x^{n}\frac{d^{n}y}{dx^{n}} + \sum_{j=1}^{n} q_{j}\left(\frac{1}{x}\right) x^{n-j}\frac{d^{n-j}y}{dx^{n-j}} = 0
\end{equation}
となる。ここで補題\ref{lem:falling_factorial}より、微分演算子 $D_{x} = x\frac{d}{dx}$ を用いると $x^{k}\frac{d^{k}}{dx^{k}} = D_{x}(D_{x}-1)\cdots(D_{x}-k+1)$ と表せるため、上式は次のような演算子の形で書ける。
\[
\left[ D_{x}(D_{x}-1)\cdots(D_{x}-n+1) + \sum_{j=1}^{n}q_{j}\left(\frac{1}{x}\right) D_{x}(D_{x}-1)\cdots(D_{x}-n+j+1) \right] y = 0
\]

ここで局所座標 $u = \frac{1}{x}$ を導入する。連鎖律より $\frac{d}{dx} = -u^{2}\frac{d}{du}$ であるから、オイラー作用素は
\[
D_{x} = x\frac{d}{dx} = \frac{1}{u}\left(-u^{2}\frac{d}{du}\right) = -u\frac{d}{du} = -D_{u}
\]
と極めてシンプルに変換される。ここで、2変数多項式 $f(u; t)$ を
\[
f(u; t) = t(t-1)\cdots(t-n+1) + \sum_{j=1}^{n} q_{j}(u) t(t-1)\cdots(t-n+j+1)
\]
と定義すれば、以下の命題が成り立つ。

\begin{proposition}
\eqref{eq:linear-ode-n}は、$x=\infty$ （すなわち $u=0$）のまわりでは、オイラー作用素 $D_{u} = u\frac{d}{du}$ を用いて
\[
f(u; -D_{u})y = 0
\]
と簡潔に表現できる。
\end{proposition}

\begin{Q}[【核心】$x=\frac{1}{u}$ という変換で生じるはずの「連鎖律による複雑な低階項の混ざり合い」は、どこへ消えたのか？]
\textbf{A.} オイラー作用素 $D = x\frac{d}{dx}$ の持つ「スケール不変性」に吸収された。
通常の微分演算子 $\frac{d}{dx}$ は座標変換に対して複雑に振る舞うが、オイラー作用素は $x \to x^{a}$ という変換に対して $D_{x} \to \frac{1}{a} D_{u}$ と振る舞う。特に $x = u^{-1}$ の反転変換に対しては単に符号が反転するだけ（$D_{x} = -D_{u}$）であり、微分の階数が混ざることは一切ない。方程式全体をあらかじめ $D_{x}$ の多項式として組み立てておくことで、座標反転の代数的な複雑さを「$-1$ を代入するだけ」という操作に還元できたのである。
\end{Q}

作用素 $f(u; -D_{u})$ を展開し、フックスの定理を適用できる標準形
\[
u^{n}\frac{d^{n}y}{du^{n}} + \sum_{k=1}^{n} r_{k}(u) u^{n-k}\frac{d^{n-k}y}{du^{n-k}} = 0
\]
に書き直すことを考える。
ここで重要なのは、作用素 $(-D_{u})(-D_{u}-1)\cdots(-D_{u}-m+1)$ の最高階の微分が $(-1)^{m} u^{m} \frac{d^{m}}{du^{m}}$ となり、それ以下の階数の微分は定数係数の線形結合になるという事実である。

\begin{proposition}\label{prop:infinity_fuchs_q}
\eqref{eq:linear-ode-n}に対して、特異点 $x=\infty$ が確定特異点になるための必要十分条件は、各 $j = 1, 2, \dots, n$ について $q_{j}(u)$ が $u=0$ において正則になることである。
またこのとき、$x=\infty$ における特性指数は、決定方程式 $f(0; -s) = 0$ の根 $s$ として得られる。
\end{proposition}

\begin{Q}[【核心】作用素を展開した後の係数 $r_{k}(u)$ の正則性と、元の $q_{j}(u)$ の正則性が完全に「同値」になるのはなぜか？]
\textbf{A.} 両者の関係が、「最高階の微分から順番に決まる、三角行列の構造（帰納的構造）」を持っているからである。
作用素 $f(u; -D_{u})$ における第 $j$ 項 $q_{j}(u) (-D_{u})\cdots(-D_{u}-n+j+1)$ は、最大で $n-j$ 階の微分までしか持たない。
方程式全体を最高階の係数 $(-1)^{n}$ で割って標準形にしたとき、$(n-k)$ 階微分の係数 $r_{k}(u)$ を構成するのは以下の3つのみである。
\begin{enumerate}
    \item 第 $k$ 項から出る最高階の寄与： $(-1)^{-k} q_{k}(u)$
    \item 第 $j < k$ 項から出る低階の寄与： $q_{1}(u), \dots, q_{k-1}(u)$ の定数倍
    \item 先頭の作用素（$q$ が付かない項）から出る低階の寄与： 単なる定数
\end{enumerate}
つまり、$r_{k}(u) = (-1)^{-k} q_{k}(u) + \sum_{j=1}^{k-1} c_{k,j} q_{j}(u) + c_{k,0}$ という形になる。
\begin{itemize}
    \item $k=1$ のとき： $r_{1}(u) = -q_{1}(u) + c_{1,0}$ より、$r_{1}$ が正則 $\iff$ $q_{1}$ が正則。
    \item $k=2$ のとき： $r_{2}(u) = q_{2}(u) + c_{2,1}q_{1}(u) + c_{2,0}$ となり、$q_{1}$ はすでに正則であるため、$r_{2}$ が正則 $\iff$ $q_{2}$ が正則。
\end{itemize}
このように、上から帰納的に連鎖していくため、すべての $r_{k}(u)$ が正則であること（確定特異点であること）と、すべての $q_{j}(u)$ が正則であることは、代数的に完全に同値となるのである。
\end{Q}

\begin{Q}[【核心】決定方程式が $f(0; s)=0$ ではなく $f(0; -s)=0$ となるのはなぜか。また、この根 $s$ は「$x=\infty$ における特性指数」とどう対応するのか？]
\textbf{A.} 演算子が $D_{x} = -D_{u}$ と反転した結果、固有値も反転するからである。
解を $u$ のフロベニウス級数 $y \sim u^{s}$ と仮定したとき、$D_{u} (u^{s}) = s u^{s}$ であるため、作用素 $f(u; -D_{u})$ を適用した最低次項は $f(0; -s) = 0$ となる。
一方、$x=\infty$ における特性指数は通常、$x$ の級数 $y \sim x^{\rho}$ として定義される。$x^{\rho} = (1/u)^{\rho} = u^{-\rho}$ であるため、$s = -\rho$ の関係がある。
つまり、特性指数 $\rho$ を直接求める決定方程式は $f(0; \rho) = 0$ となり、有限の特異点と全く同じ形の多項式 $f$ に $\rho$ を代入して解けばよい、という美しい結論が得られる。
\end{Q}

\begin{example}
ガウスの超幾何微分方程式\eqref{eq:gauss_hypergeometric}において、局所座標 $u=\frac{1}{x}$ とオイラー作用素 $D_{u} = u\frac{d}{du}$ を用いると、無限遠点 $x=\infty$ の周りでの方程式は
\[
f(u; -D_{u})y = 0, \quad f(u; -D_{u}) = D_{u}(D_{u}+1) - \frac{cu-(a+b+1)}{u-1}D_{u} - \frac{ab}{u-1}
\]
と表される。このとき、$x=\infty$ での特性指数は $s=a, b$ となる。
\end{example}

\begin{proof}
方程式\eqref{eq:gauss_hypergeometric}を、オイラー作用素 $D_{x} = x\frac{d}{dx}$ （$x^{2}\frac{d^{2}}{dx^{2}} = D_{x}(D_{x}-1)$ および $x\frac{d}{dx} = D_{x}$）を用いて書き換えるため、全体を $x$ で割ると、
\[
(1-x) \frac{1}{x} D_{x}(D_{x}-1)y + \left( \frac{c}{x} - (a+b+1) \right) D_{x}y - aby = 0
\]
となる。したがって、演算子 $f(u; D_{x})$ は（$u=1/x$ として）
\[
f(u; D_{x}) = (u-1)D_{x}(D_{x}-1) + (cu-(a+b+1))D_{x} - ab
\]
と書ける。ここに $D_{x} = -D_{u}$ を代入すると、
\begin{align*}
f(u; -D_{u}) &= (u-1)(-D_{u})(-D_{u}-1) + (cu-(a+b+1))(-D_{u}) - ab \\
&= (u-1)D_{u}(D_{u}+1) - (cu-(a+b+1))D_{u} - ab
\end{align*}
全体を最高階の係数 $(u-1)$ で割って標準形にすると、
\[
D_{u}(D_{u}+1) - \frac{cu-(a+b+1)}{u-1}D_{u} - \frac{ab}{u-1}
\]
となる。
決定方程式は $u=0$ とし、作用素 $D_{u}$ を $s$ に置き換える（$f(0; -s) = 0$）ことで得られる。
\begin{align*}
s(s+1) - \frac{-(a+b+1)}{-1}s - \frac{ab}{-1} &= s(s+1) - (a+b+1)s + ab \\
&= s^{2} - (a+b)s + ab \\
&= (s-a)(s-b) = 0
\end{align*}
したがって、$x=\infty$ における特性指数は $s=a, b$ である。
\end{proof}

\begin{Q}[【核心】ガウスの超幾何関数 $F(a,b,c; x)$ のパラメータ $a, b, c$ は、一見すると不規則に並んでいるように見えるが、これらは微分方程式の特異点とどのような関係にあるのか？]
\textbf{A.} パラメータ $a, b, c$ は、まさに「各特異点における特性指数」そのものを指定するための変数である。
ガウスの微分方程式は $x=0, 1, \infty$ の3点のみを確定特異点に持つ。それぞれの特性指数を計算すると、
\begin{itemize}
    \item $x=0$ での特性指数： $0, 1-c$
    \item $x=1$ での特性指数： $0, c-a-b$
    \item $x=\infty$ での特性指数： $a, b$
\end{itemize}
となる。つまり、$\infty$ での特性指数がそのまま関数の基本パラメータ $a, b$ として採用されているのである。
（※ベルンハルト・リーマンは、この「3つの特異点とそれぞれの特性指数」を指定するだけで、微分方程式全体が大域的に一意に定まってしまうことを証明した。超幾何関数とは、パラメータの代数的な操作の産物ではなく、「リーマン球面上に3つの穴を開け、そこでの局所的な振る舞い（特性指数）を指定することで定義される幾何学的な関数」なのである。）
\end{Q}
したがって$x=\infty$における特性指数$s_{1}^{(\infty)},\cdots,s_{n}^{(\infty)}$に対して、解と係数の関係より
\begin{equation}\label{eq:sum-of-infinity-exponents}
\sum_{j=1}^{n} s_{j}^{(\infty)} = -\frac{1}{2}n(n-1) + q_{1}(0)
\end{equation}
が成り立つ。
%#endregion
%#region
ところで、有限の確定特異点 $x=\xi_{j}$ （$j=1, \dots, m$）の条件から、微分方程式の第2項の係数 $p_{1}(x)$ は有理関数であり、各特異点での留数を $a_{j}$ とすると、部分分数分解により
\[
p_{1}(x) = \sum_{j=1}^{m} \frac{a_{j}}{x-\xi_{j}} + (x\text{の多項式})
\]
と表せる。

\begin{Q}[【核心】有理関数 $p_{1}(x)$ に含まれるはずの「$x$の多項式」の部分は、なぜ完全に $0$ になると断言できるのか？]
\textbf{A.} 無限遠点 $x=\infty$ も確定特異点であるという条件が、多項式部分を禁止するからである。
命題\ref{prop:infinity_fuchs_q}より、$x=\infty$ が確定特異点であるためには $x p_{1}(x)$ が $x=\infty$ で正則、すなわち $x \to \infty$ において $p_{1}(x) = O(x^{-1})$ とならなければならない。もし $p_{1}(x)$ に定数項や $x$ の1次以上の項が含まれていれば、$x \to \infty$ で発散または $O(1)$ となり、この条件に反してしまう。したがって、多項式部分は恒等的に $0$ でなければならない。（※これはリュービルの定理の帰結でもある）。
\end{Q}

したがって、$p_{1}(x)$ は厳密に
\[
p_{1}(x) = \sum_{j=1}^{m} \frac{a_{j}}{x-\xi_{j}}
\]
となる。
各有限特異点 $x=\xi_{l}$ の局所座標 $u = x-\xi_{l}$ において、$q_{1}^{(l)}(0) = \lim_{x \to \xi_{l}} (x-\xi_{l})p_{1}(x) = a_{l}$ である。よって、$x=\xi_{l}$ における $n$ 個の特性指数 $s_{1}^{(l)}, \dots, s_{n}^{(l)}$ に対して、\eqref{eq:sum-of-exponents}より
\[
\sum_{j=1}^{n} s_{j}^{(l)} = \frac{1}{2}n(n-1) - a_{l}
\]
が成り立つ。
一方、$x=\infty$ における特性指数の和を求める。局所座標 $u=1/x$ を用いると、
\[
q_{1}^{(\infty)}(0) = \left. \frac{1}{u}p_{1}\left(\frac{1}{u}\right) \right|_{u=0} = \left. \sum_{j=1}^{m} a_{j}\frac{1}{1-u\xi_{j}} \right|_{u=0} = \sum_{j=1}^{m} a_{j}
\]
となる。

\begin{Q}[【核心】有限特異点での特性指数の和は $\frac{1}{2}n(n-1) - a_{l}$ であったのに対し、無限遠点での和はなぜ符号が反転して $-\frac{1}{2}n(n-1) + \sum a_{j}$ となるのか？]
\textbf{A.} 無限遠点における決定方程式が $f(0; -s) = 0$ であり、$s$ ではなく $-s$ が代入されることで「上昇階乗冪」が現れるからである。
有限特異点の決定方程式には下降階乗冪 $s(s-1)\cdots(s-n+1) = s^{n} - \frac{1}{2}n(n-1)s^{n-1} + \cdots$ が現れた。
一方、無限遠点では $(-s)(-s-1)\cdots(-s-n+1) = (-1)^{n} s(s+1)\cdots(s+n-1)$ となり、これを展開すると $(-1)^{n} (s^{n} + \frac{1}{2}n(n-1)s^{n-1} + \cdots)$ という上昇階乗冪になる。
全体を $(-1)^{n}$ で割って最高次の係数を $1$ にすると、$s^{n-1}$ の係数が $-\frac{1}{2}n(n-1)$ から $+\frac{1}{2}n(n-1)$ に反転していることがわかる。したがって、解と係数の関係を用いた際にも和の符号が反転し、$\sum s^{(\infty)} = -\frac{1}{2}n(n-1) + q_{1}^{(\infty)}(0)$ となるのである。
\end{Q}

これより、$x=\infty$ における特性指数の和は
\[
\sum_{j=1}^{n}s_{j}^{(\infty)} = -\frac{1}{2}n(n-1) + \sum_{j=1}^{m} a_{j}
\]
となる。
有限特異点と無限遠点のすべての特性指数の和をとると、$a_{j}$ の項が見事に相殺され、以下の極めて重要な大域的関係式が成立する。

\begin{proposition}\label{prop:fuchs_relation}
リーマン球面上の特異点がすべて確定特異点であるようなフックス型微分方程式\eqref{eq:linear-ode-n}の特性指数には、
\begin{align*}
\sum_{l=1}^{m}\sum_{j=1}^{n} s_{j}^{(l)} + \sum_{j=1}^{n} s_{j}^{(\infty)} &= \sum_{l=1}^{m} \left( \frac{1}{2}n(n-1) - a_{l} \right) + \left( -\frac{1}{2}n(n-1) + \sum_{j=1}^{m} a_{j} \right) \\
&= \frac{1}{2}n(n-1)(m-1)
\end{align*}
という関係式が成立する。この方程式のパラメータを大域的に縛る拘束条件を\textbf{フックスの関係式（Fuchs' relation）}という。
\end{proposition}
%#endregion
%#region
微分方程式を完全に決定するためには、その係数に含まれる独立なパラメータの数を数え上げればよい。
$n$ 階フックス型微分方程式において、確定特異点の位置を $\Xi=\{\xi_{1}, \xi_{2}, \dots, \xi_{m}, \infty\}$ として考える。

まず、各特異点における局所的な振る舞いを決定する「特性指数」の自由度について考える。
特異点は全部で $m+1$ 個あり、各点に $n$ 個の特性指数が存在するため、特性指数の総数は $(m+1)n$ 個である。しかし、これらが満たすべき大域的な制約は命題\ref{prop:fuchs_relation}（フックスの関係式）の1つのみであるため、特性指数のうち自由に選ぶことができるパラメータの数は
\begin{equation}
E_{0}(n,m) = (m+1)n - 1
\end{equation}
個となる。

次に、微分方程式の係数 $p_{i}(x)$ 全体に含まれるパラメータの総数を数え上げる。
命題\ref{prop:fuchs_theorem_n}から、有限の特異点 $x=\xi_{l}$ における極の最大位数は $i$ であるため、$p_{i}(x)$ は $mi$ 個の未知定数 $a_{k}^{(l)}$ （$k=1,\dots,i$, $l=1,\dots,m$）を用いて、部分分数分解により
\[
p_{i}(x) = \sum_{l=1}^{m}\sum_{k=1}^{i} \frac{a_{k}^{(l)}}{(x-\xi_{l})^{k}} + (x\text{の多項式})
\]
と表される。
さらに命題\ref{prop:infinity_fuchs_q}より、無限遠点 $x=\infty$ が確定特異点であるためには $x^{i}p_{i}(x)$ が $\infty$ で正則でなければならない。これにより、$p_{i}(x)$ は $x \to \infty$ で $O(x^{-i})$ となることが要求されるため、多項式部分は $0$ となり、
\begin{equation}\label{eq:pi_fraction}
p_{i}(x) = \sum_{l=1}^{m}\sum_{k=1}^{i} \frac{a_{k}^{(l)}}{(x-\xi_{l})^{k}}
\end{equation}
に制限される。さらに、$p_{i}(x) = O(x^{-i})$ という条件から、$a_{k}^{(l)}$ たちの間に一定の拘束条件が生じる。

\begin{Q}[【核心】$p_{i}(x) = O(x^{-i})$ という無限遠点での条件は、具体的にいくつの独立なパラメータを奪う（束縛する）のか？]
\textbf{A.} 厳密に $i-1$ 個の独立な自由度を奪う。
局所座標 $u=1/x$ を用いて\eqref{eq:pi_fraction}を $u=0$ のまわりでテイラー展開する。
\[
\frac{1}{(x-\xi_{l})^{k}} = \frac{u^{k}}{(1-\xi_{l}u)^{k}} = u^{k} \left( 1 + k\xi_{l}u + \frac{k(k+1)}{2}\xi_{l}^{2}u^{2} + \cdots \right)
\]
これを $p_{i}(x)$ に代入して $u$ のべき級数として整理したとき、$u^{i}$ 以上のオーダー（$O(u^{i})$）にならなければならない。つまり、展開した際の $u^{1}, u^{2}, \dots, u^{i-1}$ の係数がすべて $0$ になる必要がある。
これは未知定数 $a_{k}^{(l)}$ に対する $i-1$ 個の線形連立方程式となる。特異点の位置 $\xi_{l}$ はすべて相異なるため、この連立方程式の係数行列（ヴァンデルモンド行列の微分型）はフルランクとなる。したがって、これら $i-1$ 個の条件式は完全に独立であり、冗長なものは一つもない。
\end{Q}

これにより、$p_{i}(x)$ に含まれる独立なパラメータの数は、元の $mi$ 個から独立な制約 $i-1$ 個を引いて
\[
mi - (i-1) = m i - i + 1
\]
個となる。
これを $i=1, 2, \dots, n$ について合計すれば、方程式全体に含まれるパラメータの総数が得られる。
\begin{align*}
\sum_{i=1}^{n} (mi - i + 1) &= (m-1)\sum_{i=1}^{n} i + \sum_{i=1}^{n} 1 \\
&= (m-1)\frac{1}{2}n(n+1) + n \\
&= \frac{1}{2}n \big( m(n+1) - (n-1) \big)
\end{align*}

以上の議論から、フックス型微分方程式のモジュライ空間の次元（パラメータ数）に関する以下の重要な命題を得る。

\begin{proposition}\label{prop:fuchs_parameters}
\eqref{eq:linear-ode-n}が $\Xi=\{\xi_{1}, \xi_{2}, \dots, \xi_{m}, \infty\}$ に確定特異点をもつフックス型微分方程式である時、この微分方程式の係数全体は、
\[
E(n,m) = \frac{1}{2}n \big( m(n+1) - (n-1) \big)
\]
個の独立なパラメータを含む。
\end{proposition}

命題\ref{prop:fuchs_parameters}より、方程式の係数に含まれる全パラメータ数は $E(n,m)$ である。一方、特性指数として自由に選べるパラメータ数は $E_{0}(n,m)$ であった。
これらの差を計算すると、
\begin{align*}
E(n,m) - E_{0}(n,m) &= \frac{1}{2}n \big( m(n+1) - (n-1) \big) - \big( (m+1)n - 1 \big) \\
&= \frac{1}{2}(n-1)(mn - n - 2)
\end{align*}
となる。$m \ge 2, n \ge 2$ のとき
\[
E(n,m) - E_{0}(n,m) \ge 0
\]
であるため、各特異点での特性指数をすべて指定したとしても、一般にはそれだけでは微分方程式を完全に決定するには至らないことがわかる。

\begin{definition}
このように、微分方程式の係数に含まれるパラメータのうち、特性指数の値に関係しない（特性指数を固定してもなお決定されない）独立なパラメータのことを\textbf{アクセサリー・パラメータ}（accessory parameter）という。その総数は $N_{acc} = \frac{1}{2}(n-1)(mn - n - 2)$ で与えられる。
\end{definition}

\begin{example}\label{ex:acc_parameter_zero}
2階フックス型微分方程式で、確定特異点が3つ（有限に $m=2$ 個、および無限遠点 $\infty$）の場合について考える。
このとき、$E(2,2) = 5$、$E_{0}(2,2) = 5$ であるため、$N_{acc} = 5 - 5 = 0$ となる。
したがって、この場合にはアクセサリー・パラメータは存在しない。
\end{example}

\begin{Q}[【核心】特異点が3つの場合（$m=2$）にアクセサリー・パラメータが $0$ になることは、解析学的に何を意味しているのか？]
\textbf{A.} 「特異点の位置」と「各特異点での特性指数（局所的な振る舞い）」を指定するだけで、微分方程式全体が大域的にただ1つに決定されてしまう（剛性を持つ）ことを意味している。
特異点3つの2階方程式といえば「ガウスの超幾何微分方程式」である。超幾何関数が持つ数々の美しい変換則や積分表示は、この「アクセサリー・パラメータを持たない（特性指数だけで方程式が一意にロックされる）」という強烈な代数的性質の恩恵に他ならない。
\end{Q}

\begin{Q}[【核心】特異点を1つ増やして4つ（$m=3$）にした場合、何が起きるのか？また、それがパンルヴェ方程式とどう繋がるのか？]
\textbf{A.} $n=2, m=3$ のとき、$N_{acc} = \frac{1}{2}(1)(6 - 2 - 2) = 1$ となり、アクセサリー・パラメータが初めて1つ出現する。これが「ホイン（Heun）の微分方程式」である。
特性指数をすべて固定しても、未知のパラメータが1つ残るため、解のモノドロミー群はこのパラメータに依存して変化してしまう。
岡本和夫の理論（モノドロミー保存変形）は、特異点の位置 $t$ を動かしたときに「モノドロミー群が変化しないように」このアクセサリー・パラメータをどう時間発展させればよいか、という非線形力学系を構築することである。その際、岡本はこの1次元のパラメータを直接扱うのではなく、あえて「見かけの特異点 $q$」と「その留数に由来する運動量 $p$」という新しい変数を導入することで、この変形方程式を美しいハミルトン系（パンルヴェVI型方程式）として見事に定式化したのである。
\end{Q}
%#endregion
%#region
ここで、\eqref{eq:linear-ode-n}の各特異点における特性指数を一覧として書き並べ、方程式の大域的な構造を視覚的に捉える記法を導入する。

\begin{definition}[リーマン図式]
    確定特異点 $x = \xi_1, \xi_2, \dots, \xi_m, \infty$ を持つフックス型微分方程式\eqref{eq:linear-ode-n}に対し、各特異点での特性指数を列挙した表
    \[
    \left\{
    \begin{matrix}
        x = \xi_1 & x = \xi_2 & \cdots & x = \xi_m & x = \infty \\
        s_1^{(1)} & s_1^{(2)} & \cdots & s_1^{(m)} & s_1^{(\infty)} \\
        s_2^{(1)} & s_2^{(2)} & \cdots & s_2^{(m)} & s_2^{(\infty)} \\
        \vdots    & \vdots    &        & \vdots    & \vdots \\
        s_n^{(1)} & s_n^{(2)} & \cdots & s_n^{(m)} & s_n^{(\infty)}
    \end{matrix}
    \right\}
    \]
    を、その微分方程式の\textbf{リーマン図式}（Riemann scheme）という。（※記号 $P$ を冠してリーマンのP関数と呼ばれることもある。）
\end{definition}

\begin{example}
    ガウスの超幾何微分方程式\eqref{eq:gauss_hypergeometric}のリーマン図式は、以下のようになる。
    \[
    \left\{
    \begin{matrix}
        x = 0 & x = 1 & x = \infty \\
        0 & 0 & a \\
        1-c & c-a-b & b \\
    \end{matrix}
    \right\}
    \]
\end{example}

\begin{proof}
    ガウスの超幾何微分方程式\eqref{eq:gauss_hypergeometric}は、$x=0, 1, \infty$ の3点に確定特異点を持つフックス型微分方程式である。
    それぞれの特異点における特性指数は、これまでの計算により
    \begin{itemize}
        \item $x=0$ での特性指数： $0, 1-c$
        \item $x=\infty$ での特性指数： $a, b$
    \end{itemize}
    であることが既にわかっている。残る $x=1$ での特性指数を求める。
    局所座標を $u = x-1$ とすると $x = u+1$ であり、$x(1-x) = (u+1)(-u) = -u(u+1)$ となる。方程式全体を $-u(u+1)$ で割って標準形にすると、
    \[
    \frac{d^{2}y}{du^{2}} - \frac{c - (a+b+1)(u+1)}{u(u+1)} \frac{dy}{du} + \frac{ab}{u(u+1)} y = 0
    \]
    となる。フックスの定理の条件に従い、$q_{1}(u), q_{2}(u)$ を計算して $u=0$ を代入する。
    \begin{align*}
        q_{1}(0) &= \left. u \left( - \frac{c - (a+b+1)(u+1)}{u(u+1)} \right) \right|_{u=0} = \left. - \frac{c - (a+b+1)(u+1)}{u+1} \right|_{u=0} = -c + a + b + 1 \\
        q_{2}(0) &= \left. u^{2} \left( \frac{ab}{u(u+1)} \right) \right|_{u=0} = \left. u \frac{ab}{u+1} \right|_{u=0} = 0
    \end{align*}
    したがって、$x=1$（すなわち $u=0$）における決定方程式は $s(s-1) + q_{1}(0)s + q_{2}(0) = 0$ より、
    \[
    s(s-1) + (a+b+1-c)s = s^{2} + (a+b-c)s = s(s - (c-a-b)) = 0
    \]
    となり、特性指数は $0, c-a-b$ であることが示された。
\end{proof}

\begin{Q}[【核心】この超幾何微分方程式のリーマン図式に対して、先ほど証明した「フックスの関係式」は本当に成り立っているか？]
\textbf{A.} 完璧に成り立っている。
特異点の数は $m+1 = 3$（有限が $m=2$）、階数は $n=2$ である。
フックスの関係式（命題\ref{prop:fuchs_relation}）の右辺は、$\frac{1}{2}n(n-1)(m-1) = \frac{1}{2}(2)(1)(1) = 1$ となる。
一方、左辺（リーマン図式に現れる全6個の特性指数の総和）を計算すると、
\[
0 + (1-c) + 0 + (c-a-b) + a + b = 1
\]
となり、パラメータ $a, b, c$ の値にいかなる依存もせず、見事に両辺が $1$ で一致する。空間のトポロジーが特性指数を大域的に縛っていることの美しい実例である。
\end{Q}

\begin{Q}[【核心】なぜわざわざ、特性指数をこのような「表（リーマン図式）」の形式で記述するのか？]
\textbf{A.} 前段で示した通り、$x=0, 1, \infty$ の3点のみを確定特異点に持つ2階方程式（$m=2, n=2$）は「アクセサリー・パラメータを持たない」からである。
アクセサリー・パラメータが $0$ 個であるということは、「特異点の位置」と「各特異点の特性指数」さえ決めれば、元の微分方程式が一意に完全に復元できる（剛性を持つ）ことを意味する。つまり、このリーマン図式は単なるパラメータのメモ書きではなく、これ自体が「ガウスの超幾何関数」という対象を完全に過不足なく定義する「関数記号（リーマンのP関数）」そのものとして機能するのである。
\end{Q}
%#endregion
%#region
ここで、フックス型微分方程式\eqref{eq:linear-ode-n}の特性指数を操作する基本的な変換について述べる。
任意の定数 $\alpha_{1}, \dots, \alpha_{m}$ を用い、未知関数 $y$ を
\begin{equation}\label{eq:gauge_transform}
y = \left( \prod_{l=1}^{m} (x-\xi_{l})^{\alpha_{l}} \right) z
\end{equation}
によって新しい未知関数 $z$ に変数変換しても、微分方程式はフックス型のままである。
\begin{Q}[【核心】$y = \left( \prod_{l=1}^{m} (x-\xi_{l})^{\alpha_{l}} \right) z$ という変換を行っても、微分方程式が「フックス型（確定特異点のみを持つ）」であるという性質は本当に壊れないのか？]
\textbf{A.} 本当である。その理由は、変換因子 $\Phi(x) = \prod_{l=1}^{m} (x-\xi_{l})^{\alpha_{l}}$ の「対数微分」が、各特異点において高々1位の極しか生み出さないからである。
簡単なため2階のフックス型方程式 $y'' + p_{1}y' + p_{2}y = 0$ で確認する。
$y = \Phi z$ とおいて微分すると、
\begin{align*}
y' &= \Phi z' + \Phi' z \\
y'' &= \Phi z'' + 2\Phi' z' + \Phi'' z
\end{align*}
となる。これらを元の方程式に代入して $\Phi$ で割ると、$z$ についての微分方程式は
\[
z'' + \left( p_{1} + 2\frac{\Phi'}{\Phi} \right) z' + \left( p_{2} + p_{1}\frac{\Phi'}{\Phi} + \frac{\Phi''}{\Phi} \right) z = 0
\]
となる。新しい係数を $\tilde{p}_{1}(x), \tilde{p}_{2}(x)$ とおく。
ここで、$\Phi(x)$ の対数微分 $\frac{\Phi'}{\Phi}$ を計算すると、
\[
\frac{\Phi'(x)}{\Phi(x)} = \frac{d}{dx} \log \Phi(x) = \frac{d}{dx} \sum_{l=1}^{m} \alpha_{l} \log(x-\xi_{l}) = \sum_{l=1}^{m} \frac{\alpha_{l}}{x-\xi_{l}}
\]
となり、これは各 $x=\xi_{l}$ において「ちょうど1位の極」を持つ有理関数となる。
さらに、$\frac{\Phi''}{\Phi} = \left(\frac{\Phi'}{\Phi}\right)' + \left(\frac{\Phi'}{\Phi}\right)^{2}$ であるため、1位の極の微分（2位の極）と平方（2位の極）の和となり、これも「高々2位の極」しか持たない。
元の $p_{1}$ が1位の極、$p_{2}$ が2位の極を持つため、足し合わせた新しい係数 $\tilde{p}_{1}$ は依然として1位の極を、$\tilde{p}_{2}$ は依然として2位の極を持つにとどまる。
無限遠点に関しても同様のオーダー計算が成り立つため、このゲージ変換はフックス型の条件（極の位数制限）を一切破壊しないのである。
\end{Q}
このとき、変換後の微分方程式（$z$ の方程式）のリーマン図式は、元の図式から以下のように特性指数がシフトする。
\begin{itemize}
    \item 有限の特異点 $x=\xi_{l}$ における特性指数は、$s_{j}^{(l)} - \alpha_{l}$ となる。
    \item 無限遠点 $x=\infty$ における特性指数は、$s_{j}^{(\infty)} + \sum_{l=1}^{m} \alpha_{l}$ となる。
\end{itemize}

\begin{Q}[【核心】有限特異点では特性指数が $-\alpha_{l}$ と引き算されるのに、なぜ無限遠点の特性指数だけ $+\sum \alpha_{l}$ と足し算になるのか？]
\textbf{A.} 無限遠点の局所座標 $u=1/x$ で変換式を展開すれば明らかになる。
変換の因子を $x \to \infty$ （すなわち $u \to 0$）で展開すると、
\[
\prod_{l=1}^{m} (x-\xi_{l})^{\alpha_{l}} = \prod_{l=1}^{m} \left( \frac{1}{u} - \xi_{l} \right)^{\alpha_{l}} = u^{-\sum \alpha_{l}} \prod_{l=1}^{m} (1 - \xi_{l}u)^{\alpha_{l}} \sim u^{-\sum \alpha_{l}}
\]
となる。
解 $y$ が無限遠点で $u^{s^{(\infty)}}$ のオーダーで振る舞うとすれば、$u^{s^{(\infty)}} \sim u^{-\sum \alpha_{l}} z$ となるため、新しい未知関数 $z$ は $u^{s^{(\infty)} + \sum \alpha_{l}}$ のオーダーで振る舞うことになる。
（※これによって、有限特異点でのシフトの合計 $-n\sum \alpha_{l}$ と無限遠点でのシフトの合計 $+n\sum \alpha_{l}$ が見事に相殺し、フックスの関係式が保たれるという極めて美しい代数幾何学的構造になっている。）
\end{Q}

この特性指数のシフトを利用することで、方程式を標準形に帰着させることができる。
各有限特異点において $\alpha_{l} = s_{1}^{(l)}$ と選べば、変換後の有限特異点における特性指数のうち1つを
\[
s_{1}^{(1)} - \alpha_{1} = s_{1}^{(2)} - \alpha_{2} = \cdots = s_{1}^{(m)} - \alpha_{m} = 0
\]
と、すべて $0$ に揃えることができる。

特に $n=2, m=2$ （確定特異点が全体で3つ）のときは、アクセサリー・パラメータを持たない（$N_{acc}=0$）。
さらに、複素平面上の任意の相異なる3点は、一次分数変換（メビウス変換）によって必ず $x=0, 1, \infty$ の3点に写すことができる。
したがって、特異点を $0, 1, \infty$ に移動させた後、上記の\eqref{eq:gauge_transform}の変換を行って $x=0, 1$ の特性指数の1つを $0$ に揃えれば、リーマン図式は以下の非常にシンプルな形に帰着される。
\[
    \left\{
    \begin{matrix}
        x = 0 & x = 1 & x = \infty \\
        0 & 0 & s_{1}^{(\infty)} + s_{1}^{(1)} + s_{1}^{(2)} \\
        s_{2}^{(1)} - s_{1}^{(1)} & s_{2}^{(2)} - s_{1}^{(2)} & s_{2}^{(\infty)} + s_{1}^{(1)} + s_{1}^{(2)} \\
    \end{matrix}
    \right\}
\]
この図式の形は、パラメータを適当に $a, b, c$ と置き直せば、まさにガウスの超幾何微分方程式のリーマン図式と完全に一致する。

\begin{Q}[【核心】これらの一連の変換によって、「ガウスの超幾何微分方程式」の数学的な位置づけはどう変わるか？]
\textbf{A.} 「数ある特殊関数の1つ」から、「3つの確定特異点を持つすべての2階線形微分方程式の、普遍的な標準形（Universal normal form）」へと昇華される。
アクセサリー・パラメータが $0$ であるため、リーマン図式が一致すれば、背後にある微分方程式も完全に一致する（剛性）。つまり、物理学や幾何学のあらゆる問題で「3つの確定特異点を持つ方程式」が現れた瞬間、それはメビウス変換と冪関数の掛け算という「自明な操作」を除いて、すべてガウスの超幾何関数に他ならないことが保証されるのである。（※これはベルンハルト・リーマンがその驚異的な洞察力で到達した結論である。）
\end{Q}
%#endregion

\subsection{みかけの特異点}
%#region
\eqref{eq:linear-ode-n}の解の基本系 $\vec{\varphi}(x)$ をとる。確定特異点の集合 $\Xi$ を通らない閉曲線 $\gamma$ について解析接続を考えると、回路行列 $\Gamma(\gamma)$ を用いて
\[
\vec{\varphi}^{\gamma}(x) = \vec{\varphi}(x)\Gamma(\gamma)
\]
とかける。ここでもう一つの解の基本系 $\vec{\psi}(x)$ をとっておくと、可逆行列 $C$ を用いて
\[
\vec{\varphi}(x) = \vec{\psi}(x)C
\]
となる。また、解析接続と線型操作は交換できるので、
\[
\vec{\varphi}^{\gamma}(x) = \vec{\psi}^{\gamma}(x)C
\]
よって、$\vec{\psi}(x)$ の回路行列は、
\[
\vec{\psi}^{\gamma}(x) = \vec{\varphi}^{\gamma}(x)C^{-1} = \vec{\varphi}(x)\Gamma(\gamma)C^{-1} = \vec{\psi}(x)C\Gamma(\gamma)C^{-1}
\]
より、$C\Gamma(\gamma)C^{-1}$ となる。

\begin{definition}[リーマンデータ]\label{def:riemann_data}
    フックス型微分方程式 \eqref{eq:linear-ode-n} に対して、定義域 $\mathbb{P}^{1}(\mathbb{C})$、確定特異点の集合 $\Xi$、およびモノドロミー $\hat{\rho}$ の3つ組
    \[
    (\mathbb{P}^{1}(\mathbb{C}), \Xi, \hat{\rho})
    \]
    のことを\textbf{リーマンデータ}という。
\end{definition}

\begin{Q}[モノドロミーって何？]
\textbf{A.} $\cdot$ 閉曲線に対して回路行列を対応させる写像のようなもの。\\
\phantom{\textbf{A.}} $\cdot$ これは基本群 $\pi_1(X)$ から行列群への表現とみなすことができ、微分方程式の解の「多価性」を調べるための中心的な道具となる。(定義\ref{def:monodromy}を参照)\\
\end{Q}
%#endregion
%#region
ここで、モノドロミー群を決定するために必要なパラメータの数を考えたい。
基点 $x_{0}$ をうまく取ることによって、全ての有限の特異点が異なる偏角を持つようにでき、特に $\arg(\xi_{1}-x_{0}) < \arg(\xi_{2}-x_{0}) < \cdots < \arg(\xi_{m}-x_{0})$ となるように並べることができる。
ここで、それぞれの特異点 $\xi_{l}$ を反時計回りに一周だけして基点に戻ってくる閉道を $\gamma_{l}$ とし、$x=\infty$ を反時計回り（無限遠点から見て反時計回り）に一周して戻ってくる閉道を $\gamma_{\infty}$ とすると、空間の位相的な構造から
\[
[\gamma_{\infty}] = ([\gamma_{1}] \cdot [\gamma_{2}] \cdots [\gamma_{m}])^{-1} = [\gamma_{m}]^{-1} \cdots [\gamma_{2}]^{-1} \cdot [\gamma_{1}]^{-1}
\]
というホモトピー同値の関係が成り立つ。
よって、各閉道に沿った解析接続から定まるモノドロミー行列を $\Gamma(\gamma)$ とすれば、
\[
\Gamma(\gamma_{\infty}) = \Gamma(\gamma_{m})^{-1} \cdots \Gamma(\gamma_{2})^{-1} \Gamma(\gamma_{1})^{-1}
\]
となり、無限遠点でのモノドロミー行列は、有限特異点のモノドロミー行列の組から完全に決定される。
したがって、モノドロミー群の生成元全体を決定するには、$m$ 個の $n \times n$ 行列 $\Gamma(\gamma_{1}), \dots, \Gamma(\gamma_{m})$、すなわち $mn^{2}$ 個のパラメータを決めれば良いことがわかる。

さらに、モノドロミー表現は解の基本系の取り替え（$GL(n, \mathbb{C})$ による全体の相似変換）という冗長な自由度を持つ。
相似変換に用いる正則行列が持つ $n^{2}$ 個の自由度のうち、変換に寄与しない定数倍（群の中心 $Z(GL(n, \mathbb{C}))$）の $1$ 次元分を除いた「実質的な相似変換の自由度」は $n^{2}-1$ である。この分を差し引いて、次の命題を得る。

\begin{proposition}\label{prop:monodromy_parameters}
    フックス型微分方程式\eqref{eq:linear-ode-n}のモノドロミー表現のモジュライ空間は、
    \[
    M(n,m):=mn^{2} - (n^{2}-1) = (m-1)n^{2} + 1
    \]
    個の独立なパラメータを含む。
\end{proposition}

\begin{Q}[【核心】微分方程式自体のパラメータ数 $E(n,m)$ と、モノドロミー表現のパラメータ数 $(m-1)n^{2} + 1$ を比較すると、常に微分方程式の方がパラメータが多くなる（自由度が高い）。これは何を意味しているのか？]
\textbf{A.} 「同じモノドロミー群を持つフックス型微分方程式が、無数に存在する」ということである。
実際に差を計算してみると、特異点の数 $m$ や階数 $n$ が大きくなるにつれて、方程式の自由度 $E(n,m)$ の方が圧倒的に大きくなる。（※唯一の例外が、両者が一致するガウスの超幾何方程式 $E(2,2)=5$ である）。
「与えられたモノドロミー表現を持つフックス型方程式を構成せよ」という問題をヒルベルトの第21問題（リーマン・ヒルベルト問題）と呼ぶ。パラメータ数にギャップがあるため、モノドロミーを指定するだけでは方程式は一意に定まらず、余剰な自由度（アクセサリー・パラメータ）が生じてしまう。
そこで、この余剰な自由度をあえて「見かけの特異点」として幾何学的に配置し、本質的な特異点 $t$ を動かしたときに「モノドロミー群が一定に保たれるように」見かけの特異点を時間発展させる、という壮大なアイデアが生まれた。これが岡本和夫らによるモノドロミー保存変形の理論であり、パンルヴェ方程式のハミルトニアン構造の起源なのである。
\end{Q}

ここまで見てきたように、一般にパラメータの数について
\[
E_{0}(n,m) \le E(n,m) \le M(n,m)
\]
が成り立つ。実際、$m \ge 2, n \ge 2$ の場合、右側の不等式について評価すると、
\begin{align}
M(n,m)-E(n,m) &= (m-1)n^{2} + 1-\frac{1}{2}n \big( m(n+1) - (n-1) \big) \notag \\
&= \frac{1}{2}n \big\{ (m-1)(n-1)-2 \big\} + 1 \ge 0
\end{align}
となる。ここで、整数 $m \ge 2, n \ge 2$ のもとで上式がちょうど $0$ となる（等号が成立する）条件を探ると、$n=2, m=2$ のケースのみであることがわかる。このとき、
\[
E_{0}(2,2)=E(2,2)=M(2,2)=5
\]
となり、この極めて特殊な状況が、ガウスの超幾何微分方程式によって実現されている。

\begin{Q}[【核心】$E(n,m)=M(n,m)$ が成り立つ（$n=2, m=2$）ということは、微分方程式論において何を意味しているのか？なぜそれがガウスの超幾何微分方程式に結びつくのか？]
\textbf{A.} $E$ は「方程式を決定するために必要なパラメータ数」、$M$ は「見かけの特異点などを伴うアクセサリー・パラメータの数」を意味する。$M - E = 0$ となることは、\textbf{「余分なアクセサリー・パラメータが一切存在せず、特異点の位置とRiemann scheme（特性指数）だけで方程式がグローバルにただ一つに決定される」}という強固な性質（剛性, rigidity）を示している。

$m=2$ は有限確定特異点が2つ（通常 $x=0, 1$）、これに $\infty$ を加えた計3つの確定特異点を持つ状況を指す。階数 $n=2$ でこの条件を満たすものは、まさにガウスの超幾何微分方程式に他ならない。計算上の等号成立条件が、歴史的な名方程式の存在条件と見事に一致しているのである。
\end{Q}
%#endregion
%#region
また、$E(n,m)-E_{0}(n,m)$ は方程式を決定するための「アクセサリー・パラメータ」の個数を表しているが、さらに
\[
K(n,m)=M(n,m)-E(n,m)
\]
が何を意味しているのかを考察する。

ここで、「モノドロミー表現（大域的な性質）を変えることなく、微分方程式の確定特異点の数（局所的な性質）を増やすことは可能か？」という問いを立てる。
式 \eqref{eq:linear-ode-n} のリーマンデータが
\[
(\mathbb{P}^{1}(\mathbb{C}),\Xi,\hat{\rho}), \quad \Xi=\{\xi_{1},\cdots,\xi_{m},\xi_{m+1}\}
\]
で与えられているとし、ここに新しい特異点の集合 $\Lambda=\{\lambda_{1},\cdots,\lambda_{g}\}$ が追加されたとする。
これらの追加された特異点がモノドロミーを変化させないためには、$\Lambda$ の各点の周りでの解が1価（分岐を持たない）でなければならない。すなわち、各 $x=\lambda_{i}$ において以下の2条件を満たす必要がある。
\begin{itemize}
    \item[(a)] 特性指数が全て整数である。（任意の解に無理数や虚数の複素べきが現れない）
    \item[(b)] 非対数的特異点である。（フロベニウス法で解を構成する際、任意の解に対数関数 $\log(x-\lambda_{i})$ が現れない）
\end{itemize}

\begin{definition}[みかけの特異点]
    上の(a), (b)の性質を同時に満たす確定特異点 $x=\lambda_{i}$ を、\textbf{みかけの特異点 (apparent singularity)} という。
\end{definition}

では、リーマンデータとして $(\mathbb{P}^{1}(\mathbb{C}),\Xi\cup\Lambda,\hat{\rho})$ を持つ微分方程式 \eqref{eq:linear-ode-n} のパラメータ数を数え上げる。

\begin{Q}[【核心】みかけの特異点を $g$ 個追加したとき、微分方程式のパラメータ数はどう変化するのか？また、その物理的・幾何学的な意味は何か？]
\textbf{A.} 結論から言えば、パラメータの数は元の $E(n,m)$ からちょうど $g$ 個だけ増え、$E(n,m)+g$ となる。

みかけの特異点が1つ（$x=\lambda$）増えたときのパラメータの増減をトレースすると以下のようになる。
\begin{itemize}
    \item \textbf{特異点追加によるベースの増加}: 新たな確定特異点が1つ増えるため、微分方程式の係数に与える自由度は本来 $\frac{1}{2}n(n+1)$ 個増える（$E(n,m+1)-E(n,m)$ に相当）。
    \item \textbf{条件(1)による減少}: 特性指数 $n$ 個が全て「整数」に固定される。連続な複素パラメータ（自由度）が離散的な値に縛られるため、自由度は $n$ 個減少する。
    \item \textbf{条件(2)による減少}: 対数項を消去するための適合条件（フロベニウス法における各特性指数の差からの共鳴条件）が課される。これは独立な代数方程式を $\frac{1}{2}n(n-1)$ 個要求するため、自由度がさらに $\frac{1}{2}n(n-1)$ 個減少する。
    \item \textbf{特異点の位置という新たな自由度}: 本来の特異点 $\Xi$ は固定されているが、みかけの特異点 $\lambda$ の「位置」自体は微分方程式の係数（有理関数）の極として自由に動かせるため、自由度が $+1$ される。
\end{itemize}
これらを足し合わせると、
\[
\frac{1}{2}n(n+1) - n - \frac{1}{2}n(n-1) + 1 = 1
\]
となり、みかけの特異点1つにつきパラメータはちょうど1つ（特異点の位置 $\lambda$ の分だけ）増えることがわかる。
\end{Q}

\begin{proposition}
    微分方程式 \eqref{eq:linear-ode-n} が、本来の確定特異点 $\Xi$ の他に、$g$ 個のみかけの特異点 $\Lambda$ を持つとすると、その微分方程式を決定するパラメータの総数は
    \[
    E(n,m)+g
    \]
    で与えられる。
\end{proposition}

\begin{Q}[【核心】予備知識なしに微分方程式の係数 $p_i(x)$ の形（有理関数）だけを見たとき、本来の確定特異点 $\xi$ とみかけの特異点 $\lambda$ は区別できるのか？]
\textbf{A.} 一見しただけでは全く区別できないが、「局所的な解析」を行えば完全に区別できる。

微分方程式の係数 $p_i(x)$ の分母の極として現れる点という視点では、$\xi$ も $\lambda$ も全く同じ「確定特異点」として振る舞う。しかし、ある特異点 $x=a$ のまわりでフロベニウス法を用いて級数解を構成しようとしたとき、以下の2つの厳しいテストをクリアした場合にのみ、その点 $a$ は「みかけの特異点である」と事後的に判定される。
\begin{enumerate}
    \item 決定方程式（indicial equation）を解くと、特性指数が「全て整数」になる。
    \item 特性指数の差が整数であることに起因する「対数項（$\log$）が発生するリスク」を、係数同士の絶妙な関係式（適合条件）によって完全に回避している。
\end{enumerate}
すなわち、みかけの特異点とは「係数の分母がゼロになるため見た目は特異点だが、解の分岐（モノドロミー）を一切引き起こさないように、係数パラメータが奇跡的なバランスで調整された点」なのである。
\end{Q}

\begin{Q}[【核心】みかけの特異点 $\lambda$ を追加した際のパラメータ増減を計算する際、なぜ最初は「特異点を固定して定義した」はずの公式 $E(n,m+1)$ を使って計算を出発しているのか？定義と矛盾していないか？]
\textbf{A.} 全く矛盾していない。これは「一旦、位置を固定した普通の確定特異点として追加してから、みかけの特異点へとクラスチェンジさせる」という段階的な思考実験（計算テクニック）を行っているからである。

具体的には、以下の3ステップの論理でパラメータを数え上げている。
\begin{enumerate}
    \item \textbf{枠組みの拡張（代数的な自由度の追加）}:
    まず、点 $x=\lambda$ を「位置が固定されたただの確定特異点」とみなして係数の分母に追加する。この時点では単なる極の追加であるため、微分方程式の係数が持つ代数的な自由度は $E(n,m+1) - E(n,m) = \frac{1}{2}n(n+1)$ 個増える。
    
    \item \textbf{条件の付加（自由度の剥奪）}:
    次に、この点 $x=\lambda$ が「みかけの特異点」になるための厳しい条件（特性指数が整数、対数項なし）を後から課す。これにより、ステップ1で追加された自由度のうち、ちょうど $\frac{1}{2}n(n+1)$ 個のパラメータが「条件を満たすための調整役」として完全に束縛される。結果として、係数由来の自由度の増加は「相殺されてゼロ」になる。
    
    \item \textbf{パラダイムシフト（位置のパラメータ化）}:
    最後に、条件を満たしたことで解の分岐（モノドロミー）に一切影響を与えなくなった点 $x=\lambda$ に対して、「位置を固定しなければならない」という足かせを外す。係数パラメータの自由度が増えなかった代わりに、「$\lambda$ という座標そのもの」が新たな動的パラメータ（$+1$）として復活する。
\end{enumerate}
すなわち、$E(n,m+1)$ は、計算の土俵に新しい極を乗せるための「一時的な踏み台」として機能しているのである。
\end{Q}

\begin{Q}[【核心】「みかけの特異点 $\lambda$ の位置がパラメータになる」という言葉の真意は何か？単に1つの微分方程式の中で、特異点を物理的に動かしているという意味か？]
\textbf{A.} そうではない。本質的な意味は、\textbf{「同じモノドロミー（大域的性質）を共有する微分方程式が、$\lambda$ の取り得る値の分だけ無数に存在し、ひとつの『族（family）』をなしている」}ということである。

本来の確定特異点 $\Xi$ とモノドロミー表現 $\hat{\rho}$ を完全に固定したとする。もしみかけの特異点が存在しなければ、方程式はただ1つに決定される（剛性）。しかし、みかけの特異点 $\lambda$ が存在すると、その位置 $\lambda$ を別の値 $\lambda'$ に変えるごとに、「係数の形は異なるが、モノドロミーは全く同じ別の微分方程式」を次々と作ることができる。

つまり、「$\lambda$ を動かす」というのは、\textbf{「同じモノドロミーを持つ方程式の空間（モジュライ空間）の中を連続的に移動し、別の方程式へ乗り換えていくこと」}を意味する。$\lambda$ は、その空間内の特定の方程式を指し示す「座標（インデックス）」として機能しているのである。このように、モノドロミーを不変に保ったまま微分方程式を連続的に変形させる概念を\textbf{等モノドロミー変形 (isomonodromic deformation)} と呼ぶ。
\end{Q}

ここで、改めて冒頭で定義した
\[
K(n,m)=M(n,m)-E(n,m)
\]
という量が何を意味しているのかを総括する。

\begin{Q}[【核心】モノドロミーのパラメータ数 $M$ と、微分方程式のパラメータ数 $E$ の差である $K(n,m)$ は、本質的に何を表している数なのか？]
\textbf{A.} 結論から言えば、$K(n,m)$ は「任意のモノドロミーを実現するために、\textbf{最低限追加しなければならない『みかけの特異点』の個数}」を表している。

微分方程式の理論における究極の目標の一つは、「与えられたモノドロミー表現 $\hat{\rho}$ から、それを持つフックス型微分方程式を逆算して構成すること」である（リーマン・ヒルベルト問題）。
しかし、一般に $M(n,m) > E(n,m)$ である場合、微分方程式側の自由度（パラメータ）が圧倒的に足りず、望みのモノドロミーを実現する方程式が存在しない。

そこで、モノドロミーを一切変えずに微分方程式側の自由度だけを増やす「助っ人」として、みかけの特異点 $\Lambda=\{\lambda_{1},\cdots,\lambda_{g}\}$ を導入する。みかけの特異点を $g$ 個追加すると、微分方程式のパラメータ数は $E(n,m) + g$ となる。
モノドロミーの自由度と微分方程式の自由度が完全に一致し、対応がつく（モジュライ空間の次元が等しくなる）ための条件は、
\[
E(n,m) + g = M(n,m)
\]
すなわち、
\[
g = M(n,m) - E(n,m) = K(n,m)
\]
となる。
つまり、一般のフックス型微分方程式の分類・構成において、我々は最初に本来の確定特異点 $\Xi$ を決めた後、「次元の不足分 $K(n,m)$ を補うために、同数のみかけの特異点 $\Lambda$ を導入する」という手続きを必然的に踏むことになるのである。
\end{Q}

\begin{Q}[【発展】次元の不足を補うためにみかけの特異点を追加するならば、最低限の $K(n,m)$ 個よりも多く（$g > K(n,m)$ 個）追加しても、与えられたモノドロミーを実現することは可能か？また、その場合に何が起こるのか？]
\textbf{A.} 実現することは可能であるが、微分方程式側に「冗長な自由度」が発生し、モノドロミーと微分方程式の1対1の美しい対応が崩れてしまう。

これを「微分方程式の係数の空間（次元 $E+g$）」から「モノドロミー表現の空間（次元 $M$）」への写像として捉えると、みかけの特異点の個数 $g$ によって以下のような状況の違いが生じる。

\begin{itemize}
    \item \textbf{$g < K$ の場合（自由度不足）}: 
    $E+g < M$ となり、写像の像がターゲット全体を覆えない。つまり、実現不可能なモノドロミーが多数存在する。
    
    \item \textbf{$g = K$ の場合（ジャスト・フィット）}:
    $E+g = M$ となり、両空間の次元が一致する。一般のモノドロミーに対して、それを実現する微分方程式が（本質的に）ただ1つ、あるいは有限個に定まる。無駄のない「剛性」が保たれた理想的な状態である。
    
    \item \textbf{$g > K$ の場合（自由度過剰）}:
    $E+g > M$ となるため、任意のモノドロミーを実現することは可能である（全射的）。しかし、次元が $g-K$ だけ余るため、写像は単射にはならない。結果として、「全く同じモノドロミー表現を持つのに、係数（みかけの特異点の配置など）が連続的に異なる微分方程式の族（次元 $g-K$ のファイバー）」が無数に発生してしまう。
\end{itemize}

のちにパンルヴェ方程式の理論（岡本和夫のハミルトン理論）において、みかけの特異点の位置 $q$ と、その点での微分方程式の係数に関するある量 $p$ が、ハミルトン力学系における「正準変数（座標と運動量）」として扱われる。力学系として非退化な（決定論的に未来が1つに決まる）相空間を構成するためには、この過剰な冗長性を排除し、$g=K$ という「必要最低限のモジュライ空間」を考えることが不可欠となるのである。
\end{Q}
%#endregion
%#region
また、みかけの特異点 $x=\lambda$ における特性指数が全て整数 $J_{1}, \cdots, J_{n}$ であるとき、未知関数を
\begin{equation}\label{eq:gauge_transform_local}
z = (x-\lambda)^{s} y
\end{equation}
と変換（ゲージ変換）すると、新しい関数 $z$ が満たす微分方程式の $x=\lambda$ における特性指数は、一律に $s$ だけシフトして
\[
s+J_{1}, \cdots, s+J_{n}
\]
となる。元の解 $y$ は対数項を持たずローラン展開可能（非対数的特異点）であったため、それに $(x-\lambda)^{s}$ を掛けただけの $z$ も当然対数項を持たない。
このとき、$z$ の微分方程式について点 $x=\lambda$ の周りを1周する閉経路 $\gamma$ に沿った解析接続を考えると、モノドロミー行列はスカラー行列
\[
\Gamma(\gamma) = e^{2\pi is} I_{n}
\]
となる。逆に、ある特異点でのモノドロミー行列がこのスカラー行列の形で書けるならば、逆変換 $y = (x-\lambda)^{-s} z$ を施すことで特性指数を全て整数にでき、かつ対数項を持たない（非対数的である）ため、それはみかけの特異点に帰着される。

\begin{definition}[みかけの特異点の拡張]
解が多価関数となる場合（すなわち $s \notin \mathbb{Z}$）であっても、適当な複素数 $s$ によって \eqref{eq:gauge_transform_local} の変換を行い、特性指数を全て整数にでき、かつ非対数的となるような特異点のことを、広い意味で \textbf{みかけの特異点} と呼ぶ。
\end{definition}

さて、$x=\lambda$ がみかけの特異点であるならば、変換 \eqref{eq:gauge_transform_local} で適切に $s$ を選ぶことで、すべての特性指数を「非負の整数」に揃えることができる。このとき、微分方程式の任意の解は $x=\lambda$ の近傍で正則（テイラー展開可能）となる。

\begin{Q}[【核心】任意の解が正則（特異性を持たない）になるように変換できたとしても、微分方程式の係数にとっては $x=\lambda$ が依然として極（特異点）になり得るのはなぜか？分子と分母が打ち消し合って極が消えることはないのか？]
\textbf{A.} 解が正則であっても、解のロンスキアン $W(x)$ が $x=\lambda$ でゼロになれば、必ず方程式の係数に極が発生する。

$n$ 階線形微分方程式 $y^{(n)} + p_1(x) y^{(n-1)} + \dots + p_n(x) y = 0$ において、係数 $p_1(x)$ と基本解のロンスキアン $W(x)$ の間には、アーベルの公式（リウヴィルの公式）により
\[
p_1(x) = -\frac{W'(x)}{W(x)}
\]
という関係がある。
解がすべて正則であるとき、$W(x)$ も正則関数となる。もし $W(\lambda) = 0$ ならば、$x=\lambda$ は $W(x)$ の $m$ 位のゼロ点（$m \ge 1$）となり、$W(x) = c(x-\lambda)^m + \cdots$ と展開できる。これをアーベルの公式に代入すると、
\[
p_1(x) = -\frac{cm(x-\lambda)^{m-1} + \cdots}{c(x-\lambda)^m + \cdots} = -\frac{m}{x-\lambda} + (\text{正則な項})
\]
となり、$p_1(x)$ は必ず $x=\lambda$ において留数 $-m$ の1位の極を持つ。したがって、分子との打ち消し合いで極が消滅することは原理的にあり得ず、微分方程式にとっては確実に特異点となる。

これが特異点ではなく「完全な正則点（係数も極を持たない点）」に還元されるための必要十分条件は、$W(\lambda) \neq 0$ となることである。
みかけの特異点においてこれを満たすのは、特性指数がちょうど
\[
0, 1, 2, \cdots, n-1
\]
という「抜けのない連続した整数」になる場合のみである。このとき、適切に基底を取り直すことで、各解のテイラー展開の先頭項を $1, (x-\lambda), \dots, (x-\lambda)^{n-1}$ とできる。これらの $x=\lambda$ における微分値を並べたロンスキアン行列は、対角成分に $0!, 1!, \dots, (n-1)!$ が並ぶ非退化な三角行列となる。したがって、行列式 $W(\lambda)$ は対角成分の積となり $W(\lambda) \neq 0$ が保証され、係数の極が完全に消滅する。
\end{Q}
%#endregion

\subsection{連立微分方程式系と単独高階微分方程式}

%#region
これまでの単独の高階微分方程式から視点を変え、以下の連立1階線形微分方程式系について考えたい。
\begin{equation}\label{eq:system_ode}
    \frac{d\vec{y}}{dx} = A(x)\vec{y}
\end{equation}
ここで、$A(x)$ は有理関数を成分とする $n \times n$ 行列であり、$\vec{y}$ は $n$ 次元の未知ベクトル関数である。

\begin{definition}
    方程式 \eqref{eq:system_ode} の任意の解ベクトル $\vec{\varphi}(x)$ の各成分 $\varphi_{i}(x)$ が、$x=\xi$ において高々多項式オーダーの増大度しか持たない（すなわち確定特異点である）とき、方程式 \eqref{eq:system_ode} は $x=\xi$ に \textbf{確定特異点} をもつという。
\end{definition}

\begin{definition}
    方程式 \eqref{eq:system_ode} の特異点が全て確定特異点であるとき、\eqref{eq:system_ode} を \textbf{フックス型微分方程式系 (Fuchsian system)} という。
\end{definition}

無限遠点 $x=\infty$ も特異点として許容し、確定特異点であるとする。また、解はベクトル値関数であるが、解析接続によるモノドロミー行列は単独方程式のときと同様に定義できるため、リーマンデータ $(\mathbb{P}^{1}(\mathbb{C}), \Xi, \hat{\rho})$ を考えることができる。

いま、有限の特異点に対しては $u=x-\xi$、無限遠点に対しては $u=\frac{1}{x}$ という局所座標を導入し、方程式 \eqref{eq:system_ode} を
\begin{equation}\label{eq:system_ode_local}
    u\frac{d\vec{y}}{du} = B(u)\vec{y}
\end{equation}
と書き換える。このとき、以下の十分条件が成り立つ。

\begin{proposition}\label{prop:fuchs_condition_system}
    \eqref{eq:system_ode_local} において、行列 $B(u)$ が $u=0$ の近傍で正則（各成分が極を持たない）ならば、$u=0$ （すなわち元の $x=\xi$ または $\infty$）は \eqref{eq:system_ode} の確定特異点である。
\end{proposition}

\begin{proof}
    $u=0$ を頂点とする任意の角領域（セクター）内で原点に向かう経路を考え、$u = r e^{i\theta}$ （$r \to +0$, $\theta$ は固定）と極座標表示する。
    このとき、微分演算子は $u \frac{d}{du} = r \frac{d}{dr}$ となるため、方程式 \eqref{eq:system_ode_local} は
    \[
    r \frac{d\vec{y}}{dr} = B(re^{i\theta})\vec{y}
    \]
    と書ける。
    仮定より $B(u)$ は $u=0$ の近傍で正則であるため、その成分は連続関数であり、原点を含むある閉円板 $|u| \le r_0$ において有界となる。すなわち、ある正の定数 $M$ が存在して、行列の作用素ノルムについて $\|B(u)\| \le M$ が成り立つ。
    
    ベクトル $\vec{y}(r)$ のノルム $\|\vec{y}\|$ について、微分不等式を評価する。
    \begin{align*}
        r \frac{d}{dr} \|\vec{y}\| &\le \left\| r \frac{d\vec{y}}{dr} \right\| \\
        &= \| B(re^{i\theta})\vec{y} \| \\
        &\le \| B(re^{i\theta}) \| \|\vec{y}\| \\
        &\le M \|\vec{y}\|
    \end{align*}
    両辺を $r \|\vec{y}\|$ ($>0$) で割ると、
    \[
    \frac{1}{\|\vec{y}\|} \frac{d\|\vec{y}\|}{dr} \le \frac{M}{r}
    \]
    これをある点 $r_1$ から $r$ ($0 < r < r_1 \le r_0$) まで積分する。経路の向きに注意して $- \int_r^{r_1}$ として評価すると、
    \[
    \log \|\vec{y}(r_1)\| - \log \|\vec{y}(r)\| \le M (\log r_1 - \log r)
    \]
    移項して整理すると、
    \[
    \log \|\vec{y}(r)\| \ge \log \|\vec{y}(r_1)\| - M \log(r_1 / r) \implies \|\vec{y}(r)\| \ge C r^M
    \]
    となる。（※注：特異点へ向かう極限は $r \to 0$ であり、マイナス乗の評価が必要なため、原点から遠ざかる方向へ不等式を逆に解く。）
    正確には、$\frac{d}{dr}(\log \|\vec{y}\|) \ge - \frac{M}{r}$ （絶対値の微分の性質から下限も評価可能）より、
    \[
    \|\vec{y}(r)\| \le C' r^{-M} = C' |u|^{-M}
    \]
    という上からの評価が得られる（$C'$ は積分定数に依存する正の定数）。
    これにより、解ベクトル $\vec{y}(u)$ の各成分は、$u \to 0$ において高々 $u$ の多項式（マイナス乗）のオーダーでしか増大しないことが示された。これは確定特異点の定義そのものであり、証明が完了する。
\end{proof}

\begin{Q}[【核心】上記の証明において、解の増大度が $\|\vec{y}(u)\| \le C' |u|^{-M}$ で抑えられたことが、なぜ「確定特異点であること」の直接的な証明になっているのか？]
\textbf{A.} 確定特異点（regular singular point）の現代的な定義が、「特異点へ任意のセクター内から近づくとき、解の増大度が高々多項式オーダー（$|u|^{-N}$ の形）に制限されること」だからである。

もし $u=0$ が不確定特異点（真性特異点のような激しい振る舞いをする点）であれば、解の成分には $e^{1/u}$ のような指数関数的な発散項が含まれる。このような関数は $u \to 0$ で $|u|^{-N}$ といういかなる多項式の逆数をもってしても上から抑え込むことができない。
証明の中で $B(u)$ の正則性（有界性 $\|B\| \le M$）を用いたことで、指数の肩に乗る発散のスピードが対数積分 $\int \frac{M}{r} dr = M \log r$ に制限され、結果として解の増大度が多項式オーダーに落ち着くことが完全に保証されたのである。
\end{Q}

この命題は、「$A(x)$ が $x=\xi$ で1位の極を持つ（単純極である）ならば、確定特異点である」ということを意味している。

\begin{Q}[【核心】上記の命題の逆、すなわち「$x=\xi$ が確定特異点ならば、必ず $B(u)$ は正則になる（$A(x)$ は高々1位の極しか持たない）」は成立するか？成り立たないとすれば、どのような反例があるか？]
\textbf{A.} 逆は全く成り立たない。驚くべきことに、単独 $n$ 階微分方程式を連立系に書き直した「コンパニオン行列」自体が、完璧な反例となっている。

単独のフックス型微分方程式 \eqref{eq:linear-ode-n} を考える。
\[
\vec{y}=
\begin{pmatrix}
    y_{1} \\ y_{2} \\ \vdots \\ y_{n}
\end{pmatrix}
=
\begin{pmatrix}
    y \\ \frac{dy}{dx} \\ \vdots \\ \frac{d^{n-1}y}{dx^{n-1}}
\end{pmatrix}
\]
とおくと、方程式は \eqref{eq:system_ode} の形になり、その係数行列 $A(x)$ は
\[
A(x)=
\begin{pmatrix}
    0 & 1 & 0 &\cdots & 0 \\
    0 & 0 & 1 & \cdots & 0\\
    \vdots&\vdots&\vdots& \ddots &\vdots\\
    0 & 0 & 0 & \cdots & 1 \\
    -p_{n}(x) & -p_{n-1}(x) & -p_{n-2}(x) & \cdots & -p_{1}(x)
\end{pmatrix}
\]
となる。元の単独方程式において $x=\xi$ が確定特異点であるとき、フックスの定理（命題\ref{prop:fuchs_theorem_n}）より $p_{k}(x)$ は $k$ 位の極を持つ。したがって、行列 $A(x)$ の一番左下の成分 $-p_{n}(x)$ は「$n$ 位の極」を持つことになる。
解ベクトル $\vec{y}$ の成分は全て確定特異点の性質を満たしているにもかかわらず、行列 $A(x)$ は1位の極どころか $n$ 位の極を持っており、$u A(u)$ は全く正則にならないのである。
\end{Q}

\begin{Q}[【核心】単独 $n$ 階方程式と連立1階系では、理論を展開する上でどのような利点と欠点の違いがあるか？]
\textbf{A.} 最大の違いは「表現の一意性（ゲージ変換の自由度）」にある。

\textbf{単独方程式の利点}は、方程式の「表示が一意」に定まることである。解の空間が決まれば方程式の係数 $p_k(x)$ はロンスキアンを通じて一意に決まり、フックスの定理によって特異点における「極の位数」が厳密に統制されるため、外見から特異点の性質を判定しやすい。

一方、\textbf{連立系の利点}は、行列式や固有値といった線形代数の強力なツールを直接適用できることである。しかし欠点として、$\vec{z} = T(x)\vec{y}$ （$T(x)$ は有理関数行列）というゲージ変換（基底の取り替え）を行うと、方程式は $\vec{z}' = (T'T^{-1} + TAT^{-1})\vec{z}$ となり、解の確定特異点としての性質（モノドロミーなど）を全く変えないまま、\textbf{係数行列の極の位数をいくらでも人為的に上げ下げできてしまう}。
この「表示の非一意性」こそが、連立系において見かけの極の位数に騙されず、本質的な特異点の性質を見抜くための理論（Poincaré rankの簡約化など）を必要とする理由である。
\end{Q}
%#endregion
%#region
$x$ の有理関数を成分とする可逆な（行列式が恒等的に0でない）$n$ 次正方行列 $T(x)$ をとって、未知関数を
\begin{equation}\label{eq:gauge_transform_x}
    \vec{y} = T(x)\vec{z}
\end{equation}
と変換（ゲージ変換）すると、方程式 \eqref{eq:system_ode} は微分の連鎖律により
\begin{equation}
    \frac{d\vec{z}}{dx} = C(x)\vec{z}, \quad \left( C(x) = T(x)^{-1}A(x)T(x) - T(x)^{-1}\frac{dT(x)}{dx} \right)
\end{equation}
と書ける。このとき、「$x=\xi$ が確定特異点ならば、局所的に $T(x)$ をうまく取ることで、$C(x)$ が $x=\xi$ において1位の極（単純極）を持つようにできる」という事実が知られている。前述の命題と合わせると、局所的には確定特異点と1位の極は同値な概念として扱える。

これを踏まえて、$A_{l} \ (l=1,\cdots,m)$ を $x$ に依らない定数行列として、特異点における留数を明示した以下の微分方程式系を考える。
\begin{equation}\label{eq:schlesinger}
    \frac{d\vec{y}}{dx} = \sum_{l=1}^{m}\frac{A_{l}}{x-\xi_{l}}\vec{y}
\end{equation}
これは $\Xi=\{\xi_{1},\cdots,\xi_{m},\infty\}$ を特異点にもち、すべての有限特異点で係数行列が1位の極を持つため、十分条件（命題\ref{prop:fuchs_condition_system}）を満たす確定特異点となる。無限遠点 $x=\infty$ についても局所座標を用いて計算すると確定特異点であることがわかるため、これは大域的にもフックス型微分方程式系である。

\begin{definition}[シュレージンガー型微分方程式]
    \eqref{eq:schlesinger} のように、係数行列が「1位の極の和」として表されるフックス型微分方程式系を、\textbf{シュレージンガー型微分方程式 (Schlesinger system)} という。
\end{definition}

\begin{Q}[【核心】「局所的」にはゲージ変換で確定特異点を1位の極に変形できるが、では「大域的」に、任意のフックス型微分方程式系を1つのゲージ変換でシュレージンガー型 \eqref{eq:schlesinger} に変形することは可能なのか？]
\textbf{A.} 一般には不可能である。（ここにリーマン・ヒルベルト問題の罠がある）

ある1つの特異点 $x=\xi$ の近傍だけであれば、有理型関数を用いた局所的なゲージ変換によって、確定特異点を必ず1位の極に還元できる。
しかし、1つの全域的な有理関数行列 $T(x)$ を用いて、\textbf{「すべての特異点を同時に1位の極に落としつつ、空間に余計な特異点を一切増やさない」}という都合の良い大域的な変換ができるとは限らない。

もしこれが常に可能であれば、任意のモノドロミー表現から必ずシュレージンガー型方程式が作れることになり、リーマン・ヒルベルト問題は無条件で肯定的に解けることになってしまう。長年そのように信じられていたが、1989年にボリブルフ (Bolibruch) が「シュレージンガー型の系では絶対に実現できないモノドロミー表現が存在する」という反例を構成した。
つまり、一般のフックス型系を無理やりシュレージンガー型の美しい形に書き換えようとすると、大域的な次元の辻褄を合わせるために、どうしても新たな「みかけの特異点」を系に発生させざるを得なくなるのである。
\end{Q}

シュレージンガー型微分方程式 \eqref{eq:schlesinger} に対して、$x$ に依存しない定数行列を用いた大域的なゲージ変換 $T(x) := G \in GL(n, \mathbb{C})$ を考えると、各特異点の留数行列は
\[
A_{l} \to G^{-1}A_{l}G
\]
と変換される。この定数ゲージ変換の自由度を用いることで、どれか1つの留数行列を標準形（ジョルダン標準形や対角化など）に固定することができる。
ここで、無限遠点 $x=\infty$ における留数行列 $A_{\infty}$ について考える。コンパクトなリーマン面（ここではリーマン球面 $\mathbb{P}^{1}(\mathbb{C})$）上の有理型微分形式の留数の総和は必ず $0$ になるという留数定理（あるいは $u=1/x$ による局所座標での直接計算）より、
\[
A_{\infty} = -\sum_{l=1}^{m}A_{l}
\]
が成り立つ。我々は慣例として、この $A_{\infty}$ をゲージ変換 $G$ によって標準形に固定することにする。

では、このシュレージンガー型方程式に含まれる「本質的な独立パラメータの数」を数え上げる。

\begin{Q}[【核心】シュレージンガー型微分方程式 \eqref{eq:schlesinger} が持つ本質的なパラメータの総数はいくつか？また、その計算の根拠は何か？]
\textbf{A.} パラメータの総数はちょうど $(m-1)n^{2}+1$ 個である。

計算の根拠は以下の通りである。
\begin{itemize}
    \item \textbf{元のパラメータ数}: 方程式を決定するのは $m$ 個の $n \times n$ 行列 $A_{1}, \cdots, A_{m}$ であるため、最初は $mn^{2}$ 個の自由度がある。（$A_{\infty}$ はこれらの和で決まるため独立ではない）
    \item \textbf{ゲージ変換による削減}: 定数行列 $G$ による共役変換 $G^{-1}(\cdot)G$ で標準形に固定できるため、$G$ の持つ自由度だけパラメータが減る。$G$ は $n \times n$ の可逆行列なので基本的には $n^{2}$ 個の自由度を持つ。
    \item \textbf{スカラー行列の除外}: しかし、$G = cI$（$c$ は非ゼロの定数、$I$ は単位行列）という「スカラー行列」を用いた場合、$c^{-1}I A_l cI = A_l$ となり、行列を全く変化させない。つまり、変換として有効に働く自由度はスカラー行列の分（1次元）を除いた $n^{2}-1$ 個である（射影一般線形群 $PGL(n, \mathbb{C})$ の次元）。
\end{itemize}
したがって、本質的なパラメータ数は $mn^{2}$ から有効なゲージ自由度 $n^{2}-1$ を引いて、
\[
mn^{2} - (n^{2} - 1) = (m-1)n^{2} + 1
\]
となる。
\end{Q}

\begin{Q}[【核心】このパラメータ数 $(m-1)n^{2}+1$ は、これまでの議論とどう結びつくのか？]
\textbf{A.} 驚くべきことに、この数は $\Xi=\{\xi_{1},\cdots,\xi_{m},\infty\}$ を特異点にもつ $n$ 階フックス型微分方程式の「モノドロミー表現のモジュライ空間の次元 $M(n,m)$」と完全に一致する。

単独のフックス型微分方程式では、パラメータ数 $E(n,m)$ が $M(n,m)$ に全く届かず、次元の不足を補うために「みかけの特異点」を多数導入する必要があった。
しかし、連立系である「シュレージンガー型微分方程式」を考えた瞬間、みかけの特異点を一切追加することなく、最初からモノドロミーと同じ自由度 $M(n,m)$ を獲得できている。この見事な次元の一致こそが、シュレージンガー系がリーマン・ヒルベルト問題（与えられたモノドロミーを実現する方程式の構成）を解くための「最も自然で強力な枠組み」として期待され、研究されてきた最大の理由である。
\end{Q}
%#endregion
%#region
\begin{definition}[可約な系]
    \eqref{eq:system_ode} の係数行列 $A(x)$ が、ある整数 $r$ （$1 \le r \le n-1$）に対して、$r$ 次正方行列 $A_{11}(x)$ と $n-r$ 次正方行列 $A_{22}(x)$ を用いて
    \[
    A(x)=
    \begin{pmatrix}
    A_{11}(x) & 0\\
    A_{21}(x) & A_{22}(x)
    \end{pmatrix}
    \]
    というブロック下三角行列の形で表せるとき、この微分方程式系 \eqref{eq:system_ode} は \textbf{可約 (reducible)} であるという。
\end{definition}

\begin{Q}[【核心】係数行列が可約（ブロック下三角）であるとき、微分方程式を解く上でどのような具体的なメリットがあるのか？]
\textbf{A.} $n$ 次元の巨大な連立方程式を、「$n-1$ 次元以下のより小さな部分系のカスケード（数珠つなぎ）解法」に完全に帰着できることである。

未知ベクトルを $\vec{y} = \begin{pmatrix} \vec{y}_{1} \\ \vec{y}_{2} \end{pmatrix}$ と分割して代入すると、系は以下の2つのステップに分解される。
\begin{enumerate}
    \item $\frac{d\vec{y}_{1}}{dx} = A_{11}(x)\vec{y}_{1}$
    \item $\frac{d\vec{y}_{2}}{dx} = A_{22}(x)\vec{y}_{2} + A_{21}(x)\vec{y}_{1}$
\end{enumerate}
まず、ステップ1の $r$ 次元斉次方程式を解いて $\vec{y}_{1}$ を完全に決定する。次に、その結果をステップ2の第2項（非斉次項）に代入し、定数変化法などを用いて $n-r$ 次元の非斉次方程式を解く。
いずれのステップも最大で $\max(r, n-r) \le n-1$ 次元の系を解けばよくなるため、方程式の難易度が劇的に下がるのである。
\end{Q}

\begin{Q}[【核心】単独 $n$ 階微分方程式を連立系に直した「コンパニオン行列（同伴行列）」は、そのままの形では絶対に可約にならないのはなぜか？]
\textbf{A.} コンパニオン行列の「一つ上の対角成分（超対角成分）」がすべて $1$ で埋まっているからである。

コンパニオン行列をどのように $r$ 行と $n-r$ 行でブロック分割しようとしても、右上ブロック（$A_{12}(x)$ に相当する部分）の境界付近には必ず「超対角成分の $1$」が少なくとも1つ紛れ込んでしまう。したがって、右上ブロックを完全に零行列 $0$ にすることは物理的に不可能である。
これは、「基底の取り方（表現）を工夫しない限り、単独方程式の自然な連立化は、自動的には小さな系に分解してくれない」という単独系の強固な構造を視覚的に表している。（※ただし、元の単独微分方程式がより低階の微分作用素に因数分解できるような特殊なケースに限り、適切なゲージ変換 $T(x)$ を挟むことで、系全体を可約な形に変換することは可能である。）
\end{Q}
%#endregion
%#region
\begin{definition}[単独高階化]
    連立微分方程式系を、特定の成分（例えば $y_1$）について $x$ で繰り返し微分し、代入と消去を繰り返すことで、元の成分を未知関数とする1つの $n$ 階単独微分方程式に帰着させる操作を \textbf{単独高階化} という。（※代数的には「巡回ベクトル (cyclic vector)」をとる操作に相当する。）
\end{definition}

既約（簡約不能）な $n$ 次連立微分方程式系に対して単独高階化を行うと、同値な $n$ 階単独微分方程式が得られ、後者に対してこれまでの単独方程式の議論を適用することができる。
また、元の連立系がフックス型（すべての特異点が確定特異点）であれば、得られた単独方程式もフックス型になる。なぜなら、「解が高々多項式オーダーの増大度しか持たない」という確定特異点の解析的な性質は、成分の足し算や微分といった線形代数的な操作を行っても破壊されないからである。

例えば、特異点 $\Xi=\{\xi_{1},\cdots,\xi_{m},\infty\}$ を持つシュレージンガー型方程式 \eqref{eq:schlesinger} について単独高階化を行ったとする。元の本来の確定特異点の位置は当然変わらない。しかしここで、方程式に含まれる「パラメータの数」について重大なパラドックスが生じる。

前述の通り、本来の特異点のみを持つ単独系のパラメータ数は
\[
E(n,m) = \frac{1}{2}n \big( m(n+1) - (n-1) \big)
\]
であった。一方で、シュレージンガー型微分方程式（連立系）のパラメータ数は、モノドロミー表現の次元と等しい
\[
M(n,m) = (m-1)n^{2} + 1
\]
であった。両者は全く等価な微分方程式であるはずなのに、パラメータの数が一致しない。

\begin{Q}[【核心】単独高階化を行った結果、$K(n,m) = M(n,m) - E(n,m)$ 個分のパラメータが空中に消え去ってしまったように見える。この「消えたパラメータ」はどこへ行ったのか？]
\textbf{A.} 消えたわけではない。連立系から単独系へ「翻訳」する過程で生じる歪みによって、ちょうど $K(n,m)$ 個の\textbf{「みかけの特異点」として微分方程式上に突如として出現した}のである。

連立系からある1つの変数 $y_1$ についての単独方程式を導出する際、$y_1$ の1階微分、2階微分、$\dots$ と計算を進め、それらを独立な基底として元のベクトル $\vec{y}$ を逆算して消去する作業を行う。
この「逆算」は、代数的にはある変換行列 $P(x)$ の逆行列 $P(x)^{-1}$ を掛ける操作（あるいはロンスキアンのような量での割り算）に他ならない。この変換行列 $P(x)$ は元の連立系の係数（有理関数）から作られるため、$P(x)$ の行列式が $0$ になる点 $x=\lambda_j$ では割り算ができず、得られた単独方程式の係数に「新たな極」が発生してしまう。

しかし、この極は「方程式の表示を変えたこと」に起因するだけの人工的な特異点であるため、解のモノドロミー（分岐）には一切影響を与えない。すなわち、これこそが「みかけの特異点」の正体である。
一般の位置にある巡回ベクトルを選んで単独高階化を行った場合、この割り算によって生じるみかけの特異点の個数は、驚くべきことに次元の不足分 $K(n,m)$ 個と完全に一致する。
連立系における「係数行列の成分の自由度」が、単独系に変換された瞬間に「みかけの特異点の位置 $\lambda_1, \dots, \lambda_K$ という空間座標」に姿を変えて、モノドロミーの自由度を裏から支えているのである。
\end{Q}

以下、$n=2$（2階）の微分方程式系について詳しく見てみる。係数行列と未知関数を
\[
A(x)=
\begin{pmatrix}
    a_{11}(x) & a_{12}(x)\\
    a_{21}(x) & a_{22}(x)\\
\end{pmatrix},
\quad 
\vec{y}=
\begin{pmatrix}
    y_{1}\\
    y_{2}\\
\end{pmatrix}
\]
とすると、成分ごとの方程式は
\[
\begin{cases}
\frac{dy_{1}}{dx}=a_{11}y_{1}+a_{12}y_{2}\\
\frac{dy_{2}}{dx}=a_{21}y_{1}+a_{22}y_{2}
\end{cases}
\]
となる。第1式から $y_2$ を $y_1$ で表し（$y_2 = \frac{1}{a_{12}}(\frac{dy_{1}}{dx}-a_{11}y_{1})$）、これを $x$ で微分して第2式に代入することで、単独高階化を行う。
\begin{align*}
\frac{d^{2}y_{1}}{dx^{2}} &= a_{11}\frac{dy_{1}}{dx} + \frac{da_{11}}{dx}y_{1} + a_{12}\frac{dy_{2}}{dx} + \frac{da_{12}}{dx}y_{2}\\
&= a_{11}\frac{dy_{1}}{dx} + \frac{da_{11}}{dx}y_{1} + a_{12}(a_{21}y_{1}+a_{22}y_{2}) + \frac{da_{12}}{dx}\frac{1}{a_{12}}\left(\frac{dy_{1}}{dx}-a_{11}y_{1}\right)\\
&= a_{11}\frac{dy_{1}}{dx} + \frac{da_{11}}{dx}y_{1} + a_{12}a_{21}y_{1} + a_{22}\left(\frac{dy_{1}}{dx}-a_{11}y_{1}\right) + \frac{da_{12}}{dx}\frac{1}{a_{12}}\left(\frac{dy_{1}}{dx}-a_{11}y_{1}\right)\\
&= \left(a_{11} + a_{22} + \frac{1}{a_{12}}\frac{da_{12}}{dx}\right)\frac{dy_{1}}{dx} + \left(a_{12}a_{21} - a_{11}a_{22} + \frac{da_{11}}{dx} - \frac{a_{11}}{a_{12}}\frac{da_{12}}{dx}\right)y_{1}\\
&= \left(\text{Tr} A(x) + \frac{d}{dx}\log a_{12}\right)\frac{dy_{1}}{dx} + \left(-\det A(x) + \frac{da_{11}}{dx} - a_{11}\frac{d}{dx}\log a_{12}\right)y_{1}
\end{align*}
これを標準的な単独2階線形微分方程式の形 $\frac{d^{2}y_{1}}{dx^{2}} + p(x)\frac{dy_{1}}{dx} + q(x)y_{1} = 0$ に合わせるため、以下のように置く。
\begin{align*}
p(x) &= -\left(\text{Tr} A(x) + \frac{d}{dx}\log a_{12}\right)\\
q(x) &= \det A(x) - \frac{da_{11}}{dx} + a_{11}\frac{d}{dx}\log a_{12}
\end{align*}

ここで、特異点の数が $m=3$ （無限遠点を含めて4つ）であるシュレージンガー型方程式
\begin{equation}\label{eq:schlesinger_m3}
\frac{d\vec{y}}{dx} = A(x)\vec{y}, \quad A(x) = \frac{A_{0}}{x} + \frac{A_{1}}{x-1} + \frac{A_{t}}{x-t}
\end{equation}
を考察する。本来の確定特異点は一次分数変換の自由度を用いて $\Xi=\{0, 1, t, \infty\}$ と固定してよい。（係数行列 $A_0, A_1, A_t$ が零行列に退化しない限り、これらは特異点として残る。）
さらに、無限遠点における留数行列 $A_{\infty} = -(A_{0} + A_{1} + A_{t})$ は、事前のゲージ変換によって対角化され、
\[
A_{\infty} =
\begin{pmatrix}
s_{\infty} & 0\\
0 & s_{\infty}+\kappa_{\infty}
\end{pmatrix}
\]
となっているものと仮定する。

単独方程式の係数 $p(x), q(x)$ には対数微分 $\frac{d}{dx}\log a_{12} = \frac{a_{12}'}{a_{12}}$ が含まれているため、元の連立系の特異点 $\{0, 1, t, \infty\}$ に加えて、\textbf{「$a_{12}(x)$ の零点」が単独方程式の新たな特異点として出現する。}
$A(x)$ を通分すると $A(x) = \frac{\tilde{A}(x)}{x(x-1)(x-t)}$ となり、分子 $\tilde{A}(x)$ は
\begin{align*}
\tilde{A}(x) &= A_{0}(x-1)(x-t) + A_{1}x(x-t) + A_{t}x(x-1)\\
&= (A_{0}+A_{1}+A_{t})x^{2} - \big((t+1)A_{0}+tA_{1}+A_{t}\big)x + tA_{0}
\end{align*}
となる。(1,2)成分に注目すると、$A_{\infty}$ が対角行列であるため、$(A_{0}+A_{1}+A_{t})_{12} = -(A_{\infty})_{12} = 0$ となる。これにより $x^{2}$ の項が消滅し、分子は $x$ の1次式になる。定数 $X$ と $\lambda$ を用いて因数分解すると、
\[
a_{12}(x) = X\frac{x-\lambda}{x(x-1)(x-t)}
\]
となる。よって、$x=\lambda$ が新しく現れた特異点（$a_{12}$ の1位の零点）である。

\begin{Q}[【核心】新しく現れた特異点 $x=\lambda$ は、本当に「みかけの特異点」の厳しい条件（特性指数が整数差であり、かつ対数項を持たない）をクリアしているのか？]
\textbf{A.} 完全にクリアしている。

まず特性指数を調べる。$a_{12}(x)$ は $x=\lambda$ で1位の零点を持つため、対数微分 $\frac{d}{dx}\log a_{12}$ は $x=\lambda$ に留数 $1$ の極を持つ。したがって、
\begin{align*}
(x-\lambda)p(x)\big|_{x=\lambda} &= -1\\
(x-\lambda)^{2}q(x)\big|_{x=\lambda} &= 0 \quad (\text{$q(x)$ の極は高々1位であるため})
\end{align*}
となる。決定方程式（indicial equation）は $s(s-1) + s \cdot (-1) + 0 = s(s-2) = 0$ となり、特性指数は $s=0, 2$ となる。すべて整数であり、その差は2である。

次に非対数的であること（対数項の非存在）の証明だが、これは計算するまでもなく明らかである。元の連立方程式 \eqref{eq:schlesinger_m3} の特異点は $\{0, 1, t, \infty\}$ のみであり、$x=\lambda$ において係数行列 $A(x)$ は完全に正則である。微分方程式の基本定理より、正則点での解ベクトル $\vec{y}(x)$ は当然正則である。
単独高階化は単なる式変形であるため、得られた単独方程式の解 $y_1(x)$ は元の解の第1成分に他ならず、$x=\lambda$ でテイラー展開可能（正則）でなければならない。正則関数に対数項 $\log(x-\lambda)$ が入り込む余地はないため、$x=\lambda$ は決定方程式の特性指数の差にかかわらず、確実に非対数的特異点となる。
\end{Q}

\begin{proposition}
    シュレージンガー型方程式 \eqref{eq:schlesinger_m3}（$m=3, n=2$）の単独高階化は、本来の特異点 $\Xi=\{0,1,t,\infty\}$ の他に、ただ1つの点 $x=\lambda$ をみかけの特異点として持つ。またその特性指数は $0, 2$ である。
\end{proposition}

この結果をパラメータ（モジュライ空間の次元）の観点から振り返る。
\begin{align*}
E(2,3) &= \frac{1}{2} \cdot 2 \cdot (3(2+1)-(2-1)) = 8\\
M(2,3) &= (3-1) \cdot 2^{2} + 1 = 9\\
K(2,3) &= M(2,3) - E(2,3) = 9 - 8 = 1
\end{align*}
不足していたパラメータの数 $g = K(2,3) = 1$ が、まさに単独高階化によって生じた「みかけの特異点の位置 $\lambda$」の数 1 と完全に符号している。
%#endregion


\subsection{リーマンの問題}

%#region
これまで、フックス型微分方程式に対して、確定特異点の集合 $\Xi$ とそのモノドロミー表現 $\hat{\rho}$ をまとめた組 $(\mathbb{P}^{1}(\mathbb{C}), \Xi, \hat{\rho})$ をリーマンデータと呼んできた。ここでは、これとは「逆」の対応を問う以下の問題を考える。

\begin{itemize}
    \item[] \textbf{問題} \\
    任意の3つの組 $(\mathbb{P}^{1}(\mathbb{C}), \Xi, \hat{\rho})$ が与えられたとき、これをリーマンデータとして持つようなフックス型微分方程式は存在するだろうか？
\end{itemize}

\begin{definition}[リーマン・ヒルベルトの問題]
    上の「与えられたモノドロミー表現を実現する微分方程式の構成問題」を、\textbf{リーマンの問題} または \textbf{リーマン・ヒルベルトの問題 (Riemann-Hilbert problem)} という。ダフィット・ヒルベルトが1900年に提唱した「ヒルベルトの23の問題」における第21問題としても知られる。
\end{definition}

この問題を幾何学的に捉え直すために、パラメータ空間（モジュライ空間）を考える。
（特異点の位置 $\Xi$ を固定した）フックス型微分方程式の係数がなす空間を $\mathcal{E}$、モノドロミー表現の同値類がなす空間を $\mathcal{M}$ とする。
微分方程式が一つ与えられれば、そこから（基底の取り替えの自由度を除いて）モノドロミー表現が一意に定まるため、空間 $\mathcal{E}$ から空間 $\mathcal{M}$ への自然な写像 $\gamma$ を定義できる。

\begin{definition}[モノドロミー写像]
    この、微分方程式からモノドロミー表現を対応させる写像 $\gamma : \mathcal{E} \to \mathcal{M}$ を、\textbf{モノドロミー写像 (monodromy map)} という。
\end{definition}

\begin{Q}[【核心】リーマン・ヒルベルトの問題を、このモノドロミー写像 $\gamma$ の言葉で言い換えるとどうなるか？また、これまでのパラメータ数 $E(n,m)$ や $M(n,m)$ の議論は、この問題に対してどのような結論をもたらすか？]
\textbf{A.} リーマン・ヒルベルトの問題は、\textbf{「モノドロミー写像 $\gamma : \mathcal{E} \to \mathcal{M}$ は全射（onto）であるか？」}という写像の性質を問う問題に完全に翻訳される。

もし $\gamma$ が全射であれば、任意のモノドロミー表現（$\mathcal{M}$ の元）に対して、それを写像先とする微分方程式（$\mathcal{E}$ の元）が必ず存在することになる。
しかし、これまでの議論で見たように、両空間の次元（パラメータの数）を比較すると、一般に
\[
\dim \mathcal{E} = E(n,m) < M(n,m) = \dim \mathcal{M}
\]
となる。定義域 $\mathcal{E}$ の次元が値域 $\mathcal{M}$ の次元よりも真に小さいため、写像 $\gamma$ が全射になることはあり得ない。

つまり、「本来の確定特異点 $\Xi$ のみを持つ微分方程式の空間 $\mathcal{E}$ だけでは、リーマン・ヒルベルトの問題は一般には（否定的に）解けない」という強烈な事実が、単純な次元の比較からエレガントに証明されるのである。そして、この「次元の不足分」を補い、写像を全射にして問題を肯定的に解決するための代償として導入されたのが、「みかけの特異点」に他ならない。
\end{Q}

ここで、シュレージンガー方程式系の空間 $\mathcal{E}_{Sch}$ を考えれば、その次元は
\[
\dim \mathcal{E}_{Sch} = \dim \mathcal{M} = (m-1)n^{2}+1
\]
となり、モノドロミー表現のモジュライ空間の次元と完全に一致するため、リーマン・ヒルベルトの問題を考えるに足る十分なパラメータを持っているように見える。
しかし、シュレージンガー方程式系の空間からモノドロミー表現の空間へのモノドロミー写像 $\gamma : \mathcal{E}_{Sch} \to \mathcal{M}$ は、一般には全射（surjective）にならないことが知られている（ボリブルフの反例）。

この「次元が一致しているにもかかわらず、連立系では実現できないモノドロミーが存在する」という困難を解決するため、再び単独高階微分方程式の視点に立ち戻り、十分な数のみかけの特異点
\[
g \ge K(n,m) = M(n,m) - E(n,m)
\]
をあえて追加した微分方程式の空間 $\mathcal{E}_g$ を定義し、この空間上での構成問題を考える。

\begin{Q}[【核心】シュレージンガー型（連立系）で次元が一致しても失敗したのに、なぜ「みかけの特異点を追加した単独方程式」に望みを託すのか？]
\textbf{A.} みかけの特異点の位置 $\{\lambda_j\}$ が、連立系では到達できなかったモノドロミー表現の空間の「欠落した領域」を補完する新たな幾何学的自由度として機能するからである。

ボリブルフの反例は、特定のモノドロミー表現に対して、それを実現する「1位の極のみを持つ連立系（シュレージンガー型）」が存在しないことを示した。
しかし、高階単独微分方程式においてみかけの特異点 $g$ 個を許容すれば、方程式の自由度は $E(n,m) + g$ となり、少なくとも $g = K(n,m)$ と取れば次元の不足は解消される。単独高階化の議論で見た通り、連立系の「係数パラメータ」を単独系の「特異点の座標パラメータ」へ変換して扱うこの枠組みは、モノドロミー表現を網羅する（全射性を確保する）上で、連立系よりも柔軟な構造を提供のである。
\end{Q}

\begin{Q}[【核心】$g$ を $K(n,m)$ よりも大きく取る（$g > K$）ことには、どのような数学的メリットがあるのか？]
\textbf{A.} リーマン・ヒルベルト問題を「確実に」肯定的に解決するためである。
$g = K(n,m)$ はあくまで次元を一致させるための「必要最低限」の数であるが、モノドロミー写像がその次元において常に全射である保証はない。しかし、みかけの特異点の数 $g$ を十分に大きく取れば、モノドロミー表現を構成する際の代数的な制約をより多くの自由度で吸収できるようになり、任意のリーマンデータに対してそれを実現する微分方程式の存在を保証しやすくなる。
\end{Q}
%#endregion
%#region
微分方程式の具体的な形からは一旦離れ、純粋に幾何学的・代数的なデータのみを抽出する。
$\Xi=\{\xi_{1},\cdots,\xi_{m},\xi_{m+1}\}$ とし、$\mathbb{X}=\mathbb{P}^{1}(\mathbb{C})\setminus \Xi$ とする。基本群 $\pi_1(\mathbb{X})$ から一般線形群への準同型（表現）
\[
\rho:\pi_{1}(\mathbb{X})\to GL(n,\mathbb{C})
\]
が与えられたとき、適当な正則行列 $P$ による相似変換 $\rho'(\gamma) = P^{-1}\rho(\gamma)P$ で移り合う表現の同値類 $\hat{\rho}$ が定義できる。

\begin{definition}[リーマンデータの純粋化]
    微分方程式の存在とは独立に、幾何学的に定まる組 $(\mathbb{P}^{1}(\mathbb{C}),\Xi,\hat{\rho})$ のことを \textbf{リーマンデータ} と呼ぶ。
\end{definition}

さて、この表現 $\rho$ は、ベクトル空間 $\mathbb{C}^n$ に対する群 $\pi_1(\mathbb{X})$ の作用とみなすことができる。ここで、微分方程式の解の解析接続の性質から、作用の「向き」について極めて重要な注意が必要となる。

\begin{Q}[【核心】表現 $\rho$ による作用を考えるとき、$\mathbb{C}^n$ は $\pi_1(\mathbb{X})$ の「左加群」と「右加群」のどちらとして扱うのが自然か？]
\textbf{A.} 微分方程式のモノドロミーの自然な構造に従うならば、\textbf{「右 $\pi_1(\mathbb{X})$-加群」}として扱うのが最も理にかなっている。

この理由は、「経路の積の順序」と「行列の積の順序」の衝突にある。
基本群における経路の積 $\gamma_1 \circ \gamma_2$ を「$\gamma_1$ を辿った後に $\gamma_2$ を辿る」と定義する。
微分方程式の解の基底を並べた基本行列 $Y(x)$ を解析接続すると、モノドロミー行列 $M_\gamma$ は $Y(x)$ の「右から」掛かる。
したがって、$\gamma_1$ の後に $\gamma_2$ を辿った場合、基本行列の変換は
\[
Y \xrightarrow{\gamma_1} Y M_{\gamma_1} \xrightarrow{\gamma_2} (Y M_{\gamma_2}) M_{\gamma_1} = Y (M_{\gamma_2} M_{\gamma_1})
\]
となり、表現行列の積の順序が逆転し、$M_{\gamma_1 \circ \gamma_2} = M_{\gamma_2} M_{\gamma_1}$ となる（すなわち表現が反準同型になる）。

この順序のねじれを解消するためには、ベクトルを列ベクトルではなく「行ベクトル $\vec{v}^T$」とし、行列を右から作用させる
\[
\vec{v}^T \cdot \gamma := \vec{v}^T \rho(\gamma)
\]
という定義を採用するのが自然である。このとき、$(\vec{v}^T \cdot \gamma_1) \cdot \gamma_2 = (\vec{v}^T \rho(\gamma_1)) \rho(\gamma_2) = \vec{v}^T (\rho(\gamma_1)\rho(\gamma_2))$ となり、群の積と作用の順序が完璧に整合する。ゆえに、$\mathbb{C}^n$ は「非可換な群環 $\mathbb{C}[\pi_1(\mathbb{X})]$ 上の右加群（ベクトル空間の非可換版）」とみなすことができるのである。

（※注：文献によっては $\mathbb{C}^n$ を無理やり「左加群」として扱うものもあるが、その場合は「経路の積 $\gamma_1 \circ \gamma_2$ を関数の合成のように右から順に辿るように定義し直す」か、あるいは「表現 $\rho(\gamma)$ をモノドロミー行列の逆行列 $M_\gamma^{-1}$ で定義する」といった、規約（コンベンション）の意図的な書き換えが裏で行われていることに注意が必要である。）
\end{Q}

\begin{definition}[表現の既約性]
    この右 $\pi_{1}(\mathbb{X})$-加群 $\mathbb{C}^{n}$ において、任意の $\gamma \in \pi_{1}(\mathbb{X})$ の作用に対して不変となる（すなわち $W \rho(\gamma) \subseteq W$ となる）真の部分空間 $W \subsetneq \mathbb{C}^{n}$ が $\{0\}$ に限られるとき、表現 $\rho$ は \textbf{既約 (irreducible)} であるという。
    このとき、その同値類 $\hat{\rho}$ も \textbf{既約} であるという。既約でない表現（$\{0\}$ 以外の不変部分空間を持つ表現）を一般に \textbf{可約 (reducible)} という。
\end{definition}

\begin{Q}[【核心】表現 $\rho$ が可約であるとき、その具体的な行列表現はどうなるか？]
\textbf{A.} 適切な基底を取ることで、すべての $\rho(\gamma)$ が共通の\textbf{ブロック行列（行ベクトルへの右作用の場合はブロック下三角行列）}になる。

表現 $\rho$ が可約であるとする。定義より、$0$ ではない $n'$ 次元（$0 < n' < n$）の不変部分空間 $W$ が存在する。$W$ の基底を最初に取り、それにベクトルを付け加えて $\mathbb{C}^{n}$ 全体の基底とする。
行ベクトルに対する右作用 $\vec{v}^T \mapsto \vec{v}^T \rho(\gamma)$ において $W$ が閉じているためには、行列 $\rho(\gamma)$ は
\[
\rho(\gamma_{l})=
\begin{pmatrix}
A_{l} & 0\\
B_{l} & C_{l}
\end{pmatrix}
\]
というブロック下三角の形にならなければならない。（ここで、$A_{l}$ は $n'$ 次正方行列、$C_{l}$ は $(n-n')$ 次正方行列）。
このとき、左上のブロック $\rho'(\gamma_{l}):=A_{l}$ を取り出すと、これも準同型の性質を満たすため、$\pi_{1}(\mathbb{X})$ の $\mathbb{C}^{n'}$ 上の新たな表現（部分表現）となる。
\end{Q}

%#endregion
%#region
これまでの議論を踏まえ、リーマンデータ $(\mathbb{P}^{1}(\mathbb{C}),\Xi,\hat{\rho})$ が与えられているとき、ここに新たな特異点の集合 $\Lambda=\{\lambda_{1},\cdots,\lambda_{g}\}$ を追加した新しいリーマンデータ
\[
(\mathbb{P}^{1}(\mathbb{C}),\Xi\cup\Lambda,\hat{\rho'})
\]
について考えたい。
ここで、新しい表現類 $\hat{\rho'}$ は以下の条件を満たす代表元 $\rho'$ を持つとする。
元の確定特異点 $\xi_{1},\cdots,\xi_{m+1}$ を反時計回りに1周する道を $\gamma_{1},\cdots,\gamma_{m+1}$、新しく追加した特異点 $\lambda_{1},\cdots,\lambda_{g}$ を1周する道を $\gamma'_{1},\cdots,\gamma'_{g}$ とする。このとき、元の表現 $\rho \in \hat{\rho}$ が存在して、任意の $i, j$ に対して
\begin{align*}
\rho'(\gamma_{i}) &= \rho(\gamma_{i}) \\
\rho'(\gamma'_{j}) &= I_{n} \quad (\text{単位行列})
\end{align*}
を満たすものとする。

\begin{Q}[【核心】新しく追加した点に対する表現を $\rho'(\gamma'_{j}) = I_{n}$ と設定することは、代数的にどのような意味を持つか？また、これは元の表現と矛盾しないか？]
\textbf{A.} これは、解析的な「みかけの特異点」を、代数的な「表現の核（自明な作用）」として完璧に翻訳した設定であり、数学的に一切の矛盾を生じない。

解析において、みかけの特異点とは「その周りを回っても解が分岐しない（元の解にそのまま戻る）」点であった。表現論の言葉では、ベクトルに右から作用させたとき $\vec{v}^T \rho'(\gamma'_j) = \vec{v}^T$ となること、すなわち作用する行列が単位行列 $I_n$ であることと完全に同値である。

また、穴の開いた球面（$\mathbb{P}^{1}(\mathbb{C})\setminus (\Xi\cup\Lambda)$）の基本群は、すべての特異点を1周するループの積が自明になるという関係式
\[
\gamma_{1} \cdots \gamma_{m+1} \gamma'_{1} \cdots \gamma'_{g} = 1
\]
を持つ。これに新しい表現 $\rho'$ を適用すると、
\[
\rho'(\gamma_{1}) \cdots \rho'(\gamma_{m+1}) \rho'(\gamma'_{1}) \cdots \rho'(\gamma'_{g}) = \rho(\gamma_{1}) \cdots \rho(\gamma_{m+1}) I_{n} \cdots I_{n} = I_{n}
\]
となり、元の表現 $\rho$ が持っていた関係式 $\prod \rho(\gamma_i) = I_n$ に完全に帰着する。したがって、$\rho'(\gamma'_j) = I_n$ と置くことは、元のモノドロミー表現の代数的な構造を一切破壊することなく、空間の穴（特異点）だけを増やす極めて自然な操作なのである。
\end{Q}

\begin{definition}[リーマンデータの拡大]
    上記の条件を満たすような新しいデータ $(\mathbb{P}^{1}(\mathbb{C}),\Xi\cup\Lambda,\hat{\rho'})$ が存在するとき、これを元のリーマンデータ $(\mathbb{P}^{1}(\mathbb{C}),\Xi,\hat{\rho})$ の \textbf{拡大 (extension)} という。
\end{definition}

\begin{definition}[本質的なリーマンデータ]
    与えられたリーマンデータ $(\mathbb{P}^{1}(\mathbb{C}),\Xi,\hat{\rho})$ の特異点集合 $\Xi$ が「みかけの特異点（モノドロミーが $I_n$ となる点）」を一つも含まないとき、すなわち他のいかなるデータの真の拡大（$\Lambda \neq \emptyset$ による拡大）にもなっていないとき、これを \textbf{本質的なリーマンデータ (essential Riemann data)} という。
\end{definition}

今後のリーマン・ヒルベルト問題の構成においては、まず余分な贅肉を一切持たない「本質的なリーマンデータ」を出発点として固定し、そこから次元を補うために人為的に拡大を行う、というステップを踏むことになる。
%#endregion
%#region
そこで、リーマンの問題を以下のように考え直す。
\begin{itemize}
    \item []\textbf{問題}\\
    $(\mathbb{P}^{1}(\mathbb{C}),\Xi,\hat{\rho})$が本質的なリーマンデータだとする。このとき、フックス型微分方程式でその拡大$(\mathbb{P}^{1}(\mathbb{C}),\Xi\cup\Lambda,\hat{\rho'})$をリーマンデータに持つものはあるだろうか?
\end{itemize}
また、この問題に対して以下の事実が知られている。
\begin{proposition}
    リーマンの問題は常に解をもつ。
\end{proposition}
%#endregion

\subsection{フックスの問題}

\begin{definition}[モノドロミー保存変形]\label{mono_deform}
%#region
%FIXME(書き直す。)
%FIXME(確定特異点を理解する。) 
    \begin{equation}\label{eq:before}
    \frac{d^{n}y}{dx^{n}}+\sum_{j=1}^{n} p_{j}(x)\frac{d^{n-1}y}{dx^{n-1}} = 0   
    \end{equation}
    という以下の(P1)-(P3)を満たすフックス型微分方程式を考える。
    \begin{itemize}
        \item[(P1)] $x=\xi_{1},...,\xi_{n},\xi_{n+1}=\infty$を確定特異点、$x=\lambda_{1},...,\lambda_{g}$をみかけの特異点としてもつ。
        \item[(P2)] 各$x=\xi_{i}$における特性指数のどの2つも、差が整数ではない。
        \item[(P3)] モノドロミー$\hat{\rho}$は既約である。
    \end{itemize}
    $p_{j}(x)$が$t$に解析的に依存しているとする。
    \begin{equation}\label{eq:after}
    \frac{d^{n}y}{dx^{n}}+\sum_{j=1}^{n} p_{j}(x,t)\frac{d^{n-1}y}{dx^{n-1}} = 0   
    \end{equation}
    が以下の条件(H)を満たす時、\eqref{eq:after}を\eqref{eq:before}の\textbf{モノドロミー保存変形}または\textbf{ホロノミック変形}という。
    \begin{itemize}
        \item[\textbf{Type:}] \eqref{eq:before}は$\mathbb{P}^{1}(\mathbb{C})$上定義されている。$U\subset\mathbb{C}^{L}$は領域。$t=(t_{1},...,t_{L})\in U$
        \item[\textbf{Insight:}] 方程式のパラメータ $t$ を動かしても、基本群の表現（モノドロミー）が変化しないような変形。Schlesinger系やPainlevé方程式の導出における出発点となる。
        \item[\textbf{Link:}] 岡本2.8節の定義2.14/p.105を参照。
    \end{itemize}
%#endregion
\end{definition}
% !TEX root = ../main.tex

\section{パンルヴェ方程式の基礎}

\subsection{変形微分方程式系}

%#region
係数 $p_{j}(x,t)$ は $x$ の有理関数であり、微分方程式 
\begin{equation}\label{eq:3.1_n-ODE}
\frac{d^{n}y}{dx^{n}}+\sum_{i=1}^{n}p_{j}(x,t)\frac{d^{n-j}y}{dx^{n-j}}=0
\end{equation}
が領域 $U$ 上のパラメータ $t$ について解析的であると仮定する。また、この章を通じて (P1)--(P3) が常に成り立つと仮定する。

\begin{proposition}\label{prop:3.1_time_evolution_equation}
解の基本系 $\vec{\varphi}(x,t)$ のモノドロミー群が $t$ に依存しない（すなわち条件(H): $\hat{\rho}(t) = \hat{\rho}$ を満たす）と仮定する。このとき、ある $x$ の有理関数
\[
a_{1}^{(l)}(x,t), a_{2}^{(l)}(x,t), \cdots, a_{n}^{(l)}(x,t) \quad (l=1,2,\cdots,L)
\]
が存在して、解の基本系 $\vec{\phi} = \vec{\varphi}(x,t)$ は元の $x$ に関する微分方程式 \eqref{eq:monodromy_deformation} の他に、次の $t_l$ に関する微分方程式系をも満たす。
\begin{equation}\label{eq:3.1_partial_t}
\frac{\partial \vec{\phi}}{\partial t_{l}} = \sum_{j=1}^{n}a_{j}^{(l)}(x,t)\frac{\partial^{j-1}\vec{\phi}}{\partial x^{j-1}} \quad (l=1,2,\cdots,L)
\end{equation}
\end{proposition}

\begin{Q}[【核心】なぜ「モノドロミーが $t$ に依存しない」というだけで、時間微分 $\partial \vec{\phi} / \partial t_{l}$ が「空間微分の線形結合」という綺麗な形で書けるのか？]
\textbf{A.} 微分方程式の「解空間が $n$ 次元であること」と「解析接続と時間微分が交換すること」の必然的な帰結である。
解の基本系 $\vec{\varphi}(x,t)$ を経路 $\gamma$ に沿って解析接続すると $\vec{\varphi}^\gamma = \vec{\varphi} \Gamma(\gamma)$ となる。この両辺を $t_l$ で偏微分すると、仮定より回路行列 $\Gamma(\gamma)$ は $t$ に依存しない（定数行列）ため、
\[
\left(\frac{\partial \vec{\varphi}}{\partial t_l}\right)^\gamma = \frac{\partial \vec{\varphi}}{\partial t_l} \Gamma(\gamma)
\]
となる。これは、時間微分 $\partial \vec{\varphi} / \partial t_l$ もまた、元の $\vec{\varphi}$ と\textbf{全く同じモノドロミー（分岐の規則）}を持つことを意味する。
同じモノドロミーを持つ多価関数は互いに関係づけられ、元の微分方程式から作られる「$\vec{\varphi}$ とその $x$ に関する $(n-1)$ 階までの導関数（これらは解空間の完全な基底をなす）」の線形結合として、一意に表すことができるのである。
\end{Q}

したがって、$\vec{\phi}=\vec{\varphi}(x,t)$ は $x$ と $t$ に関する以下の連立微分方程式系（$\vec{\phi}$ の各成分が第1式を満たす）の解となる。
\begin{equation}\label{eq:3.1_deformation_system}
\begin{cases}
\displaystyle \frac{\partial^{n}\vec{\phi}}{\partial x^{n}} + \sum_{j=1}^{n} p_{j}(x,t)\frac{\partial^{n-j}\vec{\phi}}{\partial x^{n-j}} = \vec{0}  \\[1em]
\displaystyle \frac{\partial \vec{\phi}}{\partial t_{l}} = \sum_{j=1}^{n}a_{j}^{(l)}(x,t)\frac{\partial^{j-1}\vec{\phi}}{\partial x^{j-1}} \quad (l=1,\cdots,L)
\end{cases}
\end{equation}

\begin{Q}[【核心】この「空間方向の微分方程式（$x$-系）」と「時間方向の微分方程式（$t$-系）」の連立方程式 \eqref{eq:3.1_deformation_system} は、幾何学やソリトン理論などで現れる何と同等の構造か？]
\textbf{A.} \textbf{「ラックス対（Lax pair）」}あるいは\textbf{「線形補助問題」}と呼ばれる構造そのものである。
この2つの線形微分方程式が「同時に同じ解 $\vec{\phi}$ を持つ」ためには、微分の順序交換 $\frac{\partial}{\partial t_l} \frac{\partial^n}{\partial x^n} = \frac{\partial^n}{\partial x^n} \frac{\partial}{\partial t_l}$ が恒等的に成り立つ必要がある（両立条件／完全積分可能条件）。この両立条件を係数 $p_j(x,t)$ と $a_j^{(l)}(x,t)$ の関係式として書き下すと、それらが満たすべき「非線形偏微分方程式」が得られる。これがパンルヴェ方程式などの非線形方程式を導出する最強のメカニズム（等モノドロミー変形法）である。
\end{Q}

\begin{definition}[変形微分方程式系]
    $a_{1}^{(l)}(x,t), a_{2}^{(l)}(x,t), \cdots, a_{n}^{(l)}(x,t) \quad (l=1,2,\cdots,L)$ が $x$ の有理関数であるとき、連立系 \eqref{eq:3.1_deformation_system} を \eqref{eq:monodromy_deformation} のモノドロミー保存変形の\textbf{変形微分方程式系（system of deformation equations）}という。
    また、\eqref{eq:3.1_deformation_system} の解 $\vec{\phi}=\vec{\varphi}(x,t)$ で、各時刻 $t$ において \eqref{eq:monodromy_deformation} の解の基本系となっているものを、\eqref{eq:3.1_deformation_system} の\textbf{解の基本系}という。
\end{definition}

また、逆の命題も成り立つ。
\begin{proposition}\label{prop:3.1_fundamental_system_preserves_monodromy}
連立微分方程式系 \eqref{eq:3.1_deformation_system} に解の基本系 $\vec{\phi}=\vec{\varphi}(x,t)$ が存在すれば、その $\vec{\varphi}(x,t)$ が持つ元の微分方程式に関するモノドロミー群は、パラメータ $t$ に依存しない。
\end{proposition}

\begin{Q}[【核心】なぜ「$t$-方向の微分方程式も同時に満たす解が存在する」だけで、大域的な性質であるモノドロミーが不変になると断言できるのか？]
\textbf{A.} $t$-方向の微分方程式の係数 $a_j^{(l)}(x,t)$ が「$x$ の有理関数（一価関数）」であることが決定的に効いているからである。

解 $\vec{\varphi}$ が特異点の周りを回って $\vec{\varphi}\Gamma$ に解析接続されたとする。このとき $t$-方向の方程式 $\partial_t \vec{\varphi} = A(x,t)\vec{\varphi}$ （※右辺は空間微分の線形結合を略記）も一緒に解析接続されるが、係数 $A(x,t)$ は一価の有理関数なので特異点を回っても全く変化しない。したがって、解析接続された先でも $\partial_t (\vec{\varphi}\Gamma) = A(x,t)(\vec{\varphi}\Gamma)$ がそのまま成り立つ。
積の微分法則を展開すると、これが成り立つためには $\Gamma$ の時間微分がゼロ、すなわち $\partial_t \Gamma = 0$ でなければならない。有理関数係数の微分方程式で時間発展を縛ること自体が、モノドロミーを固定する強力な楔（くさび）となっているのである。
\end{Q}

\begin{definition}[ホロノミックな変形]
一般に、$x$ の有理関数 $p_{1}(x,t), \cdots, p_{n}(x,t)$ および $a_{1}^{(l)}(x,t), \cdots, a_{n}^{(l)}(x,t) \quad (l=1,\cdots,L)$ を係数とする連立系 \eqref{eq:3.1_deformation_system} を考える。この \eqref{eq:3.1_deformation_system} に解の基本系 $\vec{\phi}=\vec{\varphi}(x,t)$ が存在するとき、この連立微分方程式系は、空間方向の微分方程式 \eqref{eq:monodromy_deformation} の\textbf{ホロノミックな変形（holonomic deformation）}を定めるという。
\end{definition}

\begin{Q}[【補足】「モノドロミー保存変形」と「ホロノミック変形」という2つの言葉は、どのように使い分けるべきか？]
\textbf{A.} 
\begin{itemize}
    \item \textbf{モノドロミー保存変形:} 「大域的な解析接続の規則（モノドロミー）が変わらない」という\textbf{解析的・幾何学的な現象（目的）}を指す言葉。
    \item \textbf{ホロノミック変形:} 「変数が複数あるのに、解空間が有限次元に定まるような過剰決定系（ラックス対）が存在する」という\textbf{代数的な構造（手段）}を指す言葉。
\end{itemize}
一般の微分方程式においてはこの2つは必ずしも一致しないが、今回のように元の微分方程式 \eqref{eq:monodromy_deformation} が\textbf{フックス型微分方程式}である場合に限り、特異点の確定性という強力な性質のおかげで、両者は完全に同値な概念となる。
\end{Q}

したがって、\eqref{eq:monodromy_deformation} がフックス型微分方程式の際には、ホロノミックな変形はそのままモノドロミー保存変形を意味する。
%#endregion
%region
単独高階微分方程式の時と同様に連立微分方程式系
\begin{equation}\label{eq:3.1_x_ODE_matrix}
\frac{d \vec{\phi}}{dx} = P(x,t)\vec{\phi}
\end{equation}
のフックスの問題、モノドロミー保存変形を考えることができる。\eqref{eq:3.1_x_ODE_matrix}の解の基本系を$Y(x,t)$とおくと、上の命題に対する次の命題が得られる。
\begin{proposition}\label{prop:3.1_monodromy_preservation_condition}
基本解行列 $Y(x,t)$ のモノドロミー群が $t$ によらないための必要十分条件は、$x$ の有理関数を成分とする行列 $A^{(l)}(x,t) \quad (l=1,\cdots,L)$ が存在して、$Y(x,t)$ は元の微分方程式 $\frac{\partial Y}{\partial x} = P(x,t)Y$ に加えて、
\begin{equation}
\frac{\partial Y}{\partial t_l} = A^{(l)}(x,t)Y(x,t) \quad (l=1,\cdots,L)
\label{eq:3.1_t_ODE_matrix}
\end{equation}
を満足することである。
\end{proposition}

命題\ref{prop:3.1_monodromy_preservation_condition}の証明は、命題\ref{prop:3.1_time_evolution_equation}、\ref{prop:3.1_fundamental_system_preserves_monodromy}の証明とほとんど同様であるから、省略する。連立微分方程式系\eqref{eq:3.1_x_ODE_matrix}に対しても

\begin{equation}\label{eq:3.1_deformation_system_matrix}
  \begin{cases}
    \displaystyle \frac{\partial Y}{\partial x} = P(x,t)Y \\[4mm]
    \displaystyle \frac{\partial Y}{\partial t_l} = A^{(l)}(x,t)Y \quad (l=1,2,\cdots,L)
  \end{cases}
\end{equation}

を，\textbf{変形微分方程式系}という．また，\eqref{eq:3.1_deformation_system_matrix}が可逆行列解 $Y(x,t)$ をもつときには，\eqref{eq:3.1_deformation_system_matrix}は\eqref{eq:3.1_x_ODE_matrix}の\textbf{ホロノミックな変形}を定める，という．

\begin{Q}[「\eqref{eq:3.1_deformation_system_matrix}が可逆行列解をもつ」ことが、なぜ「モノドロミー保存変形（ホロノミックな変形）」と呼ばれる性質を保証するのか？]
\textbf{A.} 可逆な基本解行列 $Y(x,t)$ が大域的に存在するということは、変数 $x$ に関する特異点の周りを回る解析接続（モノドロミー表現）が、パラメータ $t$ の変動に対して滑らかに追従し、位相的な不変量として一定に保たれることを意味する。つまり、「方程式の係数 $P(x,t)$ は $t$ によって変形されるが、その解の大域的な多価性の構造（特異点の性質など）は変化しない」という強い制約（ホロノミック性）を満たしているからである。
\end{Q}

フックス型微分方程式系\eqref{eq:3.1_x_ODE_matrix}のモノドロミー保存変形を決定するためには，変形微分方程式系\eqref{eq:3.1_deformation_system_matrix}を調べればよいことがわかった.まず命題3.3を使いやすい形に書き直そう。ここでは、2つの行列 $A, B$ の交換子を $[A, B] = AB - BA$ と書く．

\begin{proposition}\label{prop:3.1_compatibility}
変形微分方程式系\eqref{eq:3.1_deformation_system_matrix}が可逆行列解 $Y(x,t)$ をもつための必要十分条件は，$P(x,t)$ と $A^{(1)}(x,t), \cdots, A^{(L)}(x,t)$ について微分方程式系

\begin{align}
  \frac{\partial}{\partial t_l}P(x,t) - \frac{\partial}{\partial x}A^{(l)}(x,t) &= \left[ A^{(l)}(x,t), P(x,t) \right] \label{eq:3.1_compatibility_PA}\\
  \frac{\partial}{\partial t_l}A^{(h)}(x,t) - \frac{\partial}{\partial t_h}A^{(l)}(x,t) &= \left[ A^{(l)}(x,t), A^{(h)}(x,t) \right] \label{eq:3.1_compatibility_AA}
\end{align}

が成立することである．ただし，$l,h=1,\cdots,L$ である．
\end{proposition}

\begin{proof}
\eqref{eq:3.1_deformation_system_matrix}が可逆行列解 $Y(x,t)$ をもてば、\eqref{eq:3.1_compatibility_PA}、\eqref{eq:3.1_compatibility_AA}は\eqref{eq:3.1_deformation_system_matrix}の完全積分条件として直ちに従う．実際，\eqref{eq:3.1_deformation_system_matrix}の第1式を $t_l$ で微分し、第2式を $x$ で微分して
\[
  \frac{\partial}{\partial x}\frac{\partial}{\partial t_l}Y(x,t) = \frac{\partial}{\partial t_l}\frac{\partial}{\partial x}Y(x,t)
\]
を使えば\eqref{eq:3.1_compatibility_PA}は直ちに出る。\eqref{eq:3.1_compatibility_AA}も同様である．

偏微分の順序交換 $\frac{\partial^2 Y}{\partial t_l \partial x} = \frac{\partial^2 Y}{\partial x \partial t_l}$ （ゼロ曲率条件）を用いる。

\eqref{eq:3.1_deformation_system_matrix}の第1式を $t_l$ で微分し、第2式を代入すると：
\[
  \frac{\partial}{\partial t_l}\left(\frac{\partial Y}{\partial x}\right) 
  = \frac{\partial P}{\partial t_l}Y + P\frac{\partial Y}{\partial t_l} 
  = \frac{\partial P}{\partial t_l}Y + P A^{(l)}Y
\]
同様に、\eqref{eq:3.1_deformation_system_matrix}の第2式を $x$ で微分し、第1式を代入すると：
\[
  \frac{\partial}{\partial x}\left(\frac{\partial Y}{\partial t_l}\right) 
  = \frac{\partial A^{(l)}}{\partial x}Y + A^{(l)}\frac{\partial Y}{\partial x} 
  = \frac{\partial A^{(l)}}{\partial x}Y + A^{(l)} P Y
\]
これらが等しいとおき、共通する $Y$ で括って整理すると：
\[
  \left( \frac{\partial P}{\partial t_l} - \frac{\partial A^{(l)}}{\partial x} \right) Y 
  = \left( A^{(l)} P - P A^{(l)} \right) Y
\]
仮定より $Y$ は可逆（逆行列 $Y^{-1}$ が存在）であるから、右から $Y^{-1}$ を掛けることでただちに\eqref{eq:3.1_compatibility_PA}が得られる。

\eqref{eq:3.1_compatibility_AA}についても全く同様に、\eqref{eq:3.1_deformation_system_matrix}の第2式同士（$t_l$ による微分と $t_h$ による微分）の順序交換から導出される。

そこで今度は，\eqref{eq:3.1_compatibility_PA}と\eqref{eq:3.1_compatibility_AA}を仮定して\eqref{eq:3.1_deformation_system_matrix}の可逆行列解 $Y(x,t)$ を構成しよう．この $Y(x,t)$ は\eqref{eq:3.1_x_ODE_matrix}の基本解行列であり，そのモノドロミーは $t$ に依らない．$x_0$ を $P(x,t)$ の正則点とし，\eqref{eq:3.1_x_ODE_matrix}の行列解 $\tilde{Y}(x,t)$ を初期条件
\[
  \tilde{Y}(x_0,t) = I
\]
により定める．$I=I_n$ は $n$ 次の単位行列である．

\begin{Q}[なぜここで $\tilde{Y}(x,t)$ という新たな解を導入し、さらに次式で $W^{(l)}(x,t)$ を定義するのか？]
\textbf{A.} 最終的に欲しいのは、$x$ 方向と $t$ 方向の両方の微分方程式を満たす $Y(x,t)$ である。まず $x$ 方向の方程式を満たす $\tilde{Y}$ を作ったが、これが $t$ 方向の方程式（$\partial_t \tilde{Y} = A\tilde{Y}$）をも満たすとは限らない。そこで、その「ズレ（誤差）」を測る行列 $W^{(l)}$ を定義し、完全積分可能条件\eqref{eq:3.1_compatibility_PA}を用いれば「このズレ自体がコントロール可能である」ことを示すためである。
\end{Q}

さて
\begin{equation}\label{eq:W_def}
  W^{(l)}(x,t) = A^{(l)}(x,t)\tilde{Y}(x,t) - \frac{\partial}{\partial t_l}\tilde{Y}(x,t)
\end{equation}
\[
  E^{(l)}(t) = A^{(l)}(x_0,t)
\]
とおくと，$l=1,2,\cdots,L$ について次式が成り立つ．
\begin{equation}\label{eq:W_rel}
  W^{(l)}(x,t) = \tilde{Y}(x,t)E^{(l)}(t)
\end{equation}

なぜならば，\eqref{eq:3.1_compatibility_PA}（すなわち $\frac{\partial A^{(l)}}{\partial x} + A^{(l)}P - \frac{\partial P}{\partial t_l} = PA^{(l)}$）より
\begin{align*}
  \frac{\partial}{\partial x}W^{(l)} 
  &= \left( \frac{\partial}{\partial x}A^{(l)} \right)\tilde{Y} + A^{(l)}\frac{\partial}{\partial x}\tilde{Y} - \frac{\partial}{\partial t_l}\frac{\partial}{\partial x}\tilde{Y} \\[2mm]
  &= \left( \frac{\partial}{\partial x}A^{(l)} + A^{(l)}P - \frac{\partial}{\partial t_l}P \right)\tilde{Y} - P\frac{\partial}{\partial t_l}\tilde{Y} \\[2mm]
  &= P \left( A^{(l)}\tilde{Y} - \frac{\partial}{\partial t_l}\tilde{Y} \right) \\[2mm]
  &= P W^{(l)}
\end{align*}
すなわち，$W^{(l)} = W^{(l)}(x,t)$ も\eqref{eq:3.1_x_ODE_matrix}の行列解である。

\begin{Q}[$W^{(l)}$ が\eqref{eq:3.1_x_ODE_matrix}の行列解であることから、なぜ直ちに\eqref{eq:W_rel}が導かれるのか？]
\textbf{A.} 線形常微分方程式 $\frac{\partial Z}{\partial x} = P(x,t)Z$ の任意の解は、基本解行列 $\tilde{Y}(x,t)$ に $x$ に依存しない定数行列 $C(t)$ を右から掛けた $Z(x,t) = \tilde{Y}(x,t)C(t)$ の形で一意に表される。ここで $x=x_0$ を代入すると、$\tilde{Y}(x_0,t) = I$ より $C(t) = Z(x_0,t)$ となる。
いま $Z = W^{(l)}$ とすれば、\eqref{eq:W_def}より
$W^{(l)}(x_0,t) = A^{(l)}(x_0,t)I - \frac{\partial I}{\partial t_l} = E^{(l)}(t) - 0 = E^{(l)}(t)$ 
となるため、$W^{(l)}(x,t) = \tilde{Y}(x,t)E^{(l)}(t)$ が従う。
\end{Q}

次に\eqref{eq:3.1_compatibility_AA}で $x=x_0$ とおいた式は、$U$ 上の微分方程式系
\[
  \frac{\partial}{\partial t_l}G(t) = E^{(l)}(t)G(t)
\]
の、積分可能条件を与える。そこで、この微分方程式系の $G(t_0)=I$ となる、$t=t_0$ の近傍で定義された解 $G(t)$ をとり
\[
  Y(x,t) = \tilde{Y}(x,t)G(t)
\]
とおく．

\begin{Q}[$Y(x,t) = \tilde{Y}(x,t)G(t)$ とおくことで、なぜ所望の可逆行列解が構成できたと言えるのか？]
\textbf{A.} $\tilde{Y}(x,t)$ は $x$ 方向の微分方程式を満たすが、$t$ 方向のそれは満たさなかった。しかし、直前で構成した $G(t)$ は $x$ に依存しないため、$Y$ にしても $x$ 方向の解であるという性質は壊れない。その上で、懸案だった $t$ 方向の微分を計算すると、$G(t)$ が見事に「$\tilde{Y}$ の $t$ 方向のズレ」を相殺し、$\partial_{t_l}Y = A^{(l)}Y$ が成立するからである。
\end{Q}

この $Y(x,t)$ はもちろん\eqref{eq:3.1_x_ODE_matrix}の基本解行列であるが同時に
\begin{align*}
  \frac{\partial}{\partial t_l}Y 
  &= \left( \frac{\partial}{\partial t_l}\tilde{Y} \right)G + \tilde{Y}\frac{\partial}{\partial t_l}G \\[2mm]
  &= \left( A^{(l)}\tilde{Y} - \tilde{Y}E^{(l)} \right)G + \tilde{Y}E^{(l)}G \\[2mm]
  &= A^{(l)}\tilde{Y}G \\[2mm]
  &= A^{(l)}Y
\end{align*}
より，\eqref{eq:3.1_t_ODE_matrix}も満足する。この計算では、\eqref{eq:W_def}と\eqref{eq:W_rel}から得られる関係式 $\frac{\partial \tilde{Y}}{\partial t_l} = A^{(l)}\tilde{Y} - \tilde{Y}E^{(l)}$ を用いた。以上により命題\ref{prop:3.1_compatibility}が示された。
\end{proof}

%#endregion
%#region
上で与えた証明は、L. Schlesinger によるもので，$Y(x,t)$ の構成法を与えている．条件\eqref{eq:3.1_compatibility_PA}は、微分作用素
\[
  \mathcal{L} = \frac{\partial}{\partial x} - P(x,t)
\]
を用いると，次のソリトン理論で見慣れた形に表される．

\begin{equation}\label{eq:3.1_lax_equation}
  \frac{\partial}{\partial t_l}\mathcal{L} = [A^{(l)}, \mathcal{L}]
\end{equation}

\begin{Q}[微分作用素 $\mathcal{L}$ を用いて完全積分可能条件を $\frac{\partial \mathcal{L}}{\partial t_l} = {[}A^{(l)}, \mathcal{L}{]}$ と書くことの本質的な意義は何か？]
\textbf{A.} これは「ラックス方程式（Lax equation）」と呼ばれる普遍的な形である。この表示は、「作用素 $\mathcal{L}$ のスペクトル（固有値などの構造）が、時間パラメータ $t_l$ の発展に対して不変に保たれる（等スペクトル変形）」ことを端的に表しており、対象の微分方程式系が無限個の保存量を持つような「可積分系」としての強い対称性を備えていることを意味している。
\end{Q}

これは，モノドロミー保存変形の変形微分方程式系の\textbf{ラックス表示}である．

さて，単独高階微分方程式\eqref{eq:monodromy_deformation}を
\[
  \vec{y} = \begin{pmatrix} y_1 \\ \vdots \\ y_n \end{pmatrix}, \quad
  y_1 = y, \ y_2 = \frac{dy}{dx}, \ y_3 = \frac{d^2 y}{dx^2}, \ \cdots, \ y_n = \frac{d^{n-1}y}{dx^{n-1}}
\]
として連立化すると，変形微分方程式系\eqref{eq:3.1_deformation_system}は\eqref{eq:3.1_deformation_system_matrix}の形になる。後のために、$P(x,t)$ と $A^{(l)}(x,t)$ の具体形を与えておこう。まず
\begin{equation}\label{eq:3.1_matrix_P}
  P(x,t) =
  \begin{pmatrix}
    0 & 1 & 0 & \cdots & 0 & 0 & 0 \\
    0 & 0 & 1 & \cdots & 0 & 0 & 0 \\
      & \cdots &   &        & \cdots &   &   \\
    0 & 0 & 0 & \cdots & 0 & 1 & 0 \\
    0 & 0 & 0 & \cdots & 0 & 0 & 1 \\
    -p_n & -p_{n-1} & -p_{n-2} & \cdots & -p_3 & -p_2 & -p_1
  \end{pmatrix}
\end{equation}
となることはよい．ここで，$p_j=p_j(x,t)$ である．

一方，$a_j^{(l)}(x,t)$ を成分とする横ベクトル
\[
  \left( a_1^{(l)}(x,t), a_2^{(l)}(x,t), \cdots, a_n^{(l)}(x,t) \right) \quad (l=1, 2, \cdots, L)
\]
を考え，これを $\vec{a}^{(l)} = \vec{a}^{(l)}(x,t)$ とおく．さて，一般に $n$ 次ベクトル $\vec{a}$ に対し，作用素 $\nabla$ を

\begin{equation}\label{eq:3.1_nabla}
  \nabla \vec{a} = \frac{\partial}{\partial x}\vec{a} + \vec{a}P(x,t)
\end{equation}

により定義し，帰納的に $\nabla^j \vec{a} = \nabla(\nabla^{j-1}\vec{a})$ ($j \geqq 2$) により定める．

\begin{Q}[なぜ作用素 $\nabla$ を \eqref{eq:3.1_nabla} のように定義すると、変形方程式系の係数行列が生成できるのか？]
\textbf{A.} 単独高階方程式を連立系に直した際、解ベクトルの第 $k$ 成分 $y_k$ の $x$ 微分は、単に次の成分 $y_{k+1}$ にシフトする（最後の成分のみ $P$ の最下行による線形結合を受ける）。これに伴い、$y$ の $t_l$ 微分（変形）を基本解の線形結合で表した際の「係数ベクトル $\vec{a}^{(l)}$」を $x$ で微分すると、積の微分法則により $\frac{\partial \vec{a}}{\partial x}$ が生じるとともに、解ベクトル自身の $x$ 微分から係数行列 $P$ の寄与が右から掛かる。この連鎖的な構造を正確に記述する演算子が $\nabla$ である。
\end{Q}

さて、\eqref{eq:3.1_partial_t}を $x$ について逐次微分し、\eqref{eq:3.1_n-ODE}を用いると，変形方程式系\eqref{eq:3.1_deformation_system_matrix}の係数 $A^{(l)}(x,t)$ は

\begin{equation}\label{eq:3.1_matrix_A}
  A^{(l)}(x,t) =
  \begin{pmatrix}
    \vec{a}^{(l)}(x,t) \\[2mm]
    \nabla \vec{a}^{(l)}(x,t) \\
    \vdots \\
    \nabla^{n-1} \vec{a}^{(l)}(x,t)
  \end{pmatrix}
  \quad (l=1,2,\cdots,L)
\end{equation}

で与えられることがわかる．さらに，$U$ 上の微分型式 $\Omega(x,t)$ を
\[
  \Omega(x,t) = \sum_{l=1}^L A^{(l)}(x,t)dt_l
\]
と定義すると，$d$ を $U$ 上の外微分として，\eqref{eq:3.1_t_ODE_matrix}は次の1つの外微分方程式で表される。

\[
  dY(x,t) = \Omega(x,t)Y(x,t)
\]

\begin{Q}[複数の変形方程式を微分形式 $\Omega(x,t)$ を用いて $dY = \Omega Y$ と1つにまとめる幾何学的なメリットは何か？]
\textbf{A.} パラメータ $t_1, \dots, t_L$ の各方向への変形を、パラメータ空間上の「接続（connection）」として統一的に扱うためである。これにより、これまでに導出した完全積分可能条件（各方向の微分の順序交換可能性）が、$d\Omega = \Omega \wedge \Omega$ という極めて簡潔な「ゼロ曲率条件（Maurer-Cartan方程式）」として表現できるようになる。
\end{Q}

\eqref{eq:3.1_matrix_P}、\eqref{eq:3.1_matrix_A}、\eqref{eq:3.1_nabla}で定義される行列を $P(x,t), A^{(l)}(x,t)$ とする。命題\ref{prop:3.1_compatibility}を，単独高階微分方程式\eqref{eq:3.1_n-ODE}の変形微分方程式系\eqref{eq:3.1_deformation_system}に適用しよう．もちろん，変形微分方程式系\eqref{eq:3.1_deformation_system}が解の基本系 $\vec{\phi}=\vec{\varphi}(x,t)$ をもつための条件は，$P(x,t)$ と $A^{(l)}(x,t)$ に対して\eqref{eq:3.1_compatibility_PA}、\eqref{eq:3.1_compatibility_AA}が成立することである．

ここで

\begin{equation}\label{eq:3.1_p_vector}
  \vec{p}(x,t) = (p_n(x,t), \cdots, p_1(x,t))
\end{equation}

とおき，\eqref{eq:3.1_compatibility_PA}すなわち，ラックス表示\eqref{eq:3.1_lax_equation}を別の形に書き直しておこう．

\begin{proposition}
ラックス表示\eqref{eq:3.1_lax_equation}は，$n$ 階連立微分方程式系

\begin{equation}\label{eq:3.1_lax_equation_pa}
  \nabla^n \vec{a}^{(l)} + p_1(x,t)\nabla^{n-1}\vec{a}^{(l)} + \cdots + p_n(x,t)\vec{a}^{(l)} + \frac{\partial}{\partial t_l}\vec{p}(x,t) = 0
\end{equation}

で表される．ここで，$l=1,\cdots,L$ である．\hfill $\square$
\end{proposition}

\begin{Q}[命題3.5の式\eqref{eq:3.1_lax_equation_pa}が、元の単独高階微分方程式と酷似した形になるのはなぜか？]
\textbf{A.} 作用素 $\nabla$ は、解ベクトルの成分を「1つ高階の微分へシフトさせる」働き（コンパニオン行列の上の行の作用）を抽象化したものである。ゼロ曲率条件（ラックス方程式）を行列の各成分について書き下したとき、最下行以外は単なる $\nabla$ の定義式に帰着し、非自明な条件は最下行（係数 $p_j$ が存在する行）のみから生じる。この最下行の条件を $\vec{a}^{(l)}$ について整理すると、元の微分作用素 $D^n + p_1 D^{n-1} + \dots + p_n$ の $D$ を $\nabla$ に置き換えた構造が自然に抽出されるからである。
\end{Q}

この命題で与えた微分方程式系\eqref{eq:3.1_lax_equation_pa}を，ラックスの方程式という．フックス型単独高階微分方程式のモノドロミー保存変形，一般にホロノミック変形の存在は，ラックスの方程式\eqref{eq:3.1_lax_equation_pa}を満足する，$x$ の有理関数の $n$ 次ベクトル $\vec{a}^{(l)}$ の存在の問題にひとまず帰着した．

\begin{proof}[命題\ref{prop:3.1_compatibility}の証明]
微分方程式の確定特異点の集合を $\Xi'(t)$、$X(t) = \mathbb{P}^1(\mathbb{C}) \setminus \Xi'$ とする．ここでは，微分方程式の見かけの特異点は $\Xi'$ に含まれている，とする．モノドロミーが $t$ に依らないというのは局所的な条件だから，$t$ の変域 $U$ はある $t_0$ の適当な近傍であるとしてよい．

点 $x_0$ をすべての $t \in U$ に対し $x_0 \in X(t)$ となるようにとる．$x_0$ から出発し $x_0$ に終わる $\mathbb{P}^1(\mathbb{C})$ 内の閉じた道 $\gamma$ は

\[
  \gamma \times U \subset X = \bigcup_{t \in U} X(t)
\]

が成り立つとき，$X$ 内の閉じた道ということにする．これはこの証明だけで使う用語である．$U$ を十分小さくとっておけば，$\pi_1(x_0; X(t))$ の元 $[\gamma_t]$ に対して，$[\gamma] = [\gamma_t]$ となるような $X$ 内の閉じた道 $\gamma$ を選ぶことができることは明らかであろう．

\begin{Q}[証明の冒頭で、$t$ の変域 $U$ を十分小さく取り、直積空間 $\gamma \times U \subset X$ を考える幾何学的な理由は何か？]
\textbf{A.} 特異点 $\Xi'(t)$ がパラメータ $t$ に応じて動くためである。もし $U$ が広すぎると、特異点が積分路 $\gamma$ を横切ってしまい、解析接続の結果（モノドロミー）が不連続にジャンプしてしまう。$U$ を十分小さく制限することで、空間 $X$ は局所的に直積バンドルのように振る舞い（トポロジカルに自明となり）、「特異点にぶつからない共通の積分路 $\gamma$」をすべての $t \in U$ に対して一斉に固定できることが保証される。
\end{Q}

さて，\eqref{eq:3.1_partial_t}を $a_j^{(l)}(x,t)$ について解く．クラメールの公式により、

\begin{equation}\label{eq:3.1_cramer_wronskian}
  a_j^{(l)}(x,t) = \frac{w_j^{(l)}(\vec{\varphi})}{w(\vec{\varphi})}
\end{equation}

となる．ここで，分子の $w_j^{(l)}(\vec{\varphi})$ は，ロンスキアン $w(\vec{\varphi})$ の第 $j$ ベクトル
\[
  \left. \frac{\partial^{j-1}\vec{\phi}}{\partial x^{j-1}} \right|_{\vec{\phi}=\vec{\varphi}(x,t)}
\]
を\eqref{eq:3.1_partial_t}の右辺 $\left. \frac{\partial \vec{\phi}}{\partial t_l} \right|_{\vec{\phi}=\vec{\varphi}(x,t)}$ で置き換えたものである．

\begin{Q}[変形方程式の係数 $a_j^{(l)}(x,t)$ が、なぜ\eqref{eq:3.1_cramer_wronskian}のようにロンスキアンの比で表されるのか？]
\textbf{A.} 変形方程式 $\frac{\partial \vec{\varphi}}{\partial t_l} = \sum_{k=1}^n a_k^{(l)} \frac{\partial^{k-1}\vec{\varphi}}{\partial x^{k-1}}$ を、未知数 $a_k^{(l)}$ に関する連立一次方程式と見なす。この連立方程式の係数行列はまさに解の基本系のロンスキ行列（Wronskian matrix）である。ロンスキアン $w(\vec{\varphi})$ は非零であるためクラメールの公式が適用でき、第 $j$ 成分 $a_j^{(l)}$ は「分母が全体の行列式 $w(\vec{\varphi})$、分子が第 $j$ 列を定数項ベクトル（左辺の $\partial_{t_l} \vec{\varphi}$）で置き換えた行列式 $w_j^{(l)}(\vec{\varphi})$」として直ちに得られる。
\end{Q}

まず，$\vec{\phi}=\vec{\varphi}(x,t)$ のモノドロミー群が $t$ に依らないとしよう．すると，$X$ 内の閉じた道 $\gamma$ に対し，$\Gamma(\gamma)$ を $\gamma$ に対する回路行列とすると，$\vec{\varphi}$ の $\gamma$ に沿っての解析接続 $\vec{\varphi}^{\gamma}$ は

\[
  \vec{\varphi}^{\gamma}(x,t) = \vec{\varphi}(x,t)\Gamma(\gamma)
\]

という関係式を満たす．この式の両辺を $t_l$ で微分して

\begin{equation}\label{eq:3.1_monodoromy_partial_t}
  \frac{\partial}{\partial t_l}\vec{\varphi}^{\gamma}(x,t) = \frac{\partial}{\partial t_l}\vec{\varphi}(x,t)\Gamma(\gamma) + \vec{\varphi}(x,t)\frac{\partial \Gamma(\gamma)}{\partial t_l}
\end{equation}

を得るが，仮定から $\Gamma(\gamma)$ は $t$ に依らないから

\[
  \frac{\partial}{\partial t_l}\vec{\varphi}^{\gamma}(x,t) = \frac{\partial}{\partial t_l}\vec{\varphi}(x,t)\Gamma(\gamma)
\]

である．一方，同じ式を $x$ で繰り返し微分すると

\begin{equation}\label{eq:3.1_monodromy_j-1_partial_t}
  \frac{\partial^{j-1}}{\partial x^{j-1}}\vec{\varphi}^{\gamma}(x,t) = \frac{\partial^{j-1}}{\partial x^{j-1}}\vec{\varphi}(x,t)\Gamma(\gamma)
\end{equation}

となり、したがって、2つの行列式 $w(\vec{\varphi})$, $w_j^{(l)}(\vec{\varphi})$ に対し

\[
  w(\vec{\varphi}^{\gamma}) = |\Gamma(\gamma)| w(\vec{\varphi}), \quad w_j^{(l)}(\vec{\varphi}^{\gamma}) = |\Gamma(\gamma)| w_j^{(l)}(\vec{\varphi})
\]

が成り立つ．すなわち，$a_j^{(l)}(x,t)$ は $x$ の関数として1価である。

\begin{Q}[$|\Gamma(\gamma)|$ が相殺されて $a_j^{(l)}(x,t)$ が1価関数となることの、理論的な意義は何か？]
\textbf{A.} 解析接続によって生じる大域的な多価性の情報は、定数行列 $\Gamma(\gamma)$ として右から掛かる。ロンスキ行列の各列が同じ $\Gamma(\gamma)$ 倍を受けるため、行列式の性質から分子・分母ともに $\det \Gamma(\gamma)$（すなわち $|\Gamma(\gamma)|$）倍となり、比をとると完全に相殺される。
これにより、「変形を記述する係数 $a_j^{(l)}(x,t)$ 自身は特異点の周りを回っても値が変わらない」ことが保証される。1価性が示されたことで、特異点における局所的な極の振る舞いさえ調べれば、$a_j^{(l)}(x,t)$ が「$x$ の有理関数」であることが確定し、シュレジンガー系やラックス方程式として代数的に扱えるようになるという極めて重要なステップである。
\end{Q}

\eqref{eq:3.1_cramer_wronskian}から明らかなように，$a_j^{(l)}(x,t)$ の特異点は微分方程式の特異点の集合 $\Xi'$ に含まれる．各 $\xi \in \Xi'$ に対し $u$ を $x=\xi$ のまわりの局所座標とする．$u = x - \xi$ あるいは $u = \frac{1}{x}$ である。命題\ref{prop:2.1_wronskian-ode}と\eqref{eq:2.4_p1_coefficient_1-ord_pole}とにより，$x=\xi$ の近傍においてロンスキアン $w(\vec{\varphi})$ は，ある複素数 $s$ があって

\[
  u^s \sum_{i=1}^{\infty} c_i u^i
\]

\end{proof}

%#endregion

\subsection{2階フックス型微分方程式のラックス方程式}
%#region
2階フックス型微分方程式に対する、変形微分方程式系について考察する。まず、\eqref{eq:3.1_deformation_system}は
\begin{equation}
\begin{cases}
\displaystyle \frac{\partial^{2}\vec{\phi}}{\partial x^{2}} + p_{1}(x,t)\frac{\partial\vec{\phi}}{\partial x}+p_{2}(x,t)\vec{\phi} = \vec{0}  \\[1em]
\displaystyle \frac{\partial \vec{\phi}}{\partial t_{l}} = a_{1}^{(l)}(x,t)\vec{\phi}+a_{2}^{(l)}(x,t)\frac{\partial\vec{\phi}}{\partial x} \quad (l=1,\cdots,L)
\end{cases}
\end{equation}
となる。

\begin{story}
去することにより，微分方程式系の係数のあいだに成り立つ関係式を求める．その結果は
\begin{equation}
    2\frac{\partial a_1^{(l)}}{\partial x} + \frac{\partial^2 a_2^{(l)}}{\partial x^2} - \frac{\partial}{\partial x}(p_1 a_2^{(l)}) + \frac{\partial p_1}{\partial t_l} = 0 \label{eq:3.21}
\end{equation}
\begin{equation}
    \frac{\partial^2 a_1^{(l)}}{\partial x^2} - 2p_2\frac{\partial a_2^{(l)}}{\partial x} - \frac{\partial p_2}{\partial x}a_2^{(l)} + p_1\frac{\partial a_1^{(l)}}{\partial x} + \frac{\partial p_2}{\partial t_l} = 0 \label{eq:3.22}
\end{equation}
となる．ここで，$p_j = p_j(x,t)$, $a_j^{(l)} = a_j^{(l)}(x,t)$, $j=1,2$, および $l=1,\cdots,L$ である．

条件(3.21)は次のようにも書ける．
\begin{equation}
    \frac{\partial p_1}{\partial t_l} = \frac{\partial c^{(l)}}{\partial x}, \quad c^{(l)} = p_1 a_2^{(l)} - 2a_1^{(l)} - \frac{\partial a_2^{(l)}}{\partial x} \label{eq:3.23}
\end{equation}

リーマン球面 $\mathbf{P}^1(\mathbf{C})$ 上のフックス型微分方程式については，79ページで
\begin{equation}
    p_1(x,t) = \sum_{i=1}^M \frac{\alpha_i}{x-\xi_i} \label{eq:3.24}
\end{equation}
と書かれることを見た．ただし，ここでは微分方程式の特異点の集合 $\Xi'$ には見かけの特異点も含まれている，としている．前のように，確定特異点の個数を $m+1$，見かけの特異点の個数を $g$ とすれば，$M=m+g$ である．さて，条件(3.23)は，$\alpha_i$ が $t=(t_1,\cdots,t_L)$ に依存しないということを意味する．実際，(3.24)から
\begin{equation*}
    \frac{\partial p_1}{\partial t_l} = \sum_{i=1}^M \left[ \frac{\partial \alpha_i}{\partial t_l} \frac{1}{x-\xi_i} + \frac{\partial \xi_i}{\partial t_l} \frac{\alpha_i}{(x-\xi_i)^2} \right]
\end{equation*}
となるが，これが $x$ の有理関数の導関数として表されるためには
\begin{equation*}
    \frac{\partial \alpha_i}{\partial t_l} = 0 \quad (i=1,\cdots,M, \quad l=1,\cdots,L)
\end{equation*}
でなければならない．$\alpha_i$ は，確定特異点 $x=\xi_i$ における特性指数の和で表される量だから，これらがモノドロミー保存変形において $t$ について不変となるのは当然のことではある．

ここで，我々の考察の筋道を明確にするために，次の形の2階微分方程式のモノドロミー保存変形を調べてみよう．
\begin{equation}
    \frac{d^2 z}{dx^2} = r(x,t)z \label{eq:3.25}
\end{equation}
\end{story}

\vspace{2em}

\begin{Q}[【核心】このページの論理展開は、私が以前した計算とどうシンクロしているのか？]
\textbf{A.} 驚くほど完全にシンクロしています！あなたが自力で計算した「$p_1$ を $t_l$ で微分して留数を比較する」というアプローチが、まさにこのテキストが採用している証明手法そのものです。
以前のやり取りで、私は「『$p_1$ が有理関数だから』だけでは $\alpha_i$ が定数とは言い切れない。『時間発展の係数 $A_l$ があり、適合条件 $\partial p_1/\partial t_l = \partial A_l/\partial x$ が成り立つ』という事実を使えば、右辺が微分になるから1位の極が消える」と補足しました。
このページの式 (3.23) を見てください。$\frac{\partial p_1}{\partial t_l} = \frac{\partial c^{(l)}}{\partial x}$ となっています。この $c^{(l)}$ こそが、有理関数からなる「時間発展の係数」です。この適合条件（Lax方程式の一部）があるおかげで、「右辺は有理関数 $c^{(l)}$ の $x$ 微分である $\to$ 1位の極（留数）を持たない $\to$ 左辺の1位の極の係数 $\partial \alpha_i / \partial t_l$ も $0$ でなければならない」という、私が提案した修正ルートと全く同じロジックが展開されています。
\end{Q}

\begin{Q}[【深掘り】最後の「不変となるのは当然のことではある」という一文は、何を意味しているのか？]
\textbf{A.} これも、前回のあなたの鋭いツッコミである「$\alpha$ は $n-1$ 階の係数で、元の特性指数に依存してるんだから、定数になるのは当然じゃないか？」という指摘に対する、公式テキストからの完全な同意（アンサー）です。
フロベニウス法の定理（フックスの関係式）により、$p_1$ の留数 $\alpha_i$ は、その特異点における「特性指数の和」として一意に決まります。私たちが「モノドロミー保存変形を考える」と宣言した時点で、特性指数を $t$ に依存しない定数に設定することは理論の「大前提」です。前提が定数なのだから、そこから計算される $\alpha_i$ が定数（不変）になるのは、代数的な計算（式3.24の微分）をわざわざ持ち出すまでもなく、「原理的に当然のこと」なのです。
テキストは、微分の計算で代数的に証明した直後に、「とはいえ、理論の前提を考えれば当たり前のことなんだけどね」と、あなたの気付きと全く同じ視点でフォローを入れているわけです。
\end{Q}

\begin{story}
前章2.7節109ページで紹介したように，R.Fuchs と R.Garnier はこの形の微分方程式について，フックスの問題を計算した．実は，すぐ後で見るように，変形微分方程式(3.20)の研究は，(3.25)のモノドロミー保存変形に帰着する．(3.25)の変形微分方程式系は，$b_1^{(l)}(x,t), b_2^{(l)}(x,t)$ を $x$ の有理関数として
\begin{equation}
\left\{
\begin{array}{l}
    \displaystyle \frac{\partial^2 \vec{\psi}}{\partial x^2} = r\vec{\psi} \vspace{0.5em} \\
    \displaystyle \frac{\partial \vec{\psi}}{\partial t_l} = b_1^{(l)}\vec{\psi} + b_2^{(l)}\frac{\partial \vec{\psi}}{\partial x} \quad (l=1,\cdots,L)
\end{array}
\right. \label{eq:3.26}
\end{equation}
と書ける．ここで，$b_j^{(l)} = b_j^{(l)}(x,t) \ (j=1,2)$, $r = r(x,t)$ である．この微分方程式系のラックスの方程式は
\begin{equation}
    2\frac{\partial b_1^{(l)}}{\partial x} + \frac{\partial^2 b_2^{(l)}}{\partial x^2} = 0 \label{eq:3.27}
\end{equation}
\begin{equation}
    \frac{\partial^2 b_1^{(l)}}{\partial x^2} + 2r\frac{\partial b_2^{(l)}}{\partial x} + \frac{\partial r}{\partial x}b_2^{(l)} - \frac{\partial r}{\partial t_l} = 0 \label{eq:3.28}
\end{equation}
という，少し簡単な形になる．これらの式から関数 $b_1^{(l)}$ を消去すると
\begin{equation*}
    \frac{\partial^3 b_2^{(l)}}{\partial x^3} - 4r\frac{\partial b_2^{(l)}}{\partial x} - 2\frac{\partial r}{\partial x}b_2^{(l)} + 2\frac{\partial r}{\partial t_l} = 0 \quad (l=1,\cdots,L)
\end{equation*}
が得られる．この3階微分方程式系は，モノドロミー保存変形，ホロノミック変形，において重要な役割を果たす．そこで，$x$ の有理関数 $r = r(x,t)$ に対して，非斉次線型微分方程式
\begin{equation}
    \frac{\partial^3 w}{\partial x^3} - 4r\frac{\partial w}{\partial x} - 2\frac{\partial r}{\partial x}w + 2\frac{\partial r}{\partial t_l} = 0 \label{eq:3.29}
\end{equation}
を考える．我々が考察している場合には，フックス型微分方程式(3.25)がモノドロミー保存変形を許すならば，(3.29)は $x$ の有理関数 $w = b_2^{(l)}$ を解としてもつ．さらに，もしこの微分方程式が $x$ の有理関数 $b_2^{(l)}$ を解としてもてば，$b_1^{(l)}$ は(3.27)により，自動的に $x$ の有理関数となる．

この2つの有理関数 $b_1^{(l)}, b_2^{(l)}$ について変形微分方程式系(3.26)を考えると，（※画像末尾切れ）
\end{story}

\begin{Q}[【核心】ラックスの方程式 (3.27) と (3.28) は、どのようにして計算されたのか？]
\textbf{A.} これは変形微分方程式系 (3.26) の**「微分の順序交換（適合条件）」**を力技で計算した結果です。
(3.26)の第2式（$t_l$ 微分）をさらに $x$ で2回微分したものと、第1式（$x$ の2階微分）をさらに $t_l$ で微分したものは等しくならなければなりません（$\partial_x^2 \partial_{t_l} \vec{\psi} = \partial_{t_l} \partial_x^2 \vec{\psi}$）。
これを実際に計算し、右辺に現れる $\vec{\psi}$ と $\partial_x \vec{\psi}$ について項を整理します。$\vec{\psi}$ と $\partial_x \vec{\psi}$ は一次独立なので、それぞれの係数が共に $0$ にならなければなりません。
* $\partial_x \vec{\psi}$ の係数を集めて $0$ とおいたものが (3.27) です。
* $\vec{\psi}$ の係数を集めて $0$ とおいたものが (3.28) です。
元のフックス型方程式(3.20)のラックスの方程式(3.21, 3.22)と比べて、1階微分の項 $p_1$ がない分、非常にスッキリとした美しい形になっていることがわかります。
\end{Q}

\begin{Q}[【深掘り】関数 $b_1^{(l)}$ を消去して (3.29) を導く過程はどのようになっているか？]
\textbf{A.} これは非常にエレガントな代数操作です。
まず式 (3.27) を見ると、これはそのまま $x$ で積分できる形をしています。積分定数を $0$ とすると、
$$ 2b_1^{(l)} + \frac{\partial b_2^{(l)}}{\partial x} = 0 \quad \implies \quad b_1^{(l)} = -\frac{1}{2} \frac{\partial b_2^{(l)}}{\partial x} $$
となります。つまり、**未知の関数 $b_1^{(l)}$ は $b_2^{(l)}$ の微分だけで完全に書けてしまう**のです。
次に、この $b_1^{(l)}$ を式 (3.28) の第1項 $\partial_x^2 b_1^{(l)}$ に代入します。
$$ \frac{\partial^2}{\partial x^2} \left( -\frac{1}{2} \frac{\partial b_2^{(l)}}{\partial x} \right) = -\frac{1}{2} \frac{\partial^3 b_2^{(l)}}{\partial x^3} $$
これを (3.28) に代入し、全体を $-2$ 倍して符号と係数を整えると、まさに目的の3階微分方程式 (3.29) が姿を現します。
この計算により、2つの有理関数 $b_1, b_2$ を探すという問題が、**「$w = b_2^{(l)}$ というたった1つの有理関数が、作用素 $L(w) + 2\partial_{t_l}r = 0$ を満たすか？」**という極めて純粋で扱いやすい問題に帰着されたのです。
\end{Q}

\begin{story}
一方，2階線型常微分方程式
\begin{equation}
    \frac{d^2 y}{dx^2} + p_1(x,t)\frac{dy}{dx} + p_2(x,t)y = 0 \label{eq:3.33}
\end{equation}
において，1階線型微分方程式
\begin{equation}
    \frac{d\phi}{dx} + \frac{1}{2}p_1(x,t)\phi = 0 \label{eq:3.34}
\end{equation}
の解 $\phi = \phi(x,t)$ を1つとり，(3.33)の未知変数を
\begin{equation}
    y = \phi(x,t)z \label{eq:3.35}
\end{equation}
と変換する．このとき，$z$ は(3.25)の形の微分方程式を満足するが，その係数 $r=r(x,t)$ はまさに(3.32)で与えられる．

(3.35)は2つの微分方程式の解の基本系の関係であり，各々が定めるモノドロミー群は関数 $\phi$ の回路行列の作用――といってもスカラー倍されるだけであるが――の違いしかない．もし(3.33)がモノドロミー保存変形を許すとすると，$\phi$ の作用がパラメータ $t$ に依存しないならば，(3.25)もモノドロミー保存変形を許す．(3.25)と(3.35)の役割を入れ替えても同じである．ここで述べたことを，次のようにまとめておこう．

\begin{proposition}[3.7]
2階フックス型微分方程式(3.33)のモノドロミー保存変形と，2階フックス型微分方程式(3.25)のモノドロミー保存変形とは，1階フックス型微分方程式(3.34)がモノドロミー保存変形を許すという条件の下で，同値である．
\end{proposition}

実は，リーマン球面 $\mathbf{P}^1(\mathbf{C})$ 上定義された微分方程式の場合，命題3.7の条件は自動的に満たされている．実際，(3.34)を解くと，(3.24)より，解 $\phi$ は定数倍を除いて関数
\begin{equation}
    \phi(x, t) = \prod_{i=1}^M (x-\xi_i)^{-\frac{1}{2}\alpha_i} \label{eq:3.36}
\end{equation}
で与えられるから，この解のモノドロミー群は，$M$ 個のスカラー
\begin{equation*}
    e^{-\pi\sqrt{-1}\alpha_i}
\end{equation*}
で生成される可換群である．これは $t$ に依存しない．
\end{story}

\begin{Q}[【核心】このページ全体が主張している「一番嬉しいこと」は何か？]
\textbf{A.} 「1階微分の項 $p_1(x,t)y'$ を含む複雑な方程式 $y$ の変形を真面目に計算する必要は全くなく、見通しの良いSL型方程式 $z'' = r(x,t)z$ の変形だけを考えれば、理論的に完全に十分である」という強烈なお墨付きです。
ゲージ変換 $y = \phi z$ を行ったとき、両者のモノドロミーのズレは $\phi$ のモノドロミーだけです。もし $\phi$ のモノドロミーがパラメータ $t$ に依存してフラフラ動いてしまうと、$y$ と $z$ の同値性は崩壊します（命題3.7の但し書き）。しかし、舞台がリーマン球面 $P^1(\mathbb{C})$ であれば、その心配は一切無用であるということが、後半で証明されています。
\end{Q}

\begin{Q}[【深掘り】なぜ $\phi$ のモノドロミーは $t$ に依存しないと言い切れるのか？]
\textbf{A.} これこそ、まさにあなたが以前の考察で見抜いた事実です！
方程式 \eqref{eq:3.34} より、$\phi$ の解は \eqref{eq:3.36} のように各特異点 $\xi_i$ の周りで指数の形になります。この特異点を1周したときのモノドロミー（位相のズレ）は、$e^{2\pi i (-\alpha_i/2)} = e^{-\pi i \alpha_i}$ というスカラー倍になります。
ここで、$\alpha_i$ は元の2階方程式における係数 $p_1$ の留数です。「元の2階方程式のモノドロミーが保存される（特性指数が定数である）」という大前提に立っている以上、解と係数の関係から留数 $\alpha_i$ も絶対に定数でなければなりません。$\alpha_i$ が定数なら、当然 $e^{-\pi i \alpha_i}$ も定数です。
したがって、「$\phi$ がモノドロミー保存変形を許す」という命題の条件は、私たちが意図的に何かを仮定しなくても、$P^1(\mathbb{C})$ の構造と特性指数の設定から「自動的に満たされる（タダで手に入る）」のです。
\end{Q}

\begin{story}
$x$ の有理関数 $a_2^{(l)}$ が(3.29)を満たすとすると，有理関数 $a_1^{(l)}$ を(3.23)で定めれば，(3.33)のモノドロミー保存変形に関する，ラックスの方程式が成立する．一方，微分方程式(3.25)の変形微分方程式系(3.26)において，この $a_2^{(l)}$ を $b_2^{(l)}$ として，上述の議論を繰り返すと，(3.25)のモノドロミー保存変形のラックスの方程式も成立する．

なお，変形方程式系(3.20)の第2式から，(3.30)と同様に，積分可能条件として
\begin{align}
    \left( \frac{\partial}{\partial t_h} - a_2^{(h)}\frac{\partial}{\partial x} \right) a_1^{(l)} &= \left( \frac{\partial}{\partial t_l} - a_2^{(l)}\frac{\partial}{\partial x} \right) a_1^{(h)} \label{eq:3.37} \\
    \left( \frac{\partial}{\partial t_h} - a_2^{(h)}\frac{\partial}{\partial x} \right) a_2^{(l)} &= \left( \frac{\partial}{\partial t_l} - a_2^{(l)}\frac{\partial}{\partial x} \right) a_2^{(h)} \label{eq:3.38}
\end{align}
が得られる．ここで，(3.23)を考慮すると，(3.37)は
\begin{equation*}
    \left( \frac{\partial}{\partial x} - p_1 \right)\left( \frac{\partial}{\partial t_h} - a_2^{(h)}\frac{\partial}{\partial x} \right) a_2^{(l)} = \left( \frac{\partial}{\partial x} - p_1 \right)\left( \frac{\partial}{\partial t_l} - a_2^{(l)}\frac{\partial}{\partial x} \right) a_2^{(h)}
\end{equation*}
となる．したがって，(3.25)の完全積分可能条件は，$a_2^{(l)}$ に関するラックスの方程式(3.29)と条件(3.38)に帰着する．

他方，この章のはじめに述べた微分方程式(3.33)に関する仮定(P2), (P3)の下では，次の結果が成り立つ．

\begin{proposition}[3.8]
非斉次微分方程式(3.29)を満たす，$x$ の有理関数は存在するとしてもただ1つである
\end{proposition}

この命題の証明の前に，これから従う事実を整理しておこう．まず，$b_2^{(l)} = a_2^{(l)}$ でなければならない．結局，(3.33)の変形微分方程式系の完全積分可能条件は，条件(3.30)を満足する(3.29)の有理関数解 $b_2^{(l)}$ の存在に集約される．
以上により，微分方程式(3.33)のモノドロミー保存変形は，変換(3.35)を通して，微分方程式(3.25)のモノドロミー保存変形に帰着された．

\begin{proof}[命題3.8の証明]
$a^{(l)}$ と $b^{(l)}$ を(3.29)の解とし，$c^{(l)} = a^{(l)} - b^{(l)}$ とおくと，$c^{(l)} = c^{(l)}(x,t)$ は
\begin{equation}
    \frac{\partial^3 \phi}{\partial x^3} - 4r\frac{\partial \phi}{\partial x} - 2\frac{\partial r}{\partial x}\phi = 0 \label{eq:3.39}
\end{equation}
\end{proof}
\end{story}
（※画像末尾切れ）

\vspace{2em}

\begin{Q}[【核心】このページが主張する「SL型方程式(3.25)への帰着」の決定打は何か？]
\textbf{A.} **「有理関数解の一意性（ただ1つしかない）」**という事実です。
元の複雑な方程式(3.33)の変形を司る有理関数を $a_2^{(l)}$ とし、シンプルなSL型方程式(3.25)の変形を司る有理関数を $b_2^{(l)}$ とします。これらは共に、同じ形をした非斉次の微分方程式（ラックスの方程式）を満たすことが計算からわかります。
ここで「解は存在すればただ1つ」という命題3.8が威力を発揮します。これにより、「$a_2^{(l)}$ と $b_2^{(l)}$ は別物かもしれない」という疑いが完全に排除され、**$b_2^{(l)} = a_2^{(l)}$ であることが確定**します。両者の変形をつかさどる「心臓部」が全く同一であることが証明されたため、私たちは心置きなく、シンプルなSL型方程式の変形だけを考えればよい（帰着された）と結論づけることができるのです。
\end{Q}

\begin{Q}[【深掘り】命題3.8の証明に登場した (3.39) 式は、何を意味しているのか？]
\textbf{A.} これは、非斉次方程式の解の差 $c^{(l)}$ が満たす「斉次（右辺が0）の微分方程式」です。
そしてこの式の左辺を見てください。$\partial_x^3 - 4r\partial_x - 2r_x$……そう、これはあなたが以前ノートにまとめていた**KdV方程式のLax対の生成子 $M_3$ （あるいは3階の微分作用素 $L$）そのもの**です！
命題3.8の証明の筋書きはこうです。「もし2つの異なる解があったら、その差 $c^{(l)}$ はこの3階の斉次方程式 (3.39) を満たす有理関数になる。しかし、仮定(P2, P3)のような特異点の配置のもとでは、この $L(\phi)=0$ を満たす有理関数は $\phi=0$ （つまり $c^{(l)}=0$）しか存在し得ない。ゆえに解は1つである」。
可積分系の奥底に潜む「KdV的な3階の作用素」が、ここで「解の一意性を担保する」という極めて重要な役割を果たして再登場する、非常にドラマチックな展開です。
\end{Q}

ここで、元の2階微分方程式から1階微分の項を消去（正規化）することを考える。未知関数を変換することで、方程式を
\begin{equation}\label{eq:2-ode-normalized}
\frac{\partial^{2}z}{\partial x^{2}} = r(x,t)z
\end{equation}
という形に帰着させたい。この変換を実現するためには、新たな関数 $\phi(x,t)$ が
\begin{equation}\label{eq:1-ode-gauge}
\frac{\partial \phi}{\partial x} + \frac{1}{2}p_{1}(x,t)\phi = 0
\end{equation}
を満たせばよいことが知られている。

\begin{proposition}\label{prop:equivalence_2_ode}
    2階フックス型微分方程式 \eqref{eq:monodromy_deformation} のモノドロミー保存変形と、正規化された2階フックス型微分方程式 \eqref{eq:2-ode-normalized} のモノドロミー保存変形とは、1階微分方程式 \eqref{eq:1-ode-gauge} がモノドロミー保存変形を許すという条件のもとで同値である。
\end{proposition}
%#endregion

\subsection{変形微分方程式系の解}
%#region
前節の2階の場合の議論は、一般の $n$ 階フックス型微分方程式にも自然に拡張できる。

\begin{Q}[【核心】なぜわざわざ $n$ 階の変形微分方程式系に拡張して一般論を展開するのか？]
\textbf{A.} $n=2$ という特定の階数に縛られずに論理を展開する方が、背後にある線型代数的な構造（ロンスキアンの不変性や、モノドロミー行列の行列式が $1$ になる仕組み）が透けて見え、見通しが圧倒的に良くなるからである。この線型方程式の一般的な変形理論を通ることで、結果としてパンルヴェ方程式のような美しい「非線形な関係」が必然的に浮かび上がってくる。
\end{Q}

一般の $n$ 階フックス型微分方程式 \eqref{eq:monodromy_deformation} を考える。
方程式から $(n-1)$ 階微分の項を消去するため、
\begin{equation}\label{eq:auxiliary-phi}
\frac{\partial \phi}{\partial x} + \frac{1}{n}p_{1}(x,t)\phi = 0
\end{equation}
で定まる関数 $\phi(x,t)$ を用いて、未知関数を
\[
y = \phi(x,t)z
\]
と変換する。これを \eqref{eq:monodromy_deformation} に代入して整理すると、$z$ についての微分方程式は
\begin{equation}\label{eq:reduced-fuchsian}
\frac{\partial^{n}z}{\partial x^{n}} + \sum_{j=2}^{n} q_{j}(x,t)\frac{\partial^{n-j}z}{\partial x^{n-j}} = 0
\end{equation}
となる。ここで、$q_{j}(x,t)$ は $\phi$ の各階微分 $\frac{\partial^i \phi}{\partial x^i} \ (i=0,1,\cdots,j)$ と $p_{i}(x,t) \ (i=1,\cdots,j)$ の多項式として表される。この \eqref{eq:reduced-fuchsian} のリーマンデータを $(\mathbb{P}^{1}(\mathbb{C}), \Xi'(t), \hat{\rho}')$ とする。

\begin{Q}[【基礎】\eqref{eq:auxiliary-phi} を満たす関数 $\phi(x,t)$ は常に存在するのか？また、その解が多価関数になったり発散したりして元の問題を壊す心配はないか？]
\textbf{A.} 常に存在し、元の問題を壊すことはない。
\eqref{eq:auxiliary-phi} は変数分離形の1階線形微分方程式であるため、単に積分を実行するだけで
\[
\phi(x,t) = \exp\left( -\frac{1}{n} \int p_1(x,t) dx \right)
\]
と具体的に求まる。$p_1(x,t)$ は有理関数であるから、その積分は有理関数と対数関数の線形結合となる。したがって、$\phi(x,t)$ は特異点以外の全域で解析的であり、（多価にはなり得るが）その分岐の規則（モノドロミー）は完全にコントロール可能である。
\end{Q}

ここで、方程式 \eqref{eq:reduced-fuchsian} の解の基本系を $\vec{\varphi}_z(x,t)$ とする。この基本系に沿って閉曲線 $\gamma$ に沿って解析接続を行うと、モノドロミー行列 $\Gamma_z(\gamma)$ が右から掛かり、
\[
W(\vec{\varphi}_z^{\gamma}) = W(\vec{\varphi}_z) \det(\Gamma_z(\gamma))
\]
となる（$W$ はロンスキアン）。
一方で、\eqref{eq:reduced-fuchsian} は $(n-1)$ 階微分の項を持たない（係数が $0$ である）ため、命題\ref{prop:2.1_wronskian-ode}のアーベルの公式より、ロンスキアンの微分は
\[
\frac{\partial W}{\partial x} = -0 \cdot W = 0
\]
となる。すなわちロンスキアン $W(\vec{\varphi}_z)$ は $x$ に依存しないゼロでない定数である。
定数は解析接続しても変化しないため $W(\vec{\varphi}_z^{\gamma}) = W(\vec{\varphi}_z)$ となり、これを先ほどの式と比較して
\[
\det(\Gamma_z(\gamma)) = 1
\]
が得られる。
よって、\eqref{eq:reduced-fuchsian} のモノドロミー群は一般線形群 $GL(n, \mathbb{C})$ ではなく、特殊線形群 $SL(n, \mathbb{C})$ の部分群となることがわかった。すなわち、表現は
\[
\rho' : \pi_1(x_{0}, \mathbb{X}') \to SL(n, \mathbb{C})
\]
となる（ただし $\mathbb{X}' = \mathbb{P}^{1}(\mathbb{C}) \setminus (\Xi'(t) \cup \Lambda'(t))$）。

\begin{definition}[SL型の微分方程式]
    上記のように、$(n-1)$ 階微分の項が存在せず、モノドロミー群が $SL(n, \mathbb{C})$ の部分群となるような微分方程式を \textbf{SL型} という。
\end{definition}

\begin{proposition}\label{prop:equivalence_n_ode}
    $n$ 階フックス型微分方程式 \eqref{eq:monodromy_deformation} のモノドロミー保存変形と、SL型に還元された $n$ 階フックス型微分方程式 \eqref{eq:reduced-fuchsian} のモノドロミー保存変形とは、1階フックス型微分方程式 \eqref{eq:auxiliary-phi} がモノドロミー保存変形を許すという条件のもとで同値である。
\end{proposition}

\begin{proof}
パラメータ $t$ に依存する任意の閉じた道 $\gamma(t)$ に沿う解析接続を考える（解析接続と $t$ 微分は交換可能であることに注意する）。
元の方程式 \eqref{eq:monodromy_deformation} の解の基本系を $\vec{\varphi}_y(x,t)$、SL型方程式 \eqref{eq:reduced-fuchsian} の解の基本系を $\vec{\varphi}_z(x,t)$ とし、それぞれのモノドロミー行列を $\Gamma_{y}(\gamma(t))$, $\Gamma_{z}(\gamma(t))$ とする。また、1階方程式 \eqref{eq:auxiliary-phi} の解 $\phi(x,t)$ が $\gamma(t)$ に沿って解析接続されたときにかかるスカラー倍（1次元モノドロミー）を $m(\gamma(t))$ とする。

変換の定義より
\[
\vec{\varphi}_y(x,t) = \phi(x,t) \vec{\varphi}_z(x,t)
\]
であるから、これを $\gamma(t)$ に沿って解析接続すると、
\begin{align*}
\vec{\varphi}_y^{\gamma(t)}(x,t) &= \phi^{\gamma(t)}(x,t) \vec{\varphi}_z^{\gamma(t)}(x,t) \\
&= \left( \phi(x,t) m(\gamma(t)) \right) \left( \vec{\varphi}_z(x,t) \Gamma_{z}(\gamma(t)) \right) \\
&= \phi(x,t) \vec{\varphi}_z(x,t) \, m(\gamma(t)) \Gamma_{z}(\gamma(t)) \\
&= \vec{\varphi}_y(x,t) \, m(\gamma(t)) \Gamma_{z}(\gamma(t))
\end{align*}
となる（スカラー $m(\gamma(t))$ は行列と可換であるため前に出せる）。
一方で、定義より $\vec{\varphi}_y^{\gamma(t)} = \vec{\varphi}_y \Gamma_{y}(\gamma(t))$ であるから、両者を比較して
\[
\Gamma_{y}(\gamma(t)) = m(\gamma(t)) \Gamma_{z}(\gamma(t))
\]
という関係式が得られる。
ここで、仮定より1階方程式 \eqref{eq:auxiliary-phi} はモノドロミー保存変形を許すため、スカラー $m(\gamma(t))$ はパラメータ $t$ に依存しない一定値 $m(\gamma)$ である。
したがって、$\Gamma_{y}$ が $t$ に依存しないことと、$\Gamma_{z}$ が $t$ に依存しないことは完全に同値となる。
ゆえに、元の方程式のモノドロミー保存変形と、SL型方程式のモノドロミー保存変形は同値である。
\end{proof}
%#endregion
%#region
\eqref{eq:monodromy_deformation}の変形微分方程式系は\eqref{eq:3.1_deformation_system}で与えられた。これを、\eqref{eq:3.1_matrix_P}、\eqref{eq:3.1_matrix_A}の表示を用いて、連立変形微分方程式系に書き直す。この微分方程式系は、基本解行列 $Y=Y(x,t)$ を用いて

\begin{equation}
  \begin{cases}
    \displaystyle \frac{\partial Y}{\partial x} = P(x,t)Y \\
    \displaystyle \frac{\partial Y}{\partial t_l} = A^{(l)}(x,t)Y \quad (l=1,2,\cdots,L)
  \end{cases}
\end{equation}

と表されるのであった。いまの場合、基本解行列としては、\eqref{eq:monodromy_deformation}の解の基本系 $\vec{\varphi}(x,t)$ のロンスキー行列 $Y=W(\vec{\varphi})$ をとればよい。

\begin{Q}[なぜ基本解行列 $Y(x,t)$ の微分方程式系\eqref{eq:3.1_deformation_system_matrix}から、交換子を含む完全積分可能条件が導かれるのか？]
\textbf{A.} $\frac{\partial^2 Y}{\partial t_l \partial x} = \frac{\partial^2 Y}{\partial x \partial t_l}$ という偏微分の順序交換可能性（系の適合条件）を要求するためである。(3.6)の各々を他方の変数で微分して等置し、$Y$ が可逆（ロンスキアンが非零）であることを用いて右から $Y^{-1}$ を掛けることで、行列 $P, A^{(l)}$ に対する非線形な微分方程式であるゼロ曲率条件（ラックス方程式）が得られる。
\end{Q}

さて、微分方程式系\eqref{eq:3.1_deformation_system_matrix}の完全積分可能条件は、命題\ref{prop:3.1_compatibility}で与えた。すなわち

\begin{align}
  \frac{\partial}{\partial t_l}P(x,t) - \frac{\partial}{\partial x}A^{(l)}(x,t) &= [A^{(l)}(x,t), P(x,t)]\\
  \frac{\partial}{\partial t_l}A^{(h)}(x,t) - \frac{\partial}{\partial t_h}A^{(l)}(x,t) &= [A^{(l)}(x,t), A^{(h)}(x,t)]
\end{align}

である。

\begin{Q}[完全積分可能条件\eqref{eq:3.1_compatibility_PA}の両辺の行列について、そのトレース（跡）をとる操作は本質的に何を意味しているか？]
\textbf{A.} 行列の交換子のトレースは常に $0$ であるため（$\mathrm{Trace}[A, P] = 0$）、右辺の複雑な非線形項が消去される。これにより、行列方程式から系全体の「行列式の対数微分」の挙動を支配する簡潔なスカラー方程式を抽出することができる。
\end{Q}

ここで、\eqref{eq:3.1_compatibility_PA}の両辺の行列についてそのトレースをとり、\eqref{eq:3.1_matrix_P}を用いると

\begin{equation}\label{eq:3.3_trace_compatibility}
  \frac{\partial}{\partial t_l}p_1(x,t) + n\frac{\partial}{\partial x}c^{(l)}(x,t) = 0
\end{equation}
\[
  c^{(l)}(x,t) = \frac{1}{n}\mathrm{Trace}A^{(l)}(x,t) \quad (l=1,2,\cdots,L)
\]

を得る。この式は、1階微分方程式系

\begin{equation}\label{eq:3.3_1-deformation_system}
  \begin{cases}
    \displaystyle \frac{\partial \phi}{\partial x} + \frac{1}{n}p_1(x,t)\phi = 0 \\[4mm]
    \displaystyle \frac{\partial \phi}{\partial t_l} - c^{(l)}(x,t)\phi = 0
  \end{cases}
\end{equation}

の積分可能条件であり，\eqref{eq:3.3_trace_compatibility}は方程式系\eqref{eq:3.3_1-deformation_system}のラックスの方程式に他ならない．

さて，$\vec{\varphi}=\vec{\varphi}(x,t)$ を\eqref{eq:3.1_n-ODE}の解の基本系で，モノドロミー群が $t$ に依存しないもの，としよう．$\vec{\varphi}(x,t)$ は，変形微分方程式系\eqref{eq:3.1_deformation_system}の解の基本系である．このとき，微分方程式\eqref{eq:3.1_partial_t}の係数 $a_j^{(l)}(x,t)$ は公式\eqref{eq:3.1_cramer_wronskian}で与えられる．

$\vec{\psi}(x,t)$ を，\eqref{eq:3.1_n-ODE}の別の解の基本系であるとすると，$n$ 次可逆行列 $G(t)$ がとれて

\begin{equation}\label{eq:3.3_linear_transform_G}
  \vec{\psi}(x,t) = \vec{\varphi}(x,t)G(t)
\end{equation}

となる．我々は $t$ について解析的な解の基本系のみを考えているから，$G(t)$ の各成分は $t$ の解析関数である．さて，$\vec{\varphi}(x,t)$ の，ある閉じた道 $\gamma$ に対する回路行列を $\Gamma(\gamma)$ とすると，$\vec{\psi}(x,t)$ の対応する回路行列は $G(t)^{-1}\Gamma(\gamma)G(t)$ である．

\begin{Q}[なぜ $\vec{\psi}$ の回路行列は $G(t)^{-1}\Gamma(\gamma)G(t)$ となり、$G(t)=g(t)G$ の形であればモノドロミー群は $t$ に依存しなくなるのか？]
\textbf{A.} $\vec{\varphi}$ を $\gamma$ に沿って解析接続すると $\vec{\varphi}\Gamma(\gamma)$ となる。したがって $\vec{\psi} = \vec{\varphi}G(t)$ を解析接続すると $(\vec{\varphi}\Gamma(\gamma))G(t)$ となるが、これを $\vec{\psi}$ 自身を基底として表し直すために $\vec{\varphi} = \vec{\psi}G(t)^{-1}$ を代入すると、$\vec{\psi} (G(t)^{-1}\Gamma(\gamma)G(t))$ となるからである。
ここで $G(t) = g(t)G$（$g(t)$ はスカラー関数、$G$ は定数行列）と仮定すると、スカラー $g(t)$ は行列と可換であるため逆行列の $g(t)^{-1}$ とキャンセルし合い、回路行列は $G^{-1}\Gamma(\gamma)G$ となる。これは $t$ を含まない定数行列であるため、モノドロミー群は $t$ に依存しないことが保証される。
\end{Q}

特に，$g(t)$ を $t$ のスカラー関数，$G$ を $t$ に依存しない定数行列として

\begin{equation}\label{eq:3.3_scalar_matrix_gG}
  G(t) = g(t)G
\end{equation}

となっているとしよう．もちろん，考えている領域で $g(t) \neq 0$ である．このとき，\eqref{eq:3.3_linear_transform_G}で定義される解の基本系 $\vec{\psi}(x,t)$ のモノドロミー群も $t$ に依存しない．さらに，$\vec{\psi}(x,t)$ に対して

\[
  b_j^{(l)}(x,t) = \frac{w_j^{(l)}(\vec{\psi})}{w(\vec{\psi})}
\]

により，$x$ の有理関数 $b_j^{(l)}(x,t)$ を定めると，$\vec{\psi}(x,t)$ は，$b_j^{(l)}(x,t)$ を係数とする，変形微分方程式系の解である．ここで，$w_j^{(l)}(\vec{\psi})$ は，ロンスキアン $w(\vec{\psi})$ の第 $j$ 列ベクトル $\frac{\partial^{j-1}\vec{\psi}}{\partial x^{j-1}}$ を $\frac{\partial\vec{\psi}}{\partial t_l}$ で置き換えたものであった．上式の右辺に\eqref{eq:3.3_linear_transform_G}と\eqref{eq:3.3_scalar_matrix_gG}を代入すると，$j=2,\cdots,n$ については

\begin{align*}
  w_j^{(l)}(\vec{\psi}) &= g(t)^n (\det G) \cdot w_j^{(l)}(\vec{\varphi}) \quad (l=1,\cdots,L) \\
  w(\vec{\psi}) &= g(t)^n (\det G) \cdot w(\vec{\varphi})
\end{align*}

が成り立ち，$j \geqq 2$ のとき $a_j^{(l)}(x,t) = b_j^{(l)}(x,t)$ である．

\begin{Q}[$w_j^{(l)}(\vec{\psi})$ の計算において、なぜ $j \geqq 2$ のときのみ綺麗な比例関係が成り立ち、係数が不変（$a_j^{(l)} = b_j^{(l)}$）となるのか？]
\textbf{A.} 行列式（ロンスキアン）の多重線形性と交代性による。
$\vec{\psi} = g(t)\vec{\varphi}G$ を $t_l$ で偏微分すると、積の微分法により
$\frac{\partial \vec{\psi}}{\partial t_l} = \frac{\partial g(t)}{\partial t_l}\vec{\varphi}G + g(t)\frac{\partial \vec{\varphi}}{\partial t_l}G$
となる。これをロンスキアンの第 $j$ 列に代入して展開すると、第1項（$\partial_{t_l}g$ を含む項）は、「第1列（$\vec{\varphi}G$）」に比例するベクトルとなる。
$j \geqq 2$ の場合、行列式の中には既に第1列として $\vec{\varphi}G$ が存在するため、この比例するベクトルの寄与は行列式の交代性により完全に $0$ になる。結果として第2項のみが残り、$g(t)^n (\det G)$ が綺麗に括り出され、比をとると約分されて元の係数と一致する。
（逆に言えば、$j=1$ のときは第1列そのものを置き換えるため、この項が消えずに残り、係数が変化することを予見させる）。
\end{Q}

一方，$j=1$ のときは

\[
  w_1^{(l)}(\vec{\psi}) = g(t)^{n-1} \frac{\partial g}{\partial t_l}(t)(\det G) \cdot w(\vec{\varphi}) + g(t)^n (\det G) \cdot w_1^{(l)}(\vec{\varphi}) \quad (l=1,\cdots,L)
\]
となる．すなわち
\[
  b_1^{(l)}(x,t) = \frac{\partial}{\partial t_l} \log g(t) + a_1^{(l)}(x,t)
\]
である．

\begin{Q}[$j=1$ のときだけ係数が不変にならず、「対数微分 $\frac{\partial}{\partial t_l} \log g(t)$」のズレが生じることの理論的意義は何か？]
\textbf{A.} $j \geqq 2$ の係数 $a_j^{(l)}$ がゲージ変換に依存しない「真の不変量」であるのに対し、$a_1^{(l)}$ はゲージ（基本解のスケール因子 $g(t)$）の取り方に依存して変動する量であることを示している。このズレの項 $\partial_{t_l} \log g(t)$ こそが、方程式系の変形に伴う全体のスケール変化を司っており、後にパンルヴェ方程式等において「$\tau$ 関数の対数微分」として見かけの特異点のダイナミクスを記述する極めて重要な役割を果たす。
\end{Q}

実は，ここで述べたことの逆が成り立つ．

\begin{proposition}\label{prop:3.3_transform_G_implies_scalar_gG}
関係\eqref{eq:3.3_linear_transform_G}で結ばれている2つの解の基本系 $\vec{\varphi}(x,t), \vec{\psi}(x,t)$ のモノドロミー群がともに $t$ に依存しないならば，\eqref{eq:3.3_scalar_matrix_gG}（すなわち $G(t)=g(t)G$）が成立する．
\end{proposition}

この命題により，変形微分方程式系の係数
\[
  a_j^{(l)}(x,t) \quad (j=2,\cdots,n, \quad l=1,\cdots,L)
\]
は，微分方程式\eqref{eq:3.1_n-ODE}だけで決まり、モノドロミー群が $t$ に依存しない解の基本系の選び方には依らない．我々は $n=2$ の場合にこの事実を，ラックスの方程式を用いて確かめたのであった．

この節の残りを使って、命題\ref{prop:3.3_transform_G_implies_scalar_gG}を証明しよう。そのために3つの命題を準備する。まず、単独高階微分方程式\eqref{eq:3.1_n-ODE}を，\eqref{eq:3.1_matrix_P}、\eqref{eq:3.1_matrix_A}により連立微分方程式系に書き直し，変形微分方程式系\eqref{eq:3.1_deformation_system_matrix}を考えよう．$\vec{\varphi}(x,t)$ と $\vec{\psi}(x,t)$ を、命題\ref{prop:3.3_transform_G_implies_scalar_gG}の条件を満たす，\eqref{eq:3.1_n-ODE}の解の基本系とする．このとき，$\vec{\varphi}(x,t)$ に対し，$Y=W(\vec{\varphi})$ の満足する変形微分方程式系の係数を $A_1^{(l)}(x,t)$，同様に $\vec{\psi}(x,t)$ から決まる変形微分方程式系の係数を $A_2^{(l)}(x,t)$ とする。この2種類の行列 $A_1^{(l)}, A_2^{(l)}$ について，\eqref{eq:3.1_compatibility_PA}と\eqref{eq:3.1_compatibility_AA}が成立する．

\[
  B^{(l)}(x,t) = A_1^{(l)}(x,t) - A_2^{(l)}(x,t)
\]
とおくと，\eqref{eq:3.1_compatibility_PA}よりこれは行列微分方程式

\begin{equation}\label{eq:3.3_ODE_matrix_B}
  \frac{\partial}{\partial x} B^{(l)}(x,t) + \left[ B^{(l)}(x,t), P(x,t) \right] = 0
\end{equation}
を満たす．

\begin{Q}[なぜ2つの係数行列の差 $B^{(l)} = A_1^{(l)} - A_2^{(l)}$ を考え、方程式(3.49)を導出したのか？]
\textbf{A.} 非線形な完全積分可能条件（ラックス方程式）から、解析しやすい「線形な同次方程式」を抽出するためである。
$A_1^{(l)}, A_2^{(l)}$ は共に $\partial_{t_l} P - \partial_x A = [A, P]$ を満たす。これらを辺々引くと、左辺の $\partial_{t_l} P$ が完全に相殺されて消去され、$- \partial_x (A_1^{(l)} - A_2^{(l)}) = [A_1^{(l)} - A_2^{(l)}, P]$ となる。これにより、$B^{(l)}$ は $t$ 微分を含まない、$x$ 方向の線形な交換子方程式(3.49)を満たすことがわかり、その代数的な構造（ランクや成分の性質）を決定しやすくなる。
\end{Q}

ここで一般に，$Z$ を未知行列とし，行列微分方程式
\begin{equation}\label{eq:3.3_ODE_matrix_Z}
  \frac{dZ}{dx} + [Z, P(x)] = 0
\end{equation}

を考えよう．いま，$Y(x)$ を微分方程式
\begin{equation}\label{eq:3.3_ODE_matrix_Y}
  \frac{dY}{dx} = P(x)Y
\end{equation}
の基本解行列，$B_0$ を $x$ には依存しない可逆行列としたとき，\eqref{eq:3.3_ODE_matrix_Z}の一般解は
\[
  Z(x) = Y(x)B_0 Y(x)^{-1}
\]
で与えられる．これは計算で簡単に確かめることができるから，詳細は省略する．

\begin{Q}[「詳細は省略する」とあるが、一般解 $Z(x) = Y(x)B_0 Y(x)^{-1}$ が行列微分方程式 $\frac{dZ}{dx} + {[}Z, P{]} = 0$ を満たすことは、具体的にどう計算すれば確かめられるか？]
\textbf{A.} 基本解行列 $Y$ が $\frac{dY}{dx} = PY$ を満たすことと、逆行列の微分の公式 $\frac{d}{dx}(Y^{-1}) = -Y^{-1}\frac{dY}{dx}Y^{-1} = -Y^{-1}P$ を用いる。積の微分法則より、
\begin{align*}
  \frac{dZ}{dx} &= \frac{dY}{dx}B_0 Y^{-1} + Y B_0 \frac{d(Y^{-1})}{dx} \\[2mm]
  &= (PY)B_0 Y^{-1} + Y B_0 (-Y^{-1}P) \\[2mm]
  &= P(Y B_0 Y^{-1}) - (Y B_0 Y^{-1})P \\[2mm]
  &= PZ - ZP = -[Z, P]
\end{align*}
となり、移項すれば $\frac{dZ}{dx} + [Z, P] = 0$ を満たすことがわかる。
\end{Q}

この結果をいま考えている場合に適用すると，次のことがわかる．

\begin{lemma}\label{lem:3.3_solution_ODE}
行列微分方程式\eqref{eq:3.3_ODE_matrix_B}の一般解は，\eqref{eq:3.3_ODE_matrix_Y}の基本解行列 $Y=Y(x,t)$ により
\[
  B^{(l)}(x,t) = Y(x,t)B_0^{(l)}(t)Y(x,t)^{-1}
\]
で与えられる．ここで $B_0^{(l)}(t)$ は，成分が $t$ の解析関数である可逆行列である．
\end{lemma}

さらに，微分方程式\eqref{eq:3.1_n-ODE}に対する仮定(P3)の下で，次のことが成り立つ．

\begin{lemma}\label{lem:3.3_scalar_matrix_B}
$B_0^{(l)}(t)$ ($l=1,\cdots,L$) はスカラー行列である．
\end{lemma}

\begin{Q}[補題3.2において、なぜ仮定(P3)を置くだけで $B_0^{(l)}(t)$ が「スカラー行列」であると強力に結論づけられるのか？]
\textbf{A.} 表現論における「シューアの補題（Schur's lemma）」が本質的な役割を果たしている。
$B^{(l)}(x,t)$ は $x$ の有理関数行列（1価）であるため、特異点の周りを回る解析接続（モノドロミー変換 $Y \mapsto Y\Gamma$）に対して不変でなければならない。これを $B^{(l)} = Y B_0^{(l)} Y^{-1}$ に代入すると、$(Y\Gamma) B_0^{(l)} (Y\Gamma)^{-1} = Y B_0^{(l)} Y^{-1}$ となり、整理すると $\Gamma B_0^{(l)} = B_0^{(l)} \Gamma$ を得る。つまり $B_0^{(l)}$ はすべてのモノドロミー行列 $\Gamma$ と可換である。
ここで仮定(P3)によりモノドロミー表現が既約であるとすれば、シューアの補題より、「すべての表現行列と可換な行列はスカラー行列に限られる」ため、$B_0^{(l)}(t)$ はスカラー行列となる。
\end{Q}

したがって，スカラー関数 $f_0^{(l)}(t)$ により，$B_0^{(l)}(t) = f_0^{(l)}(t)I_n$ と表すことができる．$I_n$ は $n$ 次単位行列である．このとき，\eqref{eq:3.3_linear_transform_G}の行列 $G(t)$ について次のことが成り立つ．

\begin{lemma}\label{lem:3.3_t_PDE_matrix_G}
行列 $G(t)$ は微分方程式系
\begin{equation}\label{eq:3.3_t_PDE_matrix_G}
  \frac{\partial}{\partial t_l}G(t) = -f_0^{(l)}(t)G(t) \quad (l=1,\cdots,L)
\end{equation}
を満たす．この方程式系は完全積分可能である．
\end{lemma}

補題\ref{lem:3.3_scalar_matrix_B}、補題\ref{lem:3.3_t_PDE_matrix_G}を仮定すれば、命題\ref{prop:3.3_transform_G_implies_scalar_gG}の証明は簡単である．

\begin{proof}[命題\ref{prop:3.3_transform_G_implies_scalar_gG}の証明]
微分方程式\eqref{eq:3.3_t_PDE_matrix_G}の完全積分可能条件から
\begin{equation}\label{eq:3.3_t-ODE_scalar_g}
  \frac{\partial}{\partial t_l}g(t) = -f_0^{(l)}(t)g(t) \quad (l=1,\cdots,L)
\end{equation}
というスカラー関数 $g(t)$ が存在する．この関数を用いて，\eqref{eq:3.3_t_PDE_matrix_G}の一般解は……

\end{proof}
%#endregion

\subsection{シュレージンガー系}

%#region
まず、$n$次連立微分方程式系\eqref{eq:3.1_x_ODE_matrix}のモノドロミー保存変形についてこれまでわかったことをまとめよう。微分方程式系\eqref{eq:3.1_x_ODE_matrix}のリーマンデータを
\[
  (\mathbb{P}^1(\mathbb{C}), \Xi', \hat{\rho}')
\]
とし、条件(P2), (P3)を仮定する。

微分方程式系\eqref{eq:3.1_x_ODE_matrix}の基本解行列 $Y=Y(x,t)$ で、そのモノドロミー群がパラメータ$t=(t_1, \cdots, t_L)$ に依存しないものが存在するならば，$x$ の有理関数を成分とする行列 $A^{(l)}(x,t)$ ($l=1,\cdots,L$) が存在して，$Y(x,t)$ は変形微分方程式系\eqref{eq:3.1_deformation_system_matrix}の解となる．この微分方程式系の完全積分可能条件は，\eqref{eq:3.1_compatibility_PA}と\eqref{eq:3.1_compatibility_AA}である。\eqref{eq:3.1_x_ODE_matrix}のモノドロミー保存変形はこの条件の研究に帰着する。前節最後に述べた命題\ref{prop:3.3_transform_G_implies_scalar_gG}、補題\ref{lem:3.3_scalar_matrix_B}、補題\ref{lem:3.3_t_PDE_matrix_G}によれば，\eqref{eq:3.1_compatibility_PA}、\eqref{eq:3.1_compatibility_AA}を満足する $x$ の有理関数を成分とする行列が2種類、$A_1^{(l)}$ と $A_2^{(l)}$、存在するときには，スカラー関数 $g(t)$ が存在して
\[
  A_2^{(l)}(x,t) = A_1^{(l)}(x,t) + \left( \frac{\partial}{\partial t_l} \log g(t) \right) I_n \quad (l=1,\cdots,L)
\]
となる．

\begin{Q}[【補足と確認】補題\ref{lem:3.3_scalar_matrix_B}、補題\ref{lem:3.3_t_PDE_matrix_G}、および命題\ref{prop:3.3_transform_G_implies_scalar_gG}の証明は、$P(x,t)$ が単独高階方程式から作られる特殊な行列（コンパニオン行列）でなくとも、本当に一般の $n \times n$ 行列で成り立つのか？ その厳密な証明は？]
\textbf{A.} \textbf{完全に一般の行列 $P(x,t)$ に対して成り立つ。}
証明のどの段階においても $P(x,t)$ の成分の具体的な形は要求されず、「既約なモノドロミー表現（シューアの補題）」と「行列の積の微分法則」という普遍的な性質のみで完結する。以下に、一般の行列に対する完全な証明プロセスを示す。

\textbf{1. 準備と線形化（補題3.1の領域）}
一般の $\frac{\partial Y}{\partial x} = P(x,t)Y$ に対して、モノドロミーが $t$ に依存しない2つの基本解行列 $Y_1, Y_2$ があるとする。これらはある可逆行列 $G(t)$ を用いて $Y_2(x,t) = Y_1(x,t)G(t)$ と結ばれる。
それぞれの変形行列 $A_k^{(l)} = \frac{\partial Y_k}{\partial t_l}Y_k^{-1} \ (k=1,2)$ は、ともにゼロ曲率条件 $\partial_{t_l} P - \partial_x A_k^{(l)} = [A_k^{(l)}, P]$ を満たす。辺々引くことで、$P$ の形に依らず
\[
  \frac{\partial}{\partial x}(A_1^{(l)} - A_2^{(l)}) + [A_1^{(l)} - A_2^{(l)}, P] = 0
\]
となる。差を $B^{(l)} = A_1^{(l)} - A_2^{(l)}$ とおくと、この微分方程式の一般解は $x$ に依存しない行列 $B_0^{(l)}(t)$ を用いて $B^{(l)}(x,t) = Y_1 B_0^{(l)}(t) Y_1^{-1}$ と書ける。

\textbf{2. 補題3.2（シューアの補題の適用）}%TODO あとでみる
$A_1^{(l)}, A_2^{(l)}$ は有理関数行列（1価）なので、その差 $B^{(l)}$ も1価である。特異点の周りの任意の閉路に対する解析接続（モノドロミー行列を $\Gamma$ とする）を考えると、$Y_1 \mapsto Y_1 \Gamma$ となる。$B^{(l)}$ は不変であるから、
\[
  Y_1 B_0^{(l)}(t) Y_1^{-1} = (Y_1 \Gamma) B_0^{(l)}(t) (Y_1 \Gamma)^{-1} = Y_1 (\Gamma B_0^{(l)}(t) \Gamma^{-1}) Y_1^{-1}
\]
両端から $Y_1^{-1}$ と $Y_1$ を掛けて $\Gamma B_0^{(l)}(t) = B_0^{(l)}(t) \Gamma$ を得る。すなわち $B_0^{(l)}(t)$ はすべてのモノドロミー行列と可換である。
ここで\textbf{仮定(P3)（モノドロミー表現が既約である）}を適用すれば、シューアの補題より直ちに $B_0^{(l)}(t) = f_0^{(l)}(t)I_n$ （スカラー行列）であることが結論づけられる。ここにも $P$ の形は一切登場しない。したがって $B^{(l)}(x,t) = f_0^{(l)}(t)I_n$ である。

\textbf{3. 補題3.3の証明（ゲージ変換の直接計算）}
一方で、$Y_2 = Y_1 G(t)$ を用いて $A_2^{(l)}$ を直接計算する。
\begin{align*}
  A_2^{(l)} &= \frac{\partial (Y_1 G)}{\partial t_l} (Y_1 G)^{-1} 
  = \left( \frac{\partial Y_1}{\partial t_l} G + Y_1 \frac{\partial G}{\partial t_l} \right) G^{-1} Y_1^{-1} \\
  &= \frac{\partial Y_1}{\partial t_l} Y_1^{-1} + Y_1 \frac{\partial G}{\partial t_l} G^{-1} Y_1^{-1} 
  = A_1^{(l)} + Y_1 \left( \frac{\partial G}{\partial t_l} G^{-1} \right) Y_1^{-1}
\end{align*}
定義より $B^{(l)} = A_1^{(l)} - A_2^{(l)}$ であるから、上式より $B^{(l)} = - Y_1 \left( \frac{\partial G}{\partial t_l} G^{-1} \right) Y_1^{-1}$ となる。

\textbf{4. 命題3.10の結論}
補題3.2の結果 $B^{(l)} = f_0^{(l)}(t)I_n$ と、ステップ3の計算結果を等置する。
\[
  - Y_1 \left( \frac{\partial G}{\partial t_l} G^{-1} \right) Y_1^{-1} = f_0^{(l)}(t)I_n
\]
両辺に左から $Y_1^{-1}$、右から $Y_1$ を掛けると（右辺はスカラー行列なので不変）、
\[
  - \frac{\partial G}{\partial t_l} G^{-1} = f_0^{(l)}(t)I_n \implies \frac{\partial G}{\partial t_l} = -f_0^{(l)}(t)G
\]
これが補題3.3である。この微分方程式系を積分すれば、$G(t) = g(t)G$ （$g(t)$ はスカラー関数、$G$ は定数行列）が得られ、命題3.10が完全に証明される。
\end{Q}

\begin{Q}[$A_1^{(l)}$ と $A_2^{(l)}$ の差が対数微分 $\left( \frac{\partial}{\partial t_l} \log g(t) \right) I_n$ となるこの美しい関係は、$P(x,t)$ が単独高階方程式から作られるコンパニオン行列である場合に限定されるのか？一般の行列 $P(x,t)$ でも成り立つか？]
\textbf{A.} \textbf{一般の行列 $P(x,t)$ でも完全に成り立つ。}
この性質は $P$ の具体的な形ではなく、「モノドロミー表現が既約である（仮定P3）」という条件からシューアの補題を経由して得られる、極めて普遍的な帰結である。

実際、一般の $\frac{\partial Y}{\partial x} = P(x,t) Y$ に対する2つの基本解行列を $Y_1, Y_2$ とすると、$Y_2 = Y_1 G(t)$ と書ける。両者が $t$ に依存しない既約なモノドロミー群を持つとき、命題\ref{prop:3.3_transform_G_implies_scalar_gG}により $G(t) = g(t)G$（$g(t)$ はスカラー関数、$G$ は定数可逆行列）となる。
これを変形行列の定義 $A_2^{(l)} = \frac{\partial Y_2}{\partial t_l}Y_2^{-1}$ に直接代入すると、積の微分法則より
\begin{align*}
  A_2^{(l)} &= \frac{\partial}{\partial t_l}\left(Y_1 g(t)G\right) \cdot \left(Y_1 g(t)G\right)^{-1} \\
  &= \left( \frac{\partial Y_1}{\partial t_l} g(t)G + Y_1 \frac{\partial g(t)}{\partial t_l} G \right) G^{-1} g(t)^{-1} Y_1^{-1} \\
  &= \frac{\partial Y_1}{\partial t_l} Y_1^{-1} + Y_1 \left( \frac{\frac{\partial g(t)}{\partial t_l}}{g(t)} I_n \right) Y_1^{-1} \\
  &= A_1^{(l)} + \left( \frac{\partial}{\partial t_l} \log g(t) \right) I_n
\end{align*}
となり、計算過程で $P(x,t)$ に関する情報は一切使用されない。したがって、一般のフックス型連立系においてもこのゲージ変換の公式は成立する。
\end{Q}

この節では，具体的な例について，\eqref{eq:3.1_x_ODE_matrix}のモノドロミー保存変形を調べよう．

まず，有理関数を成分とする行列 $A^{(l)}(x,t)$ の形を決めることを考えてみよう．$x=\xi$ を確定特異点とすると，基本解行列 $Y(x,t)$ は，条件(P2)により
\begin{equation}\label{eq:3.4_matrix_solution_representation_Phi_G}
  Y(x,t) = \Phi(x,t)(x-\xi)^A G(t)
\end{equation}
と表される．ここで，$A$ は対角行列で，$t$ に依存しない．また，$\Phi=\Phi(x,t)$ の各成分は $x=\xi$ の近傍において，$x$ と $t$ について正則，$G(t)$ は各成分が $U$ 上正則である可逆行列である．$Y=Y(x,t)$ の表示\eqref{eq:3.4_matrix_solution_representation_Phi_G}については，単独高階微分方程式の場合と同様に求めることができるから，ここでは既知として話を進める．

\begin{Q}[係数行列 $P(x,t)$ が特異点 $x=\xi$ において高々1位の極（フックス型特異点）を持つと仮定したとき、基本解行列が局所表示 $Y(x,t) = \Phi(x,t)(x-\xi)^A G(t)$ を持つことは、具体的にどのように証明（構成）されるか？]
\textbf{A.} 行列に対する「フロベニウスの方法」と「解の基底の取り換え」を用いて、以下の4つのステップで構成される。

\textbf{ステップ1：係数行列のローラン展開}
$x=\xi$ において $P(x,t)$ は1位の極を持つため、その近傍で次のように展開できる。
\[
  P(x,t) = \frac{P_{-1}(t)}{x-\xi} + \sum_{k=0}^{\infty} P_k(t) (x-\xi)^k
\]
ここで、$P_{-1}(t)$ は $x=\xi$ における留数行列である。

\textbf{ステップ2：局所的な基本解行列 $\tilde{Y}$ の仮定（Ansatz）}
$x=\xi$ の近傍におけるある1つの局所的な解行列を、特異性を持つ部分 $(x-\xi)^A$ と、正則な部分 $\Phi(x,t)$ に分離して探す。
\[
  \tilde{Y}(x,t) = \Phi(x,t)(x-\xi)^A, \quad \Phi(x,t) = \sum_{k=0}^{\infty} \Phi_k(t) (x-\xi)^k
\]
ここで、正則性を保証するために $\Phi_0(t)$ は可逆行列（簡単のため $\Phi_0(t) = I$ と選んでよい）とし、$A$ はこれから決定する定数行列とする。

\textbf{ステップ3：代入と行列 $A$ の決定}
仮定した $\tilde{Y}(x,t)$ を微分方程式 $\frac{\partial \tilde{Y}}{\partial x} = P(x,t)\tilde{Y}$ に代入する。

\begin{Q}[行列のべき乗 $(x-\xi)^A$ の $x$ による微分は、スカラーの場合と同様に $A(x-\xi)^{A-I}$ となるか？ また、指数関数をマクローリン展開して項別微分する際、この無限級数の収束性（収束半径）に問題はないのか？]
\textbf{A.} スカラーの場合と全く同じ結果となる。収束性についても問題ない。
行列の指数関数 $\exp(X) = \sum_{m=0}^{\infty} \frac{1}{m!} X^m$ は、全空間で絶対収束する（収束半径は $\infty$ である）。いま $X = A \log(x-\xi)$ とおくと、$x \neq \xi$ であり対数関数の分枝を一つ固定する限り、展開された級数は常に収束するため項別微分は完全に正当化される。

具体的な計算は以下の通りである。行列 $A$ と単位行列 $I$ が可換である（$e^X e^Y = e^{X+Y}$ が使える）ことに注意して整理する。

\begin{align*}
  (x-\xi)^A &= \exp(A \log(x-\xi)) = \sum_{m=0}^{\infty} \frac{1}{m!} A^m (\log(x-\xi))^m \\[3mm]
  \frac{d}{dx} (x-\xi)^A &= \sum_{m=1}^{\infty} \frac{1}{(m-1)!} A^m (\log(x-\xi))^{m-1} \cdot \frac{1}{x-\xi} \\[2mm]
  &= A e^{A \log(x-\xi)} \cdot \frac{1}{x-\xi} \\[2mm]
  &= A e^{A \log(x-\xi)} \cdot e^{-\log(x-\xi)} \\[2mm]
  &= A e^{(A-I) \log(x-\xi)}
\end{align*}
すなわち、$\frac{d}{dx}(x-\xi)^A = A(x-\xi)^{A-I}$ が厳密に成立する。
\end{Q}

左辺は積の微分法則より、
\begin{align*}
  \frac{\partial \tilde{Y}}{\partial x} &= \frac{\partial \Phi}{\partial x}(x-\xi)^A + \Phi(x,t) A (x-\xi)^{A-I} \\
  &= \left( \frac{\partial \Phi}{\partial x}(x-\xi) + \Phi(x,t)A \right) (x-\xi)^{A-I}
\end{align*}
右辺は、
\[
  P(x,t)\tilde{Y} = \left( P_{-1}(t) + P_0(t)(x-\xi) + \cdots \right) \Phi(x,t) (x-\xi)^{A-I}
\]
となる。両辺から $(x-\xi)^{A-I}$ を括り出し、定数項（$(x-\xi)^0$ の係数）を比較すると、
\[
  \Phi_0(t) A = P_{-1}(t) \Phi_0(t)
\]
という行列の決定方程式が得られる。$\Phi_0(t) = I$ と選んでいるため、$A = P_{-1}(t)$ となる。
さらに、テキストの条件(P2)（留数行列の固有値の差が整数でない）により、高次の係数 $\Phi_k(t)$ を決定する漸化式が常に解ける（共鳴が起きない）ことが保証され、対数項を含まない正則な $\Phi(x,t)$ が一意に構成できる。また、モノドロミー保存の仮定から $P_{-1}(t)$ の固有値は $t$ に依存しないため、$A$ は $t$ に依存しない対角行列として取ることができる。

\begin{Q}[高次の係数 $\Phi_k(t)$ を決定する漸化式とは具体的にどのような形か？また、条件(P2)が「共鳴（対数項の混入）を防ぐ」ことの厳密な数学的証明は、何を用いれば示せるか？]%TODO 詳細は後で見る
\textbf{A.} 行列のべき級数展開から得られる「シルベスター方程式（Sylvester equation）」の可解条件を用いることで厳密に証明される。

\textbf{ステップ1：漸化式の導出}
$P(x,t) = \sum_{m=-1}^\infty P_m (x-\xi)^m$ と $\Phi(x,t) = \sum_{k=0}^\infty \Phi_k (x-\xi)^k$ を $\frac{\partial \tilde{Y}}{\partial x} = P\tilde{Y}$ に代入し、両辺の $(x-\xi)^{k-1} (x-\xi)^A$ の係数を比較する。
左辺の微分からの寄与は $k\Phi_k + \Phi_k A$ であり、右辺の積からの寄与は $\sum_{j=0}^k P_{j-1}\Phi_{k-j}$ となる。$j=0$ の項（$P_{-1}\Phi_k$）を分離し整理すると、次の方程式を得る。
\[
  P_{-1}\Phi_k - \Phi_k A - k\Phi_k = -\sum_{j=1}^k P_{j-1}\Phi_{k-j}
\]

\textbf{ステップ2：シルベスター方程式としての解釈}
右辺を $R_k$ とおく（$R_k$ は $\Phi_0, \dots, \Phi_{k-1}$ までで既に決定されている既知の行列である）。この方程式は、未知行列 $X = \Phi_k$ に対する線形代数方程式
\[
  M X - X N = -R_k \quad (\text{ただし } M = P_{-1} - k I, \ N = A)
\]
と見なせる。このような形の方程式をシルベスター方程式と呼ぶ。一般論として、この方程式が一意な解を持つ必要十分条件は「行列 $M$ と行列 $N$ が共通の固有値を持たないこと」である。

\textbf{ステップ3：条件(P2)による可逆性の保証}
$\Phi_0 = I$ のとき、定数項の比較から $A = P_{-1}$ である。
$P_{-1}$ の固有値（特性指数）を $\lambda_1, \dots, \lambda_n$ とすると、行列 $M = P_{-1} - k I$ の固有値は $\lambda_i - k$ であり、行列 $N = A$ の固有値は $\lambda_j$ である。
これらが共通の固有値を持たないための条件は、任意の $i, j$ に対して
\[
  \lambda_i - k \neq \lambda_j \iff \lambda_i - \lambda_j \neq k
\]
となることである。ここで $k \geq 1$ は整数であるため、「特性指数の差が整数でない」という\textbf{条件(P2)}が満たされていれば、この条件はすべての $k \geq 1$ に対して自動的にクリアされる。
すなわち、線形作用素 $X \mapsto MX - XN$ は常に可逆（正則）となり、各ステップで $\Phi_k$ が一意に定まる。これにより、特異点において対数項 $\log(x-\xi)$ の補正を一切必要とせずに、純粋な級数解が構成できることが厳密に示された。

\textbf{ステップ4：$A$ の対角化}
(P2)により特性指数はすべて異なるため、留数行列 $P_{-1}(t)$ は常に対角化可能である。基本解の取り替え（適当な定数可逆行列によるゲージ変換）を行えば、初めから $A = P_{-1}(t)$ を対角行列として選んでおくことができる。さらにモノドロミー保存の仮定により固有値は $t$ に依存しないため、$A$ は $t$ に依存しない定数対角行列となる。
\end{Q}

\begin{Q}[行列 $A$ が対角行列でない一般の行列である場合、解の局所表示に現れる行列のべき乗 $(x-\xi)^A$ は数学的にどのように定義されるのか？]
\textbf{A.} 行列の指数関数（テイラー展開）と自然対数を用いて、次のように厳密に定義される。
\[
  (x-\xi)^A := \exp(A \log(x-\xi)) = \sum_{m=0}^{\infty} \frac{1}{m!} A^m (\log(x-\xi))^m
\]
もし $A$ が対角行列 $\mathrm{diag}(\lambda_1, \dots, \lambda_n)$ であれば、この行列のべき乗は各成分を計算するだけで済み、直感通り $\mathrm{diag}((x-\xi)^{\lambda_1}, \dots, (x-\xi)^{\lambda_n})$ に一致する。

しかし、特性指数の差が整数になるなどして $A$ が対角化不可能となり、ジョルダン標準形に非対角成分を含む場合、この定義から必然的に「対数特異性」が生じる。
例えば、2次のジョルダン細胞 $A = \begin{pmatrix} \lambda & 1 \\ 0 & \lambda \end{pmatrix}$ の場合、行列の指数関数の性質（可換な行列の和の指数法則）を用いると、
\begin{align*}
  (x-\xi)^A &= \exp\left( \begin{pmatrix} \lambda & 1 \\ 0 & \lambda \end{pmatrix} \log(x-\xi) \right) \\
  &= \exp\left( \begin{pmatrix} \lambda \log(x-\xi) & 0 \\ 0 & \lambda \log(x-\xi) \end{pmatrix} + \begin{pmatrix} 0 & \log(x-\xi) \\ 0 & 0 \end{pmatrix} \right) \\
  &= (x-\xi)^\lambda I \cdot \exp\left( \begin{pmatrix} 0 & \log(x-\xi) \\ 0 & 0 \end{pmatrix} \right)
\end{align*}
ここで、右側の行列は2乗すると零行列になる（べき零行列である）ため、指数関数の級数展開は $m=1$ の項で打ち切られ、
\begin{align*}
  (x-\xi)^A &= (x-\xi)^\lambda \left( I + \begin{pmatrix} 0 & \log(x-\xi) \\ 0 & 0 \end{pmatrix} \right) \\
  &= (x-\xi)^\lambda \begin{pmatrix} 1 & \log(x-\xi) \\ 0 & 1 \end{pmatrix}
\end{align*}
となる。このように、一般の行列に対するべき乗の定義そのものが、解に $\log(x-\xi)$ という対数項を発生させる代数的なからくりとなっている。
\end{Q}

\textbf{ステップ4：一般の基本解行列 $Y(x,t)$ への移行}
ステップ3までで、微分方程式を満たすある特定の基本解行列 $\tilde{Y}(x,t) = \Phi(x,t)(x-\xi)^A$ が得られた。
線形常微分方程式の理論より、任意の基本解行列 $Y(x,t)$ は、$\tilde{Y}(x,t)$ に $x$ に依存しない可逆な定数行列 $G(t)$ を右から掛けたものとして一意に表される。
したがって、
\[
  Y(x,t) = \tilde{Y}(x,t)G(t) = \Phi(x,t)(x-\xi)^A G(t)
\]
となり、求める局所表示\eqref{eq:3.4_matrix_solution_representation_Phi_G}が完全に導出された。
\end{Q}

\eqref{eq:3.4_matrix_solution_representation_Phi_G}を\eqref{eq:3.1_deformation_system_matrix}の第2式（$\partial_{t_l} Y = A^{(l)}Y$）に代入すると

\begin{align}\label{eq:3.4_matrix_representation_A_Phi_G}
  A^{(l)}(x,t) &= \frac{\partial Y}{\partial t_l}Y^{-1}\\
  &= \frac{\partial \Phi}{\partial t_l}\Phi^{-1} - \frac{1}{x-\xi}\frac{\partial \xi}{\partial t_l}\Phi A \Phi^{-1} + \Phi(x-\xi)^A \frac{\partial G}{\partial t_l}G^{-1}(x-\xi)^{-A}\Phi^{-1}\nonumber
\end{align}
となる．

\begin{Q}[なぜ特異点 $x=\xi$ の周りで局所表示\eqref{eq:3.4_matrix_solution_representation_Phi_G}を考え、複雑な\eqref{eq:3.4_matrix_representation_A_Phi_G}の計算を行う必要があるのか？ また、\eqref{eq:3.4_matrix_representation_A_Phi_G}は具体的にどのように導出されるか？]
\textbf{A.} \textbf{【目的】} $A^{(l)}(x,t)$ が $x$ の有理関数行列であることを利用し、各特異点での「極の位数」と「主要部（留数）」を決定するため。これらを足し合わせることで $A^{(l)}$ の全体像（大域的な形）を決定でき、それがシュレジンガー系の具体的な方程式となる。

\textbf{【導出】} 積の微分法則を用いる。$Y = \Phi (x-\xi)^A G$ と $Y^{-1} = G^{-1} (x-\xi)^{-A} \Phi^{-1}$ を $\frac{\partial Y}{\partial t_l}Y^{-1}$ に代入する。
$\frac{\partial Y}{\partial t_l}$ は3つの項の和になる。\\
1. $\Phi$ の微分: $\frac{\partial \Phi}{\partial t_l} (x-\xi)^A G$\\
2. $(x-\xi)^A$ の微分: チェインルールより $\Phi \left( - \frac{\partial \xi}{\partial t_l} A (x-\xi)^{A-I} \right) G = - \frac{1}{x-\xi} \frac{\partial \xi}{\partial t_l} \Phi A (x-\xi)^A G$\\
3. $G$ の微分: $\Phi (x-\xi)^A \frac{\partial G}{\partial t_l}$
これら各項の右から $Y^{-1}$ を掛けると、右側の $(x-\xi)^A G$ 等が相殺・変形され、それぞれ\eqref{eq:3.4_matrix_representation_A_Phi_G}の第1項、第2項（極を持つ主要部）、第3項に一致する。
\end{Q}

この表示から，次の結果が得られる．詳しい説明は省略する．

\begin{proposition}\label{prop:3.4_matrix_ODE_monodoromy_deformation_property}
(1) 行列 $A^{(l)}(x,t)$ は，微分方程式\eqref{eq:3.1_x_ODE_matrix}の正則点，すなわち $P(x,t)$ の正則点，では正則となる．\\
(2) \eqref{eq:3.1_x_ODE_matrix}の確定特異点 $x=\xi$ においては，もし $\xi$ が $t_l$ に依存しなければ，$x=\xi$ で $A^{(l)}(x,t)$ は正則となる．\\
(3) 確定特異点の位置 $\xi$ が $t_l$ に依存すれば，$x=\xi$ は $A^{(l)}(x,t)$ の1位の極である。
\end{proposition}

\begin{Q}[命題\ref{prop:3.4_matrix_ODE_monodoromy_deformation_property}の証明について：(1)正則点での振る舞いは式\eqref{eq:3.4_matrix_representation_A_Phi_G}に $P(x,t)$ が出てこないのにどう証明するのか？ また、(2)(3)において、式\eqref{eq:3.4_matrix_representation_A_Phi_G}の第1項や第3項が極を持ったり、第2項と打ち消し合ったりしないとどうして言い切れるのか？]
\textbf{A.} (1)は式\eqref{eq:3.4_matrix_representation_A_Phi_G}を使わず「常微分方程式の基本定理」から直接導かれ、(2)(3)は「行列の可換性」によって第3項の特異性が完全に消滅することで証明される。

\textbf{【(1) 正則点での正則性の証明】}
式\eqref{eq:3.4_matrix_representation_A_Phi_G}は特異点 $x=\xi$ の近傍でのみ作られた局所表示であるため、正則点では使えない。正則点 $x=x_0$ においては、定義 $A^{(l)} = \frac{\partial Y}{\partial t_l}Y^{-1}$ に立ち返る。
$x_0$ が $P(x,t)$ の正則点であれば、常微分方程式の基本定理（初期値問題の解の解析性）により、基本解行列 $Y(x,t)$ は $x=x_0$ の近傍で $x$ および $t$ について正則（解析的）であり、かつ可逆である。
したがって、その $t_l$ 微分 $\frac{\partial Y}{\partial t_l}$ も正則であり、正則な行列と正則な逆行列の積である $A^{(l)}$ もまた $x=x_0$ において正則となる。

\textbf{【(2)(3) 確定特異点での極の評価】}
確定特異点 $x=\xi$ の近傍では式\eqref{eq:3.4_matrix_representation_A_Phi_G}を用いる。各項の特異性（極の有無）を個別に評価する。
\[
  A^{(l)} = \underbrace{\frac{\partial \Phi}{\partial t_l}\Phi^{-1}}_{\text{第1項}} - \underbrace{\frac{1}{x-\xi}\frac{\partial \xi}{\partial t_l}\Phi A \Phi^{-1}}_{\text{第2項}} + \underbrace{\Phi(x-\xi)^A \frac{\partial G}{\partial t_l}G^{-1}(x-\xi)^{-A}\Phi^{-1}}_{\text{第3項}}
\]

\begin{itemize}
    \item \textbf{第1項の評価：} $\Phi(x,t)$ は $x=\xi$ で正則かつ可逆であるため、その微分も正則。よって第1項は完全に正則である。
    
    \item \textbf{第3項の評価（最大のポイント）：} 
    一見すると $(x-\xi)^{\pm A}$ があるため極や分岐を持ちそうに見えるが、実は相殺される。
    ゲージ行列 $G(t)$ はモノドロミーを不変に保つため、特異点におけるモノドロミー行列 $\exp(2\pi i A)$ と可換でなければならない。特性指数の差が整数でない（条件P2）ため、これは $G(t)$ が $A$ 自身と可換であることを意味する。
    $G(t)$ が $A$ と可換なら、$\frac{\partial G}{\partial t_l}G^{-1}$ も $A$ と可換になる。行列の指数関数の性質から、可換な行列は $(x-\xi)^A$ と順序を交換できる。
    
    \begin{Q}[前項の証明において、「モノドロミーを不変に保つこと」からなぜ可換性が生じるのか？ また、「特性指数の差が整数でない（条件P2）」ことが、なぜ $\exp(2\pi i A)$ との可換性を $A$ 自身との可換性にすり替えることを正当化するのか？]
\textbf{A.} 厳密には、$G(t)$ 自身ではなくその対数微分 $U = \frac{\partial G}{\partial t_l}G^{-1}$ が可換になる。これは「モノドロミー行列の $t$ 微分がゼロであること」と「行列成分の比較」から以下のように鮮やかに証明される。

\textbf{1. モノドロミー不変性から生じる可換性} \\
基本解行列 $Y = \Phi (x-\xi)^A G(t)$ が特異点 $x=\xi$ の周りを1周したとき、$(x-\xi)^A$ は $(x-\xi)^A \exp(2\pi i A)$ となる。したがって解析接続後の解は
\[
  Y^\gamma = \Phi (x-\xi)^A \exp(2\pi i A) G(t) = Y \cdot G(t)^{-1} \exp(2\pi i A) G(t)
\]
となる。すなわち、大域的なモノドロミー行列は $\Gamma = G(t)^{-1} \exp(2\pi i A) G(t)$ である。
仮定よりモノドロミーは $t$ に依存しない（不変である）ため、$\frac{\partial \Gamma}{\partial t_l} = 0$ である。
$\Gamma$ を $t_l$ で微分すると（$\frac{\partial G^{-1}}{\partial t_l} = -G^{-1}\frac{\partial G}{\partial t_l}G^{-1}$ に注意して）、
\[
  0 = -G^{-1}\frac{\partial G}{\partial t_l}G^{-1} \exp(2\pi i A) G + G^{-1} \exp(2\pi i A) \frac{\partial G}{\partial t_l}
\]
両辺の左から $G$、右から $G^{-1}$ を掛け、$U = \frac{\partial G}{\partial t_l}G^{-1}$ とおくと、
\[
  0 = -U \exp(2\pi i A) + \exp(2\pi i A) U \quad \implies \quad [U, \exp(2\pi i A)] = 0
\]
となり、「$U = \frac{\partial G}{\partial t_l}G^{-1}$ は局所モノドロミー $\exp(2\pi i A)$ と可換である」ことが厳密に導かれる。

\textbf{2. 条件(P2)が引き起こす奇跡（$A$ との可換性）} \\
次に、$\exp(2\pi i A)$ と可換な行列 $U$ が、なぜ $A$ 自身とも可換になるのかを示す。
$A$ は対角行列 $A = \mathrm{diag}(\lambda_1, \dots, \lambda_n)$ としてよい。すると $\exp(2\pi i A)$ も対角行列 $\mathrm{diag}(e^{2\pi i \lambda_1}, \dots, e^{2\pi i \lambda_n})$ となる。
可換性の式 $\exp(2\pi i A) U = U \exp(2\pi i A)$ の $(j, k)$ 成分を比較すると、
\[
  e^{2\pi i \lambda_j} U_{jk} = U_{jk} e^{2\pi i \lambda_k} \quad \implies \quad U_{jk} \left( e^{2\pi i \lambda_j} - e^{2\pi i \lambda_k} \right) = 0
\]
ここで\textbf{条件(P2)}が火を噴く。(P2)は「任意の $j \neq k$ に対して $\lambda_j - \lambda_k$ が整数ではない」という条件である。オイラーの公式から、差が整数でなければ $e^{2\pi i \lambda_j} \neq e^{2\pi i \lambda_k}$ である。
したがって、$j \neq k$ のときは括弧の中身が絶対に $0$ にならないため、\textbf{$U_{jk} = 0$ でなければならない}。
これは「行列 $U$ が対角行列である」ことを意味する。
$A$ も対角行列、$U$ も対角行列であるため、対角行列同士は必ず可換である。よって
\[
  [U, A] = 0 \quad \implies \quad \left[ \frac{\partial G}{\partial t_l}G^{-1}, A \right] = 0
\]
が結論づけられる。この可換性のおかげで、前項の第3項における $\frac{\partial G}{\partial t_l}G^{-1}$ と $(x-\xi)^A$ の順序交換が完全に正当化され、特異性の見事な相殺が起こるのである。
\end{Q}

    したがって、
    \begin{align*}
      \text{第3項} &= \Phi \cdot (x-\xi)^A \left( \frac{\partial G}{\partial t_l}G^{-1} \right) (x-\xi)^{-A} \cdot \Phi^{-1} \\
      &= \Phi \cdot \left( \frac{\partial G}{\partial t_l}G^{-1} \right) (x-\xi)^A (x-\xi)^{-A} \cdot \Phi^{-1} \\
      &= \Phi \left( \frac{\partial G}{\partial t_l}G^{-1} \right) \Phi^{-1}
    \end{align*}
    となり、なんと $x-\xi$ の依存性が完全に消滅する！ $\Phi$ も $G$ も正則であるため、第3項も完全に正則であることが厳密に保証される。
    
    \item \textbf{第2項の評価と結論：}
    以上より、極を生み出す可能性のある項は第2項のみに絞られた。
    もし $\xi$ が $t_l$ に依存しなければ $\frac{\partial \xi}{\partial t_l} = 0$ となり、第2項は消滅する。すべて正則な項のみとなるため、$A^{(l)}$ は正則である（(2)の証明）。
    もし $\xi$ が $t_l$ に依存すれば $\frac{\partial \xi}{\partial t_l} \neq 0$ であり、第2項は $\frac{1}{x-\xi}$ を持つ。$\Phi A \Phi^{-1}$ は正則かつ非零であるため、他の項と打ち消し合うことはなく、$A^{(l)}$ は $x=\xi$ に「ちょうど1位の極」を持つことが結論づけられる（(3)の証明）。
\end{itemize}
\end{Q}

以下、第2章第5節の結果を使って、行列型のシュレジンガー型微分方程式
\begin{equation}\label{eq:3.4_Schleginger}
  \frac{dY}{dx} = \left( \sum_{l=1}^L \frac{A_l}{x-t_l} \right) Y
\end{equation}
のモノドロミー保存変形について調べてみよう．計算の細部に立ち入ることはせず，筋道を紹介する．

変形のパラメータ $t_l$ としては，確定特異点の位置をとる．すなわち
\[
  t = (t_1, t_2, \cdots, t_L)
\]
とする．$t=(t_1, \cdots, t_L)$ の変域 $U \subset \mathbb{C}^L$ は必要に応じて小さくとっておく．また
\[
  A_\infty = - \sum_{l=1}^L A_l
\]
は，対角化可能であるとし，あらかじめ対角行列にとっておく．$Y=Y(x,t)$ を\eqref{eq:3.4_Schleginger}の基本解行列で，そのモノドロミー群は $t$ に依存しないとする．

\eqref{eq:3.4_matrix_representation_A_Phi_G}により $A^{(l)}(x,t)$ を定めれば、命題\ref{prop:3.4_matrix_ODE_monodoromy_deformation_property}により，$A^{(l)}(x,t)$ は $x=t_l$ のみで1位の極をもち、それ以外の $\mathbb{P}^1(\mathbb{C})$ 上の各点で正則である．我々は，ここで L. Schlesinger にならって，$x=\infty$ の近傍において
\[
  Y(x,t) = \left( \sum_{j=0}^\infty \Phi_j(t) x^{-j} \right) x^{-A_\infty}, \quad \Phi_0 = I_n
\]
と表されている，と仮定する．すると，\eqref{eq:3.4_matrix_representation_A_Phi_G}から
\[
  A^{(l)}(\infty, t) = 0
\]
である．

\begin{Q}[無限遠 $x=\infty$ における展開の仮定から、なぜ $A_\infty$ の固有値に依存することなく、直ちに $A^{(l)}(\infty,t) = 0$ が導かれるのか？]
\textbf{A.} 変形行列の定義 $A^{(l)} = \frac{\partial Y}{\partial t_l}Y^{-1}$ を計算する際、$Y$ と $Y^{-1}$ に含まれる $x^{\pm A_\infty}$ が「微分されることなく完璧に相殺される」からである。

具体的には以下の通りである。
前提として、無限遠における留数行列 $A_\infty = -\sum A_l$ は、モノドロミー保存の仮定によりその固有値（特性指数）が $t$ に依存しない。テキストの仮定によりあらかじめ定数対角行列として取ってあるため、$\frac{\partial A_\infty}{\partial t_l} = 0$ である。

$Y(x,t) = \Phi(x,t)x^{-A_\infty}$ とおく。ここで $\Phi(x,t) = I_n + \Phi_1(t)x^{-1} + \Phi_2(t)x^{-2} + \dots$ である。
これを $t_l$ で偏微分すると、$A_\infty$ が定数であるため、微分作用素は $\Phi$ にしかかからない。
\[
  \frac{\partial Y}{\partial t_l} = \frac{\partial \Phi}{\partial t_l} x^{-A_\infty} + \Phi \frac{\partial}{\partial t_l}(x^{-A_\infty}) = \frac{\partial \Phi}{\partial t_l} x^{-A_\infty}
\]
一方、基本解行列の逆行列は
\[
  Y^{-1} = (x^{-A_\infty})^{-1} \Phi^{-1} = x^{A_\infty} \Phi^{-1}
\]
である。これらを $A^{(l)}$ の定義式に代入すると、中央で $x^{-A_\infty}$ と $x^{A_\infty}$ が隣り合い、完全に単位行列となって消滅する。
\begin{align*}
  A^{(l)} &= \left( \frac{\partial \Phi}{\partial t_l} x^{-A_\infty} \right) \left( x^{A_\infty} \Phi^{-1} \right) \\
  &= \frac{\partial \Phi}{\partial t_l} \Phi^{-1}
\end{align*}
ここで、$\Phi_0 = I_n$ は定数であるため $\frac{\partial \Phi}{\partial t_l}$ は $O(x^{-1})$ の項から始まる。また $\Phi^{-1} = I_n + O(x^{-1})$ である。
したがって、
\[
  A^{(l)} = O(x^{-1}) \cdot (I_n + O(x^{-1})) = O(x^{-1})
\]
となり、$A_\infty$ の固有値のいかんにかかわらず全体のオーダーは厳密に $O(x^{-1})$ となる。よって $x \to \infty$ で $A^{(l)}(\infty, t) = 0$ となる。
\end{Q}

一方，$x=t_l$ における解の表示\eqref{eq:3.4_matrix_solution_representation_Phi_G}を\eqref{eq:3.4_Schleginger}に代入すると
\[
  \Phi(t_l, t) A \Phi(t_l, t)^{-1} = A_l
\]
が得られる。

\begin{Q}[$x=t_l$ における局所表示 $Y(x,t) = \Phi(x,t)(x-t_l)^A G(t)$ をシュレジンガー系に代入すると、なぜ $\Phi(t_l, t) A \Phi(t_l, t)^{-1} = A_l$ が得られるのか？]
\textbf{A.} 代入後、右から $G(t)^{-1}$ および $(x-t_l)^{-A}$ を掛けて整理し、$x=t_l$ における留数（$(x-t_l)^{-1}$ の係数）を比較することで直ちに導かれる。

具体的な計算プロセスは以下の通りである。
基本解行列 $Y = \Phi (x-t_l)^A G(t)$ を微分方程式 $\frac{\partial Y}{\partial x} = \left( \sum_{h=1}^L \frac{A_h}{x-t_h} \right) Y$ に代入する。
左辺は積の微分法則と行列のべき乗の微分 $\frac{\partial}{\partial x}(x-t_l)^A = A(x-t_l)^{A-I} = A(x-t_l)^{-1}(x-t_l)^A$ より、
\[
  \text{左辺} = \frac{\partial \Phi}{\partial x}(x-t_l)^A G(t) + \Phi A \frac{1}{x-t_l} (x-t_l)^A G(t)
\]
となる。右辺は、
\[
  \text{右辺} = \left( \sum_{h=1}^L \frac{A_h}{x-t_h} \right) \Phi (x-t_l)^A G(t)
\]
である。ここで、両辺の右側から $G(t)^{-1} (x-t_l)^{-A}$ を掛けて特異性を持つ因子を相殺すると、
\[
  \frac{\partial \Phi}{\partial x} + \Phi A \frac{1}{x-t_l} = \left( \sum_{h=1}^L \frac{A_h}{x-t_h} \right) \Phi
\]
を得る。
この等式の両辺について、$x \to t_l$ としたときの最も特異性の高い項（$(x-t_l)^{-1}$ の項、すなわち特異点における留数）を比較する。
\begin{itemize}
    \item 左辺：$\frac{\partial \Phi}{\partial x}$ は正則であるため、極を持つのは第2項のみであり、その留数は $\Phi(t_l, t) A$ である。
    \item 右辺：総和のうち $h=l$ の項 $\frac{A_l}{x-t_l}\Phi$ のみが $x=t_l$ で極を持ち、その他の $h \neq l$ の項は正則である。したがって留数は $A_l \Phi(t_l, t)$ である。
\end{itemize}
これらを等置すると、
\[
  \Phi(t_l, t) A = A_l \Phi(t_l, t)
\]
となる。$\Phi(t_l, t)$ は可逆行列であるから、右から $\Phi(t_l, t)^{-1}$ を掛けることで
\[
  \Phi(t_l, t) A \Phi(t_l, t)^{-1} = A_l
\]
が厳密に導出される。
\end{Q}

上の2つの式から、結局
\begin{equation*}
  A^{(l)}(x,t) = - \frac{A_l}{x-t_l}
\end{equation*}

\begin{Q}[テキストの「上の2つの式から、結局…」の部分にはどのような論理の飛躍があり、どうやって $A^{(l)}(x,t)$ の大域的な形を決定しているのか？]
\textbf{A.} 複素解析における強力な定理である「リウヴィルの定理（拡張版）」を暗黙に適用している。
有理関数行列 $A^{(l)}(x,t)$ について、以下の3つの局所的な情報が揃っている。
\begin{enumerate}
    \item \textbf{極の位置:} 命題3.11より、$x=t_l$ にのみ1位の極を持つ。よって $A^{(l)}(x,t) = \frac{C_l}{x-t_l} + P(x)$ （$P(x)$は多項式）と書ける。
    \item \textbf{無限遠での振る舞い:} $A^{(l)}(\infty,t) = 0$ であるため、多項式部分 $P(x)$ は恒等的に $0$ でなければならない。すなわち $A^{(l)}(x,t) = \frac{C_l}{x-t_l}$ に確定する。
    \item \textbf{留数の値 $C_l$:} 前ページの(3.56)式において $\xi = t_l$ とすると、主要部（極を持つ項）は第2項の $- \frac{\partial t_l}{\partial t_l} \Phi A \Phi^{-1} = -\Phi A \Phi^{-1}$ となる。直前の式 $\Phi A \Phi^{-1} = A_l$ を代入すれば、留数は $C_l = -A_l$ となる。
\end{enumerate}
これらをつなぎ合わせることで、大域的な形が $A^{(l)}(x,t) = - \frac{A_l}{x-t_l}$ に一意に定まるのである。
\end{Q}

となることがわかる．詳しい説明は省略したが，上で述べたような考察により次の命題が得られる．

\begin{proposition}\label{prop:3.4_Schleginger_deformation_system}
シュレジンガー型微分方程式\eqref{eq:3.4_Schleginger}の変形微分方程式は以下で与えられる．
\begin{equation*}
  \begin{cases}
    \displaystyle \frac{\partial Y}{\partial x} = \left( \sum_{l=1}^L \frac{A_l}{x-t_l} \right) Y \\[5mm]
    \displaystyle \frac{\partial Y}{\partial t_l} = \left( - \frac{A_l}{x-t_l} \right) Y
  \end{cases}
\end{equation*}
\end{proposition}

次に、命題\ref{prop:3.4_Schleginger_deformation_system}の変形微分方程式系の完全積分可能条件\eqref{eq:3.1_compatibility_PA}を具体的に計算する．まず
\[
  \sum_{h=1}^L \frac{1}{x-t_h} \frac{\partial A_h}{\partial t_l} = \left[ -\frac{A_l}{x-t_l}, \sum_{h=1}^L \frac{A_h}{x-t_h} \right]
\]
となるが，右辺を丁寧に変形すると
\[
  \sum_{h=1(h \neq l)}^L \frac{[A_h, A_l]}{t_h - t_l} \left( \frac{1}{x-t_h} - \frac{1}{x-t_l} \right)
\]
を得る．

\begin{Q}[「右辺を丁寧に変形すると」とあるが、交換子の計算と部分分数分解の具体的なプロセスはどうなっているか？]
\textbf{A.} 右辺の交換子は双線形性により和の外に出せる。各項は
$\left[ -\frac{A_l}{x-t_l}, \frac{A_h}{x-t_h} \right] = -\frac{1}{(x-t_l)(x-t_h)} [A_l, A_h] = \frac{1}{(x-t_l)(x-t_h)} [A_h, A_l]$ 
となる。ここで $h=l$ の項は自己交換子 $[A_l, A_l] = 0$ となり消滅するため、和の添字は $h \neq l$ のみとなる。
さらに、有理関数の部分を部分分数分解すると
$\frac{1}{(x-t_l)(x-t_h)} = \frac{1}{t_h - t_l} \left( \frac{1}{x-t_h} - \frac{1}{x-t_l} \right)$ 
となるため、これを代入することでテキストの展開式が厳密に得られる。
\end{Q}

ここで，$\frac{1}{x-t_h}$ の係数を比較すれば，$A_h$ に関する微分方程式系

\begin{align}
  \frac{\partial A_h}{\partial t_l} &= \frac{[A_h, A_l]}{t_h - t_l} \quad (h \neq l)\label{eq:3.4_compatibility_hl} \\
  \frac{\partial A_l}{\partial t_l} &= - \sum_{h=1(h \neq l)}^L \frac{[A_h, A_l]}{t_h - t_l}\label{eq:3.4_compatibility_ll}
\end{align}
が成立することがわかる．

\begin{Q}[左辺と右辺の係数比較から、なぜ $h \neq l$ の場合\eqref{eq:eq:3.4_compatibility_hl}と $h=l$ の場合\eqref{eq:3.4_compatibility_ll}の2つの異なる方程式が抽出されるのか？]
\textbf{A.} 左辺は $\sum_{h=1}^L \frac{\partial A_h / \partial t_l}{x-t_h}$ であり、すべての極 $x=t_k$ $(k=1,\dots,L)$ に独立した係数を持っている。
これを右辺の展開式と比較する。\\
1. $h \neq l$ の極 $\frac{1}{x-t_h}$：右辺ではシグマの中の1つの項からしか寄与がない。係数をそのまま比較して $\frac{\partial A_h}{\partial t_l} = \frac{[A_h, A_l]}{t_h - t_l}$ を得る（\eqref{eq:3.4_compatibility_hl}式）。\\
2. $h = l$ の極 $\frac{1}{x-t_l}$：右辺を見ると、$-\frac{1}{x-t_l}$ という項が「すべての $h \ (\neq l)$」に共通して存在している。つまり、右辺ではこの極に対する寄与が和全体から集中して集まってくる。これらを合算して左辺の係数 $\frac{\partial A_l}{\partial t_l}$ と等置することで、総和記号を含む\eqref{eq:3.4_compatibility_ll}式が得られる。
\end{Q}

\begin{definition}
非線型偏微分方程式系\eqref{eq:3.4_compatibility_hl}、\eqref{eq:3.4_compatibility_ll}を\textbf{シュレジンガー系}という．
\end{definition}

他方，条件\eqref{eq:3.1_compatibility_AA}は $h=l$ のときは自明であり，$h \neq l$ のときは微分方程式系
\[
  -\frac{1}{x-t_h} \frac{\partial A_h}{\partial t_l} + \frac{1}{x-t_l} \frac{\partial A_l}{\partial t_h} = \frac{[A_l, A_h]}{t_l - t_h} \left( \frac{1}{x-t_l} - \frac{1}{x-t_h} \right)
\]
が従う．この両辺の係数を等しいとおいて得られる微分方程式系は上の(3.58)と同じものである．すなわち，条件\eqref{eq:3.1_compatibility_AA}は\eqref{eq:3.1_compatibility_PA}に含まれている．

\begin{proposition}\label{prop:3.4_schleginger_compatibility}
シュレジンガー系\eqref{eq:3.4_compatibility_hl}、\eqref{eq:3.4_compatibility_ll}は完全積分可能である．
\end{proposition}

\begin{proof}
$df$ を，関数 $f(t)$ の $t=(t_1,\cdots,t_L)$ に関する微分とする．すなわち
\[
  dA_l = \sum_{h=1}^L \frac{\partial A_l}{\partial t_h} dt_h
\]
である。さらに微分1型式 $\Omega_l$ を
\[
  \Omega_l = \sum_{h=1(h \neq l)}^L \frac{[A_h, A_l]}{t_h - t_l} dt_h - \sum_{h=1(h \neq l)}^L \frac{[A_h, A_l]}{t_h - t_l} dt_l
\]
により定める．するとシュレジンガー系\eqref{eq:3.4_compatibility_hl}、\eqref{eq:3.4_compatibility_ll}は1種類のパッフ型式
\[
  dA_l = \Omega_l \quad (l=1,\cdots,L)
\]
で表される．

\begin{Q}[シュレジンガー系の完全積分可能性を示すために、なぜ微分形式 $\Omega_l$ を導入し $d\Omega_l = 0$ を確かめるというアプローチをとるのか？]%TODO ステートメントを確認
\textbf{A.} 多数の変数 $t_1, \dots, t_L$ に対する複雑な偏微分方程式系を、微分1形式を用いた「パッフ方程式（Pfaffian system） $dA_l = \Omega_l$」として幾何学的に統合するためである。
フロベニウスの定理（あるいはポアンカレの補題）により、このような方程式系が矛盾なく解を持つ（完全積分可能である）ための必要十分条件は、外微分をとった $d(dA_l) = d\Omega_l$ が $0$ になること（ゼロ曲率条件）である。直接的な微分の順序交換を計算するよりも、外微分の代数的な操作（$d^2 = 0$ や $dt_h \wedge dt_l$ の反対称性）を用いることで、はるかに見通し良く証明を記述できるからである。
\end{Q}

命題3.13に主張する完全積分可能性を示すには
\[
  d\Omega_l = 0 \quad (l=1,\cdots,L)
\]
を示せばよい．このことは，それほど簡単ではないにしても，とにかく直接計算で確かめることができる．
\end{proof}

\vspace{1em}
\noindent
\textbf{【補足：$d\Omega_l = 0$ の直接計算による証明】}

微分1型式 $\Omega_l$ を見やすくするため、次のように書き直す。
\[
  \Omega_l = \sum_{h=1 (h \neq l)}^L \frac{[A_h, A_l]}{t_h - t_l} (dt_h - dt_l)
\]
この外微分 $d\Omega_l$ を計算する。積の微分法則（ライプニッツ則）と $d(dt_h - dt_l) = 0$ より、係数の微分のみを考えればよい。
\[
  d\Omega_l = \sum_{h \neq l} d\left( \frac{[A_h, A_l]}{t_h - t_l} \right) \wedge (dt_h - dt_l)
\]
係数の微分は以下のようになる。
\[
  d\left( \frac{[A_h, A_l]}{t_h - t_l} \right) = \frac{[dA_h, A_l] + [A_h, dA_l]}{t_h - t_l} - \frac{[A_h, A_l]}{(t_h - t_l)^2}(dt_h - dt_l)
\]
これを $d\Omega_l$ の式に代入するが、第2項は $(dt_h - dt_l) \wedge (dt_h - dt_l) = 0$ となるため消滅する。したがって、調べるべき式は以下の2つの項の和に帰着する。
\begin{align*}
  d\Omega_l &= \underbrace{ \sum_{h \neq l} \frac{[dA_h, A_l]}{t_h - t_l} \wedge (dt_h - dt_l) }_{I_1} + \underbrace{ \sum_{h \neq l} \frac{[A_h, dA_l]}{t_h - t_l} \wedge (dt_h - dt_l) }_{I_2}
\end{align*}

\textbf{1. 第2項 $I_2$ の計算} \\
$I_2$ に $dA_l = \sum_{k \neq l} \frac{[A_k, A_l]}{t_k - t_l}(dt_k - dt_l)$ を代入する。
\[
  I_2 = \sum_{h \neq l} \sum_{k \neq l} \frac{[A_h, [A_k, A_l]]}{(t_h - t_l)(t_k - t_l)} (dt_k - dt_l) \wedge (dt_h - dt_l)
\]
$h=k$ の項はウェッジ積がゼロになるため消える。残る $h \neq k$ の項について、添字 $h$ と $k$ を対称化してまとめる。$(dt_k - dt_l) \wedge (dt_h - dt_l)$ をくくり出すと、係数は
\[
  \frac{1}{2} \left( \frac{[A_h, [A_k, A_l]]}{(t_h - t_l)(t_k - t_l)} - \frac{[A_k, [A_h, A_l]]}{(t_k - t_l)(t_h - t_l)} \right) = \frac{ [A_h, [A_k, A_l]] + [A_k, [A_l, A_h]] }{2(t_h - t_l)(t_k - t_l)}
\]
となる（交換子の反対称性 $[A_h, A_l] = -[A_l, A_h]$ を用いた）。ここで\textbf{ヤコビの恒等式}より
\[
  [A_h, [A_k, A_l]] + [A_k, [A_l, A_h]] = -[A_l, [A_h, A_k]] = [[A_h, A_k], A_l]
\]
となるため、
\begin{equation}
  I_2 = \frac{1}{2} \sum_{h, k \neq l (h \neq k)} \frac{[[A_h, A_k], A_l]}{(t_h - t_l)(t_k - t_l)} (dt_k - dt_l) \wedge (dt_h - dt_l)
  \label{eq:3.4_I2_result}
\end{equation}
を得る。

\textbf{2. 第1項 $I_1$ の計算} \\
次に $I_1$ に $dA_h = \sum_{k \neq h} \frac{[A_k, A_h]}{t_k - t_h}(dt_k - dt_h)$ を代入する。
\[
  I_1 = \sum_{h \neq l} \sum_{k \neq h} \frac{[[A_k, A_h], A_l]}{(t_h - t_l)(t_k - t_h)} (dt_k - dt_h) \wedge (dt_h - dt_l)
\]
関係式 $dt_k - dt_h = (dt_k - dt_l) - (dt_h - dt_l)$ を用いると、
\[
  (dt_k - dt_h) \wedge (dt_h - dt_l) = (dt_k - dt_l) \wedge (dt_h - dt_l)
\]
となる（$dt_h \wedge dt_h = 0$ のため）。また $k=l$ の項もウェッジ積がゼロになるため除外できる。
$I_2$ と同様に $h$ と $k$ で対称化する。入れ替えにより、ウェッジ積と内側の交換子 $[A_k, A_h]$ からそれぞれ $-1$ が出て打ち消し合い、係数の分母の差から次のような構造が現れる。
\begin{align*}
  \frac{1}{2} & [[A_k, A_h], A_l] \left( \frac{1}{(t_h - t_l)(t_k - t_h)} - \frac{1}{(t_k - t_l)(t_h - t_k)} \right) \\
  &= \frac{1}{2} [[A_k, A_h], A_l] \frac{1}{t_k - t_h} \left( \frac{1}{t_h - t_l} - \frac{1}{t_k - t_l} \right) \\
  &= \frac{1}{2} [[A_k, A_h], A_l] \frac{1}{t_k - t_h} \frac{t_k - t_h}{(t_h - t_l)(t_k - t_l)} \\
  &= \frac{1}{2} \frac{[[A_k, A_h], A_l]}{(t_h - t_l)(t_k - t_l)}
\end{align*}
したがって、
\begin{equation}
  I_1 = \frac{1}{2} \sum_{h, k \neq l (h \neq k)} \frac{[[A_k, A_h], A_l]}{(t_h - t_l)(t_k - t_l)} (dt_k - dt_l) \wedge (dt_h - dt_l)
  \label{eq:3.4_I1_result}
\end{equation}
となる。

\textbf{3. 結論（奇跡の相殺）} \\
式(\eqref{eq:3.4_I2_result})と式(\eqref{eq:3.4_I1_result})を足し合わせる。
\begin{align*}
  d\Omega_l &= I_1 + I_2 \\
  &= \frac{1}{2} \sum_{h, k \neq l (h \neq k)} \frac{ [[A_k, A_h], A_l] + [[A_h, A_k], A_l] }{(t_h - t_l)(t_k - t_l)} (dt_k - dt_l) \wedge (dt_h - dt_l)
\end{align*}
交換子の反対称性 $[A_k, A_h] = -[A_h, A_k]$ により、分子の括弧の中身は厳密にゼロとなる。
以上より、$d\Omega_l = 0$ が示された。 \hfill $\blacksquare$

与えられたモノドロミーを実現する線型フックス型微分方程式として，常にシュレジンガー型微分方程式(3.57)がとれるか，という問題の考察はまだ終わっていない．行列のサイズが $2 \times 2$ の場合にはこの問題の答えは肯定的である．必ずしもシュレジンガー型微分方程式で実現できるとは限らないモノドロミーはどのようなものであろうか．一般の場合には，モノドロミーを勝手に与えると，連立微分方程式系でも見かけの特異点が必要であり，確かにそのような例も知られている．この反例で与えたモノドロミーは可約である．

では既約なモノドロミーを実現する線型フックス型微分方程式に限ったら，そのモノドロミーはシュレジンガー型微分方程式で実現できるだろうか．このことが証明されるならば，単独高階微分方程式のモノドロミー保存変形の結果得られる微分方程式はシュレジンガー系に含まれることになる．

\begin{Q}[テキスト後半で語られている「与えられたモノドロミーを実現する微分方程式が常にシュレジンガー型にとれるか」という問題意識は、数学的に何を意味しているか？ また「見かけの特異点」の役割は何か？]
\textbf{A.} これは「ヒルベルトの第21問題（リーマン・ヒルベルト問題）」の核心部分である。
任意のモノドロミー表現を、見かけの特異点を持たないフックス型方程式（シュレジンガー系）だけで実現することは、一般には不可能であることが知られている（1989年のBolibrukhによる反例など）。
単独高階微分方程式を連立系に直す際などに必然的に生じる「見かけの特異点」は、モノドロミー（大域的な構造）を崩さずに、微分方程式の係数の局所的な自由度を調整する不可欠な役割を担う。可積分系の理論において、この見かけの特異点のダイナミクスこそがパンルヴェ方程式の解の挙動（$\tau$ 関数の零点）を支配するという、極めて重要な事実を示唆する記述である。
\end{Q}

少なくともわかることは，モノドロミー保存変形を許すということから，シュレジンガー型微分方程式は十分広いクラスのフックス型微分方程式である，ということである．

シュレジンガー系にはいくつかの第一積分が存在する．たとえば
\[
  \mathrm{Trace} A_l, \quad \sum_{h=1}^L A_h
\]
等である．これらは，モノドロミー保存変形の不変量であり，実際\eqref{eq:3.4_compatibility_hl}、\eqref{eq:3.4_compatibility_ll}の形からも直接見てとれることである．しかし，シュレジンガー系から種々の積分を用いて変数を減らしていく操作を実行することは得策とは言えない．何か見通しがない限り微分方程式の形は複雑になるばかりである．

\begin{Q}[保存量（第一積分）が存在するにもかかわらず、それを用いてシュレジンガー系の変数を消去し、直接解こうとしないのはなぜか？]
\textbf{A.} シュレジンガー系は行列を成分とする巨大な非線形偏微分方程式系であるため、見通しのないまま代数的な関係式で変数を消去すると、残された方程式の形が絶望的に複雑化・巨大化し、解析が不可能になってしまうからである。
そのため、一般論を深追いするのではなく、「単独高階方程式への帰着」という別の指導原理（特に有理関数 $a_n^{(l)}$ のラックス方程式に着目するルート）をとる方が、パンルヴェ方程式などの具体的な可積分系を抽出する上で理論上も実際上もはるかに有効なアプローチとなる。
\end{Q}

したがって，単独高階フックス型微分方程式のモノドロミー保存変形は別の指導原理を求めて，別個に考察するのが理論上も実際上も良い方法である．

単独高階フックス型微分方程式のモノドロミー保存変形において，重要な役割を果たすのは，$x$ の有理関数 $a_n^{(l)}(x,t)$ である．これらはラックスの方程式\eqref{eq:3.1_lax_equation_pa}を満たす、我々の対象である。パンルヴェ方程式は2階微分方程式のモノドロミー保存変形に関係している。第3章第2節で調べたように、$a_2^{(l)}(x,t)$ のラックスの方程式は3階の非斉次線型微分方程式(3.29)であった．また，単独高階微分方程式のモノドロミー保存変形は，変換(3.43)によりSL型の場合に帰着するのであった。また、与えられた微分方程式の2つの解の基本形 $\vec{\varphi}(x,t), \vec{\psi}(x,t)$ で，そのモノドロミー群が $t$ に依存しないものが存在すれば，それらは\eqref{eq:3.3_linear_transform_G}、\eqref{eq:3.3_scalar_matrix_gG}という簡単な変換で移りあう。したがって、変形微分方程式系は本質的に1つである。

以降の節において、具体的に2階フックス型微分方程式のモノドロミー保存変形を調べる。次節では、パンルヴェVI型方程式の多変数化である、ガルニエ系を説明する。2次のシュレジンガー型微分方程式のモノドロミー保存変形との関係、$n=2$ のシュレジンガー系とガルニエ系の関係等についても述べる．

パンルヴェVI型方程式と結びついたフックス型微分方程式は第2.8節で既に触れた．重複を恐れずもう一度書く．

\begin{equation}\label{eq:3.4_2-PDE_x,t}
  \frac{d^2 y}{dx^2} + p_1(x,t)\frac{dy}{dx} + p_2(x,t)y = 0
\end{equation}
\begin{equation}\label{eq:3.4_PVI_coefficient_p1}
  p_1(x,t) = \frac{1-\kappa_0}{x} + \frac{1-\kappa_1}{x-1} + \frac{1-\theta}{x-t} - \frac{1}{x-\lambda}
\end{equation}

\begin{Q}[\eqref{eq:3.4_PVI_coefficient_p1}式における係数 $p_1(x,t)$ の極（特異点） $x = 0, 1, t, \lambda$ は、モノドロミー保存変形の理論においてそれぞれどのような幾何学的・解析的な役割を担っているか？]
\textbf{A.} これらは方程式の特異点であるが、その性質は大きく3つに分類される。
\begin{itemize}
    \item \textbf{$x = 0, 1, \infty$（固定特異点）}: 位置が動かない確定特異点。方程式の土台となる空間（$\mathbb{P}^1 \setminus \{0, 1, \infty\}$）を定める。
    \item \textbf{$x = t$（動く特異点／変形パラメータ）}: 時間発展に相当する独立変数であり、この特異点の位置を動かしてもモノドロミーが不変となるように系を変形させる。
    \item \textbf{$x = \lambda$（見かけの特異点）}: 解が多価性を持たない（モノドロミーを持たない）特殊な特異点。パラメータ $t$ を動かすと、この $\lambda$ は $x$ 平面上を $t$ の関数 $\lambda(t)$ として複雑に動き回る。実は、この \textbf{$\lambda(t)$ が満たす非線形微分方程式こそが「パンルヴェVI型方程式」そのもの}である。
\end{itemize}
\end{Q}

\begin{equation}\label{eq:3.4_PVI_coefficient_p2}
  p_2(x,t) = \frac{\kappa}{x(x-1)} + \frac{\lambda(\lambda-1)\mu}{x(x-1)(x-\lambda)} - \frac{t(t-1)H}{x(x-1)(x-t)}
\end{equation}
\[
  \kappa = \chi(\chi+\kappa_\infty) = \frac{1}{4}(\kappa_0+\kappa_1+\theta-1)^2 - \frac{1}{4}\kappa_\infty^2
\]

ここで，$H$ は $\lambda, \mu, t$ の関数として次式で与えられた．
\begin{align}\label{eq:3.4_PVI_H_representation}
  t(t-1)H &= \lambda(\lambda-1)(\lambda-t)\mu^2 \nonumber \\
  &\quad - \{\kappa_0(\lambda-1)(\lambda-t) + \kappa_1\lambda(\lambda-t) + (\theta-1)\lambda(\lambda-1)\} \mu \nonumber \\
  &\quad + \kappa(\lambda-t)
\end{align}

\begin{Q}[$p_2(x,t)$ の表示式\eqref{eq:3.4_PVI_coefficient_p2}に現れるパラメータ $\mu$ と関数 $H$ は、幾何学的・力学系的にどのような意味を持っているか？]
\textbf{A.} $\lambda$ が「見かけの特異点の位置」という空間的な「座標（$q$）」に相当するのに対し、$\mu$ はそれに付随して決まる留数に関する「アクセサリー・パラメータ」であり、力学系における「正準運動量（$p$）」の役割を果たす。そして $H$ はまさにこの力学系の「ハミルトニアン（エネルギー関数）」である。
すなわち\eqref{eq:3.4_PVI_coefficient_p2}は、方程式の係数という静的な対象の中に、既に背後で躍動する力学系の変数たちが完全に埋め込まれていることを示している。
\end{Q}

我々は，フックス型微分方程式\eqref{eq:3.4_2-PDE_x,t}、\eqref{eq:3.4_PVI_coefficient_p1}、\eqref{eq:3.4_PVI_coefficient_p2}に関するフックスの問題を再考察する．目標となるのは次の定理である．

\begin{theorem}
\eqref{eq:3.4_2-PDE_x,t}、\eqref{eq:3.4_PVI_coefficient_p1}、\eqref{eq:3.4_PVI_coefficient_p2}のモノドロミー保存変形は，ハミルトン系
\begin{equation}\label{eq:3.4_PVI_Hamilton_system}
  \frac{d\lambda}{dt} = \frac{\partial H}{\partial \mu}, \quad \frac{d\mu}{dt} = -\frac{\partial H}{\partial \lambda}
\end{equation}
で与えられる．ここで，$H=H(t;\lambda,\mu)$ は\eqref{eq:3.4_PVI_H_representation}で定められる。
\end{theorem}

すなわち，\eqref{eq:3.4_2-PDE_x,t}、\eqref{eq:3.4_PVI_coefficient_p1}、\eqref{eq:3.4_PVI_coefficient_p2}の解の基本形で，そのモノドロミー群が $t$ に依存しないものが存在するための必要十分条件は、2つのアクセサリー・パラメータ $\lambda, \mu$ が非線型常微分方程式系\eqref{eq:3.4_PVI_Hamilton_system}を満足すること，である．

\eqref{eq:3.4_PVI_Hamilton_system}を\eqref{eq:3.4_PVI_H_representation}により表し，連立微分方程式から $\mu$ を消去すると，パンルヴェ方程式 $\mathrm{P}_{\mathrm{VI}}$ が得られる．ハミルトン系からパンルヴェ方程式を導く計算は初等的ではあるが筆算で確かめるのは簡単ではない．

\begin{Q}[パンルヴェVI型方程式（2階の非線形常微分方程式）を直接導出せず、なぜ定理3.1のように $\lambda$ と $\mu$ のハミルトン系\eqref{eq:3.4_PVI_Hamilton_system}を経由して定式化するアプローチをとるのか？]
\textbf{A.} 主に以下の3つの圧倒的な利点があるためである。
\begin{enumerate}
    \item \textbf{計算量の劇的な削減:} 直接 $\mu$ を消去して複雑極まりない $\mathrm{P}_{\mathrm{VI}}$ を導出・証明する古典的な手法（R. Fuchs や R. Garnier の手法）に比べ、ハミルトン系をベースにすることで証明の見通しが圧倒的に良くなる。
    \item \textbf{力学系理論の応用:} 微分方程式を「ハミルトン系」の枠組みに乗せることで、シンプレクティック幾何学や正準変換といった、古典力学で培われた強力な数学的ツールをそのまま可積分系の解析に流用できる。
    \item \textbf{対称性の可視化:} $\mathrm{P}_{\mathrm{VI}}$ が持つベックルント変換などの隠れた対称性（アフィン・ワイル群の作用など）は、$\lambda$ だけの式よりも、$\lambda$ と $\mu$ が対等に現れるハミルトニアン $H$ の多項式の形を見る方が遥かに発見・証明しやすい。
\end{enumerate}
\end{Q}

一方，この定理の主張するハミルトン系を導入することで，古典的な R. Fuchs や R. Garnier の結果は，その見通しがよくなり計算量も少なくなる．さらに非線型微分方程式自体を扱うときにも，ハミルトン系に書かれるということは，力学系の考え方が使えるなど，なかなか便利である．上の定理に述べた事実を指導原理として，ガルニエ系の計算を実行する．

\subsection{ガルニエ系}

%#region
\begin{story}
\begin{definition}[ガルニエ系を導くフックス型常微分方程式とリーマン図式]
この節以下で考察する2階フックス型線型常微分方程式の形を以下のように定める。
\begin{equation}
\frac{d^2 y}{dx^2} + p_1(x, t) \frac{dy}{dx} + p_2(x, t) y = 0 \label{eq:3.33}
\end{equation}
ここで対象となる微分方程式 \eqref{eq:3.33} のリーマン図式は、次のように与えられる。
\begin{equation}
\left\{
\begin{array}{ccccc}
x=0 & x=1 & x=t_l & x=\lambda_k & x=\infty \\
0 & 0 & 0 & 0 & \varepsilon \\
\kappa_0 & \kappa_1 & \theta_l & 2 & \varepsilon+\kappa_\infty
\end{array}
\right\} \label{eq:3.65}
\end{equation}
ただし、$k, l = 1, 2, \cdots, g$ とする。
\end{definition}

前節では、モノドロミー保存変形（ホロノミック変形）のパラメータを $t = (t_1, t_2, \cdots, t_L)$ としていた。ここでは、$L=g$ の場合を考察する。すなわち、変形のパラメータとしては確定特異点の位置 $x=t_l$ をとり、さらに、変形パラメータの次元 $g$ は、見かけの特異点の個数と一致している場合を考えるのである。当面 $t=(t_1, \cdots, t_g)$ の変域 $U \subset \mathbb{C}^g$ は必要に応じて小さくとっておく。

リーマン図式からわかるように、微分方程式 \eqref{eq:3.33}-\eqref{eq:3.65} は $\mathbf{P}^1(\mathbb{C})$ 上に $2g+3$ 個の確定特異点をもつ。
特異点の集合は
\begin{equation*}
\Xi(t) = \{0, 1, \infty, t_1, \cdots, t_g\}, \quad \Lambda(t) = \{\lambda_1(t), \cdots, \lambda_g(t)\}
\end{equation*}
の和集合 $\Xi(t) \cup \Lambda(t)$ である。我々は、確定特異点の位置
\begin{equation*}
t = (t_1, t_2, \cdots, t_g)
\end{equation*}
をパラメータとして、\eqref{eq:3.33} のモノドロミー保存変形を考察する。結局これからの考察は、見かけの特異点の位置 $\lambda_k$ が変形パラメータにどのように依存するか、ということに帰着するので、先走って $\Lambda(t)$ と書いている。
\end{story}

\begin{Q}[【核心】リーマン図式 \eqref{eq:3.65} において、$x=\lambda_k$ の特性指数が $(0, 2)$ となっているのは、解析的にどのような要請によるものか？]
\textbf{A.} $x=\lambda_k$ が「見かけの特異点（apparent singularity）」であることを保証するためである。

\begin{proof}
一般に、特性指数の差が整数（この場合は $2 - 0 = 2$）となる確定特異点では、フロベニウスの定理により局所解の片方に「対数項（$\log(x-\lambda_k)$）」が現れることが知られている。対数項が存在すると、その点の周りを一周した際に解が元の値に戻らず、非自明なモノドロミー（分岐）が生じてしまう。

しかし、係数関数 $p_1(x,t), p_2(x,t)$ に対して極の留数などに特定の条件（後に現れるハミルトニアンの制約式）を課すことで、この対数項の係数を意図的に $0$ にすることができる。
対数項が消滅すると、解はその点の近傍で単なる有理型関数（分岐を持たない解）として振る舞い、モノドロミーは完全に自明な単位行列となる。
方程式の係数としては特異点（極）を持つように見えるが、解の空間のトポロジー（モノドロミー）には何の影響も与えないため、これを「見かけの特異点」と呼ぶ。図式の特性指数 $(0, 2)$ は、まさにこの条件を準備するためのセットアップなのである。
\end{proof}
\end{Q}

\begin{Q}[【大局観】なぜ「見かけの特異点の個数」と「変形パラメータの次元」を同じ $g$ に揃える（$L=g$ とする）必然性があるのか？]
\textbf{A.} ガルニエ系が「パンルヴェ第VI方程式（$\mathrm{P_{VI}}$）の多変数化」であることを強く反映しているからである。

\begin{proof}
パラメータの数 $g=1$ の場合を考えてみる。
このとき、真の特異点は $\{0, 1, \infty, t_1\}$ の4つであり、変形パラメータは $t_1$ の1つだけである。そして見かけの特異点は $\lambda_1$ の1つとなる。
この「4つの確定特異点と1つの見かけの特異点」を持つ2階フックス型方程式のモノドロミー保存変形から導出される非線形常微分方程式こそが、パンルヴェ第VI方程式（$\mathrm{P_{VI}}$）に他ならない。

これを一般の $g$ 変数に拡張しようとしたとき、確定特異点の数 $t_1, \dots, t_g$ を増やす（変形の自由度を $g$ にする）ならば、モノドロミーを一定に保つための「調整弁（ダイナミクスの従属変数）」として、同数の見かけの特異点 $\lambda_1, \dots, \lambda_g$ を動的な変数として用意しなければならない。これが次元を一致させる物理的・幾何学的な理由であり、この多変数化された力学系全体を指して「ガルニエ系」と呼ぶのである。
\end{proof}
\end{Q}

\begin{story}

各確定特異点における特性指数 $\kappa_0, \kappa_1, \theta_l, \varepsilon, \varepsilon+\kappa_\infty$ の間には，フックスの関係式
\begin{equation}
\kappa_0 + \kappa_1 + \sum_{l=1}^g \theta_l + \kappa_\infty + 2\varepsilon = 1 \label{eq:3.66}
\end{equation}
が成り立っている．

微分方程式 \eqref{eq:3.33} に対する条件 (P1), (P2), (P3) は常に仮定する．念のために繰り返しておく．
\begin{itemize}
    \item[(P1)] $\Xi$ の各点を確定特異点，$\Lambda$ の各点を見かけの特異点としてもつ．
    \item[(P2)] 各 $x = \xi_l \in \Xi(t)$ における特性指数のどの2つも，差が整数ではない．
    \item[(P3)] モノドロミー $\hat{\rho}(t)$ は既約である．
\end{itemize}

(P2) により，どの $\kappa_0, \kappa_1, \theta_l, \kappa_\infty$ も整数ではない．また，(P3) は，微分方程式 \eqref{eq:3.33} が既約である，すなわち
\begin{equation*}
\left( \frac{d}{dx} + r_1(x,t) \right) \left( \frac{d}{dx} + r_2(x,t) \right) y = 0
\end{equation*}
という形には表せない，という仮定と同等である．ここで，$r_1(x,t), r_2(x,t)$ は $x$ の有理関数である．

以上の仮定の下で，\eqref{eq:3.33} の係数は次のように書かれる．
\begin{align}
p_1(x,t) &= \frac{1-\kappa_0}{x} + \frac{1-\kappa_1}{x-1} + \sum_{l=1}^g \frac{1-\theta_l}{x-t_l} - \sum_{k=1}^g \frac{1}{x-\lambda_k} \label{eq:3.67} \\
p_2(x,t) &= \frac{\kappa}{x(x-1)} + \sum_{k=1}^g \frac{\lambda_k(\lambda_k-1)\mu_k}{x(x-1)(x-\lambda_k)} - \sum_{l=1}^g \frac{t_l(t_l-1)H_l}{x(x-1)(x-t_l)} \label{eq:3.68}
\end{align}
\begin{equation*}
\kappa = \varepsilon(\varepsilon+\kappa_\infty) = \frac{1}{4} \left( \kappa_0 + \kappa_1 + \sum_{l=1}^g \theta_l - 1 \right)^2 - \frac{1}{4}\kappa_\infty^2
\end{equation*}

$g=1$ とすると，\eqref{eq:3.33}, \eqref{eq:3.67}, \eqref{eq:3.68} は前出の (3.60), (3.61), (3.62) になる．これは R.Fuchs の考察した，パンルヴェ方程式 $\mathrm{P_{VI}}$ に対応する場合である．

$\mathrm{P_{VI}}$ の拡張である，$t=(t_1, \cdots, t_g)$ の関数 $\lambda=(\lambda_1, \cdots, \lambda_g)$ が満たす2階非線型偏微分方程式系は，R.Garnier により1912年に求められた．これを書くために，少し記号を準備する．

\end{story}

\begin{Q}[【確認】フックスの関係式 \eqref{eq:3.66} の右辺は「2」ではなく、テキスト通り「1」で正しいのか？]
\textbf{A.} テキストの通り「1」で厳密に正しい。方程式の階数と特異点の個数を正確にカウントすると一致する。

\begin{proof}
2階のフックス型微分方程式において、$n$ 個の特異点 $x_i$ を持つ場合、すべての特性指数 $\rho_{i, 1}, \rho_{i, 2}$ の総和は「$n-2$」になるという定理がある。
今回のリーマン図式における特異点の総数 $n$ は、$0, 1, \infty$ の3個に加えて、$t_l$ が $g$ 個、$\lambda_k$ が $g$ 個の、計 $n = 2g + 3$ 個である。したがって、特性指数の総和は $(2g+3) - 2 = 2g+1$ とならなければならない。

有限の特異点での指数の和を計算する。
\begin{itemize}
    \item $x=0, 1, t_l$ の和： $\kappa_0 + \kappa_1 + \sum_{l=1}^g \theta_l$
    \item $x=\lambda_k$ の和（見かけの特異点）： $g$ 個の点それぞれで指数が $0+2=2$ なので、和は $2g$
\end{itemize}
無限遠点 $x=\infty$ での指数の和は $\varepsilon + (\varepsilon+\kappa_\infty) = 2\varepsilon + \kappa_\infty$ である。

これらをすべて足し合わせて $2g+1$ と等号で結ぶと、
\begin{equation*}
\left( \kappa_0 + \kappa_1 + \sum_{l=1}^g \theta_l \right) + 2g + \left( 2\varepsilon + \kappa_\infty \right) = 2g + 1
\end{equation*}
両辺の $2g$ が見事に相殺され、
\begin{equation*}
\kappa_0 + \kappa_1 + \sum_{l=1}^g \theta_l + \kappa_\infty + 2\varepsilon = 1
\end{equation*}
となり、テキストの式 \eqref{eq:3.66} が正しく導出される。
\end{proof}
\end{Q}

\begin{Q}[【導出】係数 $p_1(x,t)$ の形 \eqref{eq:3.67} はどのように決定されるのか？]
\textbf{A.} 確定特異点における特性方程式（決定方程式）の「解の和」から直接逆算される。

\begin{proof}
フックス型方程式の一般論として、点 $x=x_i$ が確定特異点であるとき、係数 $p_1(x)$ は $x_i$ において1位の極を持ち、その留数は $1 - (\rho_{i,1} + \rho_{i,2})$ になることが知られている（ここで $\rho$ は特性指数）。
これを用いて各特異点での項を書き下す。
\begin{itemize}
    \item $x=0$: 指数の和は $\kappa_0$。よって $\frac{1-\kappa_0}{x}$。
    \item $x=1$: 指数の和は $\kappa_1$。よって $\frac{1-\kappa_1}{x-1}$。
    \item $x=t_l$: 指数の和は $\theta_l$。よって $\frac{1-\theta_l}{x-t_l}$。
    \item $x=\lambda_k$: 指数の和は $0+2=2$。よって $\frac{1-2}{x-\lambda_k} = -\frac{1}{x-\lambda_k}$。
\end{itemize}
これらをすべて足し合わせた部分分数分解こそが、式 \eqref{eq:3.67} に他ならない。無限遠点で正しく振る舞うこと（$x \to \infty$ で $p_1(x) \sim 2/x$ となること）は、前述のフックスの関係式によって自動的に保証される。
\end{proof}
\end{Q}

\begin{Q}[【導出と意味】係数 $p_2(x,t)$ の形 \eqref{eq:3.68} はどのように導出されるか？また $\mu_k$ や $H_l$ という未知の関数はどこから来たのか？]
\textbf{A.} 特性指数の「積」の性質と、方程式が持つ未定の自由度（アクセサリー・パラメータ）を表現論的に整理した結果である。

\begin{proof}
係数 $p_2(x)$ は各特異点 $x_i$ において高々2位の極を持ち、その係数は特性指数の積 $\rho_{i,1}\rho_{i,2}$ で与えられる。
ここで最大のポイントは、見かけの特異点 $x=\lambda_k$ における指数の積が $0 \times 2 = \mathbf{0}$ になるという事実である。これにより、$p_2(x)$ は $x=\lambda_k$ において $1/(x-\lambda_k)^2$ という2位の極を持たず、1位の極しか持たないことが確定する。

方程式全体が無限遠点 $\infty$ で正しくフックス型として振る舞うため（$p_2(x) \sim O(1/x^2)$）、全体を $\frac{1}{x(x-1)}$ という共通因子でくくり、残りの部分を各特異点での部分分数に分解するのが定石である。このとき、各1位の極の留数として「未定の関数」が現れる。
\begin{itemize}
    \item 見かけの特異点 $\lambda_k$ での留数を（係数込みで） $\mu_k$ と名付ける。これは力学系における「運動量（正準共役量）」となる。
    \item 真の特異点 $t_l$ での留数を（係数込みで） $-H_l$ と名付ける。これが「ハミルトニアン」となる。
\end{itemize}
定数項の $\kappa$ は無限遠点での指数の積 $\varepsilon(\varepsilon+\kappa_\infty)$ から定まる。
このように、$\mu_k$ や $H_l$ は勝手に湧いてきたのではなく、$p_1(x)$ だけでは決めきれない $p_2(x)$ の「自由度（アクセサリー・パラメータ）」に対して、モノドロミー保存変形をハミルトン系として美しく記述できるように意図的に割り振られた記号なのである。
\end{proof}
\end{Q}

\begin{Q}[【鋭い疑問】式 \eqref{eq:3.68} の $p_2(x,t)$ において、たとえば $x=0$ で「2位の極（$1/x^2$）」が出てくるべきではないのか？なぜ分母は1乗の $x$ なのか？]
\textbf{A.} 結論から言うと、2位の極は出てこない（1位の極しか持たない）。そして、「2位の極が消滅すること」こそが、提示されたリーマン図式から導かれる必然的な帰結なのである。

\begin{proof}
フックス型微分方程式 $\frac{d^2y}{dx^2} + p_1(x)\frac{dy}{dx} + p_2(x)y = 0$ において、確定特異点 $x=x_i$ まわりの特性指数（決定方程式の解）を $\rho_1, \rho_2$ とする。
このとき、係数関数が特異点まわりでどう振る舞うべきかは、以下の関係式で完全に決まっている。
\begin{align*}
\lim_{x \to x_i} (x-x_i) p_1(x) &= 1 - (\rho_1 + \rho_2) \\
\lim_{x \to x_i} (x-x_i)^2 p_2(x) &= \rho_1 \rho_2
\end{align*}

ここで、問題の $x=0$ におけるリーマン図式 \eqref{eq:3.65} を見てみる。特性指数は $0$ と $\kappa_0$ である。
これを上の関係式に当てはめると、以下のようになる。
\begin{itemize}
    \item 指数の和： $0 + \kappa_0 = \kappa_0 \implies p_1(x)$ は $x=0$ で $\frac{1-\kappa_0}{x}$ という1位の極を持つ。（式 \eqref{eq:3.67} の第1項と完全に一致！）
    \item 指数の積： $0 \times \kappa_0 = \mathbf{0} \implies p_2(x)$ の $1/x^2$ の係数は $0$ になる。
\end{itemize}
つまり、「片方の特性指数が $0$ である」という図式の要請により、$p_2(x)$ は $x=0$ において2位の極を持てず、高々1位の極しか持たないことが数学的に確定するのである。

式 \eqref{eq:3.68} の $p_2(x,t)$ の分母をよく見ると、$x, (x-1), (x-t_l), (x-\lambda_k)$ はすべて「1乗」でしか現れていない。
これはリーマン図式 \eqref{eq:3.65} の「上段の指数がすべて $0$ である」という美しい設定のおかげで、すべての有限の特異点において2位の極が見事に消滅していることを表している。

（※逆に、無限遠点 $x=\infty$ では指数の積が $\varepsilon(\varepsilon+\kappa_\infty) = \kappa \neq 0$ となるため、$x \to \infty$ の極限をとると $p_2(x,t) \sim \frac{\kappa}{x^2}$ となり、無限遠点では正しく2位の極の振る舞いを示すことが確認できる。）
\end{proof}
\end{Q}

\begin{Q}[【最大の核心】リーマン図式や特性指数の情報を使って $p_1(x)$ を完全に決定できたように、$p_2(x)$ の1位の極の留数（$\mu_k$ や $H_l$）も代数的な計算で一意に求めることはできないのか？]
\textbf{A.} 絶対にできない。これらは方程式の「アクセサリー・パラメータ（未定の自由度）」であり、ある固定されたパラメータ $t$ における微分方程式をいくら睨んでも値は決まらない。これらを決定するためには「モノドロミー保存変形」という大域的な要請（ダイナミクス）を考える必要がある。

\begin{proof}
フックス型微分方程式において、特異点の数が増えると、リーマン図式（特異点の位置と特性指数という「局所的なデータ」）だけでは方程式の係数を一意に決定できなくなる。この決定できない余分なパラメータのことを「アクセサリー・パラメータ」と呼ぶ。
具体的には、
\begin{itemize}
    \item 確定特異点が3個（超幾何微分方程式）のときは、アクセサリー・パラメータは0個であり、係数は完全に一意に決まる。
    \item 特異点が4個以上（ホイン方程式やパンルヴェ、ガルニエ系の設定）になると、アクセサリー・パラメータが出現する。今回の式 \eqref{eq:3.68} における $\mu_k$ や $H_l$ がまさにそれである。
\end{itemize}
では、これらをどうやって決めるのか？
ここで登場するのが**「モノドロミー保存変形（ホロノミック変形）」**の概念である。
「確定特異点の位置 $t = (t_1, \dots, t_g)$ を動かしたとしても、微分方程式の大域的な解の接続関係（モノドロミー表現）が一切変化しないようにせよ」という極めて強い条件を課す。
すると、この条件を満たすために、見かけの特異点の位置 $\lambda_k(t)$ と、そこでの留数 $\mu_k(t)$ が、パラメータ $t_l$ の変化に合わせて「どのように動かなければならないか」を記述する非線形な微分方程式系が導出される。
これこそが「ガルニエ系」であり、1位の留数 $\mu_k$ や $H_l$ は、単なる代数計算ではなく、このガルニエ系という「微分方程式を解く」ことによって初めて決定される（関数として求まる）のである。
\end{proof}
\end{Q}

\begin{Q}[【大局観】モノドロミー保存変形から導かれる「ガルニエ系」において、留数 $\mu_k$ と $H_l$ はそれぞれ物理的（解析力学的）にどのような役割を果たすか？]
\textbf{A.} 解析力学における「正準変数（運動量）」と「ハミルトニアン」の役割を完璧に果たす。

\begin{proof}
先述の通り、モノドロミーを保存するように $t_l$ を動かすと、$\lambda_k$ と $\mu_k$ は $t$ の関数として変動する。驚くべきことに、この変動は次のような多変数ハミルトン系（Hamiltonian system）として極めて美しく記述されることが知られている（R. Garnier, 1912）。
\begin{equation}
\frac{\partial \lambda_k}{\partial t_l} = \frac{\partial H_l}{\partial \mu_k}, \quad \frac{\partial \mu_k}{\partial t_l} = -\frac{\partial H_l}{\partial \lambda_k} \quad (k, l = 1, \dots, g)
\end{equation}
ここで、
\begin{itemize}
    \item 見かけの特異点の位置 $\lambda_k$ は、力学系における「一般化座標（位置）」に相当する。
    \item $\lambda_k$ での留数 $\mu_k$ は、それに正準共役な「運動量」に相当する。
    \item 確定特異点 $t_l$ は「多次元的な時間」に相当する。
    \item $t_l$ での留数 $H_l$ は、その時間発展を支配する「ハミルトニアン（エネルギー）」に相当する。
\end{itemize}
つまり、式 \eqref{eq:3.68} で $p_2(x,t)$ を部分分数分解した際に現れた未知の留数たちは、単なる未定係数ではなく、「微分方程式の変形理論」を「解析力学（シンプレクティック幾何学）」の言葉で記述するための最も自然で必然的なパーツだったというわけである。
\end{proof}
\end{Q}
\begin{Q}[【核心】前節の議論から $p_2(x)$ が1位の極しか持たないことは分かった。では、未知の留数 $\mu_k, H_l$ 以外の「式の形（骨組み）」や、「$\frac{\kappa}{x(x-1)}$ という項とその係数」は、代数的に完全に導出できるということか？]
\textbf{A.} 完全に導出できる。あなたの直感通り、$p_2(x)$ の極の条件と「無限遠点での振る舞い」を組み合わせることで、この特定の表示形式が数学的に一意に強制される。

\begin{proof}
前節で確認した通り、$p_2(x)$ は $x=0, 1, t_l, \lambda_k$ において「1位の極」しか持たない有理関数である。
一方で、無限遠点 $x=\infty$ では、特性指数の積が $\kappa = \varepsilon(\varepsilon+\kappa_\infty)$ であるため、$x \to \infty$ の極限において
\begin{equation}
p_2(x) \sim \frac{\kappa}{x^2} + \mathcal{O}\left(\frac{1}{x^3}\right)
\end{equation}
という振る舞いをしなければならない（$1/x$ の項は存在してはならない）。

有理関数を部分分数分解して表現する際、この「無限遠での減衰条件」を最初から満たすような、都合の良い「基底関数」を用意するエレガントな手法がある。それが分母に $x(x-1)$ を持たせる方法である。
以下の3種類の関数を考える。
\begin{align*}
f_0(x) &= \frac{1}{x(x-1)} \quad &\sim \frac{1}{x^2} \quad &(x \to \infty) \\
f_{\lambda_k}(x) &= \frac{1}{x(x-1)(x-\lambda_k)} \quad &\sim \frac{1}{x^3} \quad &(x \to \infty) \\
f_{t_l}(x) &= \frac{1}{x(x-1)(x-t_l)} \quad &\sim \frac{1}{x^3} \quad &(x \to \infty)
\end{align*}
これらはすべて $x=0, 1$ で1位の極を持ち、下2つはそれぞれ $\lambda_k, t_l$ でも1位の極を持つため、$p_2(x)$ の極の条件をすべて表現できる完全な基底となる。

したがって、$p_2(x)$ は未定係数 $C_0, C_{\lambda_k}, C_{t_l}$ を用いて次のように必ず展開できる。
\begin{equation*}
p_2(x) = C_0 f_0(x) + \sum C_{\lambda_k} f_{\lambda_k}(x) + \sum C_{t_l} f_{t_l}(x)
\end{equation*}

ここで、両辺の $x \to \infty$ の極限をとる。
$f_{\lambda_k}(x)$ と $f_{t_l}(x)$ は $1/x^3$ のオーダーで急速に減衰するため、無限遠での $1/x^2$ の振る舞いに寄与できるのは **$f_0(x)$ だけ**である。
すなわち、$x \to \infty$ において
\begin{equation*}
p_2(x) \sim C_0 \frac{1}{x^2}
\end{equation*}
となる。これが $\frac{\kappa}{x^2}$ に一致しなければならないのだから、未知であった係数 $C_0$ は
\begin{equation*}
C_0 = \kappa
\end{equation*}
でなければならないことが、逃げ道なく確定する。これで第1項 $\frac{\kappa}{x(x-1)}$ が完全に導出された。

残りの項の係数 $C_{\lambda_k}, C_{t_l}$ は、各特異点での「留数」がシンプルな文字になるように調整されただけのものである。
例えば $x=\lambda_k$ での留数を計算すると、
\begin{equation*}
\lim_{x \to \lambda_k} (x-\lambda_k) p_2(x) = \frac{C_{\lambda_k}}{\lambda_k(\lambda_k-1)}
\end{equation*}
となる。この留数自体を「$\mu_k$」と定義したいのであれば、$C_{\lambda_k} = \lambda_k(\lambda_k-1)\mu_k$ と置けばよい。$t_l$ での留数を「$-H_l$」と定義したいのであれば、$C_{t_l} = -t_l(t_l-1)H_l$ と置けばよい。

結論として、式 \eqref{eq:3.68} の形は決して天下り的なものではなく、「局所的な極の条件」と「無限遠での漸近挙動」を繋ぎ合わせた、解析関数論における最も自然で必然的な帰結なのである。
\end{proof}
\end{Q}

\begin{Q}[【計算の検証】$x \to \infty$ における $p_2(x) \sim \kappa/x^2$ という振る舞いは、座標変換 $x=1/u$ を経た後の $u=0$ での決定方程式を正しく満たしているか？]
\textbf{A.} 完全に満たしている。座標変換に伴う $p_1$ のシフトと $p_2$ のスケーリングを正確に追うことで、方程式の解の指数がリーマン図式の通り $\varepsilon$ と $\varepsilon+\kappa_\infty$ になることが代数的に証明される。

\begin{proof}
$x = 1/u$ と変換すると、微分演算子は以下のように変換される。
\begin{equation*}
\frac{d}{dx} = -u^2 \frac{d}{du}, \quad \frac{d^2}{dx^2} = u^4 \frac{d^2}{du^2} + 2u^3 \frac{d}{du}
\end{equation*}
これを元の方程式 $y'' + p_1(x)y' + p_2(x)y = 0$ に代入し、全体を $u^4$ で割ると、$u=0$ まわりでの方程式を得る。
\begin{equation}
\frac{d^2y}{du^2} + \left( \frac{2}{u} - \frac{p_1(1/u)}{u^2} \right) \frac{dy}{du} + \frac{p_2(1/u)}{u^4} y = 0 \label{eq:u_eq}
\end{equation}

ここで、各係数の $x \to \infty$ （すなわち $u \to 0$）での漸近挙動を評価する。
式 \eqref{eq:3.67} より、$p_1(x)$ の無限遠での振る舞いは、各項を $1/x$ で展開して、
\begin{align*}
p_1(x) &\sim \frac{(1-\kappa_0) + (1-\kappa_1) + \sum (1-\theta_l) - g}{x} + \mathcal{O}(x^{-2}) \\
&= \frac{g + 2 - (\kappa_0 + \kappa_1 + \sum \theta_l) - g}{x} = \frac{2 - (\kappa_0 + \kappa_1 + \sum \theta_l)}{x}
\end{align*}
となる。これを $u$ の式に直すと $p_1(1/u) \sim \{2 - (\kappa_0 + \kappa_1 + \sum \theta_l)\}u$ である。
同様に、$p_2(x) \sim \kappa/x^2$ より $p_2(1/u) \sim \kappa u^2$ である。

これらを式 \eqref{eq:u_eq} に代入し、$u=0$ での決定方程式を構成する。
式 \eqref{eq:u_eq} の $dy/du$ の係数を $P_1(u)$、$y$ の係数を $P_2(u)$ とおくと、
\begin{itemize}
    \item $u P_1(0) = 2 - \{2 - (\kappa_0 + \kappa_1 + \sum \theta_l)\} = \kappa_0 + \kappa_1 + \sum \theta_l$
    \item $u^2 P_2(0) = \frac{\kappa u^2}{u^2} = \kappa$
\end{itemize}
となる。フックスの関係式 \eqref{eq:3.66} より、$\kappa_0 + \kappa_1 + \sum \theta_l = 1 - \kappa_\infty - 2\varepsilon$ であるから、
$u=0$ における決定方程式 $\rho(\rho-1) + uP_1(0)\rho + u^2P_2(0) = 0$ は、
\begin{equation*}
\rho^2 + (1 - \kappa_\infty - 2\varepsilon - 1)\rho + \kappa = 0 \quad \Longrightarrow \quad \rho^2 - (\kappa_\infty + 2\varepsilon)\rho + \kappa = 0
\end{equation*}
となる。ここで $\kappa = \varepsilon(\varepsilon + \kappa_\infty)$ であったことを思い出すと、この方程式は
\begin{equation*}
(\rho - \varepsilon)(\rho - (\varepsilon + \kappa_\infty)) = 0
\end{equation*}
と見事に因数分解される。
解は $\rho = \varepsilon, \varepsilon + \kappa_\infty$ となり、リーマン図式の設定と完全に一致する。したがって、$p_2(x) \sim \kappa/x^2$ という漸近挙動は、無限遠点での決定方程式を数学的に完璧に満たしている。
\end{proof}
\end{Q}

\begin{story}
\begin{align}
\mathrm{T}(x) &= x(x-1) \prod_{l=1}^g (x-t_l) \label{eq:3.69} \\
\mathrm{L}(x) &= \prod_{k=1}^g (x-\lambda_k) \label{eq:3.70}
\end{align}
とおく．$x$ の多項式 $\mathrm{T}(x), \mathrm{L}(x)$ の導関数を $\mathrm{T}'(x), \mathrm{L}'(x)$，2階導関数を $\mathrm{T}''(x), \mathrm{L}''(x)$ 等と書く．たとえば
\begin{equation*}
\mathrm{L}'(\lambda_k) = \prod_{p=1 \, (p \neq k)}^g (\lambda_k - \lambda_p)
\end{equation*}
である．次の完全積分可能非線型偏微分方程式系をガルニエ系という．
\end{story}

\begin{definition}[ガルニエ系 (Garnier system)]
\begin{align}
\frac{\partial^2 \lambda_k}{\partial t_l^2} &= \frac{1}{2} \left[ \frac{\mathrm{T}'(\lambda_k)}{\mathrm{T}(\lambda_k)} - \frac{1}{2} \frac{\mathrm{L}''(\lambda_k)}{\mathrm{L}'(\lambda_k)} \right] \left( \frac{\partial \lambda_k}{\partial t_l} \right)^2 - \left[ \frac{1}{2} \frac{\mathrm{T}''(t_l)}{\mathrm{T}'(t_l)} - \frac{\mathrm{L}'(t_l)}{\mathrm{L}(t_l)} \right] \frac{\partial \lambda_k}{\partial t_l} \nonumber \\
&\quad + \frac{1}{2} \sum_{p=1 \, (p \neq k)}^g \left[ \frac{\mathrm{T}(\lambda_k)\mathrm{L}'(\lambda_p)(\lambda_p - t_l)^2}{\mathrm{T}(\lambda_p)\mathrm{L}'(\lambda_k)(\lambda_k - t_l)^2(\lambda_k - \lambda_p)} \left( \frac{\partial \lambda_p}{\partial t_l} \right)^2 \right] \nonumber \\
&\quad - \sum_{p=1 \, (p \neq k)}^g \left[ \frac{\lambda_k - t_l}{(\lambda_p - t_l)(\lambda_p - \lambda_k)} \frac{\partial \lambda_k}{\partial t_l} \frac{\partial \lambda_p}{\partial t_l} \right] + \frac{\mathrm{L}(t_l)^2 \mathrm{T}(\lambda_k)}{2\mathrm{T}'(t_l)^2(\lambda_k - t_l)^2 \mathrm{L}'(\lambda_k)} I_{k,l} \label{eq:Garnier1}
\end{align}
\begin{equation}
\frac{\mathrm{T}'(t_l)(t_l - \lambda_k)}{\mathrm{L}(t_l)} \frac{\partial \lambda_k}{\partial t_l} - \frac{\mathrm{T}'(t_h)(t_h - \lambda_k)}{\mathrm{L}(t_h)} \frac{\partial \lambda_k}{\partial t_h} = \frac{(t_l - t_h)\mathrm{T}(\lambda_k)}{(\lambda_k - t_l)(\lambda_k - t_h)\mathrm{L}'(\lambda_k)} \label{eq:Garnier2}
\end{equation}
ここで，$h, k, l = 1, 2, \cdots, g$ であり，第1式では
\begin{align}
I_{k,l} &= \kappa_\infty^2 + \frac{\mathrm{T}'(0)}{\mathrm{L}(0)}\frac{\kappa_0^2}{\lambda_k} + \frac{\mathrm{T}'(1)}{\mathrm{L}(1)}\frac{\kappa_1^2}{\lambda_k - 1} \nonumber \\
&\quad + \sum_{h=1 \, (h \neq l)}^g \frac{\mathrm{T}'(t_h)}{\mathrm{L}(t_h)}\frac{\theta_h^2}{\lambda_k - t_h} + \frac{\mathrm{T}'(t_l)}{\mathrm{L}(t_l)}\frac{\theta_l^2}{\lambda_k - t_l} \label{eq:Garnier_I}
\end{align}
とおいた．
\end{definition}

\begin{story}
これを見ただけで，R.Garnierの計算がどんなに精緻を極めたものであったかがわかるであろう．ガルニエ系はパンルヴェ $\mathrm{VI}$ 型方程式の $g$ 変数への拡張であり，確かに $g=1$ とすると $\mathrm{P_{VI}}$ になる．

我々の目的は，線型常微分方程式(3.33), (3.67), (3.68)のモノドロミー保存変形を再考察し，ガルニエ系をハミルトン系に書き直すことである．これは，パンルヴェ $\mathrm{VI}$ 型方程式のハミルトン系表示を与える定理3.1の拡張である．

まず，考察する微分方程式の含むアクセサリー・パラメータは，$\lambda_k$ と
\end{story}

\begin{Q}[【核心】「次の完全積分可能非線型偏微分方程式系をガルニエ系という」とあるが、この恐ろしい方程式は最初から「定義（公理）」として天から降ってきたものなのか？]
\textbf{A.} もちろん違う。これは前節までに構成した「線型常微分方程式のモノドロミー保存変形条件（フロベニウスの積分可能条件）」を、力技で計算し尽くした結果得られる「導出された定理」である。

\begin{proof}
テキストの記述上は「これをガルニエ系という（定義する）」と書かれているが、数学の歴史的・論理的な立ち位置としては定理（帰結）である。
どんな天才であっても、何もない空間から \eqref{eq:Garnier1} のような複雑怪奇な非線型偏微分方程式を「定義」として思いつくことは不可能である。

Garnier が1912年に行ったことは以下の通りである。
1. 前節で定義した、特性指数の積が0になるような見かけの特異点を持つフックス型微分方程式を用意する。
2. その確定特異点 $t_l$ を動かしたときに、解のモノドロミーが変化しない（モノドロミー保存変形）という条件を課す。
3. その条件は、線型系の「変形方程式（シュレジンガー方程式）」の積分可能条件（ゼロ曲率条件）として定式化される。
4. その積分可能条件を、見かけの特異点 $\lambda_k$ の時間発展として愚直に偏微分方程式に書き下す。
その泥臭く精緻な計算の果てに「結果として出てきた」のが、この巨大な方程式系なのである。
\end{proof}
\end{Q}

\begin{Q}[【大局観】なぜ著者は、この巨大な方程式を提示した直後に「我々の目的はハミルトン系に書き直すことである」と宣言したのか？]
\textbf{A.} この2階の非線型偏微分方程式のままでは、複雑すぎて数学的な構造（対称性や解の性質）が全く見えないためである。

\begin{proof}
式 \eqref{eq:Garnier1} は $\lambda_k$ に関する2階の微分方程式として書かれている（ニュートン力学における運動方程式 $\ddot{x} = F(x, \dot{x})$ に相当する）。
しかし、解析力学の教えによれば、複雑な2階の微分方程式は、新たな変数「運動量（$\mu_k$）」を導入して1階の連立方程式（ハミルトン系）に書き直すことで、背後にある幾何学的な構造（シンプレクティック構造や正準変換の対称性）が劇的に見えやすくなる。

著者の岡本和夫氏は、パンルヴェ方程式をハミルトン系として定式化し、その対称性（アフィン・ワイル群）を明らかにした第一人者である。このテキストの真の目的は、「Garnier が残したこの恐ろしい数式群を、ハミルトニアン $H_l$ と正準変数 $(\lambda_k, \mu_k)$ を用いた洗練された幾何学の言葉へと翻訳・浄化すること」にある。そのための壮大な「フリ」として、ここでGarnierの原論文の姿をそのまま提示しているのである。
\end{proof}
\end{Q}

\begin{story}
\begin{align}
H_l &= -\operatorname{Res}_{x=t_l} p_2(x,t) \quad (l = 1, \cdots, g) \label{eq:3.71} \\
\mu_k &= \operatorname{Res}_{x=\lambda_k} p_2(x,t) \quad (k = 1, \cdots, g) \label{eq:3.72}
\end{align}
の，$3g$ 個である．一方，仮定(P1)により $x=\lambda_k$ は微分方程式の見かけの特異点であり，したがって $3g$ 個のアクセサリー・パラメータの間には $g$ 個の関係式が成り立つ．このことから，\eqref{eq:3.71} の $H_l$ は，$t=(t_1,\cdots,t_g)$, $\lambda=(\lambda_1,\cdots,\lambda_g)$, $\mu=(\mu_1,\cdots,\mu_g)$ の有理関数となることがわかる．以上のことはフロベニウスの方法により確かめることができる．実際，$H_l$ の具体形が次のようになることを次節において示す．
\begin{equation}
H_l = M_l \left[ \sum_{k=1}^g M^{k,l} \left\{ \mu_k^2 - A_{k,l}\mu_k + \frac{\kappa}{\lambda_k(\lambda_k-1)} \right\} \right] \label{eq:3.73}
\end{equation}
\begin{equation*}
A_{k,l} = \frac{\kappa_0}{\lambda_k} + \frac{\kappa_1}{\lambda_k-1} + \sum_{h=1}^g \frac{\theta_h - \delta_{hl}}{\lambda_k - t_h}
\end{equation*}
ここで，$\delta_{hl}$ はクロネッカーのデルタであり，また
\begin{align}
M_l &= - \frac{\mathrm{L}(t_l)}{\mathrm{T}'(t_l)} \label{eq:3.74} \\
M^{k,l} &= \frac{\mathrm{T}(\lambda_k)}{\mathrm{L}'(\lambda_k)(\lambda_k - t_l)} \label{eq:3.75}
\end{align}
とおいた．上の有理式 $M_l, M^{k,l}$ は我々の計算において重要な役割を果たす．

以上の準備のもとで定理を書く．この定理は，2階線型常微分方程式のホロノミックな変形を考察するときの指導原理となる．
\end{story}

\begin{theorem}[ガルニエ系のハミルトン系表示]
線型常微分方程式(3.33), (3.67), (3.68)に関するフックスの問題は，完全積分可能な多時間ハミルトン系
\begin{equation*}
G_g \quad \frac{\partial \lambda_k}{\partial t_l} = \frac{\partial H_l}{\partial \mu_k}, \quad \frac{\partial \mu_k}{\partial t_l} = - \frac{\partial H_l}{\partial \lambda_k}
\end{equation*}
によって解かれる．ここで，$H_l$ は(3.73)の有理式で与えられる．
\end{theorem}

\begin{definition}[$g$-次ガルニエ系]
ハミルトニアン(3.73)に関する多時間ハミルトン系 $G_g$ を $g$-次ガルニエ系という．
\end{definition}

\begin{story}
定理3.1は，定理3.2で $g=1$ の場合である．とくに $G_g$ は(3.64)に帰着する．
\end{story}

\begin{story}
る．多時間ハミルトン系 $G_g$ から $\mu_k$ を消去すれば，ガルニエ系，すなわち上に書いた長い2階偏微分方程式系，が得られることになる．

定理3.2 は，微分方程式(3.33)を同値な SL 型微分方程式
\begin{equation}
\frac{d^2 y}{dx^2} = r(x,t)y \label{eq:3.76}
\end{equation}
に変換し，方程式(3.76)のモノドロミー保存変形を計算することによって証明される．この微分方程式のリーマン図式は次のようになっている．
\begin{equation*}
\left\{
\begin{array}{ccccc}
x=0 & x=1 & x=t_l & x=\lambda_k & x=\infty \\
\frac{1+\kappa_0}{2} & \frac{1+\kappa_1}{2} & \frac{1+\theta_l}{2} & \frac{3}{2} & \frac{1+\kappa_\infty}{2} \\
\frac{1-\kappa_0}{2} & \frac{1-\kappa_1}{2} & \frac{1-\theta_l}{2} & -\frac{1}{2} & \frac{1-\kappa_\infty}{2}
\end{array}
\right\}
\end{equation*}

節を改めて，変数 $t=(t_1, \cdots, t_g) \in U$ についてのモノドロミー保存変形の考察を続ける．$U$ は $\mathbb{C}^g$ の適当な領域である．

この節の残りの部分では，125ページに引き続いて，トーラス $T$ 上定義された，SL 型2階線型常微分方程式
\begin{equation}
\frac{d^2 y}{dx^2} = r(x,t)y \tag{3.76}
\end{equation}
のモノドロミー保存変形について簡単に触れておこう．$T$ 上，この微分方程式が $g+1$ 個の確定特異点をもつとする．確定特異点の集合を前のように $\Xi$ と書き，$X = T \setminus \Xi$ とおく．SL 型微分方程式のモノドロミーは，基本群 $\pi(X)$ の $\mathbb{C}^2$ 上の表現類であるが，各回路行列の行列式は1である．このとき，モノドロミーの含む独立なパラメータの個数は
\begin{equation*}
\mathrm{M}^{(1)} = 3g+3
\end{equation*}
であることがわかる．一方，SL 型微分方程式(3.76)は
\begin{equation*}
\mathrm{E}^{(1)} = 2g+2
\end{equation*}
個のパラメータに依存する．したがって，微分方程式(3.76)は，$\Xi$ に加えて
\begin{equation*}
g+1 = \mathrm{M}^{(1)} - \mathrm{E}^{(1)}
\end{equation*}
\end{story}

\begin{Q}[【核心】行列式の制約がなければパラメータは $4g+5$ 個になるという計算と、テキストの $M^{(1)} = 3g+3$ の間にある「次元のギャップ」の正体は何か？]
\textbf{A.} あなたの $4g+5$ という計算は $GL(2, \mathbb{C})$ の表現空間として「完全に正解」である。ギャップが生じた理由は、ここから $SL(2, \mathbb{C})$（行列式=1）に制限する際に失われる自由度が、$g+1$ 個ではなく「$g+2$ 個」だからである。

\begin{proof}
まず、このページから考察の舞台が「球面（$\mathbb{P}^1$）」から\textbf{「トーラス（$T$）」}へと移行していることに注意する。
$g+1$ 個の特異点（穴）を持つトーラスの基本群 $\pi(X)$ は、穴の周りのループ $\gamma_1, \dots, \gamma_{g+1}$ に加えて、トーラス特有の縦横のループ $\alpha, \beta$ の、合計 \textbf{$g+3$ 個の生成元}を持つ。
これらが満たすべき基本関係式は $[\alpha, \beta]\gamma_1 \cdots \gamma_{g+1} = I$ である。

\textbf{① あなたの計算の検証（$GL(2, \mathbb{C})$ の場合）}
各生成元を $GL(2, \mathbb{C})$ （自由度4）の行列に対応させると、初期の自由度は $4(g+3) = 4g+12$ である。
関係式が1つの行列方程式（自由度4の拘束）を与え、さらに全体の基底の取り替え（共役変換）として $PGL(2, \mathbb{C})$ （自由度3）で割るため、独立なパラメータ数は
\begin{equation*}
4g+12 - 4 - 3 = \mathbf{4g+5}
\end{equation*}
となり、あなたの導き出した数字は幾何学的に完璧である。

\textbf{② $SL(2, \mathbb{C})$ への制限（行列式 $= 1$ の条件）}
ここから「すべての生成元の行列式を 1 にする」という条件を課す。生成元は $g+3$ 個あるため、$g+3$ 個の制約がつくように思える。
しかし、基本関係式 $[\alpha, \beta]\gamma_1 \cdots \gamma_{g+1} = I$ の両辺の行列式をとると、交換子 $[\alpha, \beta]$ の行列式は必ず 1 になるため、自動的に $\det(\gamma_1) \cdots \det(\gamma_{g+1}) = 1$ が成り立つ。
つまり、$g+3$ 個の行列式のうち1つは他の結果から自動的に決まるため、独立な制約（減る自由度）は $(g+3) - 1 = \mathbf{g+2 \textbf{ 個}}$ となる。

したがって、モノドロミーのパラメータ数は
\begin{equation*}
(4g+5) - (g+2) = \mathbf{3g+3}
\end{equation*}
となり、テキストの $\mathrm{M}^{(1)}$ と見事に一致する。想定した「$g+1$ 個の制限」は、トーラス特有のループ $\alpha, \beta$ に対する行列式の制限を見落としていたために生じたズレであった。
\end{proof}
\end{Q}

\begin{Q}[【核心】パラメータのカウント $\mathrm{E}^{(1)} = 2g+2$ を計算した時点では、方程式 (3.76) には「見かけの特異点」が存在していないと仮定しているのか？]
\textbf{A.} その通りである。そして、モノドロミーの次元 $\mathrm{M}^{(1)}$ との「ギャップ」を埋めるために、まさにこれから見かけの特異点を追加しようとしているのがテキストの切れた直後の文章である。

\begin{proof}
方程式 (3.76) のパラメータ数 $\mathrm{E}^{(1)} = 2g+2$ は、純粋に「$g+1$ 個の確定特異点 $\Xi$ しか持たない SL 型方程式（見かけの特異点ゼロ）」の自由度を計算したものである。
しかし前節で確認した通り、任意のモノドロミー表現（自由度 $\mathrm{M}^{(1)} = 3g+3$）を網羅するためには、方程式側の自由度が足りない。
したがって、テキストの切れた直後の文章は以下のように続く。
\begin{quote}
「したがって、微分方程式(3.76)は、$\Xi$ に加えて \textbf{$g+1 = \mathrm{M}^{(1)} - \mathrm{E}^{(1)}$ 個の見かけの特異点をもつ} と仮定しなければならない。」
\end{quote}
まさにあなたの読み通り、この次元のギャップ「$g+1$」こそが、トーラス上の SL 型方程式において導入すべき見かけの特異点の個数を数学的に要請しているのである。
\end{proof}
\end{Q}

\begin{Q}[【核心】GL型からSL型に変換した際に減る $g+1$ 個のパラメータは、「$g+1$ 個の各特異点における特性指数の和が、すべて一定（1）に固定されるから」と解釈してよいか？]
\textbf{A.} 完璧な幾何学的解釈である。あなたの直感通り、SL型（$p_1=0$）という条件は、すべての特異点での「指数の和」の自由度を奪い去る操作に他ならない。

\begin{proof}
一般の GL 型微分方程式 $y'' + p_1(x)y' + p_2(x)y = 0$ において、特異点 $x_i$ での決定方程式を考えると、2つの特性指数 $\rho_1, \rho_2$ の和は以下のようになる。
\begin{equation*}
\rho_1 + \rho_2 = 1 - \operatorname{Res}_{x=x_i} p_1(x)
\end{equation*}
ここから SL 型の方程式（$p_1(x) = 0$）に変換するということは、特異点における $p_1(x)$ の留数がすべて厳密に $0$ になることを意味する。したがって、あなたの予想通り、\textbf{すべての特異点において特性指数の和は「1」に固定される。}

この制約によって失われるパラメータがなぜ正確に「$g+1$ 個」になるのかは、トーラス上の1次微分形式 $p_1(x)dx$ の構造を見れば明らかになる。
$p_1(x)dx$ は $g+1$ 個の特異点（単純極）を持つ。
\begin{itemize}
    \item \textbf{留数の自由度（$g$ 個）:} $g+1$ 個の極での留数を自由に選べるが、「閉曲面（トーラス）上の留数の総和は必ず $0$ になる」という留数定理の制約があるため、独立に選べる留数は $(g+1) - 1 = \mathbf{g \textbf{ 個}}$ である。これらが指数の和を変動させる自由度である。
    \item \textbf{正則部分の自由度（1 個）:} 極を持たない正則な1次微分形式（トーラス上の $dx$）の自由度がトーラスの種数と同じ \textbf{1 個} だけ存在する。これはすべての指数の和を「大域的に定数シフト」する自由度に対応する。
\end{itemize}
これらを合わせた $g + 1$ 個の自由度こそが $p_1(x)$ が持っていたパラメータのすべてであり、SL 型への変換（$p_1(x) \to 0$）によってこれらが丸ごと消滅する。
あなたの「指数の和が一定になるから $g+1$ 個減る」という考察は、リーマン面の代数幾何学的な本質を見事に言い当てているのである。
\end{proof}
\end{Q}

\begin{story}
（前ページからの続き）
$g+1$ 個の見かけの特異点をもつ，と仮定し，この下でモノドロミー保存変形を考えよう．種数 $p \geqq 2$ の一般のコンパクトリーマン面 $R$ の場合と異なり，我々の考察の特徴は，$p=1$ のとき，すなわちトーラス $T$ 上の微分方程式は2重周期関数を用いてすべての量を具体的に書き下すことができる，ということである．微分方程式(3.76)の係数 $r(x, t)$ は周期 $2\omega_1, 2\omega_3$ をもつ2重周期関数とし，$\Omega$ を $2\omega_1$ と $2\omega_3$ で生成される $\mathbf{C}$ 内の周期格子であるとする．$\mathbf{C}$ における離散群 $\Omega$ の基本領域，すなわち周期平行4辺形 $F$ を1つとる．トーラス $T$ は1次元の解析的同型をもつから，$\Omega$ の各点が微分方程式(3.76)の確定特異点である，としてよい．そこで，$0 \in F$ なるように $F$ を選び，また，(3.76)のその他の確定特異点を $t_l \in F \ (l=1, \cdots, g)$ とする(図3.1)．

このようにして，$t=(t_1, \cdots, t_g) \in U$ を変形のパラメータとしてモノドロミー保存変形の変数を考察する．$U$ は $\mathbf{C}^g$ の適当な領域である．また，$\lambda_k \in F \ (k=1, \cdots, g)$ を(3.76)の見かけの特異点である，とすると，モノドロミー保存変形により，$\lambda_k$ が $t$ の関数として与えられる．この関数を定める微分方程式系は楕円関数を用いることによって書き表されるであろう．
\end{story}

\begin{Q}[【基礎】「トーラス $\boldsymbol{T}$ 上の解析」とは、複素平面 $\mathbf{C}$ の言葉で言うと、結局のところ何をすることなのか？]
\textbf{A.} 複素平面を「平行四辺形のタイル張り」とみなし、「タイル1枚分（基本領域 $\boldsymbol{F}$）の中だけで関数を考えること」に他ならない。トーラス上の関数とは、複素平面上の「2重周期関数」のことである。

\begin{proof}
幾何学的なドーナツ型（トーラス）を切り開いて平らに伸ばすと、1枚の四角形になる。逆に言えば、複素平面上に平行四辺形（基本領域 $\boldsymbol{F}$）を描き、その「上の辺と下の辺」「右の辺と左の辺」を同一視して貼り合わせるとトーラスができる。（いわゆる昔のゲームの「パックマン」のように、画面の右端に行くと左端から出てくる世界である）。

この世界において、ある点 $x$ から右に $2\omega_1$ だけ進むか、上に $2\omega_3$ だけ進むと、元の場所に戻ってくる。これを式で書くと、トーラス上で定義された関数 $f(x)$ は
\begin{equation*}
f(x + 2\omega_1) = f(x), \quad f(x + 2\omega_3) = f(x)
\end{equation*}
を満たさなければならない。このような2つの複素周期を持つ関数を「2重周期関数（または楕円関数）」と呼ぶ。
つまり、「トーラス上の微分方程式」と言われたら、幾何学的なドーナツを想像するのではなく、\textbf{「係数 $r(x,t)$ が2重周期関数であるような、普通の複素平面上の微分方程式」}と脳内変換すれば、解析学の言葉として完全に等価になるのである。テキストの図3.1は、まさにこの「タイル1枚分」の領土を描いている。
\end{proof}
\end{Q}

\begin{Q}[【大局観】著者は「種数 $p \geqq 2$ の場合と異なり、$p=1$（トーラス）のときはすべての量を具体的に書き下すことができる」と強調しているが、なぜ $p=1$ はそんなに特別なのか？]
\textbf{A.} $p=1$ の場合（楕円関数論）には、ワイエルシュトラスの $\wp$（ペー）関数という「万能のブロック」が存在し、すべての関数を極めて具体的な数式で組み立てることができるからである。

\begin{proof}
解析学において、特異点（極）を持たない正則な関数が2重周期を持つと、リウヴィルの定理により「単なる定数」になってしまう。したがって、トーラス上で意味のある関数は必ず「極」を持たなければならない。
19世紀の数学者ワイエルシュトラスは、トーラスの原点 $x=0$ に2位の極を持つ最も基本的で美しい2重周期関数 $\wp(x)$ を具体的に構成した。

驚くべきことに、\textbf{「トーラス上のあらゆる有理型関数（2重周期関数）は、$\wp(x)$ とその導関数 $\wp'(x)$ の有理式として完全に書き下せる」}ことが証明されている。
種数が $p \geqq 2$ （穴が2つ以上のプレッツェルのような図形）になると、このような「1変数のシンプルな万能関数」は存在せず、テータ関数などの多変数関数を用いた非常に抽象的で計算が困難な表現にならざるを得ない。

著者がここで喜んでいるのは、「舞台をトーラスに設定したことで、これから微分方程式の係数 $r(x,t)$ や、見かけの特異点 $\lambda_k(t)$ の動きを、正体不明の抽象的な関数ではなく、$\wp(x)$ という人類が隅々まで性質を知り尽くした具体的な関数（楕円関数）を使って、ゴリゴリと直接計算できるぞ！」という宣言なのである。
\end{proof}
\end{Q}
%#endregion

\subsection{ガルニエ系のハミルトン表示}

\begin{story}

以下この節において定理3.2の証明を行う。ただし、細かい計算には立ち入らず、筋道を示すことを第1の目標としよう。記号等は前節のものを踏襲して用いる。

まず，ハミルトニアン $H_l$ の形(3.73)を，フロベニウスの方法により決定する
\end{story}

\begin{Q}[【核心】「ハミルトニアン $H_l$ の形を、フロベニウスの方法により決定する」とは具体的にどういうことか？]
\textbf{A.} これは、前節までの「見かけの特異点 $x=\lambda_k$（特性指数 $0, 2$）」における局所的な解の振る舞いから、系全体を統制するエネルギー関数（ハミルトニアン）を逆算して導き出す、という驚くべきプロセスを指しています。
「フロベニウスの方法」とは、特異点周りで微分方程式の解を冪級数展開して求める標準的な手法です。特性指数の差が整数の場合、通常は解に「対数項（$\log(x-\lambda_k)$）」が現れます。しかし、$x=\lambda_k$ は「見かけの特異点」であるため、この対数項の係数は厳密に $0$ にならなければなりません。
この「対数項の係数が $0$ になる」という代数的な条件式を書き下すと、そこに変数 $\lambda_k$ と留数 $\mu_k$、そして変形パラメータ $t_l$ が結びついた関係式が浮かび上がります。この関係式を整理して得られる多項式（または有理関数）こそが、ハミルトニアン $H_l$ そのものになるのです。局所的な「対数項回避の条件」が、大域的な「ハミルトン力学系の生成子」を決定するというのが、岡本理論の最も美しいハイライトの一つです。
\end{Q}

\begin{Q}[【深掘り】そもそも「ガルニエ系 (Garnier system)」とは何か？]
\textbf{A.} 一言で言えば、「パンルヴェ方程式の多変数化（高次元化）」です。
パンルヴェ方程式は、独立変数が $t$ の1つだけ（方程式の特異点が例えば $0, 1, t, \infty$ の4つ）である場合のモノドロミー保存変形を記述する非線形常微分方程式でした。
ガルニエ系は、動かせる特異点（変形パラメータ）を $t_1, t_2, \dots, t_g$ と複数個に増やしたときに現れる、非線形偏微分方程式系です。変数が多次元になっても、見かけの特異点 $(\lambda_k, \mu_k)$ を正準変数とみなすことで、時間発展 $\partial/\partial t_l$ がハミルトニアン $H_l$ によって完全に統制される「多変数ハミルトン系」として綺麗に書き表すことができます。この節では、そのハミルトニアンが実際にどのような数式になるかを導出していくことになります。
\end{Q}

\begin{story}
$x=\lambda_k$ を微分方程式 (3.33) の見かけの特異点とし、局所座標 $u = x - \lambda_k$ をとる。$u=0$ の近傍で (3.33) を
\begin{equation}
    u^2 \frac{d^2 y}{du^2} + q_1(u,t)u \frac{dy}{du} + q_2(u,t)y = 0 \label{eq:3.77}
\end{equation}
と書く。確定特異点 $u=0$ における特性指数は $0, 2$ であり、これが非対数的であるための条件は、ベキ級数解
\begin{equation}
    y = 1 + y_1 u + y_2 u^2 + \cdots \label{eq:3.78}
\end{equation}
が存在することである。

\begin{proposition}[対数項を持たないための適合条件]
\eqref{eq:3.77} の係数を次のように展開する。
\begin{equation*}
    q_1(u,t) = -1 + a_k u + a'_k u^2 + \cdots, \quad q_2(u,t) = \mu_k u + b_k u^2 + b'_k u^3 + \cdots
\end{equation*}
この表示と \eqref{eq:3.78} を \eqref{eq:3.77} に代入し、未定係数法により比較すると、
\begin{equation*}
    y_1 = \mu_k, \quad j(j-2)y_j = R_j(y_1, \cdots, y_{j-1}) \quad (j=2, \cdots)
\end{equation*}
となるから、ベキ級数解 \eqref{eq:3.78} が存在するための条件は、$j=2$ の場合を考えて
\begin{equation}
    R_2(y_1) = \mu_k y_1 + a_k y_1 + b_k = \mu_k^2 + a_k \mu_k + b_k = 0 \label{eq:3.79}
\end{equation}
であることがわかる。
\end{proposition}

\end{story}

\begin{Q}[【核心】なぜ $j(j-2)y_j = R_j$ において $j=2$ のとき、$R_2(y_1) = 0$ が「対数項を持たない条件」になるのか？]
\textbf{A.} これはフロベニウス法の決定方程式の構造に由来します。特性指数が $\rho = 0, 2$ ということは、微分方程式の左辺に $u^\rho$ を代入したときの最低次数の係数（決定方程式）が $\rho(\rho-2)=0$ になるということです。
特性指数 $\rho=0$ に対応する級数解を求めようとして係数 $y_j$ を順に決めていくと、漸化式の左辺は $j(j-2)y_j$ となります。$j=1$ のときは $1 \cdot (-1) y_1 = R_1$ なので $y_1$ は一意に決まりますが、$j=2$ のステップに到達した瞬間、左辺は $2(0)y_2 = 0 \cdot y_2 = 0$ となってしまいます。
もしこのとき右辺の $R_2$ が $0$ でなければ、「$0 = (\text{0でない値})$」という矛盾が生じ、$y_2$ を決定することができません。この矛盾を回避して解を構成するためには、どうしても独立な解として $u^2 \log u$ のような「対数項」を付け加える必要が出てきます。
逆に言えば、対数項が現れない（見かけの特異点である）ためには、どうしても右辺が $R_2 = 0$ になってくれて、$0 \cdot y_2 = 0$（$y_2$ は任意定数としてよい）となる必要があるのです。この $R_2=0$ こそが、特異点が見かけ上のものであるための強烈な代数的制約（適合条件）です。
\end{Q}

\begin{story}
\begin{theorem}[ハミルトニアンを決定する連立1次方程式]
一方、元の方程式の係数を用いて $a_k, b_k$ を具体的に計算すると、
\begin{align*}
    a_k &= \frac{1-\kappa_0}{\lambda_k} + \frac{1-\kappa_1}{\lambda_k-1} + \sum_{l=1}^g \frac{1-\theta_l}{\lambda_k-t_l} - \sum_{p=1(p \neq k)}^g \frac{1}{\lambda_k-\lambda_p} \\
    b_k &= \frac{\kappa}{\lambda_k(\lambda_k-1)} + \sum_{p=1(p \neq k)}^g \frac{\lambda_p(\lambda_p-1)\mu_p}{\lambda_k(\lambda_k-1)(\lambda_k-\lambda_p)} - \frac{2\lambda_k-1}{\lambda_k(\lambda_k-1)}\mu_k - \sum_{l=1}^g \frac{t_l(t_l-1)H_l}{\lambda_k(\lambda_k-1)(\lambda_k-t_l)}
\end{align*}
である。これを \eqref{eq:3.79} に代入して整理した式は、
\begin{equation}
    \sum_{l=1}^g E_{kl} H_l = \mu_k^2 + \hat{a}_k \mu_k + \hat{b}_k \label{eq:3.80}
\end{equation}
\begin{equation*}
    E_{kl} = \frac{t_l(t_l-1)}{\lambda_k(\lambda_k-1)(\lambda_k-t_l)}
\end{equation*}
という形をしているから、これを $H_l$ に関する連立1次方程式と思って解けばよい。
\end{theorem}
\end{story}

\begin{Q}[【深掘り】この連立方程式 \eqref{eq:3.80} が解けることで、理論上どのような「嬉しいこと」があるのか？]
\textbf{A.} これは「ガルニエ系のハミルトニアン $H_l$ が、正準変数 $(\lambda_k, \mu_k)$ の有理関数として具体的に構成できる」という決定的な事実を示しています。
左辺の係数行列 $E_{kl}$ はコーシー行列の一種であり、$\lambda_k$ と $t_l$ がすべて相異なるという仮定のもとで正則（逆行列を持つ）になります。したがって、この連立1次方程式は必ず一意な解 $H_l$ を持ちます。
局所的な「対数項を消去せよ」という解析的な条件 $R_2=0$ を解きほぐしただけで、系全体の時間発展 $\partial/\partial t_l$ を支配する力学系のエネルギー関数 $H_l$ が代数的に完全に決定されてしまうという、モノドロミー保存変形の理論構造の美しさがここに凝縮されています。
\end{Q}

\begin{story}
\begin{lemma}[3.4]
行列 $E = ((E_{kl}))$ の逆行列を $F = ((F_{lk}))$ とすると、
\begin{equation}
    F_{lk} = M_l M^{k,l} \label{eq:3.81}
\end{equation}
\begin{equation*}
    M_l = - \frac{L(t_l)}{T'(t_l)}, \quad M^{k,l} = \frac{T(\lambda_k)}{L'(\lambda_k)(\lambda_k - t_l)}
\end{equation*}
となる。
\end{lemma}

補題に出てくる、$t, \lambda, \mu$ の有理関数 $M_l, M^{k,l}$ は、それぞれ (3.74), (3.75) で与えられたものである。(3.81) により、(3.80) を $H_l$ について解けば (3.73) が得られる。そのとき、関係式
\begin{equation*}
    \sum_{k=1}^g \frac{M^{k,l}}{\lambda_k(\lambda_k-1)} = 1
\end{equation*}
などを用いて計算することになるが、(3.73) を導く計算の詳細は省略する。

\begin{proof}
線型代数の演習問題であるから、簡単に証明しておこう。
$x$ の有理関数
\begin{equation}
    Z_l(x) = \frac{T(x)}{x(x-1)(x-t_l)L(x)}
\end{equation}
を考える。ここで、$T(x), L(x)$ は、(3.69), (3.70) で与えられる $x$ の多項式、すなわち
\begin{equation*}
    T(x) = x(x-1)\prod_{m=1}^g (x-t_m), \quad L(x) = \prod_{k=1}^g (x-\lambda_k)
\end{equation*}
である。$Z_l(x)$ は $x=\lambda_k \ (k=1, \cdots, g)$ に、1位の極をもつ。さらに、有理関数 $Z_l(x)$ の分母は $g$ 次多項式、分子は $g-1$ 次多項式、であることに注意してこの有理関数を部分分数展開する。実際
\begin{equation*}
    \underset{x=\lambda_k}{\mathrm{Res}} Z_l(x) = \frac{T(\lambda_k)}{\lambda_k(\lambda_k-1)(\lambda_k-t_l)L'(\lambda_k)} = \frac{M^{k,l}}{\lambda_k(\lambda_k-1)}
\end{equation*}
より、次の式が得られる。
\begin{equation}
    Z_l(x) = \sum_{k=1}^g \frac{M^{k,l}}{\lambda_k(\lambda_k-1)} \cdot \frac{1}{x-\lambda_k}
\end{equation}
ここで、まず
\begin{equation*}
    Z_l(t_h) = \sum_{k=1}^g \frac{M^{k,l}}{\lambda_k(\lambda_k-1)} \cdot \frac{1}{t_h-\lambda_k} = - \sum_{k=1}^g \frac{M^{k,l} E_{kh}}{t_h(t_h-1)}
\end{equation*}
となることに気がつく．他方，定義式から
\begin{equation*}
    Z_l(t_l) = \frac{T'(t_l)}{t_l(t_l-1)L(t_l)} = \frac{-1}{t_l(t_l-1)M_l}
\end{equation*}
\begin{equation*}
    Z_l(t_h) = 0 \quad (h \neq l)
\end{equation*}
であるから，結局
\begin{align*}
    1 &= \sum_{k=1}^g E_{kl} M_l M^{k,l} \\
    0 &= \sum_{k=1}^g E_{kh} M_l M^{k,l} \quad (h \neq l)
\end{align*}
が得られる．これは(3.81)を意味している．

(3.80), (3.81)から $H_l$ を求めれば(3.73)が得られる。
\end{proof}
\end{story}

\begin{Q}[【核心】なぜ逆行列を求めるために、突然 $Z_l(x)$ という有理関数を考えるのか？]
\textbf{A.} 行列 $E$ と $F$ が互いに逆行列であること（すなわち $\sum E_{kl}F_{lm} = \delta_{km}$）を直接計算で証明するのは、要素が複雑な分数であるため非常に困難です。
しかし、$Z_l(x)$ の定義式に $T(x)$ を代入して約分すると、実は
$$ Z_l(x) = \frac{\prod_{m \neq l}^g (x-t_m)}{L(x)} $$
となり、これは分子が $g-1$ 次、分母が $g$ 次の有理関数になります。
複素解析において「分子の次数が分母の次数より2以上小さい有理関数の、すべての極における留数の和はゼロになる」という強力な定理（留数定理）があります。この定理や、部分分数展開に $x=t_m$ などを代入して得られる「ラグランジュ補間の恒等式」を利用することで、行列の積の和 $\sum$ を「留数の和」にすり替えることができ、複雑な代数計算を一瞬でクロネッカーのデルタ $\delta_{km}$ に帰着させることができるのです。
\end{Q}

\begin{Q}[【深掘り】多項式 $T(x)$ と $L(x)$ はそれぞれ何を意味しているのか？]
\textbf{A.} 多変数（$g$ 変数）のガルニエ系において、多数の特異点を一括して扱うための「母関数（特性多項式）」です。
* $T(x) = x(x-1)\prod (x-t_m)$ は、微分方程式の「確定特異点（動かない特異点 $0, 1$ と、動くパラメータ $t_m$）」を根に持つ多項式です。
* $L(x) = \prod (x-\lambda_k)$ は、「見かけの特異点」を根に持つ多項式です。
このように特異点の位置を多項式の根としてパックすることで、ガルニエ系全体の挙動を「$T(x)$ と $L(x)$ が織りなす1変数の代数関数論」として驚くほど見通し良く記述できるようになります。これは岡本和夫先生の研究の真骨頂とも言える手法です。
\end{Q}

\begin{Q}[【核心】$Z_l(t_h)$ に $x=t_h$ を代入するだけで、なぜ逆行列の証明が終わるのか？]
\textbf{A.} これは部分分数分解と「極（分母が0になる点）」の性質を巧みに利用したマジックです。
有理関数 $Z_l(x)$ の定義式を見ると、分子に $T(x) \propto (x-t_1)\cdots(x-t_g)$ があります。したがって、$h \neq l$ のときは単に $T(t_h) = 0$ となり $Z_l(t_h) = 0$ になります。
一方、$h = l$ のときは、分母にも $(x-t_l)$ があるため $0/0$ の不定形になります。ロピタルの定理のように約分して極限をとることで、分子は微分された $T'(t_l)$ となり、$0$ ではない値が残ります。
これを、部分分数展開した側（$\sum$ の形）と比較します。すると、「行列 $E$ と、$M_l M^{k,l}$ からなる行列 $F$ の積」が、「$h=l$ のときは $1$、$h \neq l$ のときは $0$」すなわち単位行列（クロネッカーのデルタ $\delta_{hl}$）になることが、一切の泥臭い連立方程式を解くことなく、極限の評価だけで鮮やかに証明されるのです。
\end{Q}

前節までの議論により、ハミルトニアン $H_h$ は以下の連立1次方程式を満たすことがわかっている。
\begin{equation*}
    \sum_{h=1}^g E_{kh} H_h = \mu_k^2 + \hat{a}_k \mu_k + \hat{b}_k
\end{equation*}
ここで $\hat{a}_k, \hat{b}_k$ は $H_h$ を含まない項であり、具体的には次で与えられる。
\begin{align*}
    \hat{a}_k &= \frac{1-\kappa_0}{\lambda_k} + \frac{1-\kappa_1}{\lambda_k-1} + \sum_{h=1}^g \frac{1-\theta_h}{\lambda_k-t_h} - \sum_{p=1 (p \neq k)}^g \frac{1}{\lambda_k-\lambda_p} \\
    \hat{b}_k &= \frac{\kappa}{\lambda_k(\lambda_k-1)} - \frac{2\lambda_k-1}{\lambda_k(\lambda_k-1)}\mu_k + \sum_{p=1 (p \neq k)}^g \frac{\lambda_p(\lambda_p-1)\mu_p}{\lambda_k(\lambda_k-1)(\lambda_k-\lambda_p)}
\end{align*}
行列 $E$ の逆行列 $F_{lk} = M_l M^{k,l}$ を左から掛けることで、$H_l$ は次のように表される。
\begin{equation}
    H_l = M_l \sum_{k=1}^g M^{k,l} \left( \mu_k^2 + \hat{a}_k \mu_k + \hat{b}_k \right) \label{eq:Hl_base}
\end{equation}

\begin{theorem}[ハミルトニアンの最終形態]
式 \eqref{eq:Hl_base} を $\mu_k$ について整理すると、以下の式 (3.73) が得られる。
\begin{equation}
    H_l = M_l \left[ \sum_{k=1}^g M^{k,l} \left\{ \mu_k^2 - A_{k,l}\mu_k + \frac{\kappa}{\lambda_k(\lambda_k-1)} \right\} \right] \label{eq:3.73}
\end{equation}
\begin{equation*}
    A_{k,l} = \frac{\kappa_0}{\lambda_k} + \frac{\kappa_1}{\lambda_k-1} + \sum_{h=1}^g \frac{\theta_h - \delta_{hl}}{\lambda_k - t_h}
\end{equation*}
\end{theorem}

\begin{proof}[証明と係数比較]
式 \eqref{eq:Hl_base} の括弧内を展開し、$\mu_k^2$ の項、$\mu_k^0$（定数）の項、$\mu_k^1$ の項に分けて比較する。
\begin{itemize}
    \item \textbf{2次の項}: そのまま $\mu_k^2$ となり一致する。
    \item \textbf{定数項}: $\hat{b}_k$ の第1項 $\frac{\kappa}{\lambda_k(\lambda_k-1)}$ がそのまま現れるため一致する。
    \item \textbf{1次の項}: ここが非自明である。$\hat{b}_k$ の二重和の部分は、添字 $k$ と $p$ を入れ替えることで次のように変形できる。
    \begin{equation*}
        \sum_{k=1}^g M^{k,l} \sum_{p \neq k} \frac{\lambda_p(\lambda_p-1)\mu_p}{\lambda_k(\lambda_k-1)(\lambda_k-\lambda_p)} = \sum_{k=1}^g \mu_k \sum_{p \neq k} M^{p,l} \frac{\lambda_k(\lambda_k-1)}{\lambda_p(\lambda_p-1)(\lambda_p-\lambda_k)}
    \end{equation*}
    したがって、式 \eqref{eq:Hl_base} における $\mu_k$ の全体の係数（$M^{k,l}$ をくくり出したもの）を $C_k$ とすると、
    \begin{equation}
        M^{k,l} C_k = M^{k,l} \left( \hat{a}_k - \frac{2\lambda_k-1}{\lambda_k(\lambda_k-1)} \right) + \sum_{p \neq k} M^{p,l} \frac{\lambda_k(\lambda_k-1)}{\lambda_p(\lambda_p-1)(\lambda_p-\lambda_k)}
    \end{equation}
    となる。この $C_k$ がちょうど $-A_{k,l}$ に等しいことを示せばよい。
\end{itemize}
代数的な計算により、$\hat{a}_k - \frac{2\lambda_k-1}{\lambda_k(\lambda_k-1)} = - A_{k,l} + \sum_{h \neq l} \frac{1}{\lambda_k-t_h} - \sum_{p \neq k} \frac{1}{\lambda_k-\lambda_p}$ がわかる。これを代入し、$M^{k,l}C_k = M^{k,l}(-A_{k,l})$ を示すことは、以下の等式を示すことと同値である。
\begin{equation}
    \sum_{h \neq l} \frac{1}{\lambda_k-t_h} - \sum_{p \neq k} \frac{1}{\lambda_k-\lambda_p} = \sum_{p \neq k} \frac{M^{p,l} \lambda_k(\lambda_k-1)}{M^{k,l} \lambda_p(\lambda_p-1)(\lambda_k-\lambda_p)} \label{eq:identity_target}
\end{equation}
多項式 $P(x) = \prod_{m \neq l}(x-t_m)$ および $L(x) = \prod (x-\lambda_q)$ を用いると、$\frac{M^{k,l}}{\lambda_k(\lambda_k-1)} = \frac{P(\lambda_k)}{L'(\lambda_k)}$ と書けるため、式 \eqref{eq:identity_target} の右辺は
\begin{equation*}
    \sum_{p \neq k} \frac{P(\lambda_p)L'(\lambda_k)}{P(\lambda_k)L'(\lambda_p)(\lambda_k-\lambda_p)}
\end{equation*}
となる。これは、有理関数の部分分数展開（ラグランジュ補間の微分の公式）から得られる恒等式と完全に一致する。よって $C_k = -A_{k,l}$ が証明され、式 (3.73) が導出された。
\end{proof}

\begin{Q}[【核心】なぜ二重和の計算で「添字 $k$ と $p$ の入れ替え」を行う必要があるのか？]
\textbf{A.} 行列の積を計算して $H_l = \sum_k F_{lk} (\cdots)$ と書き下した段階では、和の主役は $k$ です。しかし、$\hat{b}_k$ の中にさらに別の和 $\sum_{p \neq k} \mu_p$ が含まれているため、このままでは $\mu_1, \mu_2, \dots$ がそれぞれの項に散らばってしまいます。
最終目標である式 (3.73) は「各 $k$ について $\mu_k^2$ と $\mu_k^1$ を綺麗にまとめた形」をしています。したがって、散らばった $\mu$ を主役（係数）としてまとめ直すために、和の順番を $\sum_k \sum_p$ から $\sum_p \sum_k$ へと入れ替え、形式的に $p$ を $k$ に読み替えることで、全体の $\mu_k$ の係数を一つに集約させているのです。
\end{Q}

\begin{Q}[【深掘り】証明の最後に登場した「ラグランジュ補間の微分の公式」とは何か？]
\textbf{A.} これは代数幾何や可積分系において、複雑な有理関数の和を一瞬で計算するための強力な恒等式です。
一般に、$g-1$ 次の多項式 $P(x)$ は、ラグランジュ補間により $P(x) = \sum_{p=1}^g \frac{P(\lambda_p)}{L'(\lambda_p)} \frac{L(x)}{x-\lambda_p}$ と展開できます。この両辺を $x$ で微分し、$x = \lambda_k$ を代入して整理すると、
\begin{equation*}
    \frac{P'(\lambda_k)}{P(\lambda_k)} - \frac{L''(\lambda_k)}{2L'(\lambda_k)} = \sum_{p \neq k} \frac{P(\lambda_p)L'(\lambda_k)}{P(\lambda_k)L'(\lambda_p)(\lambda_k-\lambda_p)}
\end{equation*}
という美しい関係式が自動的に得られます。左辺の対数微分 $\frac{P'}{P}$ や $\frac{L''}{L'}$ を計算すると、まさに式 \eqref{eq:identity_target} の左辺のシグマの差そのものになります。パンルヴェ理論のハミルトニアンがこれほど綺麗にまとまる背後には、こうした代数関数論の深い恒等式が隠れているのです。
\end{Q}

\begin{story}

以下，前節から考察している2階フックス型線型常微分方程式
\begin{equation}
    \frac{d^2 y}{dx^2} + p_1(x,t) \frac{dy}{dx} + p_2(x,t)y = 0 \label{eq:3.33}
\end{equation}
のモノドロミー保存変形の計算に戻ろう．係数は(3.67), (3.68)により与えられていた．我々は，$x=\lambda_k$ が見かけの特異点である，という仮定を使って，$H_l$ を他の変数の有理関数として具体的に求めたところである．

ここで，第3.2節において調べたことを思いだそう．(3.33)に対して
\begin{equation}
    r(x,t) = -p_2(x,t) + \frac{1}{4}p_1(x,t)^2 + \frac{1}{2}\frac{\partial}{\partial x}p_1(x,t) \label{eq:3.32}
\end{equation}
とし，SL型方程式
\begin{equation}
    \frac{d^2 z}{dx^2} = r(x,t)z \label{eq:3.25}
\end{equation}
を併せて考える．ここで124ページの考察を補題の形にまとめておこう．

\begin{lemma}[3.5]
(3.33)がモノドロミー保存変形を定めることと，(3.25)がそうなることは同値であり，さらにこれは，以下の偏微分方程式系が $x$ の有理関数を解としてもつことと同値である．

\begin{align}
    \frac{\partial^3}{\partial x^3}w_l - 4r(x,t)\frac{\partial}{\partial x}w_l - 2\frac{\partial r}{\partial x}(x,t)w_l + 2\frac{\partial r}{\partial t_l}(x,t) = 0 \label{eq:3.82} \\
    \left( \frac{\partial}{\partial t_h} - w_h \frac{\partial}{\partial x} \right) w_l = \left( \frac{\partial}{\partial t_l} - w_l \frac{\partial}{\partial x} \right) w_h \quad (h,l = 1,\cdots,g) \label{eq:3.83}
\end{align}

\end{lemma}
\end{story}

\begin{Q}[【深掘り】なぜわざわざ元の微分方程式 (3.33) を、1階微分の項がないSL型 (3.25) に変換して考える必要があるのか？]
\textbf{A.} モノドロミー保存変形の条件（Lax方程式などの非線形偏微分方程式系）を導出する際、**1階微分の項 $\frac{dy}{dx}$ が存在すると計算が絶望的に複雑になるから**です。
実は、未知関数 $y$ に対して $z = y \exp(\frac{1}{2}\int p_1 dx)$ という変換（ゲージ変換）を行うと、1階微分の項が完全に消去され、$z'' = r(x,t)z$ というシンプルな形（Sturm-Liouville型、あるいはシュレディンガー方程式型）に帰着させることができます。
この変換は、解の「特異点周りでの回り方（モノドロミー表現）」の本質を一切変えません。そのため、「モノドロミーが保存される」という幾何学的な条件を計算するときは、見通しの良いSL型 (3.25) を使って変形方程式を書き下すのが、可積分系理論における絶対的な定石なのです。(3.32) の $r(x,t)$ の式は、まさにそのゲージ変換から生じるポテンシャル項を表しています。
\end{Q}

\begin{story}
実際，命題3.8により(3.82), (3.83)の有理関数解 $w_l = a_2^{(l)}(x,t)$ は存在すれば唯一であり，それに対して $a_1^{(l)}(x,t)$ を(3.23)により定めると，変形方程式系
\begin{equation}
\left\{
\begin{array}{l}
    \displaystyle \frac{\partial^2 \vec{\phi}}{\partial x^2} + p_1(x,t)\frac{\partial \vec{\phi}}{\partial x} + p_2(x,t)\vec{\phi} = 0 \vspace{0.5em} \\
    \displaystyle \frac{\partial \vec{\phi}}{\partial t_l} = a_1^{(l)}(x,t)\vec{\phi} + a_2^{(l)}(x,t)\frac{\partial \vec{\phi}}{\partial x}
\end{array}
\right. \label{eq:3.20}
\end{equation}
は完全積分可能である．さらに次の命題も成り立つ．

\begin{proposition}[3.14]
偏微分方程式系(3.82)の有理関数解 $w_l$ は，(3.83)を満たす．
\end{proposition}

この命題により，モノドロミー保存変形は(3.82)の考察に帰着される．証明のために
\begin{equation*}
    L = \frac{\partial^3}{\partial x^3} - 4r(x,t)\frac{\partial}{\partial x} - 2\frac{\partial r}{\partial x}(x,t)
\end{equation*}
とおく．計算によって次の補題の成立を確認することができる．

\begin{lemma}[3.6]
任意の解析関数 $f,g$ に対して次式が成り立つ．
\begin{equation*}
    2\frac{\partial f}{\partial x}L(g) - 2\frac{\partial g}{\partial x}L(f) + f\frac{\partial}{\partial x}L(g) - g\frac{\partial}{\partial x}L(f) = L\left( f\frac{\partial g}{\partial x} - g\frac{\partial f}{\partial x} \right)
\end{equation*}
\end{lemma}

\begin{proof}[命題3.14の証明]
(3.82)に有理関数解 $w_l$ が存在するとして
\begin{equation*}
    L(w_l) + 2\frac{\partial r}{\partial t_l}(x,t) = 0
\end{equation*}
から，積分可能条件として
\begin{equation}
    0 = \frac{\partial}{\partial t_h}L(w_l) - \frac{\partial}{\partial t_l}L(w_h) = L\left( \frac{\partial w_l}{\partial t_h} - \frac{\partial w_h}{\partial t_l} \right) + L_{hl} \label{eq:3.84}
\end{equation}
が成り立っている．ここで
\begin{equation*}
    L_{hl} = -4\frac{\partial r}{\partial t_h}\frac{\partial w_l}{\partial x} + 4\frac{\partial r}{\partial t_l}\frac{\partial w_h}{\partial x} - 2w_l\frac{\partial^2 r}{\partial x \partial t_h} + 2w_h\frac{\partial^2 r}{\partial x \partial t_l}
\end{equation*}
であるが，計算して補題3.6を使うと
\begin{align*}
    L_{hl} &= 2\frac{\partial w_l}{\partial x}L(w_h) - 2\frac{\partial w_h}{\partial x}L(w_l) + w_l\frac{\partial}{\partial x}L(w_h) - w_h\frac{\partial}{\partial x}L(w_l) \\
    &= L\left( w_l\frac{\partial w_h}{\partial x} - w_h\frac{\partial w_l}{\partial x} \right)
\end{align*}
したがって(3.84)より
\begin{equation*}
    L\left( \frac{\partial w_l}{\partial t_h} - \frac{\partial w_h}{\partial t_l} + w_l\frac{\partial w_h}{\partial x} - w_h\frac{\partial w_l}{\partial x} \right) = 0
\end{equation*}
となる．命題3.8の証明で見たように，仮定(P2), (P3)によれば，3階線型常微分方程式 $L(w)=0$ は非自明な有理関数解をもたない．すなわち
\begin{equation*}
    \frac{\partial w_l}{\partial t_h} - \frac{\partial w_h}{\partial t_l} + w_l\frac{\partial w_h}{\partial x} - w_h\frac{\partial w_l}{\partial x} = 0
\end{equation*}
以上の議論は，(3.83)が(3.82)の積分可能条件に他ならないことを示している．
\end{proof}

\end{story}

\begin{Q}[【核心】作用素 $L$ と方程式(3.82)は、どこかで見たことがある形をしているが？]
\textbf{A.} お気付きの通り、以前のノートにあった**KdV方程式のLax対の生成子 $M_3$ と本質的に全く同じもの**です。
ノートでは $M_3 = -4\partial_x^3 - 6u\partial_x - 3u_x$ でした。これを定数倍でスケール調整し、$u$ をポテンシャル $r(x,t)$ に読み替えると、まさに $L = \partial_x^3 - 4r\partial_x - 2r_x$ になります。
シュレディンガー作用素 $\partial_x^2 - r(x,t)$ の固有値（あるいはモノドロミー）を保存するような変形を考えると、必ずこの3階の微分作用素 $L$ が現れます。(3.82) は $L(w_l) = -2\partial_{t_l}r$ と書けますが、これはソリトン理論における「Lenard関係式」や「KdV階層」を定義する方程式そのものです。モノドロミー保存変形（ガルニエ系）とは、本質的に「有限自由度のKdV階層」に他ならないことが、この式から明確に読み取れます。
\end{Q}

\begin{Q}[【深掘り】命題3.14は、なぜそれほどまでに重要なのか？]
\textbf{A.} モノドロミー保存変形を成立させるためには、本来 **(3.82)（ポテンシャルの時間発展）と (3.83)（異なる時間 $t_l, t_h$ のフローの可換性）の両方を同時に満たす $w_l$** を探さなければなりません。複数の時間を扱う多変数系（ガルニエ系）において、(3.83)のような非線形の連立微分方程式を解くことは通常至難の業です。
しかし命題3.14は、**「(3.82)の有理関数解を見つけさえすれば、面倒な(3.83)は自動的に満たされる」**という驚くべき事実を保証しています。つまり、問題を解くための労力が半分（あるいはそれ以下）に劇的に減るのです。
この「奇跡」を代数的に証明するための鍵となるのが、補題3.6の恒等式です。作用素 $L$ が持つ美しい代数的な対称性のおかげで、個別の時間発展(3.82)が保証されれば、系全体の無矛盾性(3.83)も自動的に担保される構造になっています。
\end{Q}

\begin{Q}[【核心】作用素 $L$ で括った中身が $=0$ になるのはなぜか？積分定数のようなものは出ないのか？]
\textbf{A.} ここがこの証明の最大の要です。一般の微分方程式 $L(W) = 0$ には当然、非自明な解（積分定数を含む関数の族）が存在します。しかしここで探しているのは「有理関数解」です。
パンルヴェ理論の枠組みでは、特異点の配置に関する一般的な仮定（P2, P3など）を置くことで、「$L(w)=0$ を満たす有理関数は $w=0$ しか存在しない（核が自明である）」という事実が代数幾何的に保証されています。
非線形偏微分方程式の可換性条件 (3.83) という極めて複雑な等式を、「線形作用素 $L$ の核が自明である」という代数的な性質に帰着させて証明し切ってしまう、極めて美しい論法です。
\end{Q}

\begin{story}
そこで，(3.32)を用いて $r(x,t)$ を計算しよう．実際，$r(x,t)$ は $x=0, x=1, x=t$ で高々2位の極をもつ有理関数である．特性指数を考慮すれば，$r(x,t)$ は

\begin{align}
    r(x,t) &= \frac{\alpha_0}{x^2} + \frac{\alpha_1}{(x-1)^2} + \frac{\alpha_\infty}{x(x-1)} \nonumber \\
    &\quad + \sum_{l=1}^g \left[ \frac{\beta_l}{(x-t_l)^2} + \frac{t_l(t_l-1)K_l}{x(x-1)(x-t_l)} \right] \label{eq:3.85} \\
    &\quad + \sum_{k=1}^g \left[ \frac{3}{4(x-\lambda_k)^2} - \frac{\lambda_k(\lambda_k-1)\nu_k}{x(x-1)(x-\lambda_k)} \right] \nonumber
\end{align}

という形をしている．他方，(3.32)に(3.67)と(3.68)を代入して，係数の関係を調べる．すると
\begin{equation*}
    \alpha_0 = \frac{1}{4}(\kappa_0^2 - 1), \quad \alpha_1 = \frac{1}{4}(\kappa_1^2 - 1), \quad \beta_l = \frac{1}{4}(\theta_l^2 - 1)
\end{equation*}
\end{story}

\begin{Q}[【深掘り】式(3.85)の見かけの特異点 $x=\lambda_k$ の項にある「$3/4$」という数字はどこから来たのか？]
\textbf{A.} これはSL型方程式（シュレディンガー方程式型）における**見かけの特異点の絶対的な刻印**です。
フックス型微分方程式をSL型にゲージ変換（$z = y \exp(\frac{1}{2}\int p_1 dx)$）すると、変換後のポテンシャル $r(x,t)$ の各特異点での2位の極の留数は、特性指数の差 $\Delta \rho$ を用いて $\frac{1}{4}((\Delta \rho)^2 - 1)$ という普遍的な形で書けることが知られています（シュワルツ微分の性質）。
確定特異点 $x=0, 1, t_l$ では特性指数の差がそれぞれ $\kappa_0, \kappa_1, \theta_l$ なので、$\alpha_0, \alpha_1, \beta_l$ はまさにこの公式通りになっています。
そして、見かけの特異点 $x=\lambda_k$ では、以前確認したように「特性指数は $0, 2$」です。したがってその差は $\Delta \rho = 2 - 0 = 2$ となります。これを公式に代入すると、
$$ \frac{1}{4}(2^2 - 1) = \frac{3}{4} $$
となります。つまりこの $\frac{3}{4}$ は、そこが「特性指数の差が $2$ である見かけの特異点である」という幾何学的な事実から代数的に必然として導かれる定数なのです。
\end{Q}



\appendix
\input{sections/99_appendix.tex}

% --- 参考文献 ---
\begin{thebibliography}{99}
    \bibitem{book:okamoto} 岡本和夫, 『パンルヴェ方程式』, 岩波書店, 1999.
    \bibitem{book:takano} 高野恭一, 『常微分方程式』, 朝倉出版, 1994.
    \bibitem{book:kimura} 木村俊房・高野恭一, 『関数論』, 朝倉出版, 1991.
    \bibitem{book:matsumoto} 松本幸夫, 『トポロジー入門』, 岩波書店, 1985.
    \bibitem{book:kobayashi} 小林俊行・大島利雄, 『リー群と表現論』, 岩波書店, 2005.
    \bibitem{book:noguchi} 野口潤次郎, 『複素解析概論』, 裳華房, 1993.
\end{thebibliography}

\end{document}